\PassOptionsToPackage{unicode=true}{hyperref} % options for packages loaded elsewhere
\PassOptionsToPackage{hyphens}{url}
%
\documentclass[]{book}
\usepackage{lmodern}
\usepackage{amssymb,amsmath}
\usepackage{ifxetex,ifluatex}
\usepackage{fixltx2e} % provides \textsubscript
\ifnum 0\ifxetex 1\fi\ifluatex 1\fi=0 % if pdftex
  \usepackage[T1]{fontenc}
  \usepackage[utf8]{inputenc}
  \usepackage{textcomp} % provides euro and other symbols
\else % if luatex or xelatex
  \usepackage{unicode-math}
  \defaultfontfeatures{Ligatures=TeX,Scale=MatchLowercase}
\fi
% use upquote if available, for straight quotes in verbatim environments
\IfFileExists{upquote.sty}{\usepackage{upquote}}{}
% use microtype if available
\IfFileExists{microtype.sty}{%
\usepackage[]{microtype}
\UseMicrotypeSet[protrusion]{basicmath} % disable protrusion for tt fonts
}{}
\IfFileExists{parskip.sty}{%
\usepackage{parskip}
}{% else
\setlength{\parindent}{0pt}
\setlength{\parskip}{6pt plus 2pt minus 1pt}
}
\usepackage{hyperref}
\hypersetup{
            pdftitle={A Student's History of Education},
            pdfauthor={Frank Pierrepont Graves},
            pdfborder={0 0 0},
            breaklinks=true}
\urlstyle{same}  % don't use monospace font for urls
\usepackage{longtable,booktabs}
% Fix footnotes in tables (requires footnote package)
\IfFileExists{footnote.sty}{\usepackage{footnote}\makesavenoteenv{longtable}}{}
\usepackage{graphicx,grffile}
\makeatletter
\def\maxwidth{\ifdim\Gin@nat@width>\linewidth\linewidth\else\Gin@nat@width\fi}
\def\maxheight{\ifdim\Gin@nat@height>\textheight\textheight\else\Gin@nat@height\fi}
\makeatother
% Scale images if necessary, so that they will not overflow the page
% margins by default, and it is still possible to overwrite the defaults
% using explicit options in \includegraphics[width, height, ...]{}
\setkeys{Gin}{width=\maxwidth,height=\maxheight,keepaspectratio}
\setlength{\emergencystretch}{3em}  % prevent overfull lines
\providecommand{\tightlist}{%
  \setlength{\itemsep}{0pt}\setlength{\parskip}{0pt}}
\setcounter{secnumdepth}{5}
% Redefines (sub)paragraphs to behave more like sections
\ifx\paragraph\undefined\else
\let\oldparagraph\paragraph
\renewcommand{\paragraph}[1]{\oldparagraph{#1}\mbox{}}
\fi
\ifx\subparagraph\undefined\else
\let\oldsubparagraph\subparagraph
\renewcommand{\subparagraph}[1]{\oldsubparagraph{#1}\mbox{}}
\fi

% set default figure placement to htbp
\makeatletter
\def\fps@figure{htbp}
\makeatother

\usepackage{etoolbox}
\makeatletter
\providecommand{\subtitle}[1]{% add subtitle to \maketitle
  \apptocmd{\@title}{\par {\large #1 \par}}{}{}
}
\makeatother
\usepackage{booktabs}
\usepackage{amsthm}
\makeatletter
\def\thm@space@setup{%
  \thm@preskip=8pt plus 2pt minus 4pt
  \thm@postskip=\thm@preskip
}
\makeatother
% https://github.com/rstudio/rmarkdown/issues/337
\let\rmarkdownfootnote\footnote%
\def\footnote{\protect\rmarkdownfootnote}

% https://github.com/rstudio/rmarkdown/pull/252
\usepackage{titling}
\setlength{\droptitle}{-2em}

\pretitle{\vspace{\droptitle}\centering\huge}
\posttitle{\par}

\preauthor{\centering\large\emph}
\postauthor{\par}

\predate{\centering\large\emph}
\postdate{\par}
\usepackage[]{natbib}
\bibliographystyle{apalike}

\title{A Student's History of Education}
\author{Frank Pierrepont Graves}
\date{1916}

\begin{document}
\maketitle

{
\setcounter{tocdepth}{1}
\tableofcontents
}
\hypertarget{a-warhorn-historicals-book}{%
\chapter*{\texorpdfstring{\href{https://historicals.warhornmedia.com/}{A Warhorn Historicals Book}}{A Warhorn Historicals Book}}\label{a-warhorn-historicals-book}}
\addcontentsline{toc}{chapter}{\href{https://historicals.warhornmedia.com/}{A Warhorn Historicals Book}}

\includegraphics[width=0.5\textwidth,height=\textheight]{images/warhornlogo.png}

Republished by Warhorn Historicals:
Making historical books available for free online in high quality, readable formats.

We hope this book is a blessing to you. If it is, please \href{https://warhornmedia.com/give}{make a one-time or recurring contribution} right now, sponsor a book from our upcoming list, or volunteer your proofreading or technical skills to help produce more content. Contact us \href{mailto:lucas@beggarsborn.com}{Lucas Weeks} to get involved.

God bless,

---The Warhorn Team

\hypertarget{historical-cover}{%
\subsubsection*{Historical Cover}\label{historical-cover}}
\addcontentsline{toc}{subsubsection}{Historical Cover}

\hypertarget{introduction}{%
\chapter*{Introduction}\label{introduction}}
\addcontentsline{toc}{chapter}{Introduction}

This is the introduction

\hypertarget{the-earliest-education}{%
\chapter{1. The Earliest Education}\label{the-earliest-education}}

\textbf{Outline}

\begin{itemize}
\tightlist
\item
  Even a brief survey of the history of education may greatly broaden one's view.
\item
  Starting with primitive man, we find that his training aims only at the necessities of life, and is acquired informally through the elders and the medicine-men.
\item
  In Oriental education, the next stage in progress, illustrated by India, a traditional knowledge is acquired through \emph{memoriter} and imitative methods.
\item
  While Oriental, Jewish education afforded greater development of individuality, but it was late in organizing schools, \emph{memoriter} in methods, and restricted in content.
\item
  Thus all education before the day of the Greeks was largely \emph{nonprogressive.}
\end{itemize}

\textbf{The Value of the History of Education.} - The History of Education from the earliest times should contribute largely to one's breadth of view and prove a study of the greatest liberal culture. A record of typical instances of the moral, aesthetic, and intellectual development of man in all lands and at all periods should certainly enlarge one's vision and enable him to appreciate more fully the part that education has played in the progress of civilization. Such cultural values may be found even in a limited survey of the world's educational development.

\textbf{It's Treatment in This Book.} - And this is all that will be undertaken here. For, while valuable as a liberal study, the History of Education finds its justification chiefly in the degree to which it functions in the professional training of a teacher, and it will be necessary in a brief treatise to omit or pass over hastily much that might be of interest and value in a more complete account of the development of civilization. Therefore, the amount of space and the perspective afforded the various peoples, epochs, and leaders must here be determined in large measure by the part they have played in the evolution of educational institutions and practices, and by the light their history sheds upon the aim, organization, content, and method of education to-day. At times, too, the history of a single epoch, state, or educational leader will be selected as a type, to the exclusion of others equally important, and treated with considerable intensiveness, instead of describing all sides of the subject with encyclopaedic monotony. Now the first historical epoch to leave a real impress upon modern practice is that of Athens at its height. Hence a mere statement of the salient features of education preceding that period is all that can be afforded in this brief survey. A detailed account of the educational processes used by savage tribes, Oriental nations, and even Judaea may prove interesting and important in other connections, but it must here be largely curtailed.

\textbf{Primitive Education.} There is little to be noted in the training of the young among primitive peoples, save that it is intended largely for the satisfaction of immediate wants - food, clothing, and shelter. Naturally no such actual institution as a school has yet been evolved, but the training is transmitted informally by the parents. The method used is simply that of example and imitation, or, more specifically, `trial and success.' But a more conscious and formal education is given at puberty through the `initiatory ceremonies' (Fig. i). In these rites the youths are definitely instructed by the older men about their relation to the spirits and the totem animals, subordination to the elders, the relations of the sexes, the sacredness of the clansman's obligations, and other traditional usages. Strict silence is enjoined upon them concerning this information, and to impress it upon their minds, and test their endurance, they are required to fast for several days and are often tortured and mutilated. As the savage does not clearly distinguish between himself and the tribe to which he belongs, there is practically no development of individuality, and since the race has not yet learned to treasure its experience in writing, he has no record of past experience and is virtually tied to the present.

\textbf{Oriental Education.} - The nations of the ancient Orient - Egypt, Babylonia, Assyria, China, India, and Persia - may be said to represent the next higher stage in civilization. Their systems of education prepare mostly for vocations, and are not sufficiently advanced to undertake a training for manhood or citizenship. But since a division of labor has now been evolved, the training has become more clearly differentiated and fits for specific occupations. In this way, class divisions, or even castes, have generally arisen in society, and the young people are educated according to the position in life they desire, or are required to fill. As an illustration of this stage of development, we may consider somewhat in detail the social environment and education of India.

\textbf{India: Its Religion and Castes} - In India, largely as a result of the debilitating climate, there was formulated about 1200 B. C. a dreamy philosophy, according to which nothing except Brahma, the one universal spirit, really exists. While men would seem to be temporarily allowed a separate existence of their own, it was held that they should remain inactive as far as possible and seek an ultimate absorption into the great Eternal Spirit. Although somewhat modified by the infusion of Buddhism, between 500 B. C. and 500 A. D., and by the British occupation of the peninsula during the nineteenth century, this mystic and static religion still dominates in India. Connected with it is the caste system, by which the people are divided into four hereditary classes. These are (1) the \emph{brahmins,} or sacerdotal class, which includes all those trained for law, medicine, teaching, and other professional occupations; (2) the warriors, or military and administrative caste; (3) the industrial group; and (4) the \emph{sudras,} or menial caste. Altogether outside the social order are the \emph{pariahs,} or outcasts. The caste system is exceedingly strict. One may fall into a lower caste, but he cannot rise, and loss of caste by one person in a family will degrade all the rest.

\textbf{The Hindu Education.} - Hence Hindu education has always endeavored to fill the pupils with the tenets of their religion, and so prepare them for absorption into the Infinite, rather than for activities in this life, and to preserve the caste system and keep all within the sphere of their occupation. The three upper castes are, therefore, supposed to gain a knowledge of certain sacred works, especially the four \emph{Vedas} or books of `knowledge,' the six \emph{Angas} on philosophical and scientific subjects, and the \emph{Code of Manu,} which is a collection of traditional customs; but few, outside the brahmin class, are ever allowed to take advantage of this opportunity. The warriors are expected to pay more attention to martial exercises, and the industrial caste to acquire through apprenticeship the arts necessary for its hereditary occupations. Sudras, pariahs, and women are generally allowed no education. Except the sudras, all the castes obtain elementary education from a study of the laws, traditions, and customs of the country through the medium of the family, and more recently through village schools held in the open air (Fig. 2). The higher education is largely carried on in brahminic colleges, called \emph{parishads,} and, as also in the case of the elementary work, the teachers have to be brahmins. Since all learning has been preserved by tradition, the chief methods of instruction are those of memorizing and imitation. Even the later texts are so written as to be easily committed, and the lines are sung aloud by the pupils until they have memorized them. Writing is learned by imitating the teacher's copy on the sand with a stick, then on palm leaves with a stylus (Fig. 2), and finally on plane leaves with ink.

\textbf{Effect of the Hindu Education.} - Hence, among the Hindus education is forbidden to ninety-five per cent of the population, and, as far as it does exist, it is a mere stuffing of the memory. It concerns itself but little with mental culture or with preparation for real living. The brahmins have handed down considerable traditional learning, grammar, phonetics, rhetoric, logic, `Arabic' notation, algebra, astronomy, and medicine, but new knowledge of any sort is barred. The Hindus still plow with sticks of wood, and their crops are harvested and threshed by devices equally primitive. They bake bricks, work metals, and weave cloth, but with the same kind of appliances that were used by their remote ancestors. Until recently, they have been greatly lacking in ambition, self-reliance, and personal responsibility, and have not yet come to any feeling of solidarity or national unity. To them prosperity and progress are foreign ideas.

\textbf{India as Typical of the Orient.} - The other countries of the ancient Orient never fixed their social classes in so hard and fast a manner, and have never included so elaborate a philosophy among the products of their culture. But India may well be considered broadly typical of the stage of development in the Orient. Certain common features appear in the education of all the nations there. In the system of each, the classes below the sacerdotal or priestly are given little intellectual education, and the women none at all, but both are trained by apprenticeship in their vocations. Actual schools, both elementary and higher, have been instituted; and the latter, except in China, are conducted at temples or priestly colleges by members of the sacerdotal class. The educational content is naturally traditional. It is, for the most part, ensured against change by being embalmed in sacred books, such as the \emph{Vedas.} The educational method consists largely in the memorizing of the text and imitation of the copy set, and little attempt is made to give a reason for the customs and traditional knowledge taught. Hence, while individuality has begun to emerge, it is suppressed by every agency possible; and, although these peoples have largely overcome the primitive enslavement to nature and the present, they are completely in bondage to the past.

\textbf{Jewish Education.} - The Jews are classed among the nations of the Orient, but they formulated loftier aims and have exerted more influence upon modern ideals in education. While their theology greatly developed in the course of their history, from the first they held to an ethical conception of God, and the chief goal of their education was the building of moral and religious character. Not until after the Babylonish captivity (586-536 B. C), however, did they establish actual schools. Before that, children were given an informal training in the traditions and observances of their religion by their parents. But they brought back from Babylon the idea of institutions for higher training and started such schools through their synagogues. In the second century B. C. the founding of elementary schools also began, and eventually the Jews made education well-nigh universal. The beneficial effect of this training is seen in the respect shown by the Jews for their women, their kind treatment of children, and their reverence for parents. The defects of their education appear in the stereotyped and formal way in which the religious material came to be interpreted, and the consequent hostility to science and art, except as they threw light on some religious festival or custom. Although appeal was made to various types of memory, systems of mnemonics devised, and other good pedagogical features suggested, their methods of instruction were largely \emph{memoriter.} The Jewish system of education, as a whole, afforded a greater development of personality than that of the other Oriental nations, and through it have been spread some of the world's most exalted religious conceptions. Nevertheless, it did not depart much from its traditions and the past, and to this extent it may be classed with the training of the primitive tribes and of the Oriental nations as predominantly \emph{non-progressive.}

Supplementary Reading

For general works, see Graves, F. P., \emph{History of Education before the Middle Ages} (Macmillan, 1909), chaps. I-XI; Monroe, P., \emph{Text-book in the History of Education} (Macmillan, 1905), chaps. I-II. A general interpretation of the evolution of education in savagery and barbarism is also given in Laurie, S. S., \emph{Pre-Christian Education} (Longmans, Green, 1909), pp.~1-207; Morgan, L. H., \emph{Ancient Society} (Holt, 1907), Part I; and Taylor, H. O., \emph{Ancient Ideals} (Macmillan, 1913), vol.~I, chaps. I-V. An illustration of primitive training of especial interest to American students is found in Spencer, F. C., \emph{Education of the Pueblo Child} (Columbia University, Department of Philosophy and Psychology, vol.~7, no. 1); and a detailed description of the puberty rites of a variety of savage tribes, in Webster, H., \emph{Primitive Secret Societies,} (Macmillan, 1908), chaps. I-V. A more complete account of the Hindu philosophy and education appears in Dutt, R. C, \emph{Civilization of India} (Dent, London), and Taylor, H. O., \emph{Ancient Ideals} (Macmillan, 1913), vol.~I, chaps. Ill and IV. A systematic statement of the Jewish training has been adapted from a German work, in Leipziger, H. M., \emph{Education of the Jews} (New York Teachers College, 1890), and a more detailed account worked out in Spiers, B., \emph{School System of the Talmud} (Stock, London, 1898).

\hypertarget{the-education-of-the-greeks}{%
\chapter{2. The Education of the Greeks}\label{the-education-of-the-greeks}}

\textbf{Outline}

\begin{itemize}
\tightlist
\item
  The Spartan training was intended to serve the state by making warriors, and little attention was paid to intellectual education.
\item
  At first the Athenian education was also mainly concerned in serving the state. For the earliest stage of the boy's education, there were schools of two types, - one for intellectual training, as well as one for physical; from fifteen to eighteen a more advanced physical training was given; and then, for two years, a preparation for military life.
\item
  After the Persian wars, the Athenians adopted ideals of education affording a larger recognition of individualism. The sophists introduced the new educational practices, and went to an extreme in their individualism.
\item
  The systematic philosophers, - Socrates, Plato, and Aristotle, tried to mediate the outworn institutional education and the extreme individualism. Socrates held that the sophistic `knowledge' was only `opinion,' and that the more universal knowledge could be reached in every person by stripping off his individualistic opinion.
\item
  But Plato maintained that only the intellectual class could attain to knowledge. For them he formulated a new course of study, in addition to that in vogue, consisting of mathematical subjects and dialectic. Aristotle held that the training for every one before seven should be bodily; up to fourteen, the irrational soul should be trained; and until twenty-one, the rational. While Plato and Aristotle had little effect upon educational practice at the time, they have since greatly influenced education.
\item
  After Aristotle, there arose individualistic schools of philosophy and formal schools of rhetoric, and out of them universities sprang up. Then Greek culture and education spread throughout the world.
\end{itemize}

\textbf{Progressive Nature of Greek Education.} - Real educational progress began with the Greeks. In their training gradually appeared considerable regard for individuality. They were the first people whose outlook seems to have been toward the future rather than the past, and they first made a serious attempt to promote human development in accordance with a remote ideal progressively revealed. As a result, they not only gave a wonderful impetus to educational practice in their own time, but ever since then the world has had constant recourse to them for inspiration and counsel. While this intellectual emancipation did not appear to any extent before its development among the Athenians in the middle of the fifth century B. C, well-planned systems of education existed in Greece several centuries before this and paved the way for the system in Athens during the Age of Pericles.

\textbf{Spartan Education: Its Aim and Early Stages.} - Among the states of ancient Greece, Sparta possessed the earliest education of which we have any extended information. Its citizens dwelt in the midst of hostile peoples they had subjugated, and this made it necessary to produce a race of hardy and patriotic warriors. Strength, courage, and obedience to the laws were held as the aim of education. The Spartan educational system was intended to serve the state, and the rights of the individual were given little or no consideration. State control began with birth. The infant was immediately inspected by a council of elders, and, if he were sickly or deformed, he was exposed to die in the mountains; but if he appeared physically promising, he was formally adopted by the state and left with his mother for rearing until seven. At that age the boys were placed in charge of a state officer and ate and slept in a kind of public barracks. Here their life became one of constant drill and discipline. In addition to hard beds, scanty clothing, and little food, they were given a graded course in gymnastics. Besides ball-playing, dancing, and the \emph{pentathlum} - running, jumping, throwing the discus, casting the javelin, and wrestling - the exercises included boxing, and even the brutal \emph{pancratium,} in which any means of overcoming one's antagonist - kicking, gouging, and biting, as well as wrestling and boxing - was permitted.

The Spartan boys, however, received only a little informal training in the way of intellectual education. They simply committed to memory and chanted the laws of Lycurgus and selections from Homer, and they listened to the conversation of the older men during the meals at the common table, and were themselves exercised in giving concise and sensible answers to questions put to test their wisdom. Every adult was also required to choose as his constant companion or `hearer' a youth to whom he might become an'inspirer.'

\textbf{Training in Youth and Manhood: Results.} - When a youth reached eighteen, he began the distinctive study of warfare. For two years he was trained in the use of arms and skirmishing, and every ten days had his courage and his physique tested by being whipped before the altar of Artemis. Then he regularly entered the army, and for ten years guarded some border fortress and lived upon the coarsest of fare. When he became thirty, he was considered a man and forced to marry at once, but even then he could visit his wife only clandestinely and was still obliged to live in common with the boys and assist in their training.

The education of women was very like that of the men. While the girls were allowed to live at home, they were given a similar physical training in the hope that they would become the mothers of sturdy sons. Thus the Spartan education was shaped entirely with reference to the welfare of the state. Their educational system served well its purpose of creating strong warriors and devoted citizens, but it failed to make for the highest manhood. Sparta developed practically no art, literature, or philosophy, and produced little that tended to promote civilization. She has left to the world little but examples of heroism and foolhardiness alike.

\textbf{Old Athenian Education: Its Aim and Early Training.} - For many centuries the Athenian education was not unlike the Spartan in promoting the welfare of the state without much consideration of individual interests. But even in early days Athens felt that the state was best served when the individual secured the most complete personal development. Hence, the Athenian boys began to receive at seven years of age two kinds of training - (1) the \emph{pentathlum} and other physical exercises in the \emph{palaestra} (Fig.3) or exercising ground, and (2) singing and playing upon the flute or lyre, and reading and writing at the \emph{didascaleum} (Fig. 4.) or music school. After the boy had learned his letters by tracing them in the sand, he was taught to copy verses and selections from well-known authors, at first upon wax-tablets with a stylus, and later upon parchment with pen and ink. It was, moreover, necessary for the pupils in singing to be taught the rhythm and melody, and to understand the poem so as to bring out its meaning. Hence the explanations and interpretations given by the teachers brought in all the learning of the times, and the moral and intellectual value of the studies must have been much greater than would be suggested by the meagerness of the course. Some moral training and discipline were also given the boy by a slave called the \emph{paedagogus,} who conducted him to school and carried his lyre and other appurtenances. This functionary was often advanced in years or incapacitated for other duties by physical disability.

\textbf{Training for the Youth.} - At fifteen the Athenian boy might take physical training of a more advanced character at one of the exercising grounds just outside Athens, which were known as \emph{gymnasia.} He was now ephebic course permitted to go wherever he wished and become acquainted with public life through first-hand contact. When eighteen the youth took the oath of loyalty to Athens, and for two years as an \emph{ephebus} or cadet continued his education with a course in military duties. The first year he spent in the neighborhood of Athens and formed part of the city garrison, but in the second year he was transferred to some fortress on the frontier. At twenty the young man became a citizen, but even then his training continued through the drama, architecture, sculpture, and art that were all about him.

\textbf{Effect of the Old Athenian Education.} - Little attention was, however, given by the Athenians to the education of woman. It was felt that her duties demanded no knowledge beyond ordinary skill in household affairs. With this exception, the Athenian education was superior to the Spartan in allowing greater opportunity for individual development and in furnishing a more rounded training. Nevertheless, until about the middle of the fifth century B.C., while differing considerably in degree from Sparta, Athens may be grouped with that country as adhering to the `old' education, where the individual was subordinated to the good of the social whole.

\textbf{Causes and Character of the New Athenian Education.} - This characterization is, of course, in contrast to Greek education in the `new' period, which is represented by Athens alone. This later type of education was probably somewhat the result of the gradual rise of democratic ideals in Athens, but a more immediate set of factors grew out of the Persian wars (492-479 B. C). This extended conflict with a powerful Oriental people, possessing a well-organized but widely different body of traditions tended to broaden the views of the Athenians greatly, and the ensuing political and commercial intercourse with a variety of dependent states and nations in the Delian League, together with social contact with the foreigners from every land that were thronging the streets of Athens, led even more directly to a reconstruction of practices and beliefs. A rapid transition in the old traditions took place and society seems for a time to have been sadly disorganized. The old was shattered, and while new ideals were being constructed, a groping ensued. Although the latitude given the individual was destined, as always, to produce progress in the long run, and was of great ultimate service to the world, more immediately a low ebb in morals at Athens resulted. Individualism ran riot. Education reflected the conditions of the period. Its ideals became more and more individualistic. The times demanded a training that would promote the happiness of the individual with little consideration for the welfare of the state as a whole. The old education seemed narrow and barren of content; and there arose a desire for all sorts of knowledge that might contribute to one's advancement, whether it increased his social usefulness or not. Skill in debate and public speaking was especially sought, because of the unusual opportunity for personal achievement in politics.

\textbf{The Sophists and Their Training.} - To meet these new demands, a set of teachers known as the \emph{sophists} came into prominence. They professed to train young men for a political career, and some of them even claimed to teach any subject whatsoever, or how to defend either side of an argument. These pretensions, together with their charging a fee for their services, contrary to Athenian custom, seriously offended the more conservative of the citizens of Athens. But many of the first sophists afforded an honest and careful training. The effect of their teaching was especially felt by the adolescents in the \emph{gymnasium} stage of education, since they were ambitious to distinguish themselves politically. The physical training that had hitherto dominated the gymnasium course gave way to a study of grammatical and rhetorical subtleties, and whenever a sophist appeared in the street, market-place, or house, the young men crowded about him to borrow from his store of experience and wisdom, and acquire his method of argument. To a less degree the same influence was felt in the lower schools and by the cadets and younger citizens. The exercises of the palaestra were no longer as rigorous, and existed for the sake of individual health and pleasure rather than for the making of citizens. The literary work of the didascaleum came to include, besides the Homeric epics, a wide range of didactic, reflective, and lyric poetry, with a superabundance of discussions. In music the old patriotic and religious songs sung to the simple Doric airs and accompanied upon the sevenstringed lyre, were replaced by rhythms of great difficulty, like the Lydian and Phrygian, and by complicated instruments of all sorts.

\textbf{Their Extreme Individualism.} - All this inroad upon the time honored curriculum shows how fully the sophists embodied the individualism of the times. Although they held no body of doctrine common to them all, they were generally at one in their position of extreme individualism. They often went so far as to insist that there could not safely be any universal criteria in knowledge or morals; that no satisfactory interpretation of life could be made for all, but that every fact and situation should be subject to the judgment of the individual. No doubt the formula attributed to Protagoras, ``Man (i. e. the individual) is the measure of all things, both of the seen and the unseen,'' would have expressed the attitude common to most of them. They but carried to its legitimate conclusion the complete reaction from the old ideal of subordination of the individual to the state.

\textbf{The Reactionaries and the Mediators.} - Meanwhile, the conservative element was making its usual attempt to adjust the unsettled conditions by suggesting a return to the old. Various schemes had been advanced, even before the sophists had come into prominence. Of these the most complete plan was that of Pythagoras (about 580-500 B.C.). By adopting an analogy from the `harmony' of the celestial bodies and from the relation of the powers in the individual to each other, he arranged a definite hierarchy in society, so that each member should have his proper place, and complete harmony and social order should ensue. As the influence of the sophists began to be felt, later representatives of the reactionary movement, such as the matchless caricaturist, Aristophanes (445-380 B.C.), began to appear and inveigh against the new conditions. But the social process can never move backward, and reconstruction on some higher plane was needed to overcome the destructive tendencies of the times. To furnish this, was the task set themselves by Socrates, Plato, and Aristotle. Like the sophists, they recognized that the traditional beliefs and sanctions, the old social order, and the former ideals and content of education, had been outlived, and that the individual could not find truth and morality through an institutional system. At the same time they felt that the extreme individualism of the sophists was too negative a basis upon which to build, and that a more socialized standard of knowledge and morality must be sought.

\textbf{The Method of Socrates.} - This mediating effort was begun by Socrates (469-399 B. C). While he started with the formula of Protagoras, he maintained that the `man' indicated thereby was not the individual, but mankind as a whole. It is not the peculiar view of any individual that represents the truth, but the knowledge that is the same for everyone. The former, which the sophists considered `knowledge,' Socrates held to be only `opinion, and declared that the reason men think so differently is because each sees but one side of the truth. He believed that everyone could get at universal knowledge by stripping off individual differences and laying bare the essentials upon which all men are agreed. He conceived it to be the mission of the philosopher or teacher to enable the individual to do this, and he endeavored to deal with the mind of all those with whom he came in contact, so that they would form valid conclusions. By his method, known as the \emph{dialectic,} or 'conversational,' he first encouraged the individual to make a definite statement of his belief, and then, through a set of clever questions, caused the person to develop his thought, until he became so involved in manifest contradictions that he was forced to admit that his view had been imperfectly formed. He thus caused the individual to see that the view he had first expressed was mere `opinion' and but a single phase of the universal truth. As Socrates further held that morality consists in right knowledge and made no distinction between the knowledge of an action and the impulse to perform it, he strove through his methods of developing knowledge to harmonize the individual welfare with that of the social group.

\textbf{Plato's System of Education for the Three Classes of Society.} - But the believers in the old traditions and institutional morality felt that Socrates was atheistic and immoral. They persuaded Athens to give him the hemlock, and thus destroyed the man who might have proved her savior. A pupil, Plato (427-347 B. C), undertook to continue his work, but his aristocratic birth and temperament caused him to underestimate the intelligence of the masses. He held that they were incapable of attaining to `knowledge' - that they possessed only `opinion.' In his most famous dialogue, \emph{The Republic,} he endeavors to show that the ideal state can exist only when the entire control of the government is entrusted to the `philosophers,' or intellectual class, who alone possess `real knowledge.' Those who are to compose the three classes of society Plato would have selected during the educational process on the basis of their ability. For all boys up to eighteen years of age he prescribes an education similar to that in vogue in the palaestra, didascaleum, and gymnasium, except that he would somewhat expurgate the literary element, and would confine the musical training to the simpler melodies and instruments. The youths who prove capable of going beyond this lower education are next to take up the cadet training between eighteen and twenty, but those who are incapable of further education are to be relegated to the industrial class. During the cadet period are to be determined those capable of going on with the higher education of philosophers, while those who here reach their limit become members of the military class.

As Athenian education did not extend beyond the twentieth year, Plato is here obliged to invent a new course of study that will enable the future philosophers to acquire the habit of speculation. This additional course, he declares, should also be graded, in order that a further test of intellectual and moral qualities may be made. Arithmetic, plane and solid geometry, music, and astronomy, are to occupy the first ten years of the course. These subjects, however, are not to be studied for calculation or practical purposes of any sort, but entirely from the standpoint of theory or the universal relations underlying them, since only thus can they furnish a capacity for abstract thought. After this, at thirty, the young men who can go no further, are to be placed in the minor offices of the state, while those who have shown themselves capable of the study of dialectic, go on with that subject for five years longer. It then becomes the duty of these highest philosophers to guide and control the state until they have reached the age of fifty, when they may be allowed to retire.

\textbf{The Weakness of Plato's System.} - Thus, where Socrates found the basis of universal truth in everyone, Plato held that only one class of people, the most intellectual, could attain to real knowledge. He, therefore, maintained that the philosophers should absolutely guide the conduct of the state, and that education should be organized with that in view. Plato's ideal state would thus become a sort of intellectual oligarchy, and in a way was a return to the old principle of subordinating the individual to society. \emph{The Republic} thus quite neglected human will as a factor in society and assumed that men can be moved about in life like pieces upon the chess board. Plato failed to see, too, that each individual really possesses all human characteristics. The workers have reason, and the philosophers have passions, and a human being is not a man unless all these functions are his. But even if his scheme had been a happy one, the treatise provided no method of evolution from current conditions, and if it were further granted that this order of things could be established at once. Plato put the ban upon all innovation or change, and so closed the door to progress.

Hence \emph{The Republic} was viewed as a visionary conception, and had no immediate effect upon education or any other institution of Athens. So in his declining years, without denying \emph{The Republic} as ideal, he wrote the more practical dialogue known as \emph{The Laws.} In it he welded elements from the educational systems of Sparta and older Athens, and reverted to traditions and ideals not dissimilar to the doctrines of Pythagoras. He replaced the philosophers with priests, an hereditary ruler, a superintendent of education, and various other officials; and the course of study reached its height with the subject of mathematics, while dialectic was not mentioned.

\textbf{His Influence upon Educational Theory and Practice.} - Thus the efforts of Socrates, as continued by Plato, to obtain the benefit of the growing individualism for society and education without disrupting them, had seemingly come to naught. Nevertheless, Plato has had considerable influence upon the thought and practice of men since the Greek period. The ideal society where everything is well managed and everyone is in the position for which nature intended him, has ever since the day of \emph{The Republic} been a favorite theme for writers, as witness More's \emph{Utopia} and the \emph{New Atlantis} of Bacon. A specific movement that shows the impress of Plato, as we shall see later, is the formulation of the more advanced studies of the mediaeval `seven liberal arts' under the name of the `quadrivium.' It is even possible that the whole conception of `liberal' studies, and so the doctrine of `formal discipline' (see p.~182), may be traced back to Plato's idea that the mathematical subjects in the course for philosophers should never be studied from a practical point of view. On the whole, Plato has been a factor in educational theory and practice that cannot be overlooked.

\textbf{Aristotle's Ideal State and Education.} - A more practical attempt to unify the new with the old in Athenian society and education was made by Aristotle (386322 B. C), the pupil of Plato. From his father, the court physician at Macedon, and from his study under Plato, Aristotle obtained an excellent scientific training, which is evident in the way he approaches his problems. It is in his \emph{Politics} especially that he discusses the ideal state and the training of a citizen. His method of investigation to determine the nature of this ideal state is inductive, and before formulating his conception of it, he makes a critical analysis of Plato's \emph{Republic} and \emph{Laws,} and analyzes the organization of many other states, both ideal and actual. He concludes that a monarchy is theoretically the best type of government, but that the form most likely to be exercised for the good of the governed is the democracy. He then considers in detail the best natural and social conditions for a state. Among these practical considerations is the proper education to make its citizens virtuous.

Since virtue is of two kinds, moral or practical, and intellectual or speculative, and the former is merely the stepping-stone to the latter, the education needed for the virtue of the state must not, like that of Sparta, be purely a training for war and practical affairs. In marking off the periods of education, Aristotle holds that ``the care of the body ought to precede that of the soul, and the training of the impulsive side of the soul ought to come next; nevertheless, the care of it must be for the sake of the reason, and the care of the body for the sake of the soul.'' The development of the body he wishes to start even before birth by having the legislator ``consider at what age his citizens should marry and who are fit to marry.'' Also he deems it necessary to sanction the usage of his time of `exposing' (see p.~13) all deformed and weakly children. However, his advice concerning the food, clothing, and exercise of children is humane and in keeping with the best modern hygiene.

The training of the body is a preparation for the formal schooling, which is to last from seven to twenty-one. This is divided into two periods by puberty, the first to be devoted to the training of the impulsive or irrational side of the soul, and the second to that of the rational side. Education, he claims, should be public, as in Sparta, for it is the business of the state to see that its citizens are all rendered virtuous. However, the industrial classes, not being citizens, have no need of education, and women are to be limited in the scope of their training. The course of study for the irrational period is largely the same as that in use at Athens, - gymnastics, music, and literary subjects, although he recommends some reforms. Gymnastics is intended for self-control and beauty of form, and the making of neither athletes nor warriors should be the object, since the training of the former exhausts the constitution, and that of the latter is brutalizing. The literary subjects, which with Aristotle includes drawing, as well as reading and writing, are not to be taught merely for utilitarian reasons. Music is to be used not so much for relaxation or intellectual enjoyment as for higher development. Since melodies that afford pleasure are connected with noble ideas, and those which give us pain are joined to debased ideas, the study of music ``cultivates the habit of forming right judgments, and of taking delight in good dispositions and noble actions.'' Another moral effect of music is that it produces \emph{katharsis} or `purification'; that is, by arousing in us pity and fear for humanity at large, it lifts us out of ourselves and affords a safe vent for our emotions.

Such was to be the training for the body and for the irrational period, but how Aristotle would have advised that the education of the rational soul be carried on can only be surmised, since the treatise breaks off suddenly at this point. It is probable that it would have included a higher training in mathematical subjects and dialectic similar to that advocated by Plato, and, from Aristotle's own predilections, he would have been likely also to add some of the physical and biological sciences.

\textbf{The Permanent Value of His Work.} - Thus Aristotle, like Plato, endeavored to work out the harmonizing of individual with social interests by the creation of an ideal state, and he similarly failed to answer the demand of the times. His work was much less visionary than \emph{The Republic,} but he did not fully recognize that the day of the small isolated states of Greece, with their narrow prescriptions for patriotism and social order, had passed forever. Hence he hoped to achieve some reform by departing but little from existing conditions and reading a philosophy into them, and this bondage to the times prevented his educational system from making any advance beyond that of Plato. But while Aristotle had little effect upon the society of the times, his works have since been considered of great value, and the methods that he formulated have been most important. He not only started, or made the first great contributions to a number of sciences, but he crystallized the laws of thought itself. Also, as instruments to assist in fashioning the various sciences, Aristotle invented a complete system of terminology, and created such pairs as `matter' and `form,' `mean' and `extreme,' and `cause' and `effect,' and such convenient expressions as `principle,' `maxim,' `habit,' and `faculty.' A more important effect of Aristotle's ideas has been that upon the formulation of doctrine in the Christian Church. After the spread of Mohammedanism, which had largely absorbed the Aristotelian principles, the Church, though at first bitterly opposing them, finally found it impossible to suppress them, and began to clothe her own doctrine in their dress. The greatest of the scholastics began to study Aristotelianism, and soon made it the effective weapon of the Church by reducing all human knowledge to a finished Aristotelian system with theology at the top.

\textbf{The Post-Aristotelian Schools of Philosophy.} - But the harmonizing attempt of Aristotle was fruitless. Like Socrates and Plato, he failed to reconcile with the old and settled order the ever-expanding movement toward individualism. Thus all efforts to control the individualistic and disintegrating tendencies of the times were in vain, and the conquest of the Greek states by Philip of Macedon (358-338 B. C) was only symptomatic of the complete collapse of corporate life and the inability to reconstruct it successfully. All possibility of social unity disappeared, and philosophy no longer considered the individual from the standpoint of membership in society. It was occupied no further with the harmonization of the individual and the state, but concerned itself with the welfare of the individual and the art of living. Individualism was completely triumphant, and education was considered simply as a means to personal development or happiness, without regard to one's fellows. The new theories of life and education were formulated by such schools of philosophy as the Epicureans, Stoics, and Skeptics, which kept themselves far removed from society. None of these `schools' could be so termed in the sense of offering an education, but rather in the modern usage of a group of adherents to certain teachings. They spent their energy, for the most part, in interpreting, elaborating, and lauding the original teachings of the founders, and with them a stereotyped dogmatism took the place of philosophy.

\textbf{The Schools of Rhetoric.} - But these schools were not the only outcome of the teaching of the sophists. Just as they came about gradually from the speculative tendencies of the sophists as developed through certain famous philosophers, there likewise grew up more directly from the sophistic efforts to train young men in rhetoric and public speaking a multitude of rhetorical schools. In these a formal study was made of oratory and the knowledge of the day. Their professed object was to make successful men of the world, and, although they at first included such reputable and influential schools as that of Isocrates (436-338 B. C), they laid little claim to teaching anything solid or profound, much less to forming any philosophic habits. They succeeded in spreading a popular education among a people that had lost all hope of a political life, but they soon degenerated into the use of narrow and formal methods. The later rhetoricians attempted to hasten oratorical training and preparation for life, by teaching their pupils ready-made speeches and dialogues, together with a general knowledge of current questions. Nevertheless, these schools flourished for several centuries and closely rivalled those of the philosophers.

\textbf{The Hellenic Universities.} - From these two classes of schools, the philosophical and the rhetorical, the fame of Athens spread rapidly, and from the fourth century B. C. onward the number of young men from all over the civilized world who came there to study steadily increased. Before the close of the century the old cadet training of Athens was united with this intellectual education, and there sprang up a regular institution or university, which the young Athenians and students from outside might attend. Before long, the Hellenic world boasted other universities, such as those at Rhodes, Pergamon, Alexandria, and Rome. Until almost 300 A. D. Athens remained the chief intellectual center of civilization, and attracted students from all parts of the Roman Empire. Gradually, however, the higher education there tended toward the study of rhetoric alone and artificiality grew apace. In consequence, Alexandria came to displace Athens as the center of culture, and her university became the leading one of the world. Here the various philosophic and religious sects gathered to study and discuss, and the abstract Greek philosophy united with the more concrete beliefs of the Orient, especially Zoroastrianism, Judaism, and Christianity. Thus there flourished here the various systems of religious philosophy known collectively as `Hellenistic,' such as Neopythagoreanism, Neomazdeism, Philonism, Gnosticism, and Neoplatonism. Considerably before this, too, there had developed at Alexandria the Ptolemaic theory of the universe. Other noted investigations, like those of Euclid in geometry, Archimedes in physics, Eratosthenes in astronomy, and Diophantus in algebra, also bore witness to the intellectual activity of this university.

\textbf{Extension of Hellenic Culture.} - It can thus be seen that the political downfall of Athens had only prepared the way for a larger intellectual influence. As Alexander extended his yoke over one Eastern country after another, he had carried with him all the culture of Greece, and within a century of his death the whole Orient was dotted with Greek gymnasia, stadia, and theaters, and saturated with Greek literature, art, philosophy, and education. Similarly Rome, which had come somewhat into contact with Greece before conquering her, had been tinctured with Greek life and learning; and, after her absorption of Macedon and Greece, she fell under the spiritual thrall of the subjugated people. The history of Greek civilization and education was so intermingled with the Roman that it can scarcely be distinguished from it. The Greek schools of philosophy and rhetoric were continued in Rome, Roman youths made up a great body of the attendance at the universities of Athens and Alexandria, and the Roman emperors did much for the support and extension of the work in these institutions. Hence from the Greeks have developed some of the most advanced intellectual and aesthetic ideas that civilization has known.

\textbf{Supplementary Reading}

Graves, \emph{Before the Middle Ages} (Macmillan, 1909), chap. XII; Monroe, \emph{Text-book} (Macmillan, 1905), chap. III. See also Laurie, \emph{Pre-Christian Education} (Longmans, Green, 1900), pp.~208-318. Davidson, T., in his \emph{Aristotle} (Scribner, 1896), develops the periods of Greek education in chronological order, and his \emph{Education of the Greek People} (Appleton, 1903) gives the social setting of its development. A most scholarly and brilliant work is Freeman, K. J., \emph{Schools of Hellas} (Macmillan, 1907), which is illustrated by vasescenes and other reproductions of Greek education. Bosanquet, B., \emph{The Education of the Young in Plato's Republic} (Cambridge University Press, 1908), Nettleship, R. L., \emph{Theory of Greek Education in Plato's Republic} (See Evelyn Abbott's \emph{Hellenica,} Longmans, Green, 1908), and Burnet, J., \emph{Aristotle on Education} (Cambridge University Press) afford a good interpretation of the theorists mentioned; while Capes, W. W., in the \emph{University Life in Ancient Athens} (Harper, 1877), and Walden, J. W., in the \emph{Universities of Ancient Greece} (Scribner, 1909), furnish a lively description of the students and professors.

\hypertarget{the-education-of-the-romans}{%
\chapter{3. The Education of the Romans}\label{the-education-of-the-romans}}

\textbf{Outline}

\begin{itemize}
\tightlist
\item
  The contribution of the Romans to progress was largely due to their absorption of Greek culture, but their primitive training had an influence in itself. This was mostly civic and practical, and was given informally in the family and the forum.
\item
  Through amalgamation with the Greek, Roman education maintained three grades of schools: (i) the elementary school or \emph{ludus,} (2) the `grammar' school, and (3) the rhetorical school. Beyond the education of these schools, a young Roman might attend a university.
\item
  Schools were gradually subsidized by the emperors, but education eventually deteriorated into a formal qualification for senatorial rank. The practical Romans, however, created a universal empire and legal system, a universal religion, and other institutions for modern society.
\end{itemize}

\textbf{Roman Education Amalgamated with Greek.} - The name of Rome is still suggestive of power and organization. These characteristics seem to have been innate; but the significance of Roman development to the history of progress and education was largely due to the fact that, in her spread over the civilized world, the Eternal City amalgamated the Greek civilization with her own. Until then her ideals of life, while effective in conquest, had been narrow and little adapted to the development of individuality or of cosmopolitanism. Unconsciously realizing the need of broader ideals, she absorbed those of Greece. But Rome could not be Hellenized without making some contributions to the result from her own genius, and for that reason it is important to learn something of Roman civilization and education, crude as they were, before they came into contact with Greek culture.

\textbf{Early Education in Rome.} - In the early days Rome was animated by intense patriotism and love for military life, and felt that each citizen was bound to merge his identity in that of the state. In the surrender of individuality they were, to be sure, not unlike the Spartans, although they believed that this subordination should be brought about voluntarily rather than by compulsion of law and state organization. But, with such a love as theirs for mere material achievement, the Athenian ideal of a full and harmonious development of one's whole nature could scarcely be expected to make any appeal. They looked not for harmony, proportion, or grace, but for stern utility. They were sedate, grave, and serious, and their education was practical, prosaic, and utilitarian.

Until the Greek institutions began to be adopted, schools did not exist in Rome, except possibly the \emph{ludus} or elementary school. During this pristine period education consisted in a practical training in Roman ideals and everyday living conducted largely through the family. In childhood the boys and girls alike were given a physical and moral training by their mother, but, as the boy grew older, he went more in the company of his father, and learned efficiency in life informally through his example and that of the older men, while the girl was taught at home by her mother. If the boy belonged to a patrician family, he might acquire much knowledge concerning Roman custom and law by hearing his father advise and aid the family \emph{clients,} or `dependents,' and by attending banquets with him. He might also receive an apprenticeship training from his parent or some other older man in the profession of soldier, advocate, or statesman. In case he was born in a less exalted station, he might learn his father's occupation at the farm or shop. The girl, whatever her social status, was trained by her mother in the domestic arts, especially in spinning and weaving wool. Through their parents children probably learned to read and write; and they committed to memory stories of Roman heroes, ballads, martial and religious songs, and the \emph{Twelve Tables} of national laws, after these had been codified (451 B. C). Physical exercise was secured largely by games, which were mostly in imitation of future occupations, and gymnastics were employed simply as training for war. The usages of home and public religion also played an important part in the education of the young Romans, especially since almost every activity in life was presided over by some deity, whom it was necessary to propitiate when engaging in it. Thus education in early Rome was practical, and, to some extent, occupational. It was intended to produce efficiency as fathers, citizens, and soldiers. It consisted in training the youths to be healthy and strong in mind and body, and sedate and simple in their habits; to reverence the gods, their parents, the laws, and institutions; and to be courageous in war, and familiar with the traditional agriculture, or the conduct of some business. It did produce a nation of warriors and loyal citizens, but it inevitably tended to make them calculating, selfish, overbearing, cruel, and rapacious. They never possessed either lofty ideals or enthusiasm. Their training was best adapted to a small state, and became unsatisfactory when they had spread over the entire Italian peninsula. The golden age of valor and stern virtue had then largely departed, and they began unconsciously to seek a more universal culture. While such a people regarded the Greeks as visionary, just as the Greeks looked upon them as barbarians, they felt instinctively that only by absorption of the Hellenic ideals could their cosmopolitan ambitions be carried out. On the other hand, it was through the organization which the Romans were able to furnish, that the great ideals formulated by the Greeks were destined to be rendered effective and to become a matter of value and concern to civilization ever since.

\textbf{The Absorption of Greek Culture.} - There was a gradual infiltration of Greek culture into Rome from very early days. This received a great impulse through the conquests of Alexander (334-323 B. C) and the absorption of Macedon by Rome (168 B.C.), but it was not until about half a century after Greece itself had become a Roman province (146 B. C.), that the Greek educational ideals and institutions can be said to have been completely absorbed by Rome. This new type of education was thus well established early in the first century B.C. It may be said to have remained almost unmodified until toward the end of the second century A. D., when political conditions at Rome became most unstable and the period of degeneracy set in. During these three centuries of Hellenized Roman education, three grades of schools resulted from the amalgamation. They were the (1) \emph{Iudus} or school of the \emph{litterator,} as the lowest school was called; (2) the `grammar' school, taught by a \emph{grammaticus} or \emph{litteratus;} and (3) the schools of rhetoric and oratory, which furnished a somewhat higher education.

\textbf{The Ludus.} - The ludus, or lowest school, may possibly have existed before the process of Hellenization even began, but if it did, it must have been intended simply to supplement the more informal training of the home. Whenever originated, it probably taught at first only reading, writing, and rudimentary calculation, as in the family, through the medium of historical anecdotes, ballads, religious songs, and the \emph{Twelve Tables.} But as the Greek influence crept in more and more, the literary content was somewhat extended. About the middle of the third century B. C, Livius Andronicus translated the \emph{Odyssey} into Latin; and a number of epics, dramas, and epigrams were soon composed after Greek models. These works, in whole or part, were introduced into the curricula of the \emph{ludi,} and by the beginning of the first century B.C., the \emph{Twelve Tables} had been displaced by the Latinized \emph{Odyssey} of Andronicus. The methods of instruction were \emph{memoriter} and imitative. The names and alphabetic order of the letters were first taught without any indication of their significance or even shape, and all possible combinations of syllables were committed before any words were learned. Reading and writing were then taught by dictation, and, in tracing the letters on wax-tablets with the stylus (Fig. 5), the hand of the pupil was at first guided by the teacher. Calculation was learned by counting on the fingers, by means of pebbles, or upon the abacus, and eventually sums were worked upon the tablets.

Methods so devoid of interest were naturally accompanied by severe discipline. The rod, lash, and whip seem to have been in frequent use, and the names ordinarily applied to schoolmasters in Latin literature are suggestive of harshness and brutality. Moreover, a fresco found at Herculaneum depicts a boy held over the shoulders of another, with the master beating the victim upon the bare back (Fig. 6). Under these circumstances, no real qualifications were required of the teacher, and his social standing was low. The Greek custom of having the boy accompanied to and from school by a slave that was otherwise incapacitated by age or physical disability soon came to be imitated by the Romans. When a special building was employed for the school, it was usually a mere booth or veranda, and the pupils sat on the floor or upon stones.

\textbf{Grammar Schools.} - The `grammar' school grew out of the increasing literary work of the \emph{Iudus.} But, while offering a more advanced course, it would seem to belong in part at least to the elementary stage of education, especially as its work was never sharply divided from that of the \emph{Iudus.} The young Roman might attend both a Greek and a Latin grammar school, but, in case he did, usually went first to the former. The curriculum in each consisted, according to Quintilian, of `the art of speaking correctly' and `the interpretation of the poets,' or, in other words, of a training in grammar and literature. `Grammar' may, however, have included some knowledge of philology and derivations, as well as drill on the parts of speech, inflections, syntax, and prosody, and practice in composition and paragraphing. The literary training was obtained by writing paraphrases of the best authors, textual and literary criticism, commentaries, and exercises in diction and verse-writing. Some other studies, like arithmetic, geometry, astronomy, geography, and music may also have been added in time, from the suggestions of Plato, but the Romans naturally gave them a practical bearing. Some gymnastics, mostly for military training, were often in the course. The methods in the grammar schools were somewhat better than those of the \emph{Indus,} but the commentary of the teacher on the text was usually taken down \emph{verbatim} by the pupil. The discipline, in consequence, was not much in advance of that of the lower schools. But the accommodations for these secondary schools were decidedly superior, and the buildings not only possessed suitable seats for the pupils and teacher, but were even adorned with paintings and sculpture.

\textbf{Rhetorical Schools.} - The `rhetorical' schools were a development of work in debate that had gradually grown up in the grammar schools. The earliest of these institutions at Rome were Greek, but by the first century B. C, there had arisen a number in which Latin was used. While they afforded a legal and forensic training, and seem more professional in spirit than the grammar schools, they were by no means narrow. The orator was for the Roman the typical man of culture and education, and he was supposed not only to have been trained in eloquence and law and history, but to possess wide learning, grace, culture, and knowledge of human emotions, sound judgment, and good memory. Besides a training in oratory, these schools furnished a linguistic and literary education of some breadth. They may be considered as belonging partly to the secondary and partly to the higher stage of education. The youths were exercised first in declamation on ethical and political subjects, which would bring in fine distinctions in Roman law and ethics, and later they were given practice in three types of speeches, - deliberative, judicial, and panegyric. Attention was given to all the various factors in making a successful oration: the matter, arrangement, style, memorizing, and delivery.

\textbf{Universities.} - When the young Roman had completed his course at a rhetorical school, he might, if he were ambitious, go to the university at Athens, Alexandria, or Rhodes for a higher training. Later, a university also sprang up at Rome, and before long these institutions spread throughout the empire. The Greek influence caused a large number of these institutions to be established in the East, but some were also located in the West. The latter gave more emphasis to practical subjects. In several instances the universities found their nucleus in one of the many libraries that were started with books brought from the sacking of Greece and Asia Minor.

\textbf{Subsidization of Education.} - Thus, through the adoption of the institutions of the Greeks, Roman education became thoroughly Hellenized. Although all the types of schools spread everywhere in the empire, there was, of course, no such thing as a real school system, except as the government gradually came to subsidize all schools. This the different emperors accomplished in various ways, - by contributing to school support, paying a salary to certain teachers, or granting them exemption from taxation and military service, or offering scholarships to a given number of pupils. As a result, schools came to be established in many cases for the purpose of getting these special privileges for the teachers, rather than for promoting education. To stop these abuses, the emperor in 425 A. D. decreed that he had the sole authority to establish schools, and that a penalty would be laid upon anyone else assuming this prerogative. In this way the schools came fully into the hands of the imperial government, and the basis for the idea of public education was laid for the first time in history.

\textbf{Decay of Education.} - Before this, however, Roman education had deteriorated. With the political and moral decay that were obvious after the second century superficial A. D., it became a mere form and mark of the aristocracy. The training in oratory was continued, because it was a necessary qualification for entering the senatorial class, but it had lost its real function, since there was no longer any occasion for oratory when the emperor dominated all the government and law. It was not intended to furnish a training of any value in life, and the careful literary preparation was more and more shirked. While the grammarians and rhetoricians were still held in high esteem, they contented themselves with mere display, and wandered from town to town more for the purpose of entertaining than of teaching. Glittering phrases, epigrams, and other artificialities took the place of instruction and argument.

\textbf{Influence of Roman Education.} - But the Roman education and civilization had left their impress upon the world. This was accomplished by the practical nature of the Romans, and by their ability to make abstract ideals concrete and embody them in institutions that have been useful to civilization and progress. Through them was created the idea of a universal empire, which has been influential throughout the world's history. Similarly, the concept of law originating with the Greek philosophers became in the hands of the Romans the great system of principles that underlies and guides all our present civilization. And it was the Roman genius for organization that institutionalized a despised religious sect and expanded it into the position of the greatest world religion. If Judaism furnished the world with exalted religious ideals, and if from Hellenism came striking intellectual and aesthetic concepts, the institutions for realizing these ideals originated with Rome.

\textbf{Supplementary Reading}

Graves, \emph{Before the Middle Ages} (Macmillan, 1909), chap. XIII; Monroe, \emph{Text-book} (Macmillan, 1905), chap. IV. Interesting brief monographs on the subject are Clarke, G., \emph{Education of Children at Rome} (Macmillan, 1896), and Wilkins, A. S., \emph{Roman Education,} (Cambridge University Press, 1905). See also the treatment in Laurie, \emph{Pre-Christian Education} (Longmans, Green, 1900), pp.~319436.

\hypertarget{the-education-of-the-early-christians}{%
\chapter{4. The Education of the Early Christians}\label{the-education-of-the-early-christians}}

\textbf{Outline}

\begin{itemize}
\tightlist
\item
  Christianity accomplished much in the reform of the degraded Roman society. The earliest education of the Christians came through their `otherworldly' life, but actual schools, called `catechumenal,' before long furnished a moral and religious training.
\item
  After the amalgamation of Christianity with Grseco-Roman philosophy, `catechetical' schools furnished a higher training. When higher education came to be utilized by the bishops for training their clergy, institutions known as `episcopal' or `cathedral' schools were founded.
\item
  Later, although opposition grew up among the Christians to the culture of Greece and Rome, its impress was found to have been left upon the doctrines and organization of Christianity.
\end{itemize}

\textbf{The Ideals of Early Christianity.} - The actual social conditions amid which the religion of Christ was born, and which it was destined to reform, were most degraded. The Roman world had become sunk in vice and corruption. The Roman virtues of patriotism, bravery, and service to the state had largely disappeared with the development of the empire, and were impotent in checking the widespread depravity. Nor could the lofty Greek thought accomplish much, since it was too intellectual and philosophic to touch the masses. The debased Eastern religions, which Rome had admitted in her easy-going skepticism, were still less productive of good. While the more philosophic forms of Judaism and the Roman development of Stoicism tended to raise the tone of morals and pave the way for Christianity, not even these forces could have accomplished a successful reform in Roman society, without the stimulus and wide appeal of the Christian teachings. Christianity was the ethical and universal religion needed as a leaven. Its truths were based on faith rather than understanding, and its appeal was to the instinctive promptings and emotions rather than to the intellect. This made it democratic and enabled it to reach the masses, for everybody can feel and have faith, even where he cannot understand.

\textbf{Early Christian Life as an Education.} - Thus it came about that, while the earliest Christians were without schools of their own and were largely illiterate, their religion itself served as an education. They were practically deprived of intellectual development, but they received moral training of a very high order. The very dishonor and unpopularity of the Christian religion, and the segregation of their Church membership, gave the Christian life itself all the effect of a species of schooling. The early Christians showed an extreme reaction to the vicious morals of the time, and endeavored to cultivate the higher ideals inculcated by the teachings of Christ. They had gathered from the statements of the Master that he would soon return and this world would come to an end. They, therefore, concerned themselves entirely with a preparation for `Jerusalem the golden' and `the life everlasting,' and the ideal of this most primitive Christian training may be described as `other-worldly.'

\textbf{Catechumenal Schools.} - Early in the second century, however, when the Church began to extend itself rapidly, it seemed necessary to insist upon some sort of formal instruction as preliminary to Church membership. It was also deemed wise to fix a period of probation after the profession of one's faith in Christ, in order that informers might not be admitted to the services, or the Church disgraced by apostasy or the lapses of those who had not well considered the step. These demands were met by the gradual institution of popular instruction in Christian principles for the Jewish and pagan proselytes, who were known as \emph{catechumens.} While some effort was made to lift the pupils of these `catechumenal' schools from the bondage of ignorance, they were primarily trained in the things needful for their souls' salvation, and the ideal of Christian education remained prevailingly `otherworldly.' The instruction was carried on in the portico or other special portion of the church; and consisted in moral and religious teachings, reading and memorizing the Scriptures, together with some training in early psalmody. The course usually lasted three years, and while some distinction was made between the general division of catechumens and those almost ready for baptism, there is little ground for supposing that the schools were divided into actual classes. The meetings in the church were held several times a week, or even every day.

\textbf{Amalgamation of Christianity with Graeco-Roman Philosophy.} - But while the Christian ideals and training were developing and crystallizing, the Greek philosophy in its Roman form was being continued and expanded. This movement has been seen to be very different from early Christianity in its general purpose. It concerned itself chiefly with life in this world. The problem it attempted to solve was how one should live so as to get the most satisfaction out of life. The Hellenized Roman schools may, therefore, be accounted as `worldly' as the Christian schools were `otherworldly' in their aim. A general feeling of this marked difference in purpose and organization between Christianity and the contemporaneous Graeco-Roman culture was destined to cause an opposition to pagan learning to spring up among the Christians. But for two or three centuries this is scarcely noticeable, especially in the Eastern empire, where it was felt that philosophy was, like Christianity, a search after truth; and, as far as it went, confirmed the Bible. There was even a tendency to unite the two movements. As the new religion spread throughout the Roman world, and was compelled to defend itself against charges of immorality, atheism, and treason, the educated converts attempted to set forth the Christian teachings in terms of Greek thought, and to solve speculative problems that had never been considered by Jesus and his disciples. The first Hellenizing Christians are known as \emph{Apologists,} since their efforts were directed toward reconciling Christianity with the Graeco-Roman philosophy. In general, they mingled Stoicism with the teachings of Jesus. Later, other Hellenistic philosophers unified Christian doctrine with the principles of Pythagoras, Plato, and Aristotle. Perhaps the most extreme of these philosophic positions within Christianity was a combination with Platonism known as \emph{Gnosticism,} which was intended to be a sort of esoteric knowledge and to show the relation of Christianity to other religions and to the universe.

\textbf{Catechetical and Episcopal or Cathedral Schools.} - In this way, during the second and third centuries, all the Christians at Alexandria, which had become the great seat of Hellenistic philosophy, had their theology tinctured with Greek thought. Before long, a sort of theological, or `catechetical' school, was gradually organized at this center, to counteract the heathen schools there and to afford higher instruction for Christian teachers and leaders. This school had no building of its own, and the students met at the teacher's house, but they were able to take advantage 01 the facilities at the University of Alexandria. In addition to a thorough training in the Bible, the pupils were allowed to study all types of Greek philosophy, except Epicureanism, the whole range of sciences, classical Greek literature, grammar, rhetoric, and other higher subjects of the pagan schools, but from a different point of view. Thus the Graeco-Roman and the Christian movements had formed an alliance in education, and in this catechetical school we find an attempted union of the `otherworldly' ideal with the `worldly.'

The best known heads of this school at Alexandria were Clement (150-215) and Origen (185-253). They were among the most noted of the Eastern Fathers in the philosophic interpretation of Christianity, and their work contributed not a little to heretical doctrine. Origen may even have been expelled for heresy. At any rate, he opened a new school of the same sort at Caesarea, where he was kindly received. Other catechetical schools sprang up rapidly at Antioch, Edessa, Nisibis, and elsewhere throughout the East. Later the accession of the followers of Nestorius, whose Hellenized theology had in 431 been proscribed by the Church at the Council of Ephesus, very greatly increased the importance of these cities as intellectual centers. In addition to the translations already there, the Nestorian Christians accumulated a larger range of the original Greek treatises on philosophy, science, and medicine.

But before this, higher training of the Hellenic type came to be regularly used by the bishops in training their clergy, and promotion in the Church began to depend upon having had this education. So higher schools of this sort were gradually instituted in every bishopric at the see city, and became known eventually as `episcopal' or `bishop's' schools, or, from their location at the bishop's church, as `cathedral' schools. These cathedral schools became the most important educational institutions of the Middle Ages. From them were derived all the schools of Western Europe, but the bishop soon became too busy to attend to them himself and was forced to commit them to various officials. Thus they developed in time into at least three types, - the `grammar' school, taught by one of the cathedral canons, known as the \emph{scholasticus;} the `song' or music school, taught by the \emph{cantor} or \emph{precentor;} and the `chorister's' school, which offered a combination of the training in the two other schools. Thus the cathedral schools virtually took the place of the old pagan schools supported by the Roman emperors.

\textbf{Influence of Graeco-Roman Culture upon Christianity.} - However, by the century after the foundation of the catechetical school at Alexandria, the Christians had begun to grow suspicious of Graeco-Roman culture and the `worldly' ideal in education. Even the Eastern or Greek Fathers of the Church appear to have cooled considerably in their attitude toward philosophy, and the Western or Latin Fathers were more pronounced in their opposition. Roman Christians could not forget the immorality of those who had been connected with this culture, nor the abuse and insults that these pagans had heaped upon them. They felt, too, that the one great mission of the Church was ethical, and that Christ's second coming was at hand, and that all philosophy and learning were somewhat impertinent.

Nevertheless, despite this growth of opposition to pagan philosophy, primitive Christianity could not endure in its simplicity after it had been in contact with the advanced intellectual concepts of the Greeks, as modified by the organizing genius of the Romans. Both Greece and Rome left a permanent impress upon Christianity; and, though dead, they yet live in the Christian Church. The influence of Greek philosophy is seen in the formulation of a system of Christian doctrine. This appears in the development of the \emph{Aposiles' Creed} during the second century, in the selection of a canon of sacred writings or \emph{New Testament} during the third century, and still more in the \emph{Nicene Creed} (325), which was not formulated until Christianity had been largely Hellenized. Similarly, the Greek tendency to attribute universal validity to their sacred writings, and the pomp, ceremonies, and mysteries of the Hellenic worship, are more or less apparent in the various ecclesiastical tenets and usages. On the other hand, the Roman concepts of administration appear in the organization of the Church, which seems to have closely paralleled the Roman civil polity. By the third century priests and bishops had largely come to be similarly located, and to correspond in control, to the Roman district and city magistrates respectively. And in 445 the recognition of the supremacy of the Bishop of Rome established a visible head of the entire Church, corresponding to the position of the emperor on the civic side.

\textbf{Rise of the Monastic Schools.} - Thus it has been seen how the two great movements of Graeco-Roman culture and Christian teaching arose independently, in time united and later separated, although, after separation, the Christian doctrines were somewhat affected by their long association with pagan philosophy. Eventually the pagan schools were suppressed by the edict of Justinian in 529 A. D., and the Christian education was left alone in the field. It then found an additional means of expression in the `monastic' schools, in which there was naturally a tendency to revert to an ascetic or `otherworldly' ideal, and to leave intellectual attainments largely out of consideration. But these monastic institutions are to be grouped with mediaevalism and belong more distinctly to the next chapter.

\textbf{Supplementary Reading}

Graves, \emph{Before the Middle Ages} (Macmillan, 1909), chap. XII; Monroe, \emph{Text-book} (Macmillan, 1905), pp.~221-243. For the moral effect of Christianity, see Lecky, W. E. H., \emph{History of European Morals} (Appleton, 1869), vol.~II, pp.~1-100. Other places in the chapter will be illumined by reading Ayer, J. C, Jr., \emph{Catechutnenal Schools} and \emph{Catechetical Schools} (Monroe Cyclopaedia of Education, vol.~I); Dill, D., \emph{Roman Society in the Last Century of the Western Empire} (Macmillan, 1899), especially book V; Hatch, E., \emph{The Influence of Greek Ideas and Usages upon the Christian Church} (Hibbert Lectures, 1888, Williams, London, 1891); Hodgson, G., \emph{Primitive Christian Education} (Clark, Edinburgh, 1906); and Leach, A. F., \emph{Bishop's Schools} and \emph{Cathedral Schools} (Monroe Cyclopaedia of Education, vol.~I).

\hypertarget{the-monastic-education}{%
\chapter{5. The Monastic Education}\label{the-monastic-education}}

\textbf{Outline}

\begin{itemize}
\tightlist
\item
  During the Middle Ages the German hordes absorbed ancient civilization under the authoritative guidance of the Church, and the chief means of leavening the barbarian lump was found in the cathedral and monastic schools.
\item
  Monasteries grew up to counteract the prevailing worldliness. To keep the monks busy, Benedict prescribed the copying of manuscripts, and this literary work rendered schools necessary. In these monastic schools were taught the `seven liberal arts' by catechetical methods.
\item
  Thus monasticism helped preserve learning and education, although it was somewhat hostile to the classics and science.
\end{itemize}

\textbf{The Middle Ages as a Period of Assimilation and Repression.} - The Middle Ages may be regarded as an era of assimilation and of repression. On the one hand, the rude German hordes, who had by the sixth century everywhere taken possession of the decadent ancient world, were enabled during this period to rise gradually to such a plane of intelligence and achievement that they could absorb the Greek, Roman, and Christian civilization, and become its carriers to modern times. On the other hand, that this absorption might take place, it was necessary that the individual should conform to the model set, and it was inevitable that a bondage to authority, convention, and institutions should ensue.

The main power in effecting this subservience on the part of mediaeval society was the Christian Church. For it was but natural during the period of assimilation that the Church, which had become completely organized and unlimited in power, should stand as the chief guide and schoolmaster of the Germanic hosts. By the decree of Justinian in 529 A. D., which closed the pagan schools and marks the beginning of the Middle Ages, Christian education was left without a rival. Hence the cathedral and monastic schools became almost the sole means of leavening the barbarian lump. Contrary to the view commonly accepted, the educational activities of the cathedral institutions were more important and general than those of the monastic schools. But the former have already been somewhat discussed, and so much relating to the course and services of the latter will also apply to them that we may now turn to a detailed description of the monastic schools.

\textbf{The Evolution and Nature of Monasticism.} - To understand these schools, it will be necessary to examine the movement out of which they arose. Monasticism grew up through the corruption in Roman society and the desire of those within the Church for a deeper reUgious life. Christianity was no longer confined to small extra-social groups meeting secretly, but was represented in all walks of society, and mingled with the world. It had become thoroughly secularized, and even the clergy had in many instances yielded to the prevailing worldliness and vice.

Under these circumstances there were Christians who felt that the only hope for salvation rested in fleeing from the world and its temptations and taking refuge in an isolated life of asceticism and devotion. This led eventually to the foundation of monasteries, in which the monks lived apart in separate cells, but met for meals, prayers, communion, and counsel. Monasticism started in Egypt, but soon spread into Syria and Palestine, and then into Greece, Italy, and Gaul. But in the West monasticism gradually adopted more active pursuits and milder discipline, and the monks turned to the cultivation of the soil and the preservation of literature.

\textbf{Benedict's `Rule' and the Multiplication of Manuscripts.} - These monastic activities were especially crystallized and promoted by the Benedictine `rule.' This was a code formulated by St.~Benedict in 529 for his monastery at Monte Cassino in Southwest Italy, and it was generally adopted by the monasteries of Western Europe. In the forty-eighth chapter of the `rule' he commanded that the monks each day engage in manual labor for at least seven hours and in systematic reading for at least two hours. The requirement of daily reading led to the collection and reproduction of manuscripts, and each monastery soon had a \emph{scriptorium,} or `writingroom,' in one end of the building (Fig. 7). Most of the works copied were of a religious nature and were limited in number, but the monks were occasionally occupied with the Latin classics, and they also became the authors of some original literature, which included histories of the Church, the monasteries, and the times, as well as works upon religious topics.

\textbf{Amalgamation of Roman and Irish Christianity.} - This preservation of learning and development of literature was especially apparent in the monasteries of England. It came about through the amalgamation at the Council of Whitby, in 664, of the Roman Church in England, with Irish Christianity, which had preserved an unusually high order of learning after its isolation. An immense enthusiasm for the Church, culture, and literature of Rome resulted from this merging of the rival organizations, and the English monasteries, such as Jarrow and Wearmouth, and cathedral schools, like York, became the great educational centers for Europe.

\textbf{The Organization of the Monastic Schools.} - The literary work of the monasteries soon led to the establishment of regular schools within their walls (Fig. 8). The course in these monastic schools may often have lasted eight or ten years, as boys of ten or even less were sometimes received, and no one could become a regular member of the order before he was eighteen. By the ninth century the schools sometimes also admitted pupils who never expected to enter the order. These latter were called \emph{externi} in distinction to the \emph{oblati,} who were preparing to become monks. Some training was also given women in convents for nuns, such as that established by the sister of Benedict.

\textbf{The `Seven Liberal Arts' as the Curriculum.} - The curriculum of the monastic schools was at first elementary and narrow. It included only reading, in order to study the Bible; writing, to copy the sacred books; and calculation, for the sake of computing Church festivals. But after a while the classical learning was gradually introduced in that dry and condensed form of the `seven liberal arts', which was also used by the cathedral schools. This mediaeval canon of studies was a gradual evolution from Graeco-Roman days. The discrimination of these liberal subjects may be said to have begun with Plato, whose educational scheme included a higher group of studies, consisting of arithmetic, geometry, music, and astronomy; and during the later days of Greece and Rome these `liberal' subjects of Plato were combined with the `practical' studies of the sophists, - grammar, rhetoric, and dialectic. These `seven liberal arts' were definitely fixed during the fifth and sixth centuries A. D., through several treatises by such writers as Martianus, Capella, Boethius, and Cassiodorus; and the grammar, rhetoric, and dialectic eventually became classed as the \emph{trivium} or lower studies, and the arithmetic, geometry, music, and astronomy as the \emph{quadrivium} or higher (Fig. 9). While this curriculum was not a broad one, the scope was much wider than would be supposed. `Grammar' was an introduction to literature, `rhetoric' included some knowledge of law and history, `dialectic' paved the way for metaphysics, `arithmetic' extended beyond mere calculation, `geometry' embraced geography and surveying, `music' covered a broad course in theory, and `astronomy' comprehended some physics and advanced mathematics.

\textbf{The Methods and Texts.} - The general method of teaching in the monastic schools was that of question and answer. As copies of the various books were scarce, the instructor often resorted to dictation, explaining the meaning as he read, and the pupils took the passage down upon tablets and committed it. The reading books preparatory to the study of literature, many of which are still extant, were generally arranged by each teacher, and careful attention was given to the etymological and literary study of the authors to be read. As to texts, the leading works upon grammar were at first the elementary work of Donatus (fourth century) and the more advanced treatise of Priscian (sixth century), but by the thirteenth century there had sprung up a series of simplified grammars, which, for the sake of memorizing, were often written in verse. As rhetoric was no longer much concerned with declamation, Cicero and Quintilian were rarely used as texts, but various mediaeval treatises upon official letters, legal documents, and forms came into use. Dialectic was studied through translations of the \emph{Organon} of Aristotle, Euclid furnished the text on geometry, the works of Boethius were generally used for arithmetic and music, and in astronomy adaptations of the treatises of Aristotle and Ptolemy became the texts.

\textbf{Effect upon Civilization of the Monastic Schools.} - Thus monasticism accomplished not a little for civilization. While the works produced in the monasteries were uncritical and superstitious, they compose most of our historical documents and sources in the Middle Ages. And, although monastic schools were decidedly hostile to classical literature as representing the temptations of the world, and at all times their rigid orthodoxy prevented every possibility of science and the development of individualism, they, together with the cathedral schools, preserved a considerable amount of GrascoRoman culture. Without the cathedral and monastic schools, the Latin and Greek manuscripts and learning could scarcely have survived and have been available at the Renaissance.

\textbf{Supplementary Reading}

Graves, \emph{History of Education during the Middle Ages and the Transition to Modern Times} (Macmillan, 1910), chaps. I-II; Monroe, \emph{Text-book} (Macmillan, 1905), pp.~243-274. For the evolution of the ascetic life, see Lecky, \emph{History of European Morals} (Appleton, 1869), vol.~II, pp.~101-274; for the development of monasticism, Taylor, H. O., \emph{The Classical Heritage of the Middle Ages} (Macmillan, 1913), chap. VII, and Wishart, A. W.j \emph{A Short History of Monks and Monasticism} (Brandt, Trenton, 1902). The contribution of Irish monasticism is shown in Healy, J., \emph{Insula Sanctorum et Doctorum} (Sealy, Dublin, 1897), and Zimmer, H., \emph{The Irish Element in Mediaval Culture} (Putnam, 1891). Succinct articles on \emph{Abbey Schools, Bishop's Schools, Church Schools,} and \emph{Cloister Schools} by Leach, A. F. (Monroe Cyclopaedia of Education, vols. I and II), furnish the most accurate ideas of monastic education as far as it is known. An account of the monastic libraries is given in Clark, J. W., \emph{Libraries in the Mediaeval and Renaissance Monasteries} (Macmillan and Bowes, Cambridge, 1894), and Putnam, G. H., \emph{Books and Their Makers during the Middle Ages} (Putnam, 1896). The best account of \emph{The Seven Liberal Arts} in English is that by Abelson, P. (Columbia University, Teachers College Contributions, No.~n, 1906).

\hypertarget{charlemagnes-revival-of-education}{%
\chapter{6. Charlemagne's Revival of Education}\label{charlemagnes-revival-of-education}}

\textbf{Outline}

\begin{itemize}
\tightlist
\item
  Learning and schools had by the eighth century been sadly disrupted, and, to restore them, Charlemagne invited Alcuin of York to become his adviser in education. Alcuin induced Charlemagne to conduct higher education at the Palace School, and to improve the cathedral, monastic, and parish schools.
\item
  Even after Alcuin retired from the active direction of education, he continued his educational influence, but he became set and narrow. A broader spirit, however, appeared in his pupils, and intellectual stagnation never again prevailed.
\end{itemize}

\textbf{Condition of Education in the Eighth Century.} - In the course of the seventh and eighth centuries mediaeval education met with considerable retrogression. The learning of the sixth century was disappearing, the copying of manuscripts had almost ceased, and the cathedral and monastic schools had been sadly disrupted. The secular clergy, monks, nobility, and others who might have been expected to be trained, at times seem even to have lost the art of writing, although the leading churchmen must generally have maintained their knowledge of ecclesiastical Latin and some acquaintance with the classical authors and various compilations of the seven liberal arts. Just before this time the Franks had succeeded in establishing a supremacy over the other barbarian tribes and had spread their rule through what is now France, Belgium, and Holland, and most of Western Germany. Under a dynasty of vigorous kings, they now drove back the Moslems, conquered the Lombards and Saxons, and subdued the Slavs and Bohemians, and finally Charlemagne (742-814) even planned to reestablish the Western Roman Empire under his sovereignty. This monarch greatly strengthened and centralized his dominions by a number of improvements in external administration, but, even before his recognition as emperor by the pope (800), he had realized that a genuine unity of his people could be brought about only through a much more effective and universal education. He had a keen sense of the unfortunate educational situation, and made every effort to improve-it. To assist him in his endeavors, in 782 he called Alcuin (735-804) from the headship of the famous cathedral school at York (see p.~56) to be his chief adviser in education.

\textbf{Higher Education at the Palace School.} - Through this noted scholar Charlemagne proceeded to revive the cathedral, monastic, and parish schools, and to increase the importance of the `Palace School.' At this latter school the great king, all his family, and many of his relatives and intellectual friends studied under the Saxon educator. Alcuin must, however, have used a more discursive and less \emph{memoriter} method with his adult students than the formal catechetical plan employed in instructing the youth. Among the subjects taught were grammar, including some study of the Latin poets and the writings of the Church Fathers, rhetoric, dialectic, arithmetic, astronomy, and theology, but Alcuin appears to have had but little command of the Greek learning. Charlemagne himself seems to have become proficient in Latin and other languages, but, in spite of strenuous efforts, he began too late in life to train his hand to write.

\textbf{Educational Improvement in the Cathedral, Monastic, and Parish Schools.} - With the cooperation of Alcuin, Charlemagne also did everything in his power to increase facilities and improve standards in the existing types of schools. In 787 he issued an educational `capitulary' or decree to the bishops and abbots, urging diligence in the pursuit of learning and the selection of teachers for this work who are able, willing, and zealous to learn themselves and to teach others." Two years later he wrote a more urgent capitulary to the bishops and abbots, in which he specified the subjects to be taught in the cathedral and monastic schools and the care to be taken in teaching them. Schools seem to have been everywhere established or revived in the various cathedrals, monasteries, and villages, and the instruction in several places became famous. All these schools came to offer at least a complete elementary course, and some added considerable work in higher education. Reading, writing, computation, singing, and the Scriptures were taught first, but, beyond this, instruction in grammar, rhetoric and dialectic was often given, and at the more noted cathedral and monastic schools the \emph{quadrivium} also appeared in the course. The schools in the villages, under the care of the parish priests, taught only the rudiments, the Lord's Prayer, the Creed, and the Psalms. Tuition was free in all schools for those intending to become monks or priests, but for the higher work a small fee was sometimes paid by the laity. It seems to have been generally intended that education should be gratuitous and open to all. A letter of the Bishop of Orleans required it of his clergy; and through a capitulary in 802 Charlemagne strove to make it compulsory.

\textbf{Alcuin's Educational Work at Tours.} - After fourteen years of strenuous service, Alcuin retired from the active headship of the educational system to the abbacy of the monastery at Tours. But even here his educational work did not cease. He soon established a model house of learning and education, whither flocked the most brilliant youths in the empire, and since they rapidly became prominent as teachers and churchmen, his influence upon the schools remained fully as marked as before. He also wrote a number of educational works, mostly on the seven liberal arts, and had a large correspondence about education with kings and the higher clergy. Alcuin, however, was by nature conservative, and with his retirement he became decidedly set and narrow. His fear of dialectic and the more advanced views of certain Irish scholars is almost ludicrous, and his repudiation of the classic poets, even his former favorite, Vergil, is pathetic.

\textbf{Rabanus Maurus, Erigena, and Others Concerned in the Revival.} - Fortunately, Alcuin's pupils, who at his death occupied practically all positions of educational importance, retained his broader spirit. This was true in particular of Rabanus Maurus (776-856), whose leadership caused the monastic school at Fulda to become the great center of learning. Rabanus wrote even more prolifically than Alcuin upon grammar, language, and theol6gy, but was not afraid to emphasize the study of classic literature or the new training in dialectic. He also greatly expanded the mathematical subjects of the curriculum, and tended to ascribe all phenomena to natural laws. Rabanus, in his turn, influenced a large number of pupils, and a further impetus was given to the movement by a cross-fertilization of Irish learning, which was also introduced, especially through the mastership of Joannes Scotus Erigena (810-876) at the Palace School.

Thus during the ninth century and the first half of the tenth there arose, through the initiative of Charlemagne and Alcuin, a marked revival in education, and for several generations the cathedral and monastic schools enthusiastically fostered education and learning. Curricula were expanded, and many famous scholars appeared. While, owing to the weakness of Charlemagne's successors and the attacks of the Northmen, learning gradually faded once more, intellectual stagnation never again prevailed. Through the revival of the great Frankish monarch, classical learning had to some extent been recalled to continental Europe from its insular asylum in the extreme West.

\textbf{Supplementary Reading}

Graves, \emph{During the Transition} (Macmillan, 1910), chap. III; Monroe, \emph{Text-book} (Macmillan, 1905), pp.~274-279. Read also Gaskoin, C. J. C, \emph{Alcuin, His Life and His Work} (Clay, London, 1904J, or West, A. F., \emph{Alcuin and the Rise of Christian Schools} (Scribner, 1892), and Mullinger, J. B., \emph{The Schools of Charles the Great} (Longmans, London, 1877).

\hypertarget{moslem-learning-and-education}{%
\chapter{7. Moslem Learning and Education}\label{moslem-learning-and-education}}

\textbf{Outline}

\begin{itemize}
\tightlist
\item
  Moslemism amalgamated in Syria with Greek philosophy and science, and the Moslem cities there became renowned for their learning.
\item
  The masses of the Moslems were suspicious of the Greek learning, however, and those who had absorbed the Hellenized philosophy were driven from the Orient into Spain, where they founded Moorish colleges.
\item
  The Moslems thus stimulated learning in the Christian schools, and introduced Aristotle once more, but, after bringing learning back, Moslemism itself reverted to its primitive stage.
\end{itemize}

\textbf{The Hellenization of Moslemism.} - One of the most important influences in awakening mediaeval Europe was the revival of learning and education that came through the advent of the Moslems. Mohammed, the founder of Moslemism, had been almost illiterate, and the \emph{Koran,} or sacred book, was a curious jumble of Judaistic, Christian, and other religious elements with which Mohammed had become acquainted during his early travels. As long as this religion was confined to the ignorant and unreflecting tribes of Arabia, it served its purpose without modification. But when it spread into Syria and came in contact with Greek philosophy, in order to appeal to the people there, it had to be interpreted in Hellenistic terms, and during the eighth, ninth, and tenth centuries, through the influence of the Nestorian scholars (see p.~46), the Mohammedans were engaged in rendering into Arabic from the Syriac, or from the original Greek, the works of the great philosophers, mathematicians, and physicians. The Mohammedan cities of Syria soon became renowned for their learning. In them arose such scholars as Avicenna (980-1037), who wrote many treatises on mathematics and philosophy, and a \emph{Canon of Medicine} that remained authoritative for five centuries. Similarly, there grew up a society called the `Brothers of Sincerity,' which in its course of study amalgamated the Moslem theology with Hellenistic philosophy.

\textbf{Hellenized Moslemism in Spain.} - But the masses of the Mohammedans were as suspicious of the Greek learning as the orthodox Christians had been, and toward the end of the eleventh century Hellenized Moslemism was driven from the Orient and found a refuge in Northern Africa and in Spain. Here the advanced Mohammedans became known as `Moors,' and their works were destined to have a pronounced influence upon the Christians. There soon appeared such scholars as Averroes (1126-1198), who became the authoritative commentator on Aristotle for several centuries; and Moorish colleges were founded at Cordova, Granada, Toledo, and elsewhere. In these institutions, while learning was still at a low ebb in the Christian schools, were taught arithmetic, geometry, trigonometry, astronomy, physics, biology, medicine, surgery, jurisprudence, logic, and metaphysics. Arabic notation was also introduced in place of the cumbersome Roman numerals and many inventions and discoveries were made.

\textbf{Effect upon Europe of the Moslem Education.} - These schools and colleges of the Moslems soon had their effect upon Christian education. Through their influence, Raymund, Archbishop of Toledo, by the middle of the twelfth century had the chief Arabic treatises on philosophy translated into Castilian by a learned Jew, and then into Latin by the monks; and Frederick II had scholars render the works of Averroes into Latin. Such translations had, however, passed through several media - Greek, Syriac, Arabic, Castilian, Latin - and could not be at all accurate. But, stimulated by this taste of Greek learning, the Christians sought a more immediate version, and a half century later when the Venetians took the city of Constantinople, the works of Aristotle were recovered in the original and translated directly into Latin. Meanwhile the orthodox Mohammedanism had been coming to the front in Spain and overwhelming the Hellenized form, and it was left to Christian schools to continue the work of the advanced Moorish institutions. Moslemism had returned to its primitive stage, but it had helped bring back learning, especially the works of Aristotle, to Christendom. As the classical learning had been restored from the West during the revival of Charlemagne, it now returned from its refuge in the East through the coming of the Moslems.

\textbf{Supplementary Reading}

Graves, \emph{During the Transition} (Macmillan, 1910), chap. V; Monroe, \emph{Text-book} (Macmillan, 1905), pp.~331-334- For a further account of Saracen education, see Coppee, H., \emph{History of the Conquest of Spain by the Arab-Moors} (Little, Brown, Boston, 1881), especially bk. X; Davidson, T., \emph{The Brothers of Sincerity} (International Journal of Ethics, July, 1898), and Draper, J. W., \emph{History of the Intellectual Development of Europe} (Harper, 1875), vol.~I, chaps. XI and XIII, and vol.~II, chaps. II and IV.

\hypertarget{educational-tendencies-of-scholasticism}{%
\chapter{8. Educational Tendencies of Scholasticism}\label{educational-tendencies-of-scholasticism}}

\textbf{Outline}

\begin{itemize}
\tightlist
\item
  Scholasticism was a peculiar method of philosophic speculation in the later mediaeval period. At first, scholastic philosophers held that faith must precede reason, but eventually reason itself tended to become the means of testing the truth.
\item
  Scholastic education was organized in the monastic and episcopal schools, and consisted in the limited learning of the times, systematized on the basis of Aristotelian deduction. Scholasticism was extreme in its discussions, but it tended to rationalize the Church doctrines.
\end{itemize}

\textbf{The Nature of Scholasticism.} - One of the movements that most tended to awaken the mediaeval mind, especially during the latter part of the Middle Ages, was the development of the Church philosophy known as `scholasticism.' This movement does not indicate any one set of doctrines, but is rather a general designation for the peculiar methods and tendencies of philosophic speculation that became prominent within the Church in the eleventh century, came to their height during the twelfth and thirteenth, and declined rapidly the following century. The name is derived from \emph{doctor scholasticus,} which was the title given during the mediaeval period to the authorized teachers in a monastic or episcopal school, for it was among these `schoolmen' that the movement started and developed. Its most striking characteristics are the narrowness of its field and the thoroughness with which it was worked.

\textbf{The History of Scholastic Development.} - The history of scholasticism belongs properly to the field of philosophy, but its influence in bringing on the Renaissance and its effect upon education make a brief consideration of its development necessary here. It began as an effort to vanquish heresy in the interest of the Church dogmas, which until late in the Middle Ages it had not generally been necessary to explain. Even then it was assumed that the Church was in possession of all final truth, which had come to it by Divine revelation, and was in harmony with reason, when fully understood. It was, therefore, the aim of the earlier schoolmen to show how these doctrines were consistent with each other and in accordance with reason. At first, as with Anselm (1033-1109), it was held that faith must precede reason, and where reason was incapable of penetrating the mysteries of revealed doctrine, it must desist from its efforts. But the conviction gradually gained ground that human reason is reliable and that truth can be reached only through investigation. Abelard (1079-1142) declared that the only justification of a doctrine is its reasonableness, that reason must precede faith, and that it is not sinful to doubt.

A new epoch for scholasticism dawned in the twelfth and thirteenth centuries through contact with the Greek philosophy of the Moors in Spain and the subsequent recovery of some original treatises of Aristotle (see p.~67). For a time the Church endeavored to suppress the great philosopher, but, failing to do so, soon utilized his works for its own defense, and even made reason identical with Aristotle, whose authority was not to be disputed. A group of most prominent schoolmen arose, and, as a result of the discussions of Aquinas (1225-1274), Duns Scotus (1274-1308), and William of Occam (1280-1347), it came to be held that truth is established by the \emph{fiat} of God, and that ecclesiastical dogmas are, consequently, not matters of reason, but purely of faith. As a result of this breach between revelation and reason, there arose two types of truth, and a tendency to choose that type which was supported by reason.

\textbf{Scholastic Education.} - The schoolmen were thus throughout attempting to rationalize the teachings of the Church, and to present them in scientific form. As an education, scholasticism aimed also at furnishing a training in dialectic and intellectual discipline that should make the student both keen and learned in the knowledge of the times. The scholastic course of study, which was given at first in the monastic and episcopal schools and later in the universities, consisted in the beliefs of the Church and the limited learning of the times arranged in a systematized form largely on the deductive basis of the Aristotelian logic. This knowledge could all be grouped under the head of philosophical theology. The best illustration of the formal and dogmatic way in which these doctrines were usually presented can be found in the \emph{Sententia} of Peter the Lombard (1100-1160) and the \emph{Summa Theologim} of Aquinas (1225-1274), which were the standard texts of the day upon theology. The work of Aquinas has four main parts, under each of which is grouped a number of problems. Every problem is concerned with some fundamental doctrine, and is further divided into several subtopics. After the problem has been stated, first the arguments and authorities for the various solutions other than the orthodox one are given and refuted in regular order, then the proper solution with its arguments is set forth, and finally, the different objections to it are answered in a similarly systematic way. Peter the Lombard's work has a like arrangement.

\textbf{Its Value and Influence.} - As a whole, the work of scholastic education has been underestimated. It has been urged that it ruined all spiritual realities by its extreme systemization of religion, that it dealt with mere abstractions, and that it indulged in over-subtle distinctions and verbal quibbles. But the scholastic arguments were not as purposeless or absurd as they seem. For example, the celebrated inquiry of Aquinas as to the number of angels that could stand on the point of a needle is simply an attempt to present the nature of the Infinite in concrete form. It is the characteristic of reasoning beings to analyze, compare, abstract, and classify, and while scholasticism may have carried its abstractions, hairsplittings, and scientific terminology to an extreme, it performed a great service for knowledge. It found a confused mass of traditional and irrational doctrines and practices, made them systematic, rational, and scientific, and greatly assisted accuracy in thinking. The discussions of the schoolmen resulted in liberating philosophy from theology, and, without intending it perhaps, scholastic education aided the cause of human reason against dogmatism and absolute authority. It greatly stimulated intellectual interests, produced the most acute and subtle minds of the age, and helped to prepare the way for the Renaissance.

\textbf{Supplementary Reading}

Graves, \emph{During the Transition} (Macmillan, 1910), chap. VI; Monroe, \emph{Text-book} (Macmillan, 1905), pp.~292-313. For a good account of all \emph{The Great Schoolmen of the Middle Ages} (Hodder, London, 1881), see the work of Townsend, W. J.; for the beginnings of scholasticism, Mullinger, J. B., \emph{The University of Cambridge} (Longmans, Green, 1888), vol.~I, pp.~47-64; for the life and influence of Abelard, Compayr6, G., \emph{Abelard} (Scribner, 1893), chap. I; McCabe, J., \emph{Abelard} (Putnam, 1901); and Rashdall, H., \emph{The Universities of Europe in the Middle Ages} (Oxford, Clarendon Press, 1895), vol.~I, chap. 11.

\hypertarget{the-mediaeval-universities}{%
\chapter{9. The Mediaeval Universities}\label{the-mediaeval-universities}}

\textbf{Outline}

\begin{itemize}
\tightlist
\item
  Universities began to spring up toward the close of the Middle Ages. Through local conditions, a course in medicine arose at Salerno; in civil and canon law at Bologna; and in theology at Paris. Bologna became the pattern for numerous universities in the South; and Paris for many in the North.
\item
  Popes and sovereigns granted privileges by charter to the various universities. The term `university' originally signified a `corporation' of students and teachers, and the students were usually grouped according to `nations.' The teaching body was divided into four or five `faculties.'
\item
  The course in arts included the seven liberal arts and portions of Aristotle; in civil and canon law, the \emph{Corpus Juris Civilis} of Justinian and the \emph{Decree} of Gratian respectively; in medicine, the treatises of Greek and other medical writers; and in theology, mostly the \emph{Sententice} of Peter the Lombard. The texts were read and explained by the lecturers, and a practical training in debate was furnished.
\item
  While the courses and methods were narrow and formal, the mediaeval university contained the germ of modern inquiry and did much to foster independence of thought and action.
\end{itemize}

\textbf{The Rise of Universities.} - A most important effect upon subsequent education came through the foundation of the mediaeval universities. These institutions grew out of the old cathedral and monastic schools, but found their models largely in the liberal and professional courses of the Moorish colleges. In general, they came into existence through the many broadening influences of the later Middle Ages. Their rise was intimately connected with the stimulus of the Moslem presentation of Greek philosophy and science, with the interest in dialectic and theological discussions, which led to the development of scholasticism, with the reaction from `otherworldliness' resulting from the ideals of chiv- alry, and with the growth of cities and wealth, and the consequent emphasis upon secular interests and knowl- edge (see chap. xi). However, while they were all more or less the product of the same factors, no two sprang from exactly the same set of causes, and special conditions played a part in the evolution of each university.

\textbf{The Foundation of Universities at Salerno, Bologna, and Paris.} - The oldest of these institutions, that at Salerno, near Naples, was simply a school of medicine, and originated through the survival of the old Greek medical works in Southwestern Italy, and through the attraction of the mineral springs and salubrity of this particular place. By the middle of the eleventh century Salerno was well known as the leading place for medical study. It was, however, never chartered as a regular university, although in 1231 Frederick II recognized it as the school of medicine for the university he had created at Naples some seven years earlier.

On the other hand, Northern Italy became known as a center for the study of Roman law. The cities here, in order to defend their independence, were led to study this subject, and endeavored to find some special charter, grant, or edict from the old Roman emperors upon which to base their claims. Several northern centers were renowned for their investigation of the Roman civil law, but early in the twelfth century Bologna became preeminent through the lectures of Irnerius. By him the entire \emph{Corpus Juris Chilis,} a compilation of Roman law made by eminent jurists in the sixth century at the command of the emperor Justinian, was collected and critically discussed. Influenced by this example, a monk of Bologna, named Gratian, undertook to codify all edicts and formulations of popes and councils in a convenient text-book. The \emph{Decree} of Gratian, which resulted, was almost immediately recognized as the authority upon the subject, and canon law came to be studied here with the same thoroughness as civil law. The university at Bologna was regularly chartered by Frederick Barbarossa in 1158, probably as a recognition of the services of its masters in support of his imperial claims, and faculties of arts, medicine, and theology were established at various times. It was thus the first real university, and its reputation soon became widespread.

Next in order of foundation came the university at Paris, which was by far the most famous of all. The special interest here, as in this part of Europe generally, was dialectic and scholasticism. The university grew out of the cathedral school at Notre Dame, which had acquired considerable reputation under the headship of William of Champeaux, Abelard, and Peter the Lombard, but it was not until 1200, after canon law and medicine had been added to the liberal arts and theology, that it received complete recognition by the charter of Philip Augustus.

\textbf{Bologna and Paris as the Models for Other Universities.} - Salerno, as we have seen, was not a real university, and it did not reproduce its type; but Bologna, and even more Paris, became the mother of universities, for many other institutions were organized after their general plans. At Bologna the students, who were usually ma- ture men, had entire charge of the government of the university. They selected the masters and determined the fees, length of term, and time of beginning. But in Paris, where the students were younger, the government was in the hands of the masters. Consequently, new foundations in the North, where Paris was the type, usually became `master-universities,' while those of the South were `student-universities.' During the thirteenth and fourteenth centuries it became fashionable for universities' the authorities, civil and ecclesiastical, to charter existing organizations or to found new institutions on one of these two plans, and by the time the Renaissance was well started about eighty universities had been established in Europe. Not all of these foundations were permanent, however, for some thirty have, in the course of time, become extinct, and those which remain are much changed in character.

\textbf{Privileges Granted to the Universities.} - From the time of the earliest official recognition of the universities, a large variety of exemptions, immunities, and other special privileges were conferred upon the organizations or upon their masters and students, by the charters of popes, emperors, kings, and municipalities. The students of the universities were in many instances taken under the immediate protection of the sovereign, and were allowed to be tried in special courts of their own, independent of civil jurisdiction, and to possess complete autonomy in all their internal affairs. Generally masters, students, and their retainers alike were relieved from all taxation and from military service. Likewise, universities were granted the right to license masters to lecture anywhere without further examination \emph{(jus ubique docendi),} and the privilege of `striking' \emph{(cessatio),} when university rights were infringed. If no redress were given in the latter case, the suspension of lectures was followed by emigration of the university to another town. This could easily be done, since none of the mediaeval universities had buildings of their own, and there was no need of expensive libraries, laboratories, and other equipment.

Through such special rights the universities obtained great power and became very independent. Soon the liberty allowed to students degenerated into recklessness and license, and they became dissipated and quarrelsome. This is especially seen in the life of the so-called `wandering students,' who migrated from university to university, begging their way, and were shiftless, rollicking, and vicious. The one compensating feature of such degeneracy was their production of jovial Latin and German songs to voice their appreciation of forbidden pleasures and their protest against restraint.

\textbf{Organization of the Universities.} - The term \emph{uttiversitas,} or `university,' did not imply originally, as often claimed since, an institution where `everything' is taught, but it was used of any legal corporation, and only in the course of time was it limited to an organization of masters and students. The phrase \emph{studium generate} was also often used of a university, to indicate a school where the students from all parts of civilization were received, and to contrast it with a \emph{studium particulare,} which was confined to pupils of a limited neighborhood. The formation of a university had been preceded by the organization of `nations,' or bodies of students grouped according to the part of Europe from which they came, but these nations soon began to combine for the sake of obtaining greater privileges and power. Every year each nation chose a `councilor,' who was to represent it and guard its interests. On the side of the masters, the university became organized into `faculties,' of which there might be at least four, - arts, law, medicine, and theology; and each faculty came to elect a `dean' as its representative. The deans and the councilors jointly elected the `rector,' or head of the university.

\textbf{Course in the Four Faculties.} - The course of study to be offered by each faculty was largely fixed by papal decree or university legislation during the thirteenth century. The course in arts, which occupied six years, included the texts on the liberal arts mentioned for the monastic schools (see pp.~56 f .)and several of the treatises of Aristotle, as rapidly as they were recovered. In the law course, \emph{Corpus Juris Civilis} was the authorized text for civil law, and the \emph{Decree} of Gratian for canon law. The faculty of medicine utilized the Greek treatises by Medicine. Hippocrates (c. 460-375 B. C.) and Galen (c. 130-200 A. D.), the \emph{Canon} of Avicenna (see p.~66), and the works of certain Jewish and Salernitan physicians. The students of theology put most of their time upon the four books of Peter the Lombard's \emph{Sententia} (Fig. 9), although the \emph{Bible} was studied incidentally.

\textbf{The Methods of Instruction.} - The training of a mediaeval student consisted not only in acquiring the subjects mentioned, but in learning to debate upon them. The acquisition of the subject-matter was accomplished through lectures, which consisted in reading and explaining the text-book under consideration (Fig. 10). Beside the text itself, the teacher would read all the explanatory notes, summaries, cross-references, and objections to the author's statements, which often quite overshadowed the original, and might even add a commentary of his own. The passage was read slowly and repeated whenever necessary. The whole exercise was carried on in Latin, which had to be learned by the student before coming to the university. The training in debate was furnished by means of formal disputations, in which one student, or group of students, was pitted against another (Fig. 11). In these contests, which also were conducted in Latin, not only were authorities cited, but the debaters might add arguments of their own. Thus, compared with the memorizing of lectures, debating afforded some acuteness and vigor of intellect, but by the close of the fifteenth century it had become no longer reputable. The aim came to be to win and to secure applause without regard to truth or consistency.

\textbf{Examinations and Degrees.} - At the close of the course, the student was examined in his ability to define and dispute; and if he passed, he was admitted to the grade of master, doctor, or professor. These degrees seem originally to have been about on a par with each other, and signified that the candidate was now ready to practice the craft of teaching. The baccalaureate was at first not a real degree, but simply permission to become a candidate for the license to teach. During the thirteenth century, however, it came to be sought as an honor by many not intending to teach, and eventually became a separate degree.

\textbf{The Value and Influence of the University Training.} - Obviously the mediaeval universities had most of the defects of their times. From a modern point of view, the content of their course of study was meager, fixed, and formal, and the methods of teaching were stereotyped and authoritative. They largely neglected the real literature of the classical age, and permitted but little that savored of investigation or thinking. Yet the universities were a product of the growing tendencies that later burst the fetters of mediaevalism. They were a great encouragement to subtlety, industry, and thoroughness, and their efforts toward philosophic speculation contained the germs of the modern spirit of inquiry and but somewhat rationality. They were even of immediate assistance in inquiry and promoting freedom of discussion and advancing democracy, and to their arbitration were often referred disputes between the civil and ecclesiastical powers. Thus they aided greatly in advancing the cause of individualism and carrying forward the torch of civilization and progress.

\textbf{Supplementary Reading}

Graves, \emph{During the Transition} (Macmillan, 1910), chap. IX; Monroe, \emph{Text-book} (Macmillan, 1905), pp.~313-327. Standard works on the universities in general are Laurie, S. S., \emph{The Rise and Early Constitution of Universities} (Appleton, 1886), and the more complete and accurate \emph{Universities of Europe in the Middle Ages} (Oxford, Clarendon Press, 1895), by Rashdall, H. For a brief source account of the privileges, courses, methods, and student life of universities, see Norton, A. O., \emph{Readings in the History of Education; Mediceval Universities} (Harvard University, 1909), or Munro, D. C, \emph{The Mediceval Student} (Longmans, Green, 1899). For the history of individual universities, see Compayre, G.,\protect\hypertarget{ch9.xmlux5cux23para.125.1.0.box.327.250.1048.219.q.50}{}{ \emph{Abelard and the Origin and Early History of Universities} (Scribner, 1893); Lyte, H. C. M., \emph{A History of the University of Oxford} (Macmillan, 1886); Mullinger, J. B., \emph{University of Cambridge} (Longmans, London, 1888); and Paulsen, F., \emph{The German Universities} (Macmillan, 1895; Scribner, 1906).}

\hypertarget{the-education-of-chivalry}{%
\chapter{10. The Education of Chivalry}\label{the-education-of-chivalry}}

\textbf{Outline}

\begin{itemize}
\tightlist
\item
  Owing to the weakness of the regular sovereignty after Charlemagne's day, the feudal system sprang up, and by the middle of the twelfth century it had developed a code of manners known as chivalry.
\item
  Out of this there arose a training for knighthood in religion, honor, and gallantry. Before becoming a knight, the boy was early trained at home, then at some castle, first as `page,' and later as `squire.'
\item
  This chivalric education produced many contradictory results, but it tended to refine the times and to counteract `otherworldliness.'
\end{itemize}

\textbf{The Development of Feudalism.} - The mediaeval education thus far described has had to do mostly with the schooling of the ecclesiastical and other select professional classes. Quite a different type of training was that given the knight. This has generally been known as the education of chivalry. Chivalry is a name for the code of manners in usage during the days of the feudal system. By this system is meant an order of society and government that gradually grew up in the Middle Ages alongside the regular political organization, and when, under the successors of Charlemagne, the monarchy became weak, tended to be substituted for it. Under Dependence feudalism small landowners and freemen lacking land had come to depend upon some powerful neighbor for protection, and even to seek from him a dependent tenure of land. Then, in time, the lords acquired a species of sovereignty over their tenants, and by the tenth century there had come to be a great social gulf between the nobility, who owned the land and lived in castles, and the peasantry, who tilled the soil and supported them. The only serious business of the former was fighting with spear, sword, or battle-axe, in their own quarrels or those of their feudal superiors. To prepare for this warfare, mock combats may occasionally have been engaged in as early as the tenth century (Fig. 12).

\textbf{The Ideals of Chivalry.} - But by the middle of the twelfth century, when the old heroic age had lapsed into an age of courtesy, with extravagant devotion to women and romantic adventure as its chief ideals, these encounters were organized into a definite species of pastime called `tournaments,' and soon degenerated into mere pageantry. Hence the rules of chivalry became fixed and formal, and the art of horsemanship and the management of the lance and spear were developed and settled. The ideals of knightly conduct and education could then be stated as `service and obedience' to God, as represented by the organized church, to one's lord, or feudal superior, and to one's lady, whose favor the knight wore in battle or tournament. The three ruling motives of chivalric education were, therefore, held to be `religion, honor, and gallantry.'

\textbf{The Three Preparatory Stages of Education.} - There were three periods in the preparatory training of a knight. First, until the child was seven or eight, he was trained in religion, politeness, and physique at home by his mother. After this he became a `valet' or `page' at the home of a nobleman, who was generally his father's feudal superior. Here he performed personal duties for his lord and lady, and his education was conducted mostly by the latter. He learned the game of chess, acquired the etiquette of love and honor, and was taught to play the harp and pipe and to sing, to read and write, and to compose in verse. Outside the castle, the pages were trained in running, wrestling, boxing, riding, and rudimentary tilting (Fig. 14). In the third stage, at fourteen or fifteen the youth passed to the grade of `squire,' and, while he still attended the lady and carved the meat or handed around the viands for the guests, his chief service was to the knight and his training came through him. He slept near him at night, groomed his horses, kept his armor and weapons in condition, and attended him at the tournament or upon the battlefield. Through this service the squire himself was practiced in all the warlike arts. Toward the close of the period the embryo knight also chose his lady-love, and learned to write verses and dance. When the squire became twenty-one, he was knighted with many religious ceremonies. After a season of fasting, the candidate entered the church in full armor and spent a night in vigil and holy meditation. In the morning he confessed, had his sword blessed upon the altar by the priest, and took an oath to defend the church, protect women, and succor the poor. He then knelt before his lord, who laid his own sword upon the candidate and dubbed him knight.

\textbf{The Effects of Chivalric Education.} - Such was the training of the knight in the `rudiments of love, war, and religion.' It contained many apparent anomalies and contradictions, and every virtue seems to have been balanced by a correlative vice. The knights were recklessly courageous in battle, but their anger was ungovernable and their cruelty extreme. A great self-respect was supposed to characterize the true knight, but this often reacted into an overweening pride. Likewise, while the knights were rated largely according to their liberality and hospitality, these virtues degenerated into a great love of display and extravagance beyond measure. Again, although great respect for womanhood was inculcated, not much consideration could be expected by the woman beneath a certain rank. Similarly, the knightly word of honor, if accompanied by certain forms, would be held sacred, but should these forms be omitted, a decided breach of faith was not uncommon. As a whole, however, the chivalric training had a beneficial effect upon the society of the times. It helped to organize the turmoil and to refine the barbarism of mediaeval Europe, and was an effective instrument in raising the position of women. Moreover, while this peculiar training was artificial and worldly, by that very tendency it did much to counteract the `otherworldly' ideal of monasticism and the general asceticism of the period. It encouraged an activity in earthly affairs and a frank enjoyment of this life, and thus helped to develop a striking characteristic of the Renaissance.

\textbf{Supplementary Reading}

Graves, \emph{During the Transition} (Macmillan, 1910), chap. VII; Monroe, \emph{Text-book} (Macmillan, 1905), pp.~284-291. Detailed descriptions of the stages of chivalric training can be found in Cornish, F. W., \emph{Chivalry} (Sonnenschein, London, 1901) (Macmillan, 1o08); Furnival, F. J., \emph{Early Education in England (Forewords to The Babees Book,} Early English Text Society, Original Series, vol.~32); and Mills, C, \emph{The History of Chivalry} (Lea and Blanchard, Philadelphia, 1844), vol.~I, chaps. I-V, and vol.~II, chap. VII. An ingenious, but uncritical reconstruction of the life of a knight in story form, is found in Gautier, L., \emph{Chivalry,} chaps. V-XX.

\hypertarget{the-burgher-gild-and-chantry-schools}{%
\chapter{11. The Burgher, Gild, and Chantry Schools}\label{the-burgher-gild-and-chantry-schools}}

\textbf{Outline}

\begin{itemize}
\tightlist
\item
  In the later Middle Ages the commerce of Europe was greatly increased. Soon the towns received a large impulse from serfs that flocked into them, and before long an influential `burgher class' arose.
\item
  There also sprang up merchant and craft gilds, which afforded an industrial training through apprenticeship, and a more formal education through `gild schools.' As the gilds merged with the town, these institutions became `burgher schools,' and afforded a practical education in reading, writing, and reckoning. Various `adventure,' `chantry,' and other schools were also absorbed by the burgher schools.
\item
  Thus these institutions came to represent the educational interests of the industrial classes, and paved the way for the civic control of education.
\end{itemize}

\textbf{The Rise of Commerce and Industry.} - A most important influence in producing a transition from the mediaeval to modern times is found in the increase of commerce during the later Middle Ages. From the Roman days down, trade had never died out in Western Europe, especially Italy, despite the injuries wrought by barbarian invasions, as the nobles had always need of luxuries, and the Church of articles of utility in its services. But the demand for vessels and transports during the Crusades, and the desire for the precious stones, silks, perfumes, drugs, spices, and porcelain from the Orient afterward, gave a tremendous impulse to commercial and industrial activity. The people of Europe began to think of what articles others outside their own little groups might want in exchange for these luxuries, and to strive to produce such commodities. They also undertook themselves to make some of the new articles, such as light and gauzy cotton and linen fabrics, silks, velvets, and tapestries. Thus the means of communication between the European states was greatly facilitated, new commercial routes and new regions were opened, geographical knowledge was increased, navigation was developed, maritime and mercantile affairs were organized, manufactures and industries were enlarged, currency was increased, and forms of credit were improved. All this tended toward a larger intellectual view and a partial dissipation of provincialism and intolerance.

\textbf{Development of Cities and the Burgher Class.} - The most noteworthy consequence of this industrial and commercial awakening was the growth of towns and cities. There was little town life in Western Europe during the Middle Ages before the twelfth century, as the old Roman towns had, through the invasions of the Germans, largely disintegrated, and but few new organ- izations had sprung up in their place. While some towns still existed in Italy and Southern France, most of the people of Europe lived in the country upon feudal es- tates. These little communities were largely isolated and independent of the rest of the world. They pro- duced among themselves all that their members needed, and little or no money was necessary for their crude forms of exchange. Their life was unbroken in its monotony, there was little opportunity for them to better their condition, and their industries were carried on in a perfunctory and wasteful fashion. But with the growth of commerce and population, these serfs began to find it more profitable to work in the towns and compensate the lord of the manor with money rather than work, and the lords, in turn, found it of advantage to accept money in lieu of services, especially as many of them had been impoverished by the Crusades. Great bodies of serfs flocked to the towns, and new centers sprang up around the manorial estates and monasteries as manufactures, trades, and commerce increased.

Feudalism thus began to be threatened as early as the twelfth century, and within a hundred years the extinction of serfdom was assured. The people soon rebelled against the rule of their lords and either expelled them altogether or secured from them for a monetary consideration a charter conferring more liberal rights and privileges. By these charters, the lord agreed to recognize the gild of merchants, and to permit the people to govern themselves. As industries, trade, and commerce continued to develop, the craftsmen and merchants grew rapidly in wealth and importance. They were soon enabled to rival the clergy in education, and the nobility in the luxury of their dwellings and living. They began to read, and books were written or adapted for their needs. The `burgher class' came to have a recognized position by the side of the clergy and nobility; and the king, in order to retain their support, was forced to take counsel with them. This development of industry and commerce, growth of town and city life, and rise of a `third estate' is one of the most noteworthy changes of the late Middle Ages.

\textbf{The Gilds and Industrial Education.} - Such a new social attitude naturally gave rise to new forms of education. An informal type of training soon sprang up in connection with the development of `gilds.' Besides the original gild of merchants, through which the town had presented a united front and gained its privileges, separate gilds for the various crafts had been established in each town. These craft gilds were the sole repositories of the traditional lore of the vocations, and became the chief channel for transmitting it. While their number and variety differed in each town, all the gilds sought to prevent anyone who had not been regularly approved and admitted to the corporation from practicing the trade he represented. In consequence of this attempt at regulation, industrial training in the craft of each gild grew up through an apprenticeship system. This was provided upon a domestic basis. The `apprentice' entered the household of his `master,' and learned the craft under his direction (Fig. 15). The time necessary for this varied greatly in different crafts. For example, in Paris it took two years to learn to become a cook, eight years an embroiderer, and ten years a goldsmith. While the apprentice received no wages during this period, he was under the protection of the gild, and might appeal to the organization against ill-treatment or defective training. At the end of his apprenticeship, he became a `journey-man' and could earn wages, but only by working for a master, and not through direct service for the public. After an examination by the gild, which might include the presentation of a `masterpiece,' or sample of his work, the journeyman eventually became a master. In other ways, the organization regulated and protected its craft. In order that journeymen and masters might not become too numerous, all masters, save those on the governing board of the gild, were forbidden to take more than one apprentice. The methods of practicing each trade and the hours to be devoted to it each day were specified, and the handiwork of each man carefully scrutinized. In many instances, the gild put its own stamp upon good work, and might often seize products that it considered defective.

\textbf{Gild Schools.} - In this way there grew up a species of industrial education, with three definite stages in its organization and with inspection at every point, more formal Before long, too, the gilds developed a more formal means of education. The existing ecclesiastical schools did not altogether meet the needs of the gilds, and they undertook the establishment of additional institutions for this purpose. Where the gilds had retained one or more priests to perform the necessary religious offices for their members, before long they also utilized these functionaries to keep a school for the benefit of their own and sometimes other children in the town. Later, endowments were furnished especially for a priest to teach school, or an amount sufficient for the purpose was paid out of the common funds of the gild. Some of these gild schools, like `Merchant Taylors'' of London, or the Grammar School at Stratford-on-Avon, where Shakespeare was educated (Fig. 16), still survive as secondary institutions. Many instances, too, are recorded where the members of a certain gild were appointed trustees of a school established by an individual, and were granted the right of appointing and dismissing the master, admitting the pupils, managing the property, and formulating statutes. In some such fashion Colet later vested the management of the famous St.~Paul's school (see p.~118) in the gild of mercers.

\textbf{Burgher Schools.} - As the gild organizations gradually merged with those of the towns, the gild schools were generally absorbed in the institutions known as `burgher' or town schools. At first these burgher schools were not very dissimilar to those established by the Church, except that they were more conveniently located, but later various types of vernacular schools arose to meet special practical demands, especially writing and reckoning. The Latin burgher schools were also somewhat practical in their course, and often admitted some pupils who desired to learn only to read, write, and reckon. Writing had become an important vocation, since printing had not yet been invented; and there was a definite demand for writers in public offices, private secretaries, letter writers for the illiterate, and teachers of writing. Reckoning grew directly out of the new commercial life, and was often taught in the writing schools. It was not taught from the standpoint of theory or discipline, as was the arithmetic in the Latin schools, but for the sake of practical calculation and bookkeeping. But even all the facilities of the regular Latin and vernacular schools of the town were not sufficient to meet the demand for a more practical education. In consequence, private `adventure' schools, taught by wandering teachers or by women, likewise often sprang up, and some teachers were even licensed by the town authorities to teach the vernacular. In most instances, however, these institutions were also combined with the burgher schools.

\textbf{Chantry Schools.} - Another type of institution that came into prominence toward the close of the Middle Ages was the `chantry school.' Schools of this sort at first arose out of bequests by wealthy persons to support foundations for priests who should `chant' masses for the repose of their souls. Since these religious duties did not absorb all the time of the priests, they were able to do some teaching. And before long, the founders of chantries themselves came to direct that the priests carrying out their will should be required to teach. Often two chantry priests were provided, one to teach a `grammar' school, and the other a `song' or vernacular school. From the first most of these chantry schools were free of all tuition charges, the priest being requested to ``teach gratis, without asking anything beyond his stipend for his pains,'' but occasionally they were gratuitous only to the children of his parishioners or to poor children whose parents or guardians asked for the privilege.

\textbf{Influence of the New Schools.} - The chantry schools likewise were often united with various other schools within a town, and became jointly known as `burgher schools.' Many new foundations of a similar nature were also made. These burgher schools were largely controlled and supported by the public authorities, although still generally taught by the priests. They came to represent the interests of the mercantile and industrial classes, and gave instruction in subjects of more practical value than had any of the schools hitherto. Such institutions sprang up everywhere during the later Middle Ages. They were often strongly opposed by the ecclesiastical authorities, who struggled hard to abolish them or bring them under control, but they continued to grow and hold their own. The number of lay teachers in them gradually increased, and thus paved the way for the tendency toward the secularization and civic control of education that appeared later on. The new schools, therefore, that arose in connection with the development of commerce and industry and the growth of towns, were one of the largest factors that led into the broadening of outlook known as the Renaissance.

\textbf{Supplementary Reading}

Graves, \emph{During the Transition} (Macmillan, 1910), chap. X; Monroe, \emph{Text-book} (Macmillan, 1905), pp.~337-339. Adams, G. B., \emph{Civilization during the Middle Ages} (Scribner, 1894), furnishes an illuminating chapter (XII) upon the \emph{Growth of Commerce and Its Results.} The development of towns and gilds in various countries of Europe is described in detail by Ashley, W. J., \emph{English Economic History and Theory} (Putnam, 1892), vol.~I, chap. II; Green, Alice S., \emph{Town Life in the Fifteenth Century} (Macmillan, 1894); Gross, C., \emph{The Gild Merchant} (Oxford, Clarendon Press, 1890); Staley, E., \emph{The Guilds of Florence} (Methuen, London, 1906); and Unwin, G., \emph{The Gilds and Companies of London} (Methuen, London, 1908; Scribner, 1909). Accounts of the new types of schools are found in Leach, A. F., \emph{English Schools at the Reformation} (Constable, 1896), chaps. 7-9; Nohle, E., \emph{History of the German School System} (Report of the U. S. Commissioner of Education, 1897-1898, vol.~I), pp.~22-26; and Watson, F., \emph{English Grammar Schools to 1660} (Cambridge University Press, 1909), chap. VII.

\hypertarget{the-humanistic-education}{%
\chapter{12. The Humanistic Education}\label{the-humanistic-education}}

\textbf{Outline}

\begin{itemize}
\tightlist
\item
  By the fourteenth century there appeared an intellectual awakening, known as the \emph{Renaissance.} It was accompanied by a `revival of learning' and an education called `humanistic'
\item
  Italy first showed evidence of the new movement. The characteristics of the Renaissance were embodied in Petrarch and Boccaccio, but little was done with the Greek classics until Chrysoloras came from Constantinople.
\item
  The tyrants of various cities often had humanistic schools started at their courts. Of these the most typical was that under Vittorino da Feltre. These schools eventually forced the universities to admit the humanities to their course. But humanism gradually degenerated into `Ciceronianism.'
\item
  Humanistic education also gradually spread to the countries north of Italy, but it there took on more of a moral color. In France, the protection of Francis I encouraged the introduction of humanism into educational institutions by various scholars. The German universities likewise began to respond to humanistic influences.
\item
  The Hieronymians first introduced the classics into the schools, and Erasmus, who was trained by them, became the leader in humanistic education. Through other humanistic schools started by Sturm and others, the `gymnasium,' the typical classical school of Germany, was evolved, and the humanistic education became fixed and formal.
\item
  In England the movement gradually developed at Oxford and Cambridge, and Colet started St.~Paul's school, which became the model for all secondary schools. Humanism in England, however, soon retrograded into a formalism, and the `grammar' and `public' schools there are little changed to-day.
\item
  The first secondary schools in the American colonies were modeled after the grammar schools of the mother country.
\end{itemize}

\textbf{The Passing of the Middle Ages.} - It can now be seen that a new spirit had crept into European civilization, and that the Middle Ages were passing. We have previously noted (pp.~53 f.) that, in order to bring the German barbarians up to the level of the past, it was necessary for the Church to set an authoritative standard and repress all variation on the part of the individual. Yet such bondage of the human spirit was unnatural, and there were periodic tendencies to rebel against the sysMediajvaiism tem. In fact, mediaevalism contained within itself the germ Sits own germ of its own emancipation. During the eighth century there came about a new political order, which culminated in Charlemagne's revival of education. While conditions were never again as desperate after this stimulus, with the disruption of Charlemagne's empire another decline set in. But by the thirteenth century a new revival, material and intellectual, had also appeared. Several developments gave evidence of the expansion within, and assisted in producing it. The broadening of horizon through contact with the Moors, the development of scholasticism, the evolution of universities, the worldly appeal of chivalry, and the growth of cities, gilds, and commerce were all helping by accumulation to dispel the mediaeval spirit.

And by the fourteenth century a new dawn had been ushered in. The period that followed was marked by a general intellectual and cultural progress that began to free men from their bondage to ecclesiasticism and to induce them to look at the world about them. The adherence to an `otherworldly' ideal, the restriction of learning, the reception of the teachings of the Church without investigation, and the conformity of the individual were by this time rapidly disappearing. Such tendencies were clearly being replaced by a genuine joy in the life of this world, a- broader field of knowledge and thought, a desire to reason and deal with all ideas more critically, and enlarged ideals of individualism. The days of mere absorption and assimilation had passed.

\textbf{The Renaissance and the Revival of Learning.} - This tremendous widening of horizon has been generally known as the \emph{Renaissance} or `new birth.' The term is used to indicate that the spirit of the Graeco-Roman development had returned, and that opportunity for expression was granted to the individual once more. But this period is also appropriately known as the `Revival of Learning.' For, while the awakening preceded and was caused by internal factors, rather than by the recovery of classical literature and learning, intellectual freedom was very greatly heightened and forwarded after a restoration of the classics once began. The only food at hand that could satisfy the awakened intelligence of the times was the literature and culture of the classical peoples. The discovery that the writings of the ancient world were filled with a genuine vitality and virility, and that the old authors had dealt with world problems in a profound and masterly fashion, and with far more vision than had ever been possible for the mediaevalists, gave rise to an eager desire and enthusiasm for the classics that went beyond all bounds. A knowledge of classical literature had never altogether disappeared, and various works had been preserved by the monks and others. To search out the manuscripts of the Latin and Greek writers, the monasteries, cathedrals, and castles were now ransacked from end to end. The manuscripts found were rapidly multiplied, and the greatest pains taken to secure the correct form of every passage. The devotees of the new movement were generally called `humanists,' and the training embodying the classics has since been termed `humanistic education.'

\textbf{Causes of the Awakening in Italy.} - While the general tendency toward an awakening was apparent throughout Western Europe, it first became evident in Italy. This was due to the fact that Italy was at the time a seat of intellectual activity resulting from several factors. It was a storm center for civic and interstate quarrels, and, as a result of this political unrest, the citizens were kept constantly on the outlook for their own safety and interests, and their wits were greatly sharpened. Even the exile, into which one civic faction or another was constantly forced, had the effect of broadening their vision and bringing out the greatest possibilities within them. Again, the commercial intercourse of the Italian cities with other countries had, for various physiographic and historic reasons, become extraordinarily active. This tended to open the minds of the Italians, break up their old conceptions, free them of prejudice, and increase their thirst for learning. Furthermore, the ghost of the classic ages still haunted its old home. A knowledge of the Latin tongue had never ceased to exist in Italy, and many manuscripts of the Latin and Greek authors had been preserved. There was only needed an intellectual awakening sufficient to shake off the thraldom to the Church and produce an appreciation of classical literature and culture, in order to bring back this spirit of the past into real pulsating life.

\textbf{The Revival of the Latin Classics.} - The earliest of the great humanists was Petrarch (1304-1374). In him we find the very embodiment of the Renaissance spirit. He completely repudiated the `otherworldly' ideal of mediaevalism, and was keenly aware of the beauties and joys of this life. He did not hesitate to attack the most hoary of traditions, nor to rely upon observation, investigation, and reason. He likewise felt a kinship with the thinkers and writers of the classic age, when independence and breadth were given more scope, and held that their works must be recovered before their spirit could be continued. This led to a tremendous enthusiasm for the Latin classics, and he spent much of his life in restoring ancient culture. He devoted himself during his extensive travels largely to collecting manuscripts of the old Latin writers, which previously had been widely scattered, and endeavoring to repair in them the ravages of time. And he inspired every one he met with a desire to gather and study the works of the classic authors. He also wrote a number of Latin works that were filled with the classic spirit. Among them were several collections of \emph{Letters,} a work of erudition \emph{On Famous Men,} and an epic poem in honor of Scipio Africanus that he called \emph{Africa.} Some of his letters were indited to Cicero, Homer, and other classical authors as if they were still living. After he had been crowned as poet laureate by the University of Rome in 1341, he spent most of his time visiting various Italian cities and spreading the humanistic spirit. Of the younger scholars and literary men influenced by him probably the most noted was Boccaccio (1313-13 75). Through Petrarch this youthful poet developed a perfect passion for the ancient writers, and devoted the rest of his life to classical culture. He obtained a wide knowledge of the Latin writers, and searched out, preserved, and had copied as many manuscripts as possible.

\textbf{The Development of Greek Scholarship.} - With all this revival of Latin literature by the \emph{coterie} of Petrarch, for some time there was little done with the Greek. That language had almost disappeared in Europe, and the greatest Greek authors were known only through Latin translations. But a knowledge of the Greek language and literature still persisted in the Eastern empire, and the humanists of Italy were, through the works of the Latin authors, constantly directed back to the writings of the Greeks. They became eager to read them in the original, and several humanists began the study of Greek. Nevertheless, Petrarch pathetically confessed: ``Homer is dumb to me, while I am most certainly deaf to him.'' And while, with the aid of his Greek teacher, Boccaccio made a translation of Homer, it showed little real appreciation of the original. Not until Chrysoloras (1350-1415) came as an envoy from the Eastern emperor and was induced in 1396 to settle in Italy and teach Greek, was any systematic training possible. During the next sixteen years this man of learning taught in the leading centers, established schools, made translations of Greek authors, and wrote a Greek grammar. From his efforts sprang a number of famous scholars, such as Vergerio, Niccolo de' Niccoli, Bruni, and Guarino da Verona and his son. These men collected or copied hundreds of volumes, started libraries and schools, made excellent translations, wrote treatises on humanistic education, and trained a number of humanists, who became distinguished later.

\textbf{The Court Schools and Vittorino da Feltre.} - A powerful support for the work of these humanists resulted from the rivalry of the Italian cities. The princes at the head of these centers were often usurpers, and depended largely upon city pride to maintain their power. To appeal to the classical enthusiasm of their people, they did everything possible to propagate the humanistic movement and make their cities illustrious. Probably the most typical examples of these humanistic tyrants are found among the Visconti at Milan and the Medici at Florence. In some instances these court circles promoted the new learning informally, but often, where a scholar had been taken into the family of a prince as private tutor, children of the neighboring aristocracy were associated and a regular school was started. `Court schools' of this sort soon existed at Florence, Venice, Padua, Pavia, Verona, Ferrara, and several other cities, but the best known of all was that organized by Vittorino da Feltre (1378-1446) at Mantua.

\textbf{The Court School at Mantua.} - Vittorino undertook this school at forty-five, when he had received the best possible education of the times in Latin, Greek, and mathematics, and had greatly distinguished himself as a teacher and a man of piety. He received into the school not only the royal princes and the scions of the leading Mantuan families, but, by special permission, the sons of his personal friends and promising boys of every degree. He dwelt with his pupils, and was most strict in his selection of masters and of attendants, that the morals of his pupils might be of the highest. Likewise, `the father of his pupils,' as Vittorino held himself to be, looked out for their food, clothing, and health, and shared in their games, interests, and pleasures. It was his intention to secure for his pupils that harmonious development of mind, body, and morals that the old Greeks had known as a `liberal education,' but he emphasized the practical and social side of the individual's efficiency, and wished to prepare his pupils for a life of activity and service rather than to create mere rhetoricians and pedants.

This he felt could be accomplished largely through a grammatical and literary study of the Greek and Roman writers. The pupils learned from the first to converse in Latin, and there were games with letters for the youngest and simple exercises to train them in clear articulation and proper accent and emphasis. Before they were ten, they were also drilled in memorizing and reciting with intelligence the easier portions of the classic authors. This elocutionary work, which was increased in length and difficulty as the boys grew older, gave them an excellent grasp of vocabulary, rhythm, and style. As they advanced, the pupils read a variety of Latin writers, and soon took up a study of the Greek authors and of the Church Fathers. The mathematical subjects were also taught with an enlarged scope, especially in their applications to drawing, mensuration, and surveying. Because of the lack of books, the teaching was carried on largely by dictation. Vittorino, however, carefully studied the ability, interests, and future career of his pupils, and selected the subjects and methods best suited to each intelligence. He thus inaugurated a thoroughly elastic course for the school. Physical and moral education were likewise insisted upon quite as fully as intellectual. Vittorino introduced especially fencing, wrestling, dancing, ball-playing, running, and leaping, in all of which he was himself an expert, but the purpose of these was to aid and stimulate the mental powers. He also by both precept and example inculcated piety, reverence, and religious observances. He believed, moreover, that truth and moral beauty could be derived not only from the Christian authors, but also, by means of expurgation, from the classic writings.

\textbf{The Relation of the Court Schools to the Universities.} - The court school at Mantua had thus a most potent influence upon the educational practice of the times, and trained a large number of distinguished ecclesiastics, statesmen, scholars, and rulers. It doubtless was broadly typical of the court schools and of the humanistic education of Italy in general. These court schools, while taking pupils very early, often retained them until they were twenty-one, and covered as much, if not more, ground than the arts course of the university. They were, in a way, competitors of the older institutions. A student might, for the sake of a degree, go from a court school to a university, but, as a rule, if what he wished were a general course, he would be satisfied with the greater prestige that came from being a pupil of one of the distinguished humanists that the court schools were generally able to retain at their head. In fact, the want of hospitality, if not actual hostility, of universities to the new learning, often stimulated the growth of court schools. In many instances where the university was especially conservative, a court school was set up by its side as a professed rival. Gradually, however, the humanistic training crept into all the universities of Italy, and the classical literature of the Greeks and Romans largely took the place of the former grammar, rhetoric, and dialectic. Before the close of the fifteenth century, Florence, Padua, Pavia, Milan, Ferrara, Rome, and other cities had admitted the humanities to their universities, and the other university seats were not long in following their example.

\textbf{Decadence of Italian Humanism.} - Toward the close of the fifteenth century, however, this liberal education of the humanists in Italy began to be fixed and formal. Until the middle of the century the ideals, content, and meaning of this training were constantly expanding, but after that time there was a gradual narrowing and hardening, and during the early years of the sixteenth century the degeneration became complete. As the subjectmatter became institutionalized, the literature of the Greeks and Romans failed more and more to be interpreted in terms of life. Emphasis was placed upon the form rather than the content of the classical writings, and grammatical drill was more and more emphasized as a means of formal discipline. Before long the course was limited largely to Cicero, and the new learning fell into that decadent state known as `Ciceronianism.' It consisted simply in an attempt to teach a perfect style with Cicero as a model, and to give one a conversational knowledge of Ciceronian Latin. The structure, metaphors, and vocabulary of all Latin writing had to be copied from the phrases of Cicero, and the literature of the day became little more than a sequence of model passages from that author.

\textbf{The Spread and Character of Humanism in the Northern Countries.} - Such was the effect of the Renaissance upon education in the country of its birth. But the humanistic training could not be confined to Italy. By the middle of the fifteenth century, with the invention of printing, the texts of the classic authors were rapidly multiplied and spread everywhere. The Renaissance and the classic literature leaped the Alps, and made their way into France, the Teutonic countries, England, and elsewhere. At first, humanistic scholars wandered into the North, soon others were invited in large numbers by patrons of learning, and, at length, students from the Northern countries thronged into Italy for instruction. Towards the close of the fifteenth century the humanists outside the peninsula became very numerous, and during the sixteenth century the movement came to its height in the Northern lands.

But the character and effects of the Renaissance and humanism in the North differed greatly from those in the country of their origin. The peoples of the North, especially those of Germanic stock, were by nature more religious than the brilliant and mercurial Italians. With them the Renaissance led less to a desire for personal development, self-realization, and individual achievement, and took on more of a social and moral color. The prime purpose of humanism became the improvement of society, morally and religiously, and the classical revival pointed the way to obtaining a new and more exalted meaning from the Scriptures. Through the revival of Greek, Northern scholars, especially the German and English, sought to get away from the ecclesiastical doctrines and traditions, and turn back to the essence of Christianity by studying the New Testament in the original. This suggested a similar insight into the Old Testament, and an interest in Hebrew was thereby aroused. In consequence, to most people in the North a renewed study of the Bible became as important a feature of humanism as an appreciation of the classics.

\textbf{The Development of Humanism in France.} - In France humanism appeared early. In 1458 a professorship of Greek was established at the University of Paris, but the humanistic movement did not amount to much in France until it was stimulated by the expeditions of Charles VIII (1494) and Louis XII (1498) into Italy. These undertakings of the monarchs did not attain the military and political objects intended, but through them France came into direct contact with humanism at its sources, and a definite impression was made upon French art, literature, and education. Even then, owing to the conservatism of the university, the new learning met at first with formidable opposition. Happily, it found an influential patron in the youthful Francis I (f. 1515-1547).

\textbf{French Humanistic Educators and Institutions.} - Under the protection of Francis, many prominent humanistic scholars and educators, like Budseus (14681540), appeared, classical manuscripts were collected, Greek and Latin authors were translated, treatises on humanistic education were produced, and the College of France, with chairs of Greek, Hebrew, and Latin, was established (1530). Humanism was also introduced into various colleges in Paris and Bordeaux by such scholars and practical teachers as Corderius (1479-1564) and Ramus (1515-1572), and many text-books and editions of the classics were published. Soon most of the schools of France responded to the new training. It would hardly be possible to consider many of them, but a brief description of the course and administration in vogue at the College of Guyenne, taken from an account of one of its teachers, may prove illuminating. This college contained ten classes in secondary work, and two years more in philosophy, which partially overlapped the faculty of arts in the university. Latin and religion were taught throughout the secondary school, and Greek, mathematics, rhetoric, and declamation could be taken in the last three or four classes. The pupils were introduced to the rudiments of Latin through the vernacular, and developmental methods and enlivening disputations were used. Probably the general conditions here were typical of the French humanistic schools everywhere during the sixteenth century.

\textbf{Humanism in the German Universities.} - Before humanism was well established in France, however, it had also spread through the Teutonic countries. By the end of the sixteenth century the German universities had begun to adopt the new learning. In 1494 Erfurt established a professorship of Poetry and Eloquence, which covered the field of classic literature, and lectures on humanistic subjects were before long given in Leipzig, Heidelberg, Tubingen, Ingoldstadt, and Vienna. Likewise, a number of new universities, Wittenberg, Marburg, Konigsberg, and Jena, were started upon a humanistic basis, and before the middle of the sixteenth century humanism prevailed in practically all of the German universities.

\textbf{The Hieronymians and Their Schools.} - The earliest factor in Germanic humanism, however, appeared in the education furnished by the Hieronymians, or Brethren of the Common Lot. For the instruction of the poor, this order had started schools, or established teachers in institutions already existing, throughout the Netherlands, Germany, and France. At first, they stressed instruction in the Bible and the vernacular, but, as the Italian influence began to be felt in the upper countries, they broadened the course by the addition of classic literature and Hebrew, and the schools soon became recognized centers of humanism and intellectual interests. The pupils that were trained there strengthened the new learning as teachers in the universities and schools throughout Germany and the Netherlands. The first educator of importance to introduce humanism into the Hieronymian training seems to have been Wessel (1420-1489). He was preeminently interested in teaching, and among his earliest pupils of distinction were Agricola (1443-1485), who had a most potent influence in introducing classics, and Reuchlin (1455-1522), who taught the classics and Hebrew at various universities, and produced a monumental grammar and lexicon upon the latter subject. An even more noteworthy teacher was Wimpfeling (1450-1528), who became professor, dean, and rector at Heidelberg. He lectured upon the classical authors and the Church Fathers, and wrote a number of treatises upon education, in which he held to the attitude of Northern humanism that all learning is vain which does not lead to the advancement of mankind. But, while a true reformer, he never broke from the Church.

\textbf{Erasmus, Leader in the Humanistic Education of the North.} - A similar attitude was held by Erasmus (1467-1531), the greatest of the humanists trained by the Hieronymians. While he was bitterly opposed to the corruption and obscurantism of ecclesiastics, he believed that the remedy lay, not in a division of the Church, but in the study of the classics and the Church Fathers, and in the general removal of ignorance. Accordingly, to advance education, he assisted in the preparation of Lily's Latin grammar, translated into Latin the Greek grammar of Theodore of Gaza, and wrote a work on Latin composition, called \emph{Be Copia Verborum et Rerum,} and an elementary text-book of Latin conversation on topics of the day, known as \emph{Colloquies.} Similarly, he produced treatises on the New Testament, and popularized the Gospels and Church Fathers through paraphrases. Even better known are the satires that he wrote in Latin to reform the abuses and foibles of his times. His \emph{Adages} and \emph{Praise of Folly} mercilessly scored the absurdities and vices of the Church and the priesthood, and in his \emph{Dialogue on Ciceronianism} he ridiculed some of the narrower tendencies into which humanism had fallen. He also made direct contributions to educational theory in his Latin treatises on \emph{The Liberal Education of Children, The Right Method of Study,} and \emph{Courteous Manners in Boys,} which are almost modern in some of their recommendations. Learning, morality, religion, and good manners, he held, must be trained together, and education must be open to everyone, according to his or her ability. It should be started in infancy by the mothers, and reading, writing, drawing, and some knowledge of familiar animals and objects taught by informal methods. At seven the boy is to be given a thorough training in the Scriptures, Church Fathers, and the classics, and the content rather than the language and form of these works is to be stressed.

\textbf{The Development of Gymnasiums: Melanchthon's Work.} - It can thus be seen what a profound effect the humanists trained in the Hieronymian schools had upon the Teutonic universities and other educational institutions. But there sprang up another set of schools, known as \emph{Gymnasien,} that was an even more typical and lasting institutional development of the Northern Renaissance. These `gymnasiums' grew largely out of the old cathedral and upper burgher schools, and were established for the benefit of the municipality, rather than for State and Church. Their development was gradual, but they were given their first definite shaping by Melanchthon (1497-1560). After a thorough humanistic training from his great-uncle, Reuchlin, and from the universities at Heidelberg and Tubingen, that scholar had become associated with Luther at the University of Wittenberg, and was requested by the Elector of Saxony in 1528 to organize the schools in his state. The `Latin Schools,' which he planned for every town and village of the electorate, were divided into three classes, and the work in Latin and religion was adapted to the grade. Not even Greek or Hebrew appeared in the course; much less the vernacular, mathematics, science, and history. Nevertheless, it was from these municipal Latin schools, when the course had been somewhat modified and expanded, that the `gymnasium' may be said to have sprung.

\textbf{Sturm at Strassburg.} - A further step in fixing the type and the first use of the term `gymnasium' are found in the case of the classical school organized by Johann Sturm (1507-1589) at Strassburg in 1538. Here during his forty-five years as rector, Sturm worked out a gymnasial course of ten classes, upon which the pupils entered at six or seven years of age. The aim of this training he held to be `piety, knowledge, and eloquence,' meaning by the last an ability to speak and write Latin readily. For `piety,' the Lutheran catechism was studied in German for three years, and in Latin for three years longer. The \emph{Sunday Sermons} were read in the fourth and fifth years, and the \emph{Letters} of Jerome also in the fifth year, while the \emph{Epistles} of St.~Paul were carefully studied from the sixth year through the rest of the course. On the `knowledge' and `eloquence' side, Latin grammar was begun immediately and the drill continued for four years, during which the pupil passed gradually from memorizing lists of words used in everyday life and reading di- alogues that embodied them to the translation of Cicero and the easier Latin poets. In the fourth year exercises in style were begun, and this was accompanied by a grammatical and literary study of Cicero, Vergil, Plautus, Terence, Martial, Horace, Sallust, and other authors, together with letter writing, declamation, disputation, and the acting of plays. Greek was begun in the fifth year, and after three years of grammatical training, Demosthenes, the dramatists, Homer, and Thucydides were undertaken.

\textbf{Formalism in the Gymnasiums.} - This training, like that of the Italian humanists, soon became set, formal, and mechanical. While other authors than Cicero were read, the object was to acquire an ability to read, write, and speak Ciceronian Latin, and words, phrases, and expressions were carefully committed. The main emphasis throughout was upon form, with little regard for content, and the Latin and Greek were largely regarded as an end in themselves. Yet the gymnasium of Sturm was an enormous success, and was soon crowded with students. His pupils became the headmasters of all the most prominent schools, and through his wide correspondence with sovereigns and educators, the course of study formulated by Sturm became a model not only for Germany, but, in a sense, for the rest of Europe. At any rate, most of the existing secondary schools in Germany, and many founded later, became gymnasiums. The majority of the Hieronymian schools soon adopted the gymnasial course. This was also the case with the \emph{Furstenschulen,} or `princes' schools, ' a type of institution started in 1543 by Duke Moritz of Saxony to train wellprepared officials for Church and State at public expense, and afterward absorbed into the gymnasial system. And the gymnasiums have to-day changed but little from Sturm's organization. Owing to the later influence of realism, the addition of mathematics, modern languages, and the natural sciences has somewhat mitigated the amount of classics prescribed, but otherwise the German gymnasiums adhere to their formal humanism as tenaciously as in the sixteenth century.

\textbf{The Humanistic Movement in England: Greek at Oxford and Cambridge.} - In its northward march the humanistic education also effected profound changes in England. By the middle of the fifteenth century many former students of Oxford began to study at various humanistic centers in Italy. But the influence of such innovators was scarcely felt until Grocyn and Linacre, who had gone to Florence about 1488, undertook to introduce Greek into education upon their return home. Grocyn (1442-1519) became the first lecturer on Greek at Oxford, but he was greatly assisted in the humanistic training by Linacre (1460-1524), although his lectureship was nominally on medicine. Among their pupils were Erasmus, More, and Colet. Humanistic education did not reach Cambridge, however, until the close of the fifteenth century, but, with the progress of the sixteenth, that university rapidly overtook her sister institution. The real development began when Erasmus, while a professor of theology at Cambridge (1510-1514), consented also to lecture upon Greek as a labor of love. Erasmus was succeeded by a number of lecturers, and in 1540 the new \emph{regius} professorship was held for four years each by the great teachers, Cheke (1514-1557) and Ascham (1515-1568).

\textbf{Humanism at the Court.} - As Cheke became private tutor to Prince Edward and Ascham to Princess Elizabeth, an Hellenic atmosphere was soon promoted in royal circles. A powerful assistance to the development of humanism was also found at. the court through the influence of More, who was especially close to Cardinal Wolsey, and so for a time to the king, Henry VIII. A number of treatises upon humanistic education were written by members of the court, like More and Vives; while Ascham produced his \emph{Scholemaster,} a well-known work on teaching Latin and Greek by `double translation.' This famous method consisted in having the child translate a passage into English, and then, after an hour, render it back into the original and have the master compare it with the text.

\textbf{Colet and His School at St.~Paul's.} - The humanistic changes in English education, however, were not limited to the universities and the court. The schools also felt the effect of the new movement, and the most important factor in bringing this about was the foundation of St.~Paul's School in 1509 by Colet. This scholar devoted most of the fortune left him by his father to establishing a humanistic school in St.~Paul's churchyard, dedicated to `the child Jesus.' The institution was thus an outgrowth of Northern humanism, and combined religious training with a study of the classics. In connection with certain Latin authors and Church Fathers, the pupils studied the catechism in English, the \emph{Latin Grammar} of Lily, who was the first headmaster of the school, and the \emph{De Copia} of Erasmus. St.~Paul's school trained a long list of brilliant scholars, literary men, clergy, and statesmen, and became the immediate model for a host of other institutions. There were in existence at the time St.~Paul's was founded some three hundred `grammar' schools of various types. These had come down from the Middle Ages, and their chief purpose had been the training of young men for the priesthood. Their curriculum was usually of the mediaeval monastic type, but they soon felt the influence of the new school. Those which survived the general dissolution of ecclesiastical foundations by Henry VIII and Edward VI were gradually remodeled on the classical basis of St.~Paul's. New schools were also established in accordance with the humanistic ideals.

\textbf{Humanism in the English `Grammar' Schools.} - But the humanism of the `grammar' schools in England, as in Italy and Germany, soon became narrow and formal. The purpose of humanistic education came to be not so much a real training in literature as a practical command of Latin as a means of communication in all lands and ages. Accordingly, the training became one of dictionaries, grammars, and phrase-books. Expressions and selections were culled from authors and treasured in notebooks, and the methods became largely \emph{memoriter} and passive. The formalism into which the schools of England had thus fallen by the seventeenth century is depicted in Brinsley's \emph{Ludus literarius: or the Grammar Schoole,} a work intended to ridicule and reform these conditions. It indicates that the training in Latin was devoted to drill in inflecting, parsing, and construing a fixed set of texts. Lily's \emph{Grammar} was memorized by the pupils, and references to it were glibly repeated, with little understanding of their meaning. All conversation was based upon some phrase-book, like the \emph{Colloquies} of Corderius, and a Latin theme had to be ground out each week.

\textbf{English `Grammar' and `Public' Schools To-day.} - Although reforms have since been made in many of these directions, the organization and the formal humanism of the English `grammar' school have been preserved in principle even to this day. Mathematics, modern languages, and sciences have been added, and a `modern side' has been established as an alternate for the old course, but the classics are still the emphasized feature, and, to a large degree, the drill methods prevail. But, while it was originally intended that the grammar schools should, by means of the endowment, be open to rich and poor alike, because of the great increase in expenses, necessary and unnecessary, there are now not many opportunities for any one in the lower classes of society to attend a grammar school. Similarly, a distinction has come to be drawn between `grammar' and `public' schools, although it is not a very clear one. In general, a `public school' has a more aristocratic and wealthier patronage. Nine `great public schools' were recognized by the Clarendon Commission in 1864, - Winchester (Fig. 17), Eton, St.~Paul's, Shrewsbury, Westminster, Rugby, Harrow, Merchant Taylors', and Charterhouse; but several other old schools and a number of the stronger foundations of Victoria's reign are generally admitted, and many others claim the dignity of the name that would not be considered eligible outside of the immediate locality.

\textbf{The `Grammar' Schools in the American Colonies.} - It was after these `grammar' schools of the mother country that the first secondary schools in America were modeled and named. In many instances the fathers of the colonies, such as Edward Hopkins, William Penn, and Roger Williams, had been educated in the grammar schools of England, and naturally sought to model the institutions in their new home after them as nearly as the different conditions would permit. The Boston Latin (Grammar) School was founded as early as 1635 (Fig. 23), and other towns of Massachusetts, - Charlestown, Ipswich, Salem, Dorchester, Newbury, Cambridge, and Roxbury, also before long established grammar schools. Similarly, towns of Connecticut, Rhode Island, New York, Pennsylvania, Virginia, and the other colonies, had in many cases founded grammar schools before the close of the century. Moreover, the legislatures of Massachusetts (1647) and Connecticut (1650) soon ordered that a `grammar' school be established in every town having one hundred families. The American grammar schools, like their prototypes, were secondary and sustained no real relation to the elementary schools. They were mostly intended to fit pupils for college, although sometimes the college had not yet been established, and thus to furnish a preliminary step to preparation for the Christian ministry. Hence their course consisted chiefly in reading the classics and the New Testament, and used among its texts Lily's \emph{Grammar} and the \emph{Colloquies} of Corderius. And while the hold of formal humanism upon secondary education was somewhat relaxed during the subsequent stages of the `academy' and the `high school,' the formal classical training was considered the only means of a liberal education until well into the nineteenth century.

\textbf{The Aim and Institutions of Humanistic Education.} - It can now be seen how far the ideals of humanism had departed from those of the mediaeval period. The `otherworldly' aim, the monastic isolation, and the scholastic discussions had given way to the interests of this life, personal and social development, and a study of the classics. In the North the movement took on rather a different color from what it did in the peninsula that gave it birth. While Northern humanism was narrower in not concerning itself so much with self-culture, personal expression, and the various opportunities of life, it had a wider vision through interesting itself in society as a whole and in endeavoring to advance morality and religion. It was democratic and social in its trend, where Italian humanism was more aristocratic and individual.

In Italy the chief educational institutions resulting from the humanistic movement were the schools that arose at the brilliant courts of the city tyrants. These institutions were sometimes connected with the universities, and gradually the universities themselves were forced to admit the new learning to the curriculum. In the North a number of new institutions - Hieronymian schools, princes' schools, gymnasiums, and grammar schools - were developed from humanism, and the existing institutions soon showed the influence of the movement, but all of them stressed moral and religious studies, as well as classical. Everywhere the curriculum of the humanistic foundations consisted mostly in the mastery of Latin and Greek, but in the North the renewal of Greek meant also a study of the New and Old Testaments and the Church Fathers. Where the Italian Renaissance re-created the liberal education of Plato and Aristotle, Cicero and Quintilian, the movement in its Northern spread found in the classical revival a means of moral and religious training. But just as humanism in Italy by the beginning of the sixteenth century had degenerated into mere Ciceronianism, so the humanistic education in the North, after about a century of development, began to grow narrow, hard, and fixed. By the middle of the sixteenth century the spirit of criticism, investigation, and intellectual activity had begun to abate, and by the opening of the seventeenth humanism had been completely formalized. In the study of the classics all emphasis was placed upon grammar, linguistics, and style; form was preferred to content; and methods became \emph{memoriter} and imitative. Humanism had largely performed its mission, and a new awakening was needed to revivify education and society in general.

\textbf{Supplementary Reading}

Graves, \emph{During the Transition} (Macmillan, 1910), chaps. XIIXIV; Monroe, \emph{Text-book} (Macmillan, 1905), chap. VI. An interesting interpretation of the Renaissance both in Italy and the North is found in Adams, G. B., \emph{Civilization during the Middle Ages} (Scribner, 1894), chap. XV. An account of the movement, including its educational aspects in Italy, is found in Burckhardt, J., \emph{Civilization of the Renaissance in Italy} (Sonnenschein, London, 1892; Macmillan), vol.~I, especially part III; Symonds, J. A., \emph{Renaissance in Italy} (Holt, Scribner), vol.~II, especially chaps. Ill-VIII; or Symonds' \emph{Short History of the Renaissance} (Holt, 1894), especially chaps. I and VII, and IX-XI. Woodward, W. H., gives us a vivid account of the educational work of \emph{Vittorino da Fcltre and Other Humanist Educators} (Cambridge University Press, 1897), and of \emph{Erasmus concerning Education} (Cambridge University Press, 1904), and of \emph{Education during the Renaissance} (Cambridge University Press, 1906) as a whole. \emph{Peter Ramus and the Educational Reformation of the Sixteenth Century} (Macmillan, 1912), by Graves, F. P., furnishes some idea of conditions in France. \emph{The Italian Renaissance in England} (Columbia University Press, 1905), especially chap. I, is succinctly described by Einstein, L.; and an account of Colet and St.~Paul's School can be found in Barnard, H., \emph{English Pedagogy,} second series, pp.~49-117.

\hypertarget{educational-influences-of-the-reformation}{%
\chapter{13. Educational Influences of the Reformation}\label{educational-influences-of-the-reformation}}

\textbf{Outline}

\begin{itemize}
\tightlist
\item
  Luther's educational positions are most fully revealed in his well-known \emph{Letter} and \emph{Sermon.} He holds that education should prepare for citizenship, and should be state-supported, and these recommendations were somewhat embodied in actual schools by his associates.
\item
  Zwingli was killed before he could greatly influence education, but the educational institutions of Calvin spread rapidly through Switzerland, France, Netherlands, Puritan England, and Scotland.
\item
  In England Henry VIII and Edward VI confiscated the property of some three hundred monastic and other ecclesiastical schools, but subsequently many of these were refounded.
\item
  The Jesuit colleges were organized to extend Catholic Christianity. The lower colleges were humanistic, and the higher taught `philosophy' and theology. The teachers were trained, and the methods, though \emph{memoriter} and emulative, were effective. The influence of the Jesuit colleges was phenomenal, but they have failed to meet new conditions.
\item
  The Port Royalists held that reason was more important than memory, but, while their `little schools' stressed vernacular, logic, and geometry, they offered nothing beyond the best elements in the education of the past.
\item
  Elementary and industrial education was given an impulse for the Catholics by the schools of the Christian Brothers. They also opened training schools for teachers, and perfected the `simultaneous' method.
\item
  Among the Protestants and some Catholics in Germany, Holland, Scotland, and certain of the American colonies, the Reformation inclined toward universal elementary education and control of the schools by the state. The secondary schools in Protestant countries also came largely under civic authorities, although the clergy still taught and inspected them; while Catholic secondary education was furnished mostly by the Jesuit colleges. In many instances the universities turned Protestant; and new universities, Protestant and Catholic, were founded.
\end{itemize}

\textbf{The Relation of the Reformation to the Renaissance.} - The series of revolts from the Catholic Church, generally known collectively as the `Reformation,' may be regarded as closely connected with the Renaissance. As shown in the last chapter, humanism in the North led to a renewed study of the Scriptures and a reform of ecclesiastical doctrines and abuses, and took on a moral and religious color. Reformers arose, like Wimpfeling and Erasmus, who, while remaining within the Church, sought to purify it of corruption and obscurantism. But the Church at first stubbornly resisted all efforts at internal reform. Its immense wealth, large numbers, and training enabled it for a long time to thwart the spirit of the age, and a condition of ecclesiastical upheaval followed. Revolts against papal authority ensued in various parts of Europe north of Italy, and were furnished support by the awakened intellectual and social conditions of the sixteenth century. The result was the establishment of a church, or rather a set of churches, outside of Catholic Christianity. While each revolt had some peculiarities of its own, there were underlying them all certain general causes that indicated their relation to the Renaissance.

\textbf{The Revolt and Educational Works of Luther.} - Even the attitude of Martin Luther (1483-1546) seems to have been bound up with the tendencies of the day. Apparently he had at first no idea of breaking from the Church, and supposed that the ninety-five theses he nailed to the church door at Wittenberg (1517) were quite consistent with Catholic allegiance. But even before this he had attacked Aristotle and scholasticism with great vigor, appealing to primitive Christianity and the right Of free thought, and thus identified himself in spirit with the Northern Renaissance. And two years later, in his contest with Eck, when he was actually led to deny the authority of both pope and council, he was evidently relying upon the humanistic and individualistic atmosphere of the times.

When once he had revolted, Luther gave much of his time to promoting the reform and education of the masses by writing. All his works, whether religious or pedagogical, were clearly intended, in a broad sense, to be educational. After his condemnation at the Diet of Worms (1521), when he had taken refuge at the Wartburg, he undertook to awaken the minds and hearts of the common people by a translation of the Greek Testament. Contrary to general opinion, a large number of translations had preceded that of Luther, and their popularity must have proved suggestive to him, but his edition was unusually close to the colloquial language of the times. A dozen years later, he had completed a translation of the entire Bible, which contributed greatly to education by getting the masses to read and reflect. For the further instruction of the people, he also followed the fashion of the day in producing two catechisms, one for adults and the other for children, together with many tracts, addresses, and letters, filled with allusions to the organization and methods of education. But the documents which most fully reveal his educational positions are his \emph{Letter to the Mayors and Aldermen of All Cities of Germany in behalf of Christian Schools} (1524), and his \emph{Sermon on the Duty of Sending Children to School} (1530).

\textbf{Luther's Ideas on Education.} - The purpose of education, Luther everywhere holds, involves the promotion of the State's welfare quite as much as that of the Church. The schools were to make good citizens as well as religious men. Educational institutions should, on that account, be maintained a.t public expense for every one, - rich and poor, high and low, boys and girls, alike, and attendance should be compelled by the civic authorities. Realizing that some pupils may find it hard to give the time to school, Luther planned that ``they should spend an hour or two a day in school, and the rest of the time in work at home, learn some trade and do whatever is desired, so that study and work may go on together.'' But he also desired a more academic course ``for the brightest pupils, who give promise of becoming accomplished teachers, preachers, and workers.'' In any case, Luther naturally believed that the chief studies should be the Bible and the catechism. But, as a Northern humanist, he recommended the ancient languages - Latin, Greek, and Hebrew - for the light they would throw on the Scriptures and the patristic writers. He likewise approved of rhetoric and dialectic, which were very valuable subjects in those days of controversy; and he made a decided advance in advocating history, natural science, vocal and instrumental music, and gymnastic exercises. History is advised, not only, as was common with the humanists, for the sake of illustrating moral truth, but also for the purpose of understanding social institutions. The study of nature was intended to reveal ``the wonders of Divine Goodness and the omnipotence of God.'' Gymnastics he considered of value both for the body and the soul, and music a means of ``driving away all care and melancholy from the heart.'' The methods he recommended were equally rational. He would utilize the natural activity of children and not attempt to repress them, and would make use of concrete examples, wherever possible. Languages he would teach less by grammar than by practice. This belief in the importance of selecting the proper content and method in education led him to rate the function of the teacher as higher, if anything, than that of the preacher.

\textbf{The Embodiment of Luther's Ideas in Schools by His Associates.} - These recommendations of Luther were largely embodied in actual institutions by his associates. The year after his \emph{Letter to the Mayors} was published, the Protestants were requested by the Count of Mansfeld to establish in Luther's native town, Eisleben, a school that should put his educational theories into practice, and this was performed by Melanchthon. The subsequent organization of Latin schools throughout the Electorate of Saxony, and the foundation of the gymnasium of Sturm at Strassburg upon the Protestant basis have already been touched upon (pp.~114 ff.). But of fully as much importance were the educational foundations of Bugenhagen (1485-1558). While engaged in reorganizing the churches in the cities and states of Northern Germany, by his general `church orders' to each, he made ample provision for schools of the Lutheran type. For instance, at Hamburg in 1520 he organized a single Latin school with a rector and seven teachers, together with a German school for boys and one for girls in every parish. Eight years afterward, the `church orders' of Brunswick provided two classical schools, two vernacular schools for boys, and four for girls, so located in the city that all children could conveniently reach a school. Within a half dozen years he made similar re- quirements for Liibeck, Minden, Gottingen, Soest, Bremen, Osnabriick, and other cities, and throughout some entire states of Germany, such as Holstein and his own native duchy of Pomerania. The educational other theories of Luther were also put into practice in a num- ber of schools taught by Trotzendorf, Neander, and other pupils of Melanchthon.

\textbf{The Revolt and Educational Ideas of Zwingli.} - The revolt under Zwingli (1484-1531) was more directly the outcome of Northern humanism than was that of Luther. Through Erasmus and others he had come to believe that there was little basis in the Bible for the traditional theology, and he carefully read the accounts himself in the original Greek and Hebrew. After he took charge of the cathedral at Zurich, he began his attack upon the dogmas and traditions of the Church, and, by securing the support of the town, managed in a fairly peaceful way to drop one form of the Church after another, until, within five years, he had abolished even the mass. Zwingli likewise made the extension of educational facilities a part of his reform. He founded a number of humanistic institutions, and introduced elementary schools into Switzerland. He also published a \emph{Brief Treatise on the Christian Education of Youth} (1523), which recommended a course of studies not unlike that of Luther, except that, from his practical temperament, he did not mention history, but did add arithmetic and surveying.

\textbf{Calvin's Revolt and His Encouragement of Education.} - While endeavoring to spread his reforms, Zwingli was slain in the prime of life. His positions were maintained by his successor in the cathedral, but the work was soon overshadowed and merged in the movement of Calvin (1509-1564). Calvin's break with the Church, like that of French Protestants generally, also began through the influence of Northern Humanism and the study of the Greek Testament. He had, however, received an excellent legal and theological education, and did not content himself with merely attacking Catholic doctrine, but was the first Protestant to formulate an elaborate system of theology. The call of Calvin to reorganize the civil and religious administration of the city of Geneva gave him an excellent opportunity for working out his theories. Although he was much engrossed in religious disputes, he established `colleges' at Geneva and elsewhere, and in other ways undertook to found schools and promote education. He succeeded, too, in persuading his former teacher, Corderius (see p.~111), to come to Switzerland, and organize, administer, and teach in the reformed colleges.

\textbf{The Colleges of Calvin.} - Corderius here wrote four books of \emph{Colloquies,} with the purpose of training boys by means of conversation on timely topics to speak Latin with facility, and from this work we can learn much of the character of the Calvinistic colleges. Clearly the ideal was the `learned piety' of Melanchthon, Sturm, and the other Northern humanists and Protestants. An attempt seems to have been made to teach Latin in such a way as to cultivate moral and religious life, and psalms were sung, public prayers offered, and selections from the Bible repeated each day. We also know that in the seven classes of a college at Geneva the pupils learned reading and grammar from the Latin catechism, and then studied Vergil, Cicero, Ovid, Caesar, Livy, and Latin composition. Greek seems to have been begun in the fourth year, and, beside classical Greek authors, the Gospels and Epistles were read. Likewise, as in the other Reformation schools, Logic and rhetoric were studied in the higher classes. The colleges of this type not only spread rapidly among Calvin's co-religionists in Switzerland and France, but, as Geneva became a city of refuge for all the oppressed, a regard for humanistic, religious, and universal education was absorbed by the persecuted Netherlanders, the English Protestants of Mary's time, and the Scotch under the leadership of Knox in the days of Mary, Queen of Scots (1505-1572).

\textbf{Henry VIII's Revolt and Its Effect upon Education.} - In England a revolt from the Church likewise occurred. This also may have been due in part to the investigative spirit of Northern humanism, but the immediate cause of the breach was the desire of Henry VIII (r. 1509-1547) to control the national Church that he might divorce his wife, and there was at first little change in doctrine. Once in ecclesiastical power, Henry began in 1536 to confiscate the monastic lands and property, and enlarged the scope of his operations until he had suppressed a large number of monastic, cathedral, collegiate, hospital, and other schools. During the reign (1547-1553) of his successor, Edward VI, the acts of suppression were extended to chantry and gild foundations, and it is estimated that, of the three hundred grammar schools that had come down in England from the Middle Ages, but few were not destroyed under Henry and Edward. Some, however, remained by the terms of the parliamentary acts of suppression, and popular sentiment caused others to be refounded. And during the reign of Elizabeth (1558-1603) and of the first two Stuart kings (1603-1649) these foundations were greatly increased out of royal funds or through the philanthropy of wealthy men. All of these schools, as we have seen (p.~118), following the example of St.~Paul's, adopted the Northern ideals of humanism and furnished a curriculum of classics and religious training. The latter became based, of course, upon the teachings of the Church of England.

\textbf{Foundation of the Society of Jesus.} - We may now turn back to the Mother Church and see what efforts she was putting forth in behalf of education during the period of Protestant revolts. Both before and after the time of Luther there were reformers inside the Church who wished to improve its practices without changing its administration, but the Cathohcs in genefarfeit it their chief duty to crush the Protestant heresy and recover the ground they had lost. This resulted in a number of religious wars, in which both sides displayed great bitterness and cruelty. But a more effective and constructive instrument in advancing the interests of Catholicism was the organization of the `Society of Jesus.' This order was founded by Ignatius de Loyola (1491-1556) in 1534. He persuaded six fellow-students at Paris to join with him in devoting themselves to the conversion of the heathen, and to strengthening the authority of the pope. Six years later, after considerable opposition, the new order was recognized by the pope and began to add rapidly to its numbers. The Jesuits have always striven first through missionary labors to extend Catholic Christianity throughout the world, and then by means of schools to hold their converts and educate all peoples to papal allegiance.

\textbf{Organization of the Jesuits.} - The organization of the Society of Jesus was outlined in its \emph{Constitution.} This fundamental document of the order received its final revision shortly after Loyola's death, but the \emph{Ratio Studiorum,} which was an expansion of Part IV of the \emph{Constitution} and described the educational administration in detail, was not finally formulated until 1599. It thus summed up the experience of the Jesuit schools during more than sixty years. The admimstration of the society has always been of a military type. Loyola had originally started upon the career of a soldier, and did not believe that any system could be effective unless it were based upon implicit obedience to one's official superiors. At the head of the order is the `general,' who is elected for life and has vast administrative powers. As the society spread, the countries that came under its control were divided into provinces, and at the head of the Jesuit interests in each of these districts is the `provincial,' who is appointed by the general for three years. In each province there are various colleges, whose presiding officer, or `rector,' is chosen for three years by the general, but is directly responsible to the provincial and reports to him. Similarly, within each college are `prefects,' immediately subordinate to the rector, but selected by the provincial; and under the inspection of the prefects are the `professors' or `preceptors.'

\textbf{The Jesuit Colleges.} - The Jesuits have never engaged in elementary education, but have required that pupils know how to read and write before being admitted to any of their schools. This may have been brought about in the first place by the fact that the number of their teachers was limited, or that the public elementary school was just coming to be regarded as of importance, and secondary education of the humanistic type was everywhere dominant. The Jesuit educational organization has, therefore, consisted of `lower colleges' with a gymnasial course, and of `upper colleges,' which are of university grade. Boys are admitted to the lower colleges at from ten to fourteen years of age, and spend five or six years there. The first three classes were at first devoted to a careful study of Latin grammar, and a little of Greek; in the fourth year a number of the Greek and Latin poets and historians were read; while the last class, to which two years were usually given, took up a rhetorical study of the classical authors. Only slight variations in the curriculum have ever been allowed since the \emph{Ratio Studiorum} was issued, until the revision in 1832. In that year work in mathematics, natural science, history and geography was added in the lower colleges, but the classics still compose the body of the course.

The full course of the upper colleges lasts seven or nine years, - the first three in `philosophy,' followed by four or six in theology. The training in `philosophy' now includes not only logic, metaphysics, psychology, ethics, and natural theology, but also work in algebra, geometry, trigonometry, analytics, calculus, and mechanics, and such natural sciences as physics, chemistry, geology, astronomy, and physiology. A successful completion of the course leads to the degree of Master of Arts. After the course in philosophy, most of the Jesuits teach in the lower colleges five or six years before going on with the work in theology. In the theological course four years are devoted to a study of the Scriptures, Hebrew, and other Oriental languages, together with Church history, canon law, and various branches of theology. After this one may elect a further training of two years, to review the work in philosophy and theology, and to prepare a thesis. After a public examination and defense of his thesis, the successful candidate is awarded the degree of Doctor of Divinity. Hence a complete Jesuit training will take from eighteen to twenty years, and a member of the order may be from thirty to thirty five years of age before completing his formal education.

\textbf{The Jesuit Methods of Teaching.} - The methods of teaching and the splendid qualification of the instructors were from the first distinctive features in the Jesuit colleges, especially when one considers how little attention up to their time had been given to the preparation of teachers. No one could teach in the lower colleges who had not passed through the course in philosophy, while professors in the universities had first to complete the theological course. Instruction was generally imparted orally, and then memorized or taken down in lecture notes. The method was the `prelection,' which meant a preliminary explanation of the passage or lectures upon the topic under consideration by the teacher. It consisted in giving, first, the general meaning of the whole passage or proposition; then, a more detailed explanation of the construction or phraseology; next, similar thoughts in other authors; fourthly, `erudition', or informational comment upon the passage; then, a study of the rhetorical figures; and finally, the moral lesson to be drawn. Obviously, with such a method, great stress would be placed upon memorizing, especially in the lower colleges. To fix subjects firmly in mind, short hours, few studies, and brief lessons were early found to be necessary. Likewise, reviews have always been frequent and systematic, and the Latin motto of the Jesuit method declares that ``repetition is the mother of learning.'' Each day begins with a review of the preceding day's work, and closes with a review of the work just accomplished. Each week ends with a repetition of all that has been covered in that time, and the last month of every year reviews the course of the year. To maintain interest in the midst of so much memorizing and reviewing, many devices to promote emulation are used. The pupils are arranged in pairs as `rivals,' whose business it is to check on the conduct and studies of each other (Fig. 18); and public `disputations' between two sides are engaged in each week.

\textbf{Value and Influence of the Jesuit Education.} - The Jesuit system, then, seems to have been in advance of that in the schools at the time of its foundation. It was organized upon a systematic and thorough basis, and was administered by a set of splendidly trained teachers through the best methods that were known in that day. The schools were interesting and pleasant, and were free to all who had the ability and desire to attend. The Jesuit teachers, too, were indefatigable and devoted to their duty. The criticism that has been offered to this educational system is based on its insistence upon absolute authority and the consequent opposition to the development of individuality. The Jesuit courses, subjects, and methods have become somewhat uniform and fixed. In the lower colleges they depend largely upon memory and appeal to interest through a system of rivalry, honors, and rewards. Such a system is likely to tend toward a reproductive attitude in the pupil.

Nevertheless, the Jesuits furnished the most effective education during the latter half of the sixteenth, the entire seventeenth, and the early part of the eighteenth centuries. The growth of their schools was phenomenal. By the death of Loyola (1556) there were already one hundred colleges, and a century and a half later they had increased to seven hundred and sixty-nine institutions, spread throughout the world. The average number of students in attendance at any of these colleges during the seventeenth century was about three hundred, and in several of the larger centers there were between one and two thousand, and the famous College of Clermont (now \emph{Lycie Louis le Grand}) at Paris is said to have run up to three thousand. At a modest estimate, there must have been some two hundred thousand students in the Jesuit colleges when they were at their height. Their graduates seem to have become prominent in every important activity of life, and included a large number of the noted authors, prelates, statesmen, and generals of the time. By the middle of the eighteenth century, however, the ideals and content of education had somewhat changed, and the Jesuits did not adapt their course to the new conditions. Moreover, the Jesuits seem to have become powerful, ambitious, and somewhat arrogant. They quarreled frequently with bishops, other monastic orders, governments, and universities. Finally, after they had been banished from France, Spain, and Portugal, in 1773 the pope himself dissolved the Society of Jesus. Forty years later the order was restored, but, owing to the development of educational ideals and organization and the increase of educational institutions, their work has never since become relatively as effective or held as important a place in education.

\textbf{The Organization of the Education of the Port Royalists.} - A type of Catholic education radically opposed to that of the Jesuits was created by a group of men belonging to the religious body known as the Jansenists. The doctrines of the Jansenists were formulated in 1621 by Cornelius Jansen, a professor in the University of Louvain. While striving to retain their place within the Church, the Jansenists opposed the prevailing doctrines of confession and penance, and adopted the rationalistic philosophy of Descartes. They also held that humanity is naturally corrupt, except as it is watched and guided, and that only a relatively few can be saved. These doctrines probably influenced a body of Jansenists that established a new departure in the way of education at the convent of Port Royal at Chevreuse. In 1643 the `Port Royalists' endeavored to remove what few children they could from the temptations of the world to a school started in this convent. Similar institutions quickly sprang up in the vicinity and then spread through Paris. To carry out their ideal of careful oversight, these schools usually took only twenty to twenty-five pupils, and each master had under him five or six boys whom he never allowed out of his immediate supervision day or night. Hence these institutions were known as `little schools.'

\textbf{The Port Royal Course and Method of Teaching.} - Since the Port Royalists held that character was of more importance than knowledge, and reason was to be developed rather than memory, these `little schools' sought to impart an education that should be sound and lasting, rather than brilliant. Unlike the Jesuits, they did not start their pupils with Latin, but with the vernacular, since this was within their comprehension. As soon as they possessed a feeling for good literature, they began the study of Latin through a minimum grammar written in French, and soon took up the Latin authors, rendering them into the vernacular. Greek literature was treated in similar fashion. To train the reason, the older pupils were also taught logic and geometry. The course of study, however, was mostly literary, and had no regard for science or investigation. Port Royal presented the best elements of the education of the past, but did not see beyond it. The methods introduced some striking innovations. The leaders in the Port Royal education departed from the alphabetic plan in teaching their pupils to read, and developed a phonetic method. The Port Royalists also refused to permit the use of emulation and prizes in their schools, but their exclusion of rivalry resulted in indifference. They were never able to secure the energy, earnestness, and pleasing environment of the Jesuit colleges. They did, however, succeed in inculcating a general spirit of piety without the formal teaching of doctrine.

\textbf{Closing of the Port Royalist Schools and Its Effects.} - In 1661 the Port Royalist schools were closed by the order of Louis XIV through the influence of the Jesuits. But this act cost the Jesuits dearly. Not only did it lose them sympathy, but it furnished the Port Royalists occasion to issue tracts against Jesuitism that have injured its repute ever since. This closing of their schools also gave the Port Royalists the opportunity of becoming educators in a larger sense by producing a great variety of writings upon their system. Later on, too, Rollin (1661-1741), who was twice elected rector of the University of Paris, summarized in his \emph{Treatise on Studies} the Port Royalist reforms wrought in that institution.

\textbf{La Salle and the Schools of the Christian Brothers.} - The Port Royalists were, however, like the Jesuits, engrossed with secondary and higher education, and gave little heed to the education of all the people in the rudiments. In fact, until toward the close of the seventeenth century, the Catholics generally did not succeed in inaugurating any effective or widespread movement toward elementary education. Numerous attempts before this were made through catechism schools and various reformers and religious orders, but teachers were scarce and often ignorant and poorly trained, and there was little progress before the organization of the Brothers of the Christian Schools through the self-sacrificing efforts of Jean Baptiste de la Salle (1651-1719). The organization sprang out of a group of five masters engaged in teaching schools for the poor in the city of Rheims in 1679, but it was not until three years later at Rheims, that La Salle completed his regulations, founded the brotherhood, and moved the members into a permanent home. The order flourished, and neighboring towns soon endeavored to secure its members as teachers in their schools for the poor. Within a year or two, four schools in and about Rheims were placed under masters trained in the house of the Christian Brothers, and a number of other institutions were soon organized in the vicinity upon the same basis.

But, being unable to supply the constant demands for his teachers that came from districts outside the towns, La Salle undertook to train boys who were sent him by the rural clergy, and were expected to return to their homes to teach after their training. To accomplish this, he established in 1684 a `seminary for schoolmasters' in a wing of the house of the brotherhood, and two other seminaries were opened in neighboring towns the following year. Four years later La Salle opened a house for the brotherhood near Paris, and the Christian Brothers were speedily requested to take charge of the schools of several parishes. Despite the jealousy and opposition of the established order of schoolmasters and of many parties in Church and State, the schools and seminaries of the Brothers greatly increased in Paris, and were rapidly extended throughout France. At Paris also La Salle started the `Christian academy,' in which drawing, geometry, and architecture were taught ambitious poor boys on Sunday, and introduced boarding colleges for higher secondary training. And these institutions likewise spread through France and the rest of Europe (Fig. 19). In 1705 La Salle retired to the estate known as Saint Yon, near Rouen, and there opened a home for the brotherhood. Here he also founded a famous boarding-school in which he trained boys for soldiery, farming, trade, and various other vocations. Before long he likewise organized in conjunction an industrial training for youthful delinquents, and both the vocational school and the `protectory' soon became models for many similar institutions in France and elsewhere.

\textbf{The Aim, Curriculum, and Method of the Christian Brothers' Schools.} - The plan of the schools of the Christian Brothers was eventually worked out and crystallized in a fixed system under the title of \emph{Conduct of Schools.} This code has not remained quite as definite and uniform as the \emph{Ratio Studiorum} of the Jesuits, for changes and revisions are permitted, and modern methods and subjects have from time to time been introduced. Considerable latitude, moreover, has been allowed to the individual houses by the Superior General at the head of the order, and by the Brothers Visitors, who have charge of the districts. The educational aim of the Christian Brothers has been preeminently religious, and the chief means of attaining this have been strict vigilance, good example, and catechetical instruction. The course has included the studies of the best schools of the time, and added other more practical subjects. Besides the rudiments - reading, writing, and arithmetic - and religious instruction and good manners, mathematics, history, botany, geography, drawing, architecture, hydrography, navigation, and other technical subjects have often been taught, and in the industrial schools a manual and vocational training has been furnished. La Salle seems to have made a great advance, too, in educational economy by perfecting and applying the `simultaneous' method, which had been practiced in a crude form by some of his forerunners. By this method is meant grading the children according to their capacity, and having those in each grade use the same book and follow the same lesson under a single master, instead of instructing each pupil individually, as was generally the custom then. Likewise, the seminaries or training schools of the Christian Brothers contributed much to the advancement of efficiency in teaching. For the first time teachers of ability and training were made possible for the elementary schools.

\textbf{Influence of the Schools of the Christian Brothers.} - The work of the Christian Brothers has met with steady growth and development. By the time of La Salle's death (1719), there had come to be twenty-seven houses of the order, with two hundred and seventy-four brothers, educating about nine thousand pupils. Before the close of the century these numbers had about quadrupled, and now they have increased nearly a hundredfold since the founder's day. During the nineteenth century these institutions were established in all the states of Europe, Asia, Northern Africa, and America. The educational system has been much modified and expanded, and now includes colleges, technical and industrial schools, academies and high schools, elementary and grammar schools, commercial schools, asylums, and protectories. Thus La Salle and his schools of the Christian Brothers have performed a great service for education in all lines, but especially in the promotion and enrichment of elementary training, which had previously been so neglected.

\textbf{Aim and Content of Education in the Reformation.} - It can now be seen that, as a result of the Reformation, the religious and theological aim of education at all stages became very prominent with Catholics and Protestants alike. In the elementary schools, beside the rudiments, the Scriptures, the Lord's prayer, the ten commandments, and the Catholic, Lutheran, Calvinist, or Anglican creed and catechism were taught, and, with the Protestants, also the hymns of the church. The courses in the secondary schools and universities contained large religious elements, as well as the formal humanism into which the Renaissance of the North had degenerated. Likewise, there was furnished in all universities a training in dialectic, rhetoric, and theology for the sake of efficient controversy with ecclesiastical opponents.

\textbf{Effect of the Reformation upon Elementary Schools.} - But while the Catholics were inclined to leave the organization of education in the hands of various religious bodies, the Protestants more often thought it wise to have its support and control administered by the princes and the state. Owing to this secular management and their position on universal education, the Protestants, with the exception of the Anglicans, who had altered but little in doctrine, were inclined to establish state school systems and hold to the duty of providing and requiring elementary education at public expense. In this way the germs of the modern tendency toward universal, free, and compulsory education began to appear, although they did not ripen until much later.

In the German states there were many illustrations of the spread of elementary education and civic control. As an immediate result of Luther's \emph{Letter to the Mayors} in 1524, the city of Magdeburg united its parish schools under one management and adopted the Protestant ideals. So, in 1525, the school at Eisleben, organized upon a Protestant basis (see p.~128), included elementary as well as secondary work. Similar ideals and organization appear in the provision for `German' schools in the `Church orders' sent out by Bugenhagen (see pp.~128 i.) to the Protestant cities and states of Northern Germany. A further step was taken in 1528 when Melanchthon drew up a plan for schools throughout the entire Electorate of Saxony. This, the first state school system in history, was followed by one in Wiirtemberg, where in 1559 Duke Christopher adopted an improvement upon the Saxon plan, which called for a religious and elementary training for the children of the common people in every village of the duchy. Brunswick in 1569, and Saxony in 1580, followed the lead of Wiirtemberg in revising their school systems. Before the middle of the next century, a number of other states of Germany, such as Weimar, Hessen-Darmstadt, Mecklenburg, Holstein, Hessen-Cassel, and Gotha modeled elementary school systems after those of Saxony and Wiirtemberg. While the Catholics did not in general maintain public elementary education, the Christian Brothers and others undertook a great work in this direction, and Duke Albrecht V of Bavaria even ordered throughout his state the establishment of `German' schools with instruction in reading, writing, and the Catholic creed. This organization of universal education continued its advance, despite the decimation and the general havoc upon finance and education wrought by the Thirty Years' War (1618-1648), and by the end of the eighteenth century practically every village throughout the German states had its \emph{Volksschule} or `people's school.' These institutions were under the direction of the pastor of each parish, and while actual conditions may often have been somewhat below the statutory level and in many cases were a wretched apology, every child not studying at a secondary school was in theory obliged, between the ages of six and thirteen, to attend one of these schools of the people (Fig. 20).

As a result of the Dutch Reformed movement, Holland also made early provision for instruction in religion, reading, and writing. The Church at various synods, and civic authorities in many statutes, recognized the need of universal training, and finally the great Synod of Dort, by a combination with the civil government, in 1618 required every parish to furnish elementary education for all. Similarly, through Knox, Scotland established elementary schools under the control of the parishes. Preliminary steps in this direction were taken by the Privy Council and the Scotch Parliament early in the seventeenth century, and in 1646 the parliament further enacted that there be ``a Schoole founded, and a Schoole master appointed in every Parish,'' and provided that if a parish should fail in this duty, the presbytery should have power to establish the school and compel the parish to maintain it. Half a century later this school system was given over more fully to the control of the State, but even then much of the old connection with the Church was apparent. These schools gave instruction in readings writing, and religion, with the Bible as text, and have done a wonderful work in raising the level of intelligence and affording an opportunity to the children of the lower classes in Scotland. England herself continued to hold to aristocratic and `selective' education, and gave little heed to the establishment of elementary schools; but the American colonies, as far as they were founded by Calvinists or Lutherans, provided early for elementary education (see p.~189). The Puritan towns of the Massachusetts colony established schools almost as soon as they were settled, and in 1647 the legislature enacted that all towns with fifty families should provide an elementary school. Connecticut followed the example three years later, and before the close of the century, similar action was taken by New Hampshire and Vermont (see pp.~197 and 199). Likewise, New Amsterdam and the villages of New Netherlands followed the example of the Mother Country and provided public schools in connection with each church through the support of the Dutch West India Company or of the civil and ecclesiastical bodies jointly (see pp.~193 f).

\textbf{Effect of the Reformation upon the Secondary Schools.} - While the development of elementary instruction and state systems of education was the most important educational outcome of the Reformation, the movement had a somewhat similar effect upon the humanistic secondary education of the time. In Protestant Germany the Latin schools and gymnasia came under the control of the princes and the State rather than the Church, and gradually became the backbone of the state school systems. But they stressed the religious element in their curriculum, and the direct management of education was simply transferred to Protestant ministers or leaders. The schools were still taught and inspected by representatives of the Church, but the form of the organization and administration of education was radically changed. In England there was a similar transfer of management to the Protestant clergy. The existence of the schools had to be authorized and their teachers licensed by the bishop, and they were at all times liable to visitation from ecclesiastical authority. The grammar schools, however, were never organized like the gymnasia, but each school remained independent of the rest and of any national combination. Nor were the Calvimstic colleges united into a national system, except where they came into Germany, when they were absorbed into the system of the gymnasia. The state system of education established by the Scotch parliament in the parishes, often gave secondary training, as well as elementary. And in America the establishment and control of the `grammar' schools, inherited from the mother country, were vested in the authorities of the state and the several towns. On the other hand, the Catholic education in all countries found its secondary schools largely in the colleges of the Jesuits, and the subordination of the individual to authority and the Church was insisted upon.

\textbf{Influence of the Reformation upon the Universities.} - In the case of the universities, many remained loyal to Catholicism and a few new Catholic foundations grew out of the Reformation. All these adhered to the principle of submission to ecclesiastical authority. But the majority of the universities in the Protestant states of Germany followed their princes when they changed from the old creed to the new. Wittenberg, through its connection with Luther and Melanchthon, was the first German university to become Protestant, but others, like Marburg, Konigsberg, Jena, Helmstadt, and Dorpat followed rapidly. Altdorf and Strassburg were developed out of gymnasia. The English universities, Oxford and Cambridge, went over to Protestantism with the national Church. In America, too, Harvard and other early colleges were closely connected with the various commonwealths and with the Calvinistic or the Anglican communion, according to the colony.

\textbf{The Lapse into Formalism.} - There came to be both in Catholic and Protestant institutions a tendency to regard the subjects taught as materials for discipline rather than as valuable for their content. The studies largely became an end in themselves and were deprived of almost all their vitality. The curriculum of the institutions became fixed and stereotyped in nature, and education lapsed into a formalism but little superior to that of the mediaeval scholastics. The methods of teaching came to stress memory more than reason. The Protestants had claimed to depend less upon uncritical and obedient acceptance of dogma than upon the constant application of reason to the Scriptures, but they soon tended to emphasize the importance of authority and the repression of the individual quite as clearly as the Catholics, who definitely held that reason is out of place and unreliable as a final guide in education and life. Hence, except for launching the great conception of state support and control of education, the Reformation accomplished but little directly making for individualism and progress, either through the Catholic awakening or the Protestant revolts. Education fell back before long into the grooves of formalism, repression, and distrust of reason. There resulted a tendency to test life and the educational preparation for living by a formulation of belief almost as much as in the days of scholasticism.

\textbf{Supplementary Reading}

Graves, \emph{During the Transition} (Macmillan, 1910), chap. XVXVI; Monroe, \emph{Text-book} (Macmillan, 1905), chap. VII. An excellent interpretative account of the Reformation is that in Adams, G. B., \emph{Civilization during the Middle Ages} (Scribner, 1894), chaps. XVI and XVII. Painter, F. V. N., furnishes a good translation of \emph{Luther on Education} (Lutheran Publication Society, Philadelphia). Richard, J. W., gives a good account of \emph{Melanchthon, the Protestant Preceptor of Germany} (Putnam, 1898), especially chaps. II-IV and VII; Watson, F., of \emph{Maturinus Corderius, the Schoolmaster of Calvin} (School Review, vol.~XII, nos. 4, 7, and 9); Graves, F. P., \emph{Ramus and the Educational Reformation of the Sixteenth Century} (Macmillan, 1912) of conditions in France; and Leach, A. F., of the dissolution acts of Henry VIII and Edward VI in \emph{English Schools at the Reformation} (Constable, London, 1896), pp.~58-122. On the side of Catholic education, one should read Schwickerath, R., \emph{Jesuit Education} (Herder, St.~Louis), chaps. Ill-VIII and XV-XVIII; Cadet, F., \emph{Port Royal Education} (Bardeen, Syracuse, 1899; George Allen and Co., London) pp.~9-119; and Wilson, Mrs.~R. F., \emph{Christian Brothers} (London, 1883), which gives an epitome of Ravelet, A., \emph{Life of La Salle.} The influence of the Reformation upon the German schools and universities, both Protestant and Catholic, is shown in Nohle E., \emph{History of the German School System} (Report of the United States Commissioner of Education, 1897-98, vol.~I), pp.~30-40; and Paulsen, F., \emph{German Education} (Scribner, 1908), pp.~79-85.

\hypertarget{early-realism-and-the-innovators}{%
\chapter{14. Early Realism and the Innovators}\label{early-realism-and-the-innovators}}

\textbf{Outline}

\begin{itemize}
\tightlist
\item
  The intellectual awakening that appeared in the Renaissance and the Reformation found another avenue for expression in early realism.
\item
  This movement had two phases: (i) humanistic realism, which emphasized the content of classical literature; and (2) social realism, which strove to adapt education to actual life. But the two phases generally occurred together, and the classification of a treatise under one head or the other is largely a matter of emphasis.
\item
  The influence of the two phases was mostly indirect, but through social realism a special training arose in the \emph{Ritterakademien} in Germany, while Milton's humanistic realism was embodied in the `academies' of England, and afterward of America.
\end{itemize}

\textbf{The Rise and Nature of Realism.} - By the seventeenth century it is obvious that humanism was everywhere losing its vitality and declining into a narrow `Ciceronianism,' and that the Reformation was hardening once more into fixed concepts and a dogmatic formalism. The awakened intellect of Europe, however, was tending to find still another mode of expression in the educational movement that is usually known as `realism.' The process of emancipating the individual from tradition and repressive authority had not altogether ceased, but it was manifesting itself mainly through a rather different channel. The movement of realism implied a search for a method by which `real things' may be known. In its most distinct and latest form, - `sense realism,' it held that real knowledge comes through the senses and reason rather than through memory and reliance on tradition, and in this way it interpreted the `real things' as being individual objects. Educational realism, therefore, concerned itself ultimately with investigation in the natural sciences; and it might well be denominated `the beginnings of the scientific movement,' were it not that such a description and the earlier neglects the earlier phases of the realistic development.

\textbf{Humanistic Realism.} - For, even before objects were regarded as the true realities, there seems to have been an effort among some later humanists to seek for the `real things' in the ideas that were represented by the written words. This broader type of humanism, in consequence, tended to break from a restriction to words and set forms and return to the interest in the content of classical literature that marked the Renaissance before its decline into formalism. It may, therefore, properly be called `humanistic realism.' With its emphasis upon content usually went a study of social and physical phenomena, *in order to throw light upon the passages under consideration. Illustrations of this humanistic realism are found in many writers of the sixteenth and seventeenth centuries. Milton (1608-1674), for example, while a remarkable classicist himself, in his \emph{Tractate of Education} objects to the usual humanistic education with ``its gramma tic flats and shallows where they stuck unreasonably to learn a few words with lamentable construction''; and says of the pupil, ``if he have not studied the solid things in them as well as the words and lexicons, he were nothing so much to be esteemed as any yeoman or tradesman competently wise in his mother dialect only.'' And he would teach the Latin writers on agriculture, and the Greek writers on natural history, geography, and medicine for the sake of the subject matter.

\textbf{Social Realism.} - But there was another phase of early realism, which often appeared in conjunction with humanistic education, and may be called `social realism.' Its adherents strove to adapt education to actual living in a real world, and to afford direct practical preparation for the opportunities and duties of life. It generally recommended as the means of education for all members of the upper social class. It sought to combine with the literary elements taught the clergy in the Middle Ages and the scholar in the Renaissance, certain remnants of the old chivalric education as the proper training for gentlemen. It held schools to be of less value as an agency for educating the young aristocrats than training through a tutor and travel. Hence an education in social realism usually included a study of heraldry, genealogy riding, fencing, and gymnastics, and involved a study of modern languages and the customs and institutions of neighboring countries.

A good illustration of this type of education is found in the educational essays of Montaigne (1533-1592). In the \emph{Education of Children} he holds that virtue comes from experience and breadth of vision rather than from reading, and declares: ``I would have travel the book my young gentleman should study with most attention; for so many humors, so many sects, so many judgments, opinions, laws, and customs, teach us to judge aright of our own, and inform our understanding to discover its imperfection and natural infirmity.'' This training, too, he feels, should be under the care of a tutor, who is to be a man of the world, one ``whose head is well tempered, rather than well filled.'' While a gentleman has need of Latin and Greek, Montaigne maintains that one should first study his own language and those of his neighbors. He also stresses physical exercise, and fears the training of boys near their mothers, who ``will not endure to see them mount an unruly horse, nor take a foil in hand against a rude fencer.''

An educational work based on social realism that has been studied even more than the \emph{Essays} of Montaigne is \emph{Some Thoughts concerning Education} by John Locke (1632-1704). Locke states the aims of education in order of their value as \emph{'Virtue, Wisdom} (i. e., worldly wisdom), \emph{Breeding,} and \emph{Learning'}; and holds that such a training can be secured by the young gentleman only through a tutor, who ``should himself be well-bred, understanding the Ways of Carriage and Measures of Civility in all the Variety of Persons, Times and Places, and keep his Pupil, as much as his Age requires, constantly to the Observation of them.'' In considering the subjectmatter of the training, he maintains that ``besides what is to be had from Study and Books, there are other \emph{Accomplishments} necessary for a Gentleman, - dancing, horseback riding, fencing and wrestling.''

\textbf{The Relations of Humanistic to Social Realism.} - Humanistic and social realism, however, constantly appear together in the works of the same author, and it is often difficult to distinguish a writer as advocating one type or the other. The differentiation seems to be largely a matter of emphasis. While one element or the other may seem to be more prominent in the treatise of a certain author, the two phases of education are largely bound up in each other. While Milton, for instance, is in the main a humanistic realist and advises an education in languages and books, he recommends that considerable time be given, toward the end of the course, to the social sciences - history, ethics, politics, economics, theology - and to such practical training as would bring one in touch with life. He also specifically advocates the experience and knowledge that would come from travel in England and abroad; and defines education as ``that which fits a man to perform justly, skillfully, and magnanimously all the offices both private and public of peace and war.'' On the other hand, Montaigne, the social realist, seems quite as strenuous in urging a more realistic humanism. In his essay, \emph{On Pedantry,} he launches most vigorous ridicule against the prevailing narrow humanistic education, with its memorizing of words and forms, and insists: ``Let the master not only examine him about the words of his lesson, but also as to the sense and meaning of them, and let him judge of the profit he has made, not by the testimony of his memory, but that of his understanding.''

And it is equally difficult to state whether humanistic or social elements prevail in Locke's \emph{Thoughts,} the \emph{Gargantua} of Rabelais (1495-1553), the \emph{Positions} of Mulcaster (1530-1611), and other treatises of the period. It is true, of course, that in certain other works written upon the training of the aristocracy, social realism is more exclusively stressed. The titles of most of these reveal their content, as can easily be seen in the case of such productions as Castiglione's \emph{The Courtier} (1528), Elyot's \emph{The Governour} (1531), Peacham's \emph{The Compleat} '\emph{Gentleman} (1622), and Brathwaite's \emph{The English Gentleman} (1630). But, in most of the early realistic works, humanistic and social elements are inextricably interwoven; and humanistic and social realism, taken together, seem to constitute a natural bridge from humanism over to sense realism.

\textbf{The Influence of the Innovators upon Education.} - There is, however, a variety of other brilliant educational suggestions in each of these early realists. All of them hold to a broader and better rounded training and more natural and informal methods than those in vogue. Mulcaster even advocates universal elementary education, the professional training of teachers, and the education of girls, and undertakes to make a naive analysis of the mind as the basis of a philosophy of education. So suggestive have the recommendations of the early realists proved to modern education that these authors are often known as the `innovators.' Yet their theories do not seem to have affected greatly the educational practice of the times. They did tend to disrupt traditionalism and the formal humanism, to bring education into touch with society and preparation for real life, and to popularize a wider content and a more informal procedure, but their influence appeared through their successors and later education rather than directly in the schools of the period. Locke, for instance, in addition to the influence he had upon Rousseau, Pestalozzi, and other reformers, must in some measure have been responsible for the great development of the physical and ethical sides of education in the public and grammar schools of England, together with the tendency of these institutions to consider such aspects of rather more importance than the purely intellectual. His plea for a tutor as the means of shaping manners and morals has also probably had its effect upon the education of the English aristocracy.

\textbf{The Ritterakademien.} - In the German states, on the other hand, there arose at the courts during the seventeenth century an actually new type of educational institution as the outgrowth of social realism. Here, in place of the old humanistic education, there was developed a special training for the young nobles in French, Italian, Spanish, and English, in such accomplishments as courtly conduct, dancing, fencing, and riding, and in philosophy, mathematics, physics, geography, statistics, law, genealogy, and heraldry. The educational institutions in which this training was embodied were known as \emph{Ritterakademien} or `academies for the nobles.' Such academies were founded at Colberg, Luneberg, Vienna, Wolffenbuttel, and many other centers before the close of the century. They originally covered the work of the gymnasia, although substituting the modern languages, sciences, and the knightly arts that have been mentioned for the Greek and Hebrew, and adding a little from the course of the university. Gradually, however, they became part of the regular secondary system.

\textbf{The Academies in England.} - Milton's suggestions were ultimately materialized in an even more influential type of school. In the \emph{Tractate} he had recommended that his ideal education be carried out in an institution to be known as an `academy.'. Such a school was to be erected `in every city throughout this land.' It should train boys from the age of twelve to twenty-one, and should provide both secondary and higher education. `Academies,' based very closely upon "this plan, were about a generation later actually organized in a number of places by the Puritans. Under the harsh Act of Uniformity (1662) two thousand non-conforming clergymen were driven from their parishes, and in many instances found school-teaching a congenial means of earning a livelihood, and at the same time of furnishing higher education to the young dissenters, who were excluded from the universities and grammar schools. The first of these academies was that established by Richard Frankland at Rathmill in 1665, and this was followed by the institutions of John Woodhouse at Sheriffhales, of Charles Morton at Newington Green, and of some thirty other educators of whom we have record at other places. These academies were largely humanistic in their realism, and, since their chief function was to fit for the ministry, they included Latin, Greek, and Hebrew in their course, but they were also rich in mathematics, natural and social sciences, modern languages, and the vernacular. The new tendency was also broadened and amplified by Locke's \emph{Thoughts} (1693), which became the great guide for the managers of the Puritan academies. In 1689, when the Act of Toleration put non-conformity upon a legal footing, the academies were allowed to be regularly incorporated.

\textbf{The Academies in America.} - Academies arose also in America. When the number of religious denominations had greatly increased and the demands upon secondary education had expanded, the `grammar schools' (see pp.~120 f.), with their narrow denominational ideals and their limitation to a classical training and college preparation, proved inadequate, and efforts were made to organize academies as a supplement. Their may have been earlier academies in America, but the first well-known suggestion of an academy was made in 1743 by Benjamin Franklin. He wished to inaugurate an education that would prepare for life, and not merely for college. Accordingly, he proposed for the youth of Pennsylvania a course in which English grammar and composition, penmanship, arithmetic, drawing, geography, history, the natural sciences, oratory, civics, and logic were to be emphasized. He would gladly have excluded Latin and other languages altogether, but for politic reasons these courses were allowed to be elective. Through the efforts of a number of leading citizens, such an academy was opened at Philadelphia in January, 1750 (although not chartered until July, 1753). During the next generation a number of similar institutions sprang up, especially in the middle and southern colonies. A great impulse was given the movement by the foundation of the two Phillips academies, - one in 1780 at Andover, Massachusetts, and the other the next year at Exeter, New Hampshire. The Dummer Grammar School was reorganized as an academy in 1782, and the movement spread rapidly throughout New England during the last two decades of the eighteenth century.

Shortly after the Revolution, owing in part to the inability or unwillingness of the towns to maintain grammar schools, and in part to the wider appeal and greater usefulness of the academies, the latter institutions quite eclipsed the former, and became for about half a century the prevailing type of secondary school in the United States. They were usually endowed institutions managed by a close corporation, but were often largely supported by subscriptions from the neighborhood, and sometimes subsidized by the state. Located in small towns or villages, they served a wide constituency and made provision for boarding, as well as day pupils. Unlike the grammar schools, they were not originally intended to prepare for the learned professions exclusively, but, as time passed, they tended more and more to become preparatory schools for the colleges, instead of finishing schools for the middle classes of society. The academies were also the first institutions of secondary education to offer opportunities to women. Many of them were co-educational, and others, frequently burdened with the name of `female seminary' were for girls exclusively. Academies for some time likewise furnished the only means of training teachers for the elementary schools, and have generally played an important part in education in the United States.

\textbf{Supplementary Reading}

Graves, \emph{During the Transition} (Macmillan, 1910), chap. XVII; and \emph{Great Educators of Three Centuries} (Macmillan, 1912), chaps. I and V; Monroe, \emph{Text-book} (Macmillan, 1905), pp.~442-460. An excellent edition of Milton's \emph{Tractate of Education} is that by Morris, E. E. (Macmillan, 1895); of Montaigne's \emph{Education of Children} that by Rector, L. E. (Appleton, 1899); of Locke's \emph{Thoughts concerning Education,} and of Mulcaster's \emph{Positions,} those by Quick, R. H. (Cambridge University Press, 1895, and Longmans, 1888, respectively); and of Rabelais' \emph{Gargantua,} that by Besant, W. (Lippincott, Foreign Classics for English Readers). The works of Castiglione, Elyot, Peacham, Brathwaite, etc., are also extant. For an account of the \emph{Ritterakademien,} see Nohle, E., \emph{History of the German School System} (Report of the U. S. Commissioner of Education, 1897-98), pp.~41 f., and Paulsen, F., \emph{German Education} (Scribner, 1908), pp.~112-116; and of the academies, Brown, E. E., \emph{The Making of Our Middle Schools} (Longmans, Green, 1902), chaps. VIII and IX.

\hypertarget{sense-realism-and-the-early-scientific-movement}{%
\chapter{15. Sense Realism and the Early Scientific Movement}\label{sense-realism-and-the-early-scientific-movement}}

\textbf{Outline}

\begin{itemize}
\tightlist
\item
  In the seventeenth century scientific investigation developed rapidly, and led theorists to introduce science into the curriculum and to advocate a study of `real things.'
\item
  Bacon undertook to formulate induction, and while he did not understand the importance of an hypothesis, he did much to rid the times of \emph{a priori} reasoning.
\item
  On the basis of sense realism, Ratich anticipated many principles of modern pedagogy, but he was unsuccessful in applying his ideas.
\item
  Comenius (1) produced texts for teaching Latin objectively, (2) crystallized his educational principles in the \emph{Great Didactic,} and (3) attempted an encyclopaedic organization of knowledge. He wished to make this knowledge part of the course at every stage of education, and, while he was not consistently inductive, he made a great advance in the use of this method.
\item
  Through sense realism, rudimentary science was introduced into the elementary schools; the \emph{Ritterakademien} and the pietist schools stressed the subject; and professorships of science were founded in the universities.
\end{itemize}

\textbf{The Development of the Sciences and Realism.} - The realistic tendency did not pause with reviving the ideas represented by the words nor with the endeavor to bring the pupil into touch with the life he was to lead. The earlier realism seems to have been simply a stage in the process of transition from the narrow and formal humanism to a realism obtained through the senses, which may be regarded as the beginning of the modern movement to develop the natural sciences. Science had started to develop as early as the time of the schoolman, Roger Bacon (1214-1294), but for three centuries it was not kindly received. Even during the Renaissance the Church had continued to oppose it bitterly, because it tended to conflict with religious dogma, although this age did not object to the revival of the classics. Accordingly, the latter subject became strongly intrenched in educational tradition, and its advocates offered the most obstinate opposition to the sciences. Its numerous representatives struggled hard to keep the sciences out of education.

However, concomitant with the growth of reason and the partial removal of the theological ban, there was developed a remarkable scientific movement, with a variety of discoveries and inventions. For more than a millennium the Greek developments in astronomy and physics had been accepted as final, but toward the close astronomy of the sixteenth and during the seventeenth century these \emph{dicta} were completely upset. The hypothesis of a solar system, which replaced the Ptolemaic interpretation, was published by Copernicus (1473-1543); Kepler (1571-1630) explained the motion of the planets by three simple laws; and, through the construction of a telescope, Galileo (1564-1642) revealed new celestial phenomena. Galileo also demonstrated that all bodies, allowing for the resistance of the air, fall at the same rate; by means of the barometer, Torricelli (1608-1647) and Boyle (1627-1691) proved the existing theories of a vacuum incorrect, and formulated important laws concerning the pressure of gases; and Guericke (1602-1686), inspired by their discoveries, succeeded in constructing an air-pump. Investigations of this kind paved the way for the formulation of the law of universal gravitation and the laws of motion by Sir Isaac Newton (1642-1727), which united the universe into a single comprehensive system and completed the foundations for modern mechanics.

Likewise, about the same time, the other great development in science among the Greeks,-anatomy and physiology, was completely revolutionized. Through the discovery of valves in the veins by means of dissection and vivisection, the hypothesis of the double circulation of the blood by Harvey (1578-1657), and the microscopic demonstrations by Malpighi (1628-1694) of the existence of capillaries connecting the veins and the arteries, the old theory of the motion of the blood through suction, which had been promulgated by Galen, was completely shattered, and a great impetus was given to investigations in anatomy and physiology. In consequence of this scientific progress, the educational theorists began to introduce science and a knowledge of real things into the curriculum. It came to be widely felt that humanism gave a knowledge only of words, books, and opinions, and did not even at its best lead to a study of real things. Hence, new methods and new books were produced, to shorten and improve the study of the classical languages, and new content was imported into the courses of study. The movement also included an attempt at a formulation of scientific principles in education and an adaptation to the nature of the child.

\textbf{Bacon and His Inductive Method.} - The new tendency, however, did not appear in education until after the time of Francis Bacon (1561-1626). The use of the scientific method by the various discoverers was largely unconscious, and it remained for Bacon to formulate method of the what he called the method of `induction,' and by advocating its use, to point the way to its development as a scientific method in education. He is, therefore, ordinarily known as the first sense realist. He reacted from deductive logic, which was currently supposed to be the sole method of Aristotle, and took his cue in formulating a new method of reasoning from the many scientific workers of his time. He made a great advance in his rejection of the contemporary method of attempting to establish the first principles of a science, and then deducing from them by means of the syllogism all the propositions which that science could contain. However, his \emph{Novum Organum,} or `new instrument,' as he called his treatise, in endeavoring to create a method whereby anyone could attain all the knowledge of which the human mind was capable, undertook far too much, and resulted in a merely mechanical procedure. Briefly stated, his plan was, after ridding the mind of individual prejudices, to observe and carefully tabulate lists of all the facts of nature, and from these discover the underlying law by comparing the cases where a certain phenomenon appears and where it does not.

But by this method neither Bacon himself nor anyone else has ever made any real contribution to science. It does not follow that, because all observed cases under certain conditions produce a particular effect, every other instance not yet observed will necessarily have the same effect. The true method of induction, which was evident even in the work of Kepler, and came to be more so in the discoveries of Harvey and Newton, stresses rather the part played by scientific imagination, as it is manifested by men of genius in the forming of an hypothesis. The modern procedure is as follows: - When certain effects are observed, of which the cause or law is unknown, the scientist frames an hypothesis (i. e., makes a conjecture) to account for them; then he tests this hypothesis, by collecting facts and comparing with these facts the conclusions to which his hypothesis would lead; and, if they correspond or agree, he holds that his hypothesis has been confirmed or verified, and maintains that he has discovered the cause or law. Nevertheless, while Bacon did not formulate the inductive method of modern science, he largely helped to rid the times of an unwise dependence upon \emph{a priori} reasoning, and he did call attention to the necessity of careful observation and experimentation, and thus opened the way for real inductive procedure. Probably no book ever made a greater revolution in modes of thinking or overthrew more prejudices than Bacon's \emph{Novum Organum.}

\textbf{Bacon's Educational Suggestions and Influence.} - Bacon was not a teacher, and his treatment of educational problems appears in brief and scattered passages. While he offers isolated suggestions concerning the mental and moral training of the young, he plans no serious modification in the existing organization of schools. He does, however, in his \emph{New Atlantis} imply an interest in promoting scientific research and higher education. In the ideal society depicted in that work, he describes an organization of scholars called `Salomon's House,' whose members in their investigations anticipate much that scientists and inventors have to-day only just begun to realize. Among these anticipations were the variation of species, the infusion of serums, vivisection, telescopes, telephones, flying-machines, submarine boats, and steamengines. From this description Bacon would seem to believe that education should be organized upon the basis of society's gradually accumulating a knowledge of nature and imparting it to all pupils at every stage. At any rate, in his \emph{Advancement of Learning,} he definitely suggests a wider course of study, more complete equipment for scientific investigation, a closer cooperation among institutions of learning, and a forwarding of `unfinished sciences.' And such a plan of \emph{pansophia,} or `universal knowledge,' was specified in the educational creed of the later sense realists, who worked out the Baconian theory of education. Hence, while not skilled or greatly interested in education himself, Bacon influenced profoundly the writing of many who were, and has done much to shape the spirit of modern practice. His method was first applied directly to education by a German known as Ratich and, in a more effective way, by Comenius, a Moravian.

\textbf{Ratich's Methods.} - Ratich (1571-1635) probably became acquainted with the sense realism of Bacon while studying in England, and, when about forty years of age, undertook to found a system of education upon it. In linguistic training, like all realists, he insisted that one ``should first study the vernacular'' as an introduction to other languages. He also held to the principle of ``one thing at a time and often repeated.'' By this he meant that, in studying a language, one should master a single book before taking up another. In his teaching at Köthen, as soon as his pupils knew their letters, they were required to learn \emph{Genesis} thoroughly for the sake of their German. Each chapter was read twice by the teacher, while the pupil followed the text with his finger. When the pupils could read the book perfectly, they were taught grammar from it as a text. The teacher pointed out the various parts of speech and made the boys find other examples, and had them decline, conjugate, and parse. In taking up Latin, a play of Terence was treated in similar fashion. Others of the principles that he used in teaching language and grammar, and especially those which applied to education in general, were even more distinctly realistic. Such, for example, were his precepts, - ``follow the order of nature'' and ``everything by experiment and induction,'' and his additional recommendation that ``nothing is to be learned by rote.'' Thus Ratich not only helped shape some of the best methods for teaching languages, but anticipated the main principles of modern pedagogy. While, owing to obtrusive failings in character and experience, he was uniformly unsuccessful in his practice, he, nevertheless, stirred up considerable thought and stimulated many treatises of others. Thus, through Comenius, who carried out his principles more fully, this German innovator, unpractical as he was, became a spiritual ancestor to Pestalozzi, Froebel, and Herbart.

\textbf{Comenius: His Training and Work.} - John Amos Comenius (1592-1670) was born at Nivnitz, Moravia, and was by religious inheritance a staunch adherent of the Moravian Church. After a course in a Latin school, he spent a couple of years in higher education at the Lutheran College of Herborn and at the University of Heidelberg. In consequence of many vicissitudes in life, he lived and wrote in a number of places, and became acquainted with the work of a variety of men engaged in educational reform and advancement. While the problems with which they were dealing were similar to his own and largely influenced his educational positions, he far surpassed them all in scope of work and greatness of repute. His educational achievements were the outgrowth of sense realism, and appear in three directions: - (1) the series of texts for learning Latin; (2) his \emph{Great Didactic;} and (3) his attempts to create an encyclopaedic organization of knowledge \emph{(pansophia).}

\textbf{His Series of Latin Texts.} - The first of the famous texts that Comenius produced to facilitate the study of Latin was issued in 1631, and has generally been known by the name of \emph{Janua Linguarum Reserata} (The Gate of Languages Unlocked). It was intended as an introductory book to the study of Latin, and consisted of an arrangement into sentences of several thousand Latin words for the most familiar objects and ideas. The Latin was printed on the right-hand side of the page, and on the left was given a translation in the vernacular. By this means the pupil obtained a grasp of all ordinary scientific knowledge and at the same time a start in his Latin vocabulary. In writing this text, Comenius may have been somewhat influenced by Ratich, a review of whose methods he had read at Herborn, but he seems to have been more specifically indebted both for his method and the felicitous name of his book to a Jesuit known as Bateus, who had written a similar work.

It was soon apparent that the \emph{Janua} would be too difficult for beginners, and two years later Comenius issued his \emph{Vestibulum} (Vestibule), as an introduction to it. While the \emph{Janua} contained all the ordinary words of the language, - some eight thousand, there were but a few hundred of the most common in the \emph{Vestibulum.} Later both of the works were several times revised, modified, and enlarged; and grammars, lexicons, and treatises were written to accompany them. He also published a third Latin reader, the \emph{Atrium} (Entrance Hall), which took the pupil one stage beyond the \emph{Janua.} We know, too, that he intended also to write a still more advanced work, to be called \emph{Sapientiae Palatium} (Palace of Wisdom). This fourth book was to consist of selections from the best Latin authors, but it was never completed. He did, however, produce as a supplementary text-book a simpler and more extensive edition of the \emph{Janua,} accompanied with pictures. Each object in the illustrations of this book was marked with a number corresponding to one in the text. This work, which he called \emph{Orbis Sensualium Pictus} (The World of Sense Objects Pictured), is the first illustrated reading book on record (Fig. 21).

\textbf{\emph{The Great Didactic.}} - But these books on teaching Latin realistically were only part of the work that Comenius contemplated. During his whole career he had in mind a definite idea of the aim of education, and of what, in consequence, he wished the organization, subject-matter, and methods to be. His ideas on the whole question of education were formulated in his \emph{Great Didactic} even before the \emph{Janua} appeared, but the work was not published until 1657. In it he strove to assimilate all that was good in the realistic movement and use it as a foundation. He developed many of the principles and methods of Ratich, Bateus, and others, but he owed a greater debt for the suggestions he took from Bacon's \emph{Advancement of Learning,} and even more from the \emph{Encyclopedia} of Alsted, one of his teachers at Herborn. In the \emph{Great Didactic} Comenius formulated an educational aim and constructed an educational organization of his own. Probably, as an outgrowth of his religious attitude, he held to `knowledge, morality, and piety' as the ideals of education, and advocated universal education for `boys and girls, both noble and ignoble, rich and poor.' His organization of education consisted of four-periods of six years each. The first period of instruction was that through infancy, or up to the age of six. It was to be given in the school of `the mother's lap,' which should exist in every house. For childhood, or from six to twelve, was to be organized the `vernacular schooL' which should appear in every hamlet and village. From that time up to eighteen comes the `Latin school' to be maintained in every city; and, finally, for youth from eighteen to twenty-four, there should be a university in every kingdom or province. Such an organization would have made education universal, and would tend to bring about the custom of education according to ability, rather than social status, which was a suggestion some three centuries in advance of the times.

\textbf{His Encyclopaedic Arrangement of Knowledge.} - The rest of the works of Comenius may be regarded as amplifications of various parts of this \emph{Great Didactic.} Besides the Janual series, which he seems to have written for the Latin school, he produced a set of texts for the vernacular school, which soon disappeared, and a handbook for the lowest work, called \emph{The School of Infancy.} But the phase of the \emph{Great Didactic} most often elaborated was the realistic one of \emph{pansophia} or universal knowledge. This principle was not only exemplified in such works as the \emph{Janua} and \emph{Orbis Pictus} and in treatises he wrote upon astronomy and physics, but in various educational institutions that he undertook to found, and it remained the ruling passion throughout his life. In the \emph{Great Didactic} he went so far as to hold that an encyclopaedic training should be given at every stage of education, - mother school, vernacular school, Latin school, and university.

But, while even in the mother school the infant was to make a beginning with geography, history, and various sciences, grammar, rhetoric, and dialectic, music, arithmetic, geometry, and astronomy, and the rudiments of economics, politics, ethics, metaphysics, and religion, his attainment was not expected to be as formidable as the names of the subjects sound. It was to consist merely in understanding simple causal, temporal, spatial, and numerical relations; in distinguishing sun, moon, and stars, hills, valleys, lakes, and rivers, and animals and plants; in learning to express oneself; and in acquiring proper habits. It was, in fact, not unlike the training of the modern kindergarten. In a similar way each succeeding stage is to enlarge the body of knowledge along all these lines. ``The different schools are not to deal with different subjects, but should treat the same subjects in different ways; throughout graduating the instruction to the age of the pupil and the knowledge that he already possesses. In the earlier schools everything is taught in a general and undefined manner, while in those which follow the information is particularized and exact.'' Moreover, beyond the university, which, like the lower schools, was to make teaching its chief function, Comenius held it to be important that somewhere in the world there should be a `didactic college' devoted to scientific investigation, in which learned men from all nations should cooperate. Such an institution would form a logical climax to his system of schools, bearing the same relation to them that the stomach does to the other members of the body by ``supplying blood, life, and strength to all.''

\textbf{The Method of Nature.} - The way in which this pansophic instruction should be given, Comenius also intended to have in full accord with sense realism. He insists that the `method of nature' must be observed and followed, and then shows how nature accomplishes all things `with certainty, ease, and thoroughness,' in what respects schools have deviated from the principles of nature, and how they can be rectified only by following her plans. These principles concerning the working of nature were laid down \emph{a priori,} but it is probable that they had been previously worked out inductively from his schoolroom experience. At times, though, they were put in the form of fanciful analogies. For example, he declares that because a bird by nature hatches her young in the spring or early part of the year, schools have erred (1) in not requiring education to begin in the springtime of life, or boyhood, and (2) in not selecting the springtime of the day, or the morning hours, for study. But it is not remarkable that, with all his realistic tendencies, Comenius did not consistently employ induction. The natural sciences were young in his day, so that he did not altogether grasp their content and method, and he had partially inherited the scholastic notion that truth cannot be fully secured through the senses or by reason. It is sufficient merit that Comenius, for the first time in history, applied anything like induction to teaching. Moreover, in the application of his general method to the specific teaching of various lines, - sciences, reading, writing, singing, languages, morality, and piety, he utilized more fully the induction of Bacon. For example, after showing the necessity for careful observation in obtaining a knowledge of the sciences, he gives nine useful precepts for their study that are clearly the inductive result of his own experience as a teacher. Likewise, he insists that, in teaching the sciences, in order to make a genuine impression upon the mind, one must deal with realities rather than books. The objects themselves, or where this is not possible, such representations of them as can be conveyed by copies, models, and pictures, must be studied. After the same principle he formulates inductive rules and methods for instruction in the other subjects.

\textbf{The Influence of Comenius upon Education.} - Thus the work of Comenius was based primarily upon sense realism, but he added many modifications and new elements of his own. He may in the fullest sense be considered the great educational theorist and practical reformer of the seventeenth century. His practical ability is especially shown in the series of Latin text-books, which far excelled the works of several contemporaries on similar lines. The \emph{Janua} was translated into a dozen European, and at least three Asiatic languages; the \emph{Orbis Pictus} proved even more popular, and went through an almost unlimited number of editions in various tongues; and the whole series became for many generations the favorite means of introducing young people to the study of Latin. But the remarkable theoretical work of Comenius had little effect upon the schools of the period, and until about the middle of the nineteenth the \emph{Great Didactic} was scarcely known. At that time, when this treatise of Comenius was brought to light by German investigators, it was discovered that the old realist of the seventeenth century had been the first to deal with education in a scientific spirit, and work out its problems practically in the schools. And the principles of Comenius were at the time unconsciously taken up by others and indirectly became the basis of modern education. His spirit appeared not only in the ideas of subsequent theorists - Francke, Rousseau, Basedow, Pestalozzi, Herbart, Froebel - but even in the actual curricula and methods of educational institutions.

\textbf{Realistic Tendencies in Elementary Schools.} - While the effect of sense realism upon the schools seems to have been slow and indirect, the movement was obvious even in the seventeenth century. In Germany there came a decided tendency throughout the elementary schools to increase instruction in the vernacular, as recommended by Ratich and Comenius, and to learn first the German grammar rather than the Latin. With this movement was joined the increase in universal and compulsory education urged by the reformers, and an introduction of elementary science, in addition to reading, writing, arithmetic, religion, and singing. At Weimar in 1619, through a pupil of Ratich, a new school system was organized; and in 1642, under the order of Duke Ernst, Andreas Reyher prepared a new course for Gotha, which afforded elementary instruction in the natural sciences, as well as the rudiments and religion. This work included teaching the children to measure with the hour-glass and sun-dial, to observe the ordinary plants and animals, and to carry on other objective studies of a simple character. Many other attempts at instruction in science were made elsewhere in the German states, both in private and public education, and the same tendency appeared in the states of Italy, and in France, Holland, and England.

\textbf{Secondary Schools.} - But the new realistic tendencies appeared also in secondary education. While in Germany it was not until the eighteenth century that there were any evidences of sense realism in the gymnasia, languages of neighboring countries and considerable science appeared in the \emph{Ritterakademien} (see p.~157) by the middle of the seventeenth, and toward the end of the century in the schools of Francke and other `pietists' at Halle were embodied all the realistic elements of Comenius. While the pietists adopted these ideas largely for their religious side, as a protest and reaction to the rationalistic \emph{Ritterakademien,} they did not hesitate also to stress the science content and the study of the vernacular. In the secondary school known as the \emph{Padagogium,} which he had started for well-to-do boys, Francke included training in the vernacular, mathematics, geography, natural science, astronomy, anatomy, and materia medica; and the \emph{Realschule,} established by his colleague, Semler, went even more fully into the vernacular, mathematics, and the sciences, pure and applied. This realistic instruction of the pietists was brought by Hecker to Berlin, where he started his famous \emph{Realschule} in 1747, and similar institutions soon spread throughout Prussia. In England, while very few of the grammar and public schools (see p.~120) as yet introduced even the elements of science into their course, the academies (see p.~157) were rich in sciences, mathematics, and the vernacular. This was also true of the academies that sprang up in America (see p.~158).

\textbf{The Universities.} - The universities were slower in responding to the movement of sense realism. As the result of its pietistic origin, however, the University of Halle was realistic almost from its beginning in 1692. Göttingen, the next institution to become hospitable to the tendency, did not start it until 1737. But soon afterward the movement became general, and by the end of the eighteenth century all the German universities - at least, all under Protestant auspices - had created professorships in the sciences. While the English universities, Oxford and Cambridge, were much slower than those of Germany in adopting the new subjects, and it was a century and a half before these institutions became known for their science, during the professorship of Isaac Newton (1669-1702) considerable was done toward making Cambridge mathematical and scientific, and in the course of the eighteenth century several chairs in the sciences were established. Besides formulating the law of gravitation, Newton lectured and wrote at Cambridge upon calculus, astronomy, optics, and the spectrum. He became one of the greatest mathematicians and physicists the world has known, and he did much to create a scientific atmosphere in other educational institutions, as well as Cambridge. America also felt the scientific impulse in its higher institutions. Some study of astronomy, botany, and physics was possible at Harvard even in the seventeenth century, and during the eighteenth Yale, Princeton, King's (afterward Columbia), Dartmouth, Union, and Pennsylvania all came to offer a little work in physics, and at times in chemistry, geology, astronomy, and biology. In his proposals for the prospective `seminary' in New York (1753), which was destined to become Columbia University, and in the actual course of the academy at Philadelphia (later the University of Pennsylvania), over which he presided, Dr.~William Smith put a most progressive program of sciences, including the rudiments of mechanics, physics, chemistry, geology, astronomy, botany, zoology, and physiology. But for half a century after this American institutions did little with the sciences as laboratory studies.

\textbf{Supplementary Reading}

Graves, \emph{During the Transition} (Macmillan, 1910), chap. XVIII; and \emph{Great Educators of Three Centuries} (Macmillan, 1912), chaps. II, IV, and VI; Monroe, \emph{Text-book} (Macmillan, 1905), pp.~461-501. The following works are standard for the authors mentioned: Adamson, J. W., \emph{Pioneers of Modern Education} (Macmillan, 1905), chap. Ill (Bacon); Barnard, H., \emph{German Teachers and Educators,} pp.~343-370 (Ratich); Fowler, T., Bacon's \emph{Novum Organum} (Oxford, Clarendon Press); Laurie, S. S., \emph{John Amos Comenius} (Bardeen, Syracuse, 1892); Monroe, W. S., \emph{Comenius} (Scribner, 1900); and Quick, R. H., \emph{Educational Reformers} (Appleton, 1896), chap. IX (Ratich) and X (Comenius). An account of sense realism is afforded by Adamson, \emph{op. cit.,} chap. I, and of its effect upon the schools by Barnard, \emph{op. cit.,} pp.~302-317, and by Paulsen, F., \emph{German Education} (Scribner, 1908), pp.~117-133.

\hypertarget{formal-discipline-in-education}{%
\chapter{16. Formal Discipline in Education}\label{formal-discipline-in-education}}

\textbf{Outline}

\begin{itemize}
\tightlist
\item
  Locke is often classed with the advocates of realism or of naturalism, but the keynote to his thought is `discipline.' This is to be obtained in intellectual training through mathematics; in moral training, through the control of desires by reason; and in physical training, through a `hardening process.'
\item
  Locke has, therefore, often been viewed as the great advocate of the theory of formal discipline, according to which certain subjects yield a general power that may be applied in any direction, and should be studied by all.
\item
  This doctrine has greatly influenced education, but in the late nineteenth century there was a decided reaction from it. Recently this extreme reaction has been modified, and a position taken with which Locke's real attitude would seem to be in harmony.
\end{itemize}

\textbf{Locke's Work and Its Various Classifications.} - Because of their relation to an important topic in modern education, the theories of John Locke (1632-1704) should receive further attention than they have yet been given. No writer on education has been more variously classified than he. We have already seen (p.~154) that the general tenor of his \emph{Thoughts concerning Education} would lead us to group him with the early realistic movement. There are also elements in this work that would seem to place him with the sense realists, and many of his ideas proved so similar and suggestive to Rousseau's thought (see p.~213), that he has sometimes been classed among the advocates of naturalism. But Locke's \emph{Thoughts,} by which his educational position is often exclusively judged, were simply a set of practical suggestions for the education of a gentleman, written for a friend as advice in bringing up his son. They make clear his general sympathy with the current educational reform, but do not bring out his main point of view. His central thought appears more definitely through the philosophical principles in his famous \emph{Essay concerning the Human Understanding,} and through the intellectual training suggested in his other educational work, \emph{Conduct of the Understanding,} which was originally an additional book and application of the \emph{Essay.}

\textbf{Locke's Disciplinary Theory in Intellectual Education.} - Probably Locke's underlying thought as to the proper method of intellectual, moral, and physical training may best be summed up in the word `discipline.' This educational attitude is a natural corollary of his philosophic position. In his \emph{Essay} he holds that ideas are not born in one, but that all knowledge comes from experience. The mind, he declares, is like `white paper, or wax,' upon which impressions from the outside world are made through our senses. When the ideas are once in mind, it is necessary to determine what they tell us in the way of truth. Hence, to train the mind to make proper discriminations, he declares in the \emph{Conduct of the Understanding} that practice and discipline are necessary. ``Would you have a man reason well, you must use him to it betimes, exercise his mind in observing the connection of ideas and following them in train.'' As to the means of effecting this mental discipline, Locke holds: ``Nothing does this better than mathematics, which therefore I think should be taught all those who have the time and opportunity, not so much to make them mathejnaticians as to make them reasonable, creatures, that having got the way of reasoning, which that study necessarily brings the mind to, they might be able to transfer it to other parts of knowledge as they shall have occasion.'' Similarly, he advises a wide range of sciences, ``to accustom our minds to all sorts of ideas and the proper ways of examining their habitudes and relations; not to make them perfect in any one of the sciences, but so to open and dispose their minds as may best make them capable of any, when they shall apply themselves to it.''

\textbf{Disciplinary Attitude in Moral and Physical Training.} - The same disciplinary conception of education underlies Locke's ideals of moral training: ``That a man is able to deny himself his own desires, cross his own inclinations, and purely follow what reason directs as best, tho' the appetite lean the other way. This power is to be got and improved by custom, made easy and familiar by an early practice.'' And even more definitely disciplinary is the well-known `hardening process,' which he recommends in physical training: ``The first thing to be taken care of is that children be not too warmly clad or covered, winter or summer. The face, when we are born, is no less tender than any other part of the body. It is use alone hardens it, and makes it more able to endure the cold.'' He likewise advises that a boy's ``feet be washed every day in cold water,'' that he ``have his shoes so thin that they might leak and let in water,'' that he ``play in the wind and sun without a hat,'' and that ``his bed be hard.''

\textbf{Origin, Significance, and Influence of the Theory of Formal Discipline.} - This emphasis upon discipline in training of every sort - intellectual, moral, physical - has often caused Locke to be regarded as the first great exponent of the educational doctrine of `formal discipline.' That theory has been so widespread and important during the past two centuries as to require consideration here. During the Middle Ages and the early period of humanism Latin was not only of cultural, but of practical utilitarian value. It was the language of the Church and of diplomacy, and in it was locked up all the learning of the times. All guidance in science, literature, philosophy, and politics that received any consideration was couched in its terms. But with the decline of ecclesiastical influence, the development of vernacular languages, and the scientific awakening in the seventeenth century (see pp.~163 f.), this utilitarian argument for the study of Latin was largely swept away. Appeal was then made in behalf of the subject to the doctrine of `formal discipline,' which was supported by the `faculty' psychology of Aristotle. It was held that the study of Latin yields results out of all proportion to the effort expended, and gives a general power that may be applied in any direction. A similar claim was before long made for Greek and mathematics. Mathematics was declared to sharpen the `faculty of reason,' while the classic languages were believed to improve the `faculty of memory.' Consequently, it gradually came to be argued by formal disciplinarians that every one should take these all-important studies, regardless of his interest, ability, or purpose in life, since he would thus best prepare himself for any field of labor. All who proved unfitted for these particular subjects have, therefore, been supposed to be not qualified for the higher duties and responsibilities, and to be unworthy of consideration in higher education.

This doctrine of formal discipline has had a tremendous effect upon each stage of education in practically every country and during every period until recently. Even the scientists and advocates of a variety of other subjects, instead of arguing for content value and particular training, have made strenuous efforts to meet this argument by pointing out the formal discipline in their own studies (see pp.~404 f.). Excellent examples of the effect of this theory upon educational institutions are found in the formal classicism of the English grammar and public schools and universities and of the German gymnasiums. While in the United States a newer and more flexible society has enabled changes to be more readily made, as late as the last decade of the nineteenth century, Greek, Latin, and mathematics largely made up the staples in many high schools, colleges, and universities, and the husks of formal grammar were often defended in elementary education upon the score of formal discipline.

\textbf{Opposition to the Disciplinary Theory and More Recent Modification.} - At the beginning of the twentieth century, however, with the abandonment of the `faculty psychology' and the development of educational theory, a decided reaction from the doctrines of formal discipline began among psychologists and common sense educators. It is now almost universally conceded that specific, rather than general, power is developed by the various studies, and no student is held to be unworthy of education or impervious to culture, simply because he is not adapted to the classics or mathematics. In consequence, the content of studies, rather than the process of acquisition has come to be emphasized, the curriculum has everywhere been broadened, and the principle of the election of subjects largely recognized. It has, however, been felt within the last half dozen years that in reacting from the old theory of formal discipline, educators went too far. While it is still held that emphasis must be laid upon the specific character of mental training, there are some generalized powers and values to be obtained. It is realized that ``a general benefit can be derived from specific training in so far as the person trained has consciously wrought out in connection with the specific training a general concept of method, based upon the specific methods used in that training'' (F. A. Hodge). Thus a student who has once realized the value of close reasoning through mathematical demonstrations is likely to develop a general concept of method, and can hardly be satisfied any longer with slovenly thinking in other fields; and the fine discriminations discovered in the classical authors, the balanced judgment used in historical method, and the accuracy required in the study of the sciences, may well be abstracted and tend to furnish a generalized ideal for other lines of endeavor.

\textbf{Locke's Real Position on Formal Discipline.} - It would seem as if this modified form of general power were all that Locke had in mind. He definitely concedes that ``learning pages of Latin by heart, no more fits the memory for retention of anything else, than the graving of one sentence in lead makes it the more capable of retaining firmly any other characters.'' And while he holds that the method of reasoning in mathematics can be transferred `to other parts of knowledge,' he declares that men who are reasonable in some things are often very unreasonable in others, and ``men who may reason well in one sort of matters today may not do so at all a year hence.'' The generalized benefits that students may obtain from mathematics are simply that it ``would show them the necessity there is, in reasoning, to separate all distinct ideas, and see the habitudes that all those concerned in the present inquiry have to one another, and to lay by those which relate not to the proposition in hand and wholly to leave them out of the reckoning. This is that which in other subjects is absolutely requisite to just reasoning.'' Thus Locke appears to be rather in harmony with modern educational theory than a thorough-going advocate of formal discipline. At any rate, it should be recognized that he did not defend, but vigorously assailed, the grammatical and linguistic grind in the English public schools. His attitude toward formal discipline seems to have sprung from his desire to root out the traditional and false, rather than to support the narrow humanistic curricula of the times.

\textbf{Supplementary Reading}

Graves, \emph{During the Transition} (Macmillan, 1910), pp.~305-311; and \emph{Great Educators} (Macmillan, 1912), chap. VI; Monroe, \emph{Textbook} (Macmillan, 1905), chap. IX. For a more extended account of Locke, read his \emph{Thoughts} and \emph{Conduct,} and Fowler, T., \emph{John Locke} (Macmillan, 1901). The literature of formal discipline is most extensive and the subject is still under discussion; but a good summary of all written up to 1911 is furnished in Heck,

W. H., \emph{Mental Discipline and Educational Values} (John Lane, New York), and later articles can be found by consulting the index of \emph{The American Psychological Review.} In a doctoral dissertation (University of Virginia), \emph{John Locke and Formal Discipline,} Hodge, F. A., makes it clear that the common interpretation of Locke as a formal disciplinarian is unfair. The most typical of the earliest opposition to the disciplinary argument is probably found in Thorndike, E. L., \emph{Educational Psychology} (Teachers College, New York, 1910), chap. VIII; the sanest discussion of the possible transfer of ideals appears in Bagley, W. C, \emph{Educative Process} (Macmillan, 1905), chap. XIII; and the reaction to the reaction is best portrayed by Angell, Pillsbury, and Judd in \emph{Educational Review,} vol.~XXXVT, pp.~1-43. Lyans, C. K., in his article upon \emph{Formal Discipline (Pedagogical Seminary,} vol.~XXI, PP- 343\textasciitilde{}393) makes a most careful analysis of the interpretations of the defenders and opponents of the theory, and gives a verythorough discussion of transfers.

\hypertarget{education-in-the-american-colonies}{%
\chapter{17. Education in the American Colonies}\label{education-in-the-american-colonies}}

\textbf{Outline}

\begin{itemize}
\tightlist
\item
  The schools of the American colonies closely resembled those of the European countries from which the colonists came, and were influenced by the various religious conceptions of education that were current in each case. In general, where the Calvinistic attitude prevailed, the colonies attempted universal education, but where the Anglican communion dominated, the aristocratic ideal of education was in evidence.
\item
  Three types of colonial school organization appeared: (1) \emph{laissez \{aire} in Virginia; (2) `parochial' in New Netherlands; and (3) governmental activity in Massachusetts. The South generally followed the same plan as Virginia, and New York (after the English occupation) and Rhode Island also developed on this basis. The other Middle and New England colonies followed the parochial and governmental patterns respectively.
\end{itemize}

American Education a Development from European. - We have hitherto had little occasion to speak of American education, except by way of anticipating certain great waves of influence and important institutions that have come into America from Europe. But we have now reached the period when the New World began to be extensively colonized, and in the rest of our study educational practices in America will become increasingly distinctive and influential. The schools of America are the offspring of European institutions, and have their roots deep in the social soil of the lands from which the colonists came. While the universal, free, and secular schools of the United States are a natural accompaniment of its republican form of government, like the new democracy itself, this development of popular education was not reached at a bound. At first the American schools resembled the institutions of the Mother Country as closely as the frontier life would permit. The seventeenth century was, therefore, for American education distinctly a period of `transplantation of schools,' with little or no conscious change; and it is only toward the middle of the next century, as new social and political conditions were evolving and the days of the Revolution were approaching, that there are evident the gradual modification of European ideals and the differentiation of American schools toward an ideal of their own.

Conditions in Europe from Which American Education Sprang. - Hence, in order to understand American education in the colonial period, we must briefly consider the social and educational conditions in Europe during the early part of the seventeenth century, when the colonists began their migrations. The thirteen American colonies were started while the fierce agitations of the Reformation period were still at their height. The settlers, for the most part, were Protestants, and many of them had emigrated in order to establish institutions - political, ecclesiastical, educational - that would conform to their own ideals, and in all cases education in the New World was given a peculiar importance by the dominant religious interests and conflicts of the old. At this time in practically all the states of Europe, educational institutions were controlled and supported

Anglicans.

by\^{}the Church and religious orders, with the assistance of private benevolence; but a few schools everywhere, and especially in Teutonic countries, were maintained by pre-Reformation craft gilds, and so had a close connection with municipalities (see p.~92). Thus the American schools at first naturally adopted the religious conception of education and religious domination, but had some acquaintance with free schools and municipal .management.

In addition to these characteristics, the religious reformers, like Luther and Calvin, generally held to the Tendency

''° \emph{\textsuperscript{J}} toward uni

idea that a system of schools should be supported, or \textsuperscript{v}.ersai educa

tion among

at least established, by the state, and that all children Caivinists, but should have an opportunity to secure an education ideals among sufficient to make them familiar with the Scriptures. If people were to be guided by the word of God, they must all be able to read it. But this view of education was not held by those for whom, as in the English Church, the Reformation was not primarily a religious and theological, but rather an ecclesiastical and political revolt. In Holland and Scotland, for example, where Calvinism prevailed, universal education was upheld by the mass of the people, but in France and England only a small minority, the Huguenots and Puritans respectively, adopted this attitude. Hence it happens that, wherever in America the influence of Puritanism, the Dutch Reformed religion, Scotch Presbyterianism, or other forms of Calvinism was felt, the nucleus of public education appeared, while in the colonies where the Anglican communion was dominant, the aristocratic idea of education prevailed and training of the masses was neglected. However, even among the Calvinists, who held that

elementary education should be universal, and that the State as well as the Church should hold itself responsible for its being furnished, the logical solution of the problem was not perceived for scores of years. In the Calvinistic colonies it was not at first believed that education should be the same in character for all or that the State should bear the expense through taxation. This distinctively American interpretation of public education did develop later, but in the beginning even the most advanced colonies to some extent placed the financial responsibility upon the parent or guardian.

Colonial School Organization: The Aristocratic Type in Virginia. - As a result of these general traditions and characteristics, there would seem to have been three chief types of school organization in the colonies. These Spra \textsuperscript{ch!ef} were (1) the \emph{laissez faire} method, current in Virginia and the South; (2) the parochial organization of New Netherlands and the Middle Colonies in general; (3) the governmental activity in Massachusetts and most of the other New England colonies. We may profitably discuss these typical organizations in order. Turning first to the aristocratic colonies of the South, we may seiecUvf\textsuperscript{mi}edu- \textsuperscript{se}l\textsuperscript{ect} Virginia, the oldest of these provinces, as represented from sentative of the type. That colony constituted the first England. attempt of England at reproducing herself in the New World, and here are found an order of society, form of government, established church, and distinction between I classes, similar to those of the Mother Country. For some time there existed a sharp line of demarcation between the gentry, or landowning class, and the masses, which 1 included the landless, indentured servants, and other dependents. In education, the colonists had brought

\textbackslash{}

with them the idea of a classical higher and secondary training for the upper classes in the semi-monastic type of university and the (Latin) grammar school (see pp.~120 f.), and but little in the way of elementary education, except private `dame' schools and the catechetical training by the clergy. There were, in addition, the family `tutorial' education, both secondary and elementary, for the children of the wealthy, and evident attempts at perpetuating the old English industrial training through apprenticeship for orphans and children of the poor. But no such institution as a public elementary school was at first known. In consequence, the educa- \^{}\^{}j"TM! tional legislation in colonial Virginia is concerned mainly legislation. with (a) the organization of a college or university, (b) individual schools of secondary grade, and (c) apprenticeship education for the poor.

During the first quarter of a century most educational efforts in Virginia were in behalf of the foundation of found VcoUege an institution of higher learning, and were aided by the king, the Anglican bishops, and the London Company. By 1619 over £2000 and a grant of ten thousand acres of land had been obtained for a University at Henrico, but this rather indefinite plan was brought to a violent end by the Indian massacre of 1622, and the funds were diverted to a school in the Bahamas. An even more fruitless endeavor to found a college was made in 1624 by Sir Edwin Palmer upon an island in the Susquehanna. During this period also there was at least one abortive attempt to establish a school by collections and gifts, and during the second quarter century of the colony there were chartered a number of secondary schools, and secondary

\textsubscript{\textgreater{}} -\_" \emph{..'■} schools.

endowed with bequests of land, money, cows, horses,

slaves, or other property. These schools, however, were local, and resembled the endowed Latin schools of England, except that they may sometimes have been obliged by circumstances to include more or less elementary instruction. In 1660 there was also a renewed attempt to establish by subscriptions a college and ``free (secondary) school for the advance of learning, education of youth, supply of the ministry and promotion of piety.'' But none of the efforts at founding schools could have been very successful, for, a decade later, when interrogated as to what kind of education existed in the colonies, Governor Berkeley made his famous reply: ``The same course that is taken in England out of towns; every man according to his ability instructing his children. \ldots{} I thank God there are no free schools, and I hope we shall not have them these hundred years; for learning has brought disobedience and heresy and sects into the world.''

However, despite these biased remarks of the testy governor, by 1692 the constant efforts to obtain an institution of learning were finally rewarded. Through the management of the Reverend James Blair, D. D., the bishop's commissary in Virginia, a charter for the College of William and Mary, a gift of £2000 and of twenty thousand acres of land, and the right to certain colonial taxes were obtained from the king, and large donations were made by the planters and additional support provided by the assembly. In fact, the college was munificently endowed for the times, and it did a great work in training the greatest scholars, statesmen, judges, military officers, and other leaders during the struggle for independence. Moreover, 'free\^{}\_schools\protect\hypertarget{ch17.xmlux5cux23para.248.1.0.box.135.239.1262.955.q.60}{}{ now greatly increased in number and their courses were much improved. But education was throughout this Apprenticeship early period regarded as a special privilege, and the the poor, masses were mostly employed in making tobacco, and other manual pursuits. For the sons of these people the only educational legislation was that provided between 1643 \textsuperscript{an}d 1748 in various acts concerning the industrial training of the poor, apprentices, wards, and orphans. In keeping with English precedents, these children were taught a trade by the masters to whom they were indentured, or trained in the flax-house established by public funds at James City. Thus, by the middle of the eighteenth century a fair provision of secondary and higher education had been voluntarily made in various localities, but as yet no real interest in common schools had been shown by the responsible classes in Virginia. Education was there predominantly \emph{'}selective' in character.}

The Parochial Schools in New Netherlands. - A second type of colonial organization of education appears in the New Netherlands, as the country between the Delaware and Connecticut rivers was called during the period of Dutch control (1621-1674). In contrast to the \emph{laissez faire} attitude of Virginia, the foundation of schools wasjDarochial. Instead of the chance endowment of schools wherever the benefactors happened to be located, a school was founded in connection with every church. This arrangement grew out of the Cal- Calvinistic

■\^{}\^{}y,\_ \emph{\textsuperscript{a} "\textsubscript{m}} conception of

vinTsHc conception of universal education, which formed universal eduan essential part of the social traditions in Holland during Holland, the seventeenth century. Long before the Dutch came to America, the parochial school, as a means of preserv

ing the Reformed faith, had become an indispensable part of church organization. But the Dutch state also had concerned itself with the facilities for education. The Reformed Dutch Church was granted the right to examine teachers, enforce subscription to the creed, and, in the case of the elementary schools at least, largely determine the appointments, but the legal support and control of education were vested in the civil authorities. Hence there early arose in New Amsterdam and the villages of New Netherlands a parochial school system and a distribution of control between Church and State prayers\textsuperscript{15}\^{} Re- very similar to that in Holland. Besides the ordinary afweu asei?' elementary branches, these parochial schools of the New benches Netherlands taught the 'true\_princjples of Christian

taught. religion,' and the catechism and prayers of the Reformed

Church. Thus the Dutch school differed from those in the Anglican colonies of the South, which stressed secondary education, in being chiefly elementary, although some attempt at conducting a Latin or `\_grammar' (see p.~120) school was also made in New Amsterdam from 1652 on. However, after the English took permanent possession of New York (1674), the parochial school of the city was limited to the support of the Reformed Church, and, as a result of its long refusal to adopt the English language, its possible influence toward the realization of universal education was completely But, with Eng- lost. While the Dutch schools of the villages generally tion, replaced retained the joint control and support of the local court conization"* and church, with a constantly increasing domination of the former, as a whole the English occupation of New York would seem to have set public education back about a hundred years. At any rate, by the eighteenth\protect\hypertarget{ch17.xmlux5cux23para.250.1.0.box.101.249.1045.206.q.70}{}{ century colonial New York seems to have fallen into the same \emph{laissezfaire} support of education that prevailed in the Southern colonies. The policy of universal education by means of parochial schools no longer existed.}

Sectarian Organization of Schools in Pennsylvania. - As a colony, Pennsylvania developed a church school organization, similar to that of the New Netherlands, except that it was carried on in connection with a number of creeds, and that the municipality was seldom a coiirdinate factor. Pennsylvania was more heterogeneous More \textsuperscript{sect}? ?°\textsuperscript{d}

\emph{\textsuperscript{J}} ° the mumci

in population than New York, as the tolerant attitude P\^{}1y \textsuperscript{not},

coordinated.

of the Quaker government had attracted a large variety of German sects, Swedes, Dutch, English, Welsh, and Scotch and Irish Presbyterians, and each was devoted to its own denominational schools. Early in the eighteenth century all Protestant religious bodies were authorized by statute to conduct schools and to receive bequests and hold land for their support. Even before this the Friends had started the `Penn Charter School,' Friends, which, while itself a secondary school, soon established elementary schools as branches\_. throughout the city upon various arrangements. In keeping with the conclusions of various `Yearly Meetings' (1722, 1746, etc.), the Friends also provided elementary, and to some extent secondary, schools in close proximity to all meetinghouses throughout the colony. Similarly, the Lutheran Lutherans, congregations, for example, each set up a school alongside of the church as early as possible. Likewise the Mennonites included in their system the famous schools Mennonites, of Christopher Dock, who in 1750 produced the first elaborate educational treatise in America. There was also some attempt at `grammar' schools (see p.~120)

Broader attempts.

and others. or secondary education, especially in the case of the well-known Moravian institutions at Bethlehem, Nazareth, and Lititz, and the Presbyterian Log College at Neshaminy, which became the cradle of Princeton, Washington and Jefferson, Hampden-Sidney, and Union Colleges.

A somewhat broader spirit was manifest in the voluntary `neighborhood' schools of Western Pennsylvania and elsewhere, in the attempts at universal education of the Connecticut colonists in the Wyoming Valley, and in the `academy' (see p.~159) set up at Philadelphia through Franklin, to train public men and teachers, and fuse the various nations in a common citizenship. But, as a whole, parochial schools exerted the greatest influence in the colony of Pennsylvania.

Town Schools in Massachusetts. - The third type of colonial school organization appeared first in Massachusetts. As compared with the \emph{laissez faire} and the parochial methods, governmental activity here prevailed. Accordingly, Massachusetts may be said to have inaugurated the first real system of public education in America. The character of the schools in this colony developed from its peculiar form of society and government. It was democratic, concentrated, and homogeneous, as

geneous society compared with the cosmopolitan and sectarian social

produced gov- \ldots{} .

ernmentai ac- structure in the Middle colomes, or the class distinctions and scattered population of the South. While there were some servants and dependents in the Massachusetts Bay Colony and a distinction was made between `freemen' and others, there were at no time rival elements that were unable to combine. The settlements were not a mere confederation, but the blending of all elements

Democratic and homo

into a single organism, where the individuality of each was merged in a new social whole. This condition was a result of the radical ingrained religious conviction that every one was a child of God, capable of becoming a vital and useful member of society, and that the community was obligated to give him training to that end Jn\_lhe\_\_home\textsubscript{J\textbar{}}\_the church, and the school.

Out of this Calvihistic attitude sprang a spirit of cooperation and helpfulness, a general participation of all townsmen in local government, and the Massachusetts type of school organization. Common schools seem to have been supported in most towns from the first by voluntaJ3L4ir\_cojripulsory subscriptions, and before the close of the first quarter of a century there had been established by the colony at large an educational system in which every citizen had a working share. Because of I this inclusiveness and unity in matters theological, the/ schools, while religious and moral, could hardly be considered sectarian. The first educational act of the colony, passed in 1642, was similar to the old English appren- \textsuperscript{Acts} \textsuperscript{of} \textsuperscript{l6}4* ticeship law in its provision for industrial education, and, while it was broadened so as to include some literary elements and a rate to procure materials was established, no school is mentioned in it. But in 1647 each \textsuperscript{and} \textsuperscript{l647} town of fifty families was required, under a penalty of £5, tomamFain an elementary school (Fig. 22), and every one of a hundred families a (Latin) `grammar' (Fig. 23) school. These schools might be supported in part by tuition fees, as well as by the town rate, and the obligation seems to have still rested on the parents to see that the children did `resort' to the school, but the germs of the present common school system in the United States

appear in the educational activity of the legislature in colonial Massachusetts. The `grammar' schools were to prepare boys for Harvard College (Fig. 24), which had been founded in 1636.

Education in the Other Colonies. - In general, the

organization of education in the remaining nine colonies

can be classed under one of the three types, described

above, but there are various modifications and some

exceptions to be noted. The \emph{laissez faire} foundation

of schools and colleges during the colonial period, which

was evident in Virginia, seems to be characteristic of

the four other colonies of the South. But the problems

were in every case a little different, and in each there

in°MaTyiand° were variations in development. Maryland, for example,

while mainly following the same random foundation of

schools as Virginia, also seriously endeavored (1696)

to support schools in every county by a general colonial

?\textsuperscript{a}south\textsuperscript{Ch}°°\textsuperscript{ls} \textsuperscript{tax}- South Carolina likewise made an unsuccessful

Carolina. attempt (1 722) at establishing a county system ofschools,

and, a decade before, it undertook to subsidize a school

financed by "\textgreater{} \textsuperscript{e}\%\textsuperscript{cn} parish. Georgia, on the other hand, until the

parliament. Revolution, had its entire budget, including the items

for education, financed by the English parliament.

Democratic And North Carolina, through a large number of Irish

tendencies in '° °

North Caro- and Scotch Presbyterians, German Protestants, and

una.

other immigrants, mostly from Pennsylvama, after 1728 began to break away from the aristocratic policy.

Moreover, after the permanent occupation (1674) by the English, New York went over to the \emph{laissez faire} plan (see p.~194). And, although in the remaining `middle' colonies, New Jersey and Delaware, something was accomplished by the parochial schools of the

Fig. 24. - The buildings of Harvard College (founded 1636) erected in 1675, 1699, and 1720.

various sects, much of the school organization there was \textsuperscript{Ra}"\textsuperscript{doi}p \textsuperscript{or}:

'° gamzation in

\emph{laissez (aire.} Likewise, Rhode Island, dominated by a New York and

. , \textsubscript{r} t11 Rhode Island.

fanatical devotion to freedom in thought and speech, failed throughout colonial days to pass any general regulations on education, like those of Massachusetts, and followed more closely the random organization of schools in Virginia. But the other New England colonies, Connecticut and New Hampshire, when it separated from Massachusetts, tended to provide schools after the Massachusetts plan. The Hartford colony of Con- Governmental

• 111. activity in

necticut i n its statutes of 1650 copied almost \emph{verbatim} New England, the phraseology used by Massachusetts in the establishment of schools. It remains for later chapters to show how the practices suggested by this type of organization have eventually overcome those of the other two, for that did not come to pass until after the colonial period.

\textbf{Supplementary Reading}

Graves, \emph{History of Education in Modern Times} (Macmillan, 1913), chap, iv; Clews, Elsie W., affords primary source material in \emph{Educational Legislation and Administration of the Colonial Governments} (Columbia University, Department of Philosophy and Psychology, No.~6). The interpretation of educational organization in \emph{Colonial Schools} used in this chapter is furnished by Monroe and Kilpatrick in the Monroe \emph{Cyclopaedia of Education} (Macmillan, 1910-14). For conditions in the various colonies, consult Dexter, E. G., \emph{History of Education in the United Stales} (Macmillan, 1904), chaps. I-VI; Jackson, G. L., \emph{The Development of School Support in Colonial Massachusetts} (Columbia University, Teachers College Contributions, No.~25, 1909); Kilpatrick, W. H., \emph{The Dutch Schools of New Netherland} (Bulletin, U. S. Bureau of Education, 1912); McCrady, E., \emph{Education in South Carolina} (Collections of the Historical Society of South Carolina, vol.~IV); Smith,\protect\hypertarget{ch17.xmlux5cux23para.257.1.0.box.257.234.1051.404.q.50}{}{ C. L., \emph{History of Education in North Carolina} (U. S. Bureau of Education, Circular of Information, no. 2, 1894); Steiner, B. C, \emph{History of Education in Connecticut} (U. S. Bureau of Education, Circular of Information, no. 2, 1893) and \emph{History of Education in Maryland} (U. S. Bureau of Education, Circular of Information, no. 2, 1894), chaps. I-IV; Stockwell, T. B., \emph{History of Public Education in Rhode Island} (Providence Press Co., Providence, 1876), pp.~281-404; and Wickersham, J. P., \emph{History of Education in Pennsylvania} (Lancaster, Pennsylvania, 1886), chaps. I-XII.}

\hypertarget{growth-of-the-democratic-ideal-in-education}{%
\chapter{18. Growth of the Democratic Ideal in Education}\label{growth-of-the-democratic-ideal-in-education}}

\textbf{Outline}

\begin{itemize}
\tightlist
\item
  During the eighteenth century, there appeared the climax to the revolt against absolutism.
\item
  This movement was directed against repression of intellect in the first half of the century, and against repression of political rights in the second half. The former phase, through Voltaire, made reason the basis of society and education, but introduced the tyranny of an intellectual few; the latter, through Rousseau, promoted an emotionalism and `naturalism' that were in keeping with the sentiments of the times.
\item
  The early treatises of Rousseau advocated a complete return to nature, but his later works somewhat modified this attitude.
\end{itemize}

\textbf{The Revolt from Absolutism.} - The ideal of universality and of state control in the education of America and other countries was greatly assisted by the climax to the general revolt against absolutism and ecclesiasticism that appeared in the eighteenth century. During this period of time most strenuous efforts were made to interpret life from a more reasonable and natural point of view and to overthrow all customs and institutions that did not square with these tests. This century marked the climax of the rebellion against authority and against the enslavement of the individual that had been manifesting itself in one form or another from the close of the Middle Ages. One revival after another - the Renaissance, the Reformation, realism, Puritanism, Pietism - had burst forth only to fade away or harden into a new formalism and authoritative standard. Yet with each effort something was really accomplished for freedom and progress, and the way was paved for the seemingly abrupt break from tradition that appears to mark the period roughly included in the eighteenth century. At this point despotism and ecclesiasticism were becoming thoroughly intolerable, and the individual tended more and more to assert his right to be an end in himself. At times all institutional barriers were swept aside, and in the French Revolution destruction went to an extreme. The logical consequence of these movements would have been complete social disintegration, had not the nineteenth century happily made conscious efforts to justify the eighteenth, and bring out the positions that were only implied in the negations of the latter. Thus the revolutionary tendencies and destruction of absolutism in the eighteenth century led to evolutionary movements and the construction of democracy in the nineteenth.

\textbf{The Two Epochs in the Eighteenth Century.} - But this revolt of the eighteenth century from absolutism in politics, religion, and thought falls naturally into two parts. During the first half of the century the movement was directed against repression in theology and intellect, and during the second half against repression in politics and the rights of man. The former tendency appears in the rationalism and skepticism of such men as Voltaire and the `encyclopedists,' while the latter becomes evident chiefly in the emotionalism and `naturalism' of Rousseau. Although these aspects of the revolutionary movement somewhat overlapped each other and had certain features in common, they should be clearly distinguished. The one prepared the way for the other by seeking to destroy existing abuses, especially of the Church, by the application of reason, but it gave no ear to the claims of the masses, and sought merely to replace the traditionalism of the clergy and monarch with the tyranny of an intellectual few. In distinction to this fuie of `reason,' `naturalism' declared that the intellect could not always be trusted as the proper monitor, but that conduct could better be\textasciitilde{} guided by the emotions as the true\textasciitilde{} expression of nature. It opposed the control of intellectual aristocracy and demanded rights for the common man.

Voltaire and the Encyclopedists. - The rationalistic and scientific tendency was chiefly developed by Diderot, Voltaire, Condillac, D'Alembert, and others interested in the production of the French \emph{Encyclopedic} Of all these `encyclopedists' the most keen and brilliant was Voltaire (1694-1278), who may well serve as the type of the whole movement. With matchless wit and literary skill, in a remarkable range of poems, epistles, epigrams, and other writings, he championed reason against the championed

. * - \textasciitilde{} . reason against

traditional institutions of State and Church. His chief traditions, object of attack was the powerful Roman Catholic Church, which seemed to him to stand seriously in the way of all liberty, individuality, and progress, and the slogan with which he often closed his letters was, - ``crush the infamous thing.'' The Protestant beliefs he likewise condemned as hysterical and irrational. While an exile in England, as the result of a quarrel with a member of the nobility, he became acquainted with the work

of Newton, Harvey, Bacon, Locke, and others (see and undertook pp. \^{}4 f.), and undertook to transplant the English English sden- scientific movemenTtcTFrance,\textasciitilde{}and make it the basis of

tific movement. \emph{-r\textasciitilde{}\textasciitilde{},}; -. 7r\textasciitilde{}r\textsuperscript{-} , ,

a new regime m society, religion, and education.

The other rationalistic writers had similar doctrines and purposes, and, although details of their ideas are hardly worthy of consideration here, most of them produced treatises upon education. In these they freely criticised the traditional school systems, and proposed

New theories of \textsubscript{new} theories of organization, content, and method, which must later have assisted to demolish the existing theory and practice in France. Thus rationalism sought to dejtroy\_\_despotism and superstition, and to establish in their place freedom fn action, social JUgjjffjJ"id-teILgious toleration. But in casting away the old, it swung to the opposite extreme and often degenerated into skepticism, anarchy, and license. In their fight against

Degenerated despotic ecclesiasticism, the rationalists often failed to

into skepticism \emph{*'}

and license. distinguish it from Christianity, and they opposed the Church because it was irrational rather than because it was not sincere. They felt that it might have a mission with the masses who were too dull and uneducated to be able to reason. So while rationalism wielded a mighty weapon against the fettering of the human intellect, it cared little about improving the condition of thelower, classes, who were sunk in poverty and ignorance, and were universally oppressed.

Rousseau and His Times. - In opposition to this intellectualistic and rationalistic attitude, Jean Jacques Rousseau (1712-1778) developed his emotionalism and `naturalism.' The social and educational positions of this reformer find a ready explanation in his antecedents\protect\hypertarget{ch18.xmlux5cux23para.264.1.0.box.129.242.1047.36.q.60}{}{ and career. From his father he inherited a mercurial}

temperament, love of pleasure, and irresponsibility, and

from his mother a morbid and emotional disposition.

His tendency toward sentimentalism, idleness, and want fnd\^{}Tt1tf\textsuperscript{sm}

oLcontrol was also s!JengBfeheaT5y\textasciitilde{}tHe"indulgent aunt \textsuperscript{contro'}

that brought him up, and by low companions during his

trade apprenticeships in the city of Geneva. At sixteen

he ran away from the city, and spent several years in

vagrancy, menial service, and dissoluteness. A love of \emph{\textsuperscript{lwe}} \textsuperscript{of} \textsuperscript{nature}

nature was impressed upon him by the wonderful scenery

of the country in which he spent his boyhood and his

years of wandering. He also learned to sympathize with \textsubscript{w}kh\textsuperscript{P}poor!

the poor and oppressed, whose condition was at this

time forced upon his attention. He received some spo- Sporadic

. education.

radic instruction, but his education was inaccurate and unsystematic.

At twenty-nine Rousseau settled down in Paris, but his days of vagabondage had left an ineffaceable stamp upon him. His sensitiveness, impulsiveness, love of nature, and sympathy for the poor were ever afterward in evidence. These characteristics blended well with a \^{}th\textsuperscript{d}tachoate body of inchoate sentiments and vague longings of this \^{}e\^{}rlod °\textsuperscript{£} period. It was the day of Louis XV and royal absolutism, when affairs in the kingdom were controlled by a small clique of idle and extravagant courtiers. A most artificial system of conduct had grown up in society. Under this veneer the degraded peasants were ground down by taxation and forced to minister to the pleasure of a vicious leisure class. But against this oppression there had gradually arisen an undefined spirit of protest and a desire to return to the original beneficent state of nature from which it was felt that man had departed. Hence

it happened that Rousseau, emotional, uncontrolled, and half-trained, was destined to bring into consciousness and give voice to the revolutionary and naturalistic ideas and tendencies of the century.

Rousseau's Works. - In 1750 he first crystallized this spirit of the age and resultant of his own experience in a

His discourses, discourse on \emph{The Progress of the Arts and Sciences.} In this he declared with much fervor and conviction, though rather illogically, that the existing oppression and corruption of society were due to the advancement of civilization. Three years later he wrote his discourse on \emph{The Origin of Inequality among Men.} Here again he held that the physical and intellectual inequalities of nature which existecTin primitive society were scarcely noticeable, but that, with the growth of civilization, most oppressive distinctions arose. This point of view in a somewhat modified form he continued in his remarkable

\emph{New Edoise,} romance, \emph{The New Heloise,} published in 1759, and three years afterward in his influential essay on political ethics,

\emph{Social Contract,} k\textsubscript{nown} \textsubscript{as} the \emph{Social Contract,} and in that most revolu

\emph{tndEmUe.} tionary treatise on education, the \emph{Emile.} The \emph{New Heloise} commends as much of primitive conditions as the crystallized institutions of society will permit. In the \emph{Social Contract,} Rousseau also finds the ideal state, not in that of nature, but in a society managed by the people, where simplicity and natural wants control, and aristocracy and artificiality do not exist. But the work that has made the name of Rousseau famous is the \emph{Emile.} This, while an outgrowth of his naturalism, assumes the modified position of the later works, and undertakes to show how education might minimize the drawbacks of civilization and bring man as near to nature\protect\hypertarget{ch18.xmlux5cux23para.266.1.0.box.129.247.1049.196.q.60}{}{ as possible. But the educational influence of the \emph{Entile} has been so far-reaching that we must turn to another chapter to study the positions of Rousseau and the effects of naturalism in education.}

\textbf{Supplementary Reading}

Graves, \emph{During the Transition} (Macmillan, 1910), pp.~311-313; \emph{History of Education in Modern Times} (Macmillan, 1913), pp.~110; and \emph{Great Educators} (Macmillan, 1912), pp.~77-85; Monroe, \emph{Text-book} (Macmillan, 1905), pp.~533-542. See also Boyd, W., \emph{The Educational Theory of Rousseau} (Longmans, Green, 1911); Morley, J., \emph{Voltaire} and \emph{Rousseau} (Macmillan).

\hypertarget{naturalism-in-education}{%
\chapter{19. Naturalism in Education}\label{naturalism-in-education}}

\textbf{Outline}

\begin{itemize}
\tightlist
\item
  Rousseau attempts in the \emph{Emile} to outline a natural education from birth to manhood. The first book takes Emile from birth to five years of age, and deals with the training of physical activities; the second, from five to twelve, treats of body and sense training; the third, from twelve to fifteen, is concerned with intellectual education in the natural sciences; the fourth, from fifteen to twenty, outlines his social and moral development; and the fifth describes the parasitic training of the girl he is to marry.
\item
  The \emph{Emile} is often inconsistent, but brilliant and suggestive; and, while anti-social, the times demanded such a radical presentation. Through it Rousseau became the progenitor of the social, scientific, and psychological movements in education.
\item
  The first attempt to put the naturalism of Rousseau into actual practice was made by Basedow. He suggested that education should be practical in content and playful in method, and he produced texts on his system, and started a school known as the `Philanthropinum.' He planned a broad course, and taught languages through conversation, games, and drawing, and other subjects by natural methods. The Philanthropinum was at first successful, and this type of school grew rapidly, but it soon became a fad.
\end{itemize}

The Influence of Rousseau's Naturalism, - The influence of Rousseau's \emph{Emile} upon education in all its aspects has been tremendous. It is shown by the library of books since written to contradict, correct, or disseminate his doctrines. During the quarter of a century following the publication of the \emph{Emile,} probably more than twice as many books upon education were published as in the preceding three-quarters of a century. This epochmaking work forced a rich harvest of educational thinking for a century after its appearance, and has affected our ideas upon education from that day to this.

Naturalistic Basis of the \emph{Emile.} - In the \emph{Emile} Rousseau aims to replace the conventional and formal education of the day with a training that should be natural and spontaneous. Under the existing \emph{rtgime} it was customary for "boys and girls to be dressed like men and women of fashion (Fig. 25), and for education to be largely one of deportment and the dancing master. On the intellectual side, education was largely traditional and consisted chiefly of a training in Latin grammar, words, The substituand \emph{memoriter} work. Rousseau scathingly criticises these urai education

\ldots{} \ldots{} for the conven

practices, and applies his naturalistic principles to an tionai type in imaginary pupil named Emile ``from the moment of his \textsuperscript{V08ue}'' birth up to the time when, having become a mature man, he will no longer need any other guide than himself." He begins the work with a restatement of his basal principle that ``everything is good as it comes from the hands of the Author of Nature; but everything degenerates in the hands of man.'' After elaborating this, he shows that we are educated by ``three kinds of teachers - nature, man\textsubscript{2} and things, and since the cooperation of the three educations is necessary for their perfection, it is to the one over which we have no control (i. e., nature) that we must direct the other two.'' Education must, therefore, conform to nature.

The Five Books of the \emph{Emile.} - Now the natural objects, through which Emile is to be educated, remain the\protect\hypertarget{ch19.xmlux5cux23para.269.1.0.box.70.233.1261.94.q.50}{}{ pulses exTM" \textsuperscript{same}, \textsuperscript{Du}t Emile himself changes from time to time. In ammed \textsuperscript{and}\ldots{} so far, therefore, as he is to be the guide of how he is to}

trained at dif-''"

ferent periods: be educated in a natural environment, his impulses must be examined at different times in his life. Hence the work is divided into five parts, four of which deal with Emile's education in the stages of infancy, childhood, boyhood, and youth respectively, and the fifth jwith the training of the girl who is to become his wife. The characteristics of the different periods in the life of Emile are marked by the different kinds of things he desires.

In the first book, which takes him from birth to five physicai\textsuperscript{1C}ac- y\textsuperscript{ears} of age, his main desire is for physical activities, tivities. \textsubscript{an}d jj\textsubscript{e} should, therefore, be placed under simple, free,

and healthful conditions, which will enable him to make the most of these. He must be removed to the country, where he will be close to nature, and farthest from the contaminating influence of civilization. His growth and training must be as spontaneous as possible. He must have nothing to do with either medicine or doctors, ``unless his life is in evident danger; for then they\_can do nothing worse than kill him.'' His natural movements must not be restrained by caps, bands, or swaddling clothes, and he should be nursed by his own mother. He should likewise be used to baths of all sorts of temperature. In fact, the child should not be forced into any fixed ways whatsoever, since with Rousseau, habit is necessarily something contrary to impulse and so unnatural. ``The only habit,'' says he, ``which the child should be allowed to form is to contract no habit whatsoever.'' His playthings should be such simple products of nature as "branches with their fruits and flowers, or a\protect\hypertarget{ch19.xmlux5cux23para.270.1.0.box.138.224.1053.42.q.70}{}{ poppy-head in which the seeds are heard to rattle."}

Language thMissimple, plain, and hence natural, should

be used with him, and he should not He hurried beyond

nature in learning to talk. He should be restricted to a

few words that express real thoughts for him.

The education of Emile during infancy is thus to be

'\_negative Land purely physical. The aim is simply to

keep his instincts and impulses, which Rousseau holds

to be good by nature, free from vice, and to afford him

the natural activity he craves. Next, in the period of

childhood, between the years of five and twelve, which i\textsuperscript{n} childhood, T' . '\_ ., ,"; \textsuperscript{llmb} \textsuperscript{and} \textsuperscript{sense}

is treated in the second book, Emile desires most to development,

exercise his legs and arms, and to touch, to see, and in other ways to sense things. This, therefore, is the time for training his li\^{}absjand senses. ``As all that enters the human understanding comes there through the senses, the first reason of man is a sensuous reason. Our first teachers of philosophy are our feet, our hands, and our eyes. \ldots{} In order to learn to think, we must then exercise our limbs, our senses, and our organs, which are the instruments of our intelligence.'' To obtain this training, Emile is to wear short, loose, and scanty clothing, go bareheaded, and have the body inured to cold and heat, and be generally subjected to a `hardening process' similar to that recommended by Locke (see p.~181). He is to learn to swim, and practice long and high jumps, leaping walls, and scaling rocks. But, what is more important, his eyes and ears are also to be exercised through natural problems in weighing, measuring, and estimating masses, heights, and distances. Drawing and constructive geometry are to be taught him, to render him more capable of observ

no geography, history, or reading,

though moral training through `natural consequences.'

ing accurately. His ear is to be rendered sensitive to harmony by learning to sing.

This body and sense training should be\^{} the nearest approach\_JxL\_\&n\_ intellectual, training aLJthis.\_period. Rousseau condemns the usual unnatural practice of requiring pupils to learn so much before they have reached the proper years. In keeping with his `negative' education, he asks rhetorically: ``Shall I venture to state at this point the most important, the most useful, rule of all education? It is not to gain time, but to lose it.'' During his childhood Emile is not to study geography, history, or languages, upon which pedagogues ordinarily depend to exhibit the attainments of their pupils, although these understand nothing of what they have memorized. ``At the age of twelve, Emile will hardly know what a book is. But I shall be told it is very necessary that he know how to read. This I grant. It is necessary that he know how to read when reading is useful to him. Until then, it serves only to annoy him.''

Incidentally, however, in order to make Emile tolerable in society, for he cannot entirely escape it, he must be given the idea of property and some ideas about conduct. But this is simply because of practical necessity, and no moral education is to be given as such, for, ``until he reaches the age of reason, he can form no idea of moral beings or social relations.'' He is to learn through `natural consequences' until he arrives at the age for understanding moral precepts. If he breaks the furniture or the windows, let him suffer the consequences that arise from his act. Do not preach to him or punish him for lying, but afterward affect not to believe him even when\protect\hypertarget{ch19.xmlux5cux23para.272.1.0.box.139.243.1048.304.q.60}{}{ he has spoken the truth. If he carelessly digs up the sprouting melons of the gardener, in order to plant beans for himself, let the gardener in turn uproot the beans, and thus cause him to learn the sacredness of property. As far as this moral training is given, then, it is to be indirect and incidental.}

However, between twelve and fifteen, after the demands of the boy's physical activities and of his senses \emph{\}\textsuperscript{n}} boyhood,

\emph{\textsuperscript{J} \textsuperscript{r} \textsuperscript{J}} , intellectual

have somewhat abated, there comes "an interval when training

\ldots{}.. . , , . , . ,, through curi

his faculties and powers are greater than his desires, osity concernwhen he displays an insistent curiosity concerning nat- 'phenomena. ural phenomena and a constant appetite for rational knowledge. This period, which is dealt with in his third book, Rousseau declares to be intended by nature itself as thfiJimje for instruction. But as not much can be learned within three years, the boy is to study only those subjects which are useful and not incomprehensible and misleading, and so is limited to the natural sciences. Later in this third book, in order that Emile may informally learn the interdependence of men and may himself become economically independent, Rousseau adds industrial experience and the acquisition of cabinetmaking to his training. The most effective method of instruction, Rousseau holds, comes through appealing fo the curiosity and interest in investigation, which are so prominent in the boy at this time. He contrasts the current methods of teaching astronomy and geography by means of globes, maps, and other misleading representations, with the more natural plan of stimulating inquiry through observing the sun when rising and setting during the different seasons, and through problems concerning the topography of the neighborhood. Emile

is taught to appreciate the value of these subjects by being lost in the forest, and endeavoring to find a way out. He learns the elements of electricity through meeting with a juggler, who attracts an artificial duck by means of a concealed magnet. He similarly discovers through experience the effect of cold and heat upon solids and liquids, and so comes to understand the thermometer and other instruments. Hence Rousseau feels that all knowledge of real value may be acquired most, clearly and naturally without the use of rivalry or textbooks. But he finds an exception to this irrational method' in one book, \emph{Robinson Crusoe,} ``where all the natural needs of man are exhibited in a manner obvious to the mind of a child, and where the means of providing for these needs are successively developed with the same facility.''

The fourth book takes Emile from the age of fifteen in youth, sex to twenty. At this period the sex interests appear and

interests, as r- - - , \textasciitilde{}, • - . \emph{\^{}\textasciitilde{},\textasciitilde{},}

basis of should Jie\_4irQp£rly\_guided and trained, especially as

Sdai training, they are the basis of social and moral relationships. Emile's first passion calls him into relations with his 'species, and he must now learn to live with others. ``We ihave formed his body, his senses, and his intelligence; it / remains to give him a heart.'' He is to become moral, affectionate, and religious. Here again Rousseau insists that the training is not to be accomplished by the formal method of precepts, but in a natural way by bringing the youth into contact with his fellowmen and appealing to his emotions. Emile is to visit infirmaries, hospitals, and prisons, and witness concrete examples of wretchedness in all stages, although not so frequently as to become hardened. That this training may not render him cynical\protect\hypertarget{ch19.xmlux5cux23para.274.1.0.box.156.222.1047.475.q.60}{}{ or hypercritical, it should be corrected by the study of history, where one sees men simply as a spectator without feeling or passion. Further, in order to deliver Emile from vanity, so common during adolescence, he is to be exposed to flatterers, spendthrifts, and sharpers, and allowed to suffer the consequences. He may at this time also be guided in his conduct by the use of fables, for ``by censuring the wrongdoer under an unknown mask, we instruct without offending him.''}

Emile at length becomes a man, and a life companion The passive

. and parasitic

must be found for him. A search should be made for a education of suitable lady, but ``in order to find her, we must know her.'' Accordingly, the last book of the Emile deals with the model Sophie and the education of woman. It is the weakest part of Rousseau's work. He entirely misinterprets the nature of women, and does not allow them any individuality of their own, but considers them as simply supplementary to the nature of men. Like men, women should be given adequate bodily training, but rather for the sake of physical charms and of producing vigorous offspring than for their own development. Their instinctive love of pleasing through dress should be made of service by teaching them sewing, embroidery, lacework, and designing. They ought to be obedient and industrious, and they ought early to be brought under restraint. Girls should also be taught singing, dancing, and other accomplishments. They should be instructed dogmatically in religion, and in ethical matters they should be largely guided by public opinion. A woman may not learn philosophy, art, or science, but shf\^{}should stnHy m\^{}n. "She must learn to penetrate their feelings through their conversation, their actions,

I their looks, and their gestures, and know how to give / them the feelings which are pleasing to her, without even 'seeming to think of them."

Estimate of the \emph{Emile.} - Such was Rousseau's notion of the naturaHndiYidualistic education for a man and the passive and repressive training suitable for a woman, and of the happiness and prosperity that were bound to ensue. To make a fair estimate of the Emile and its Defects out- influence is not easy. It is necessary to put aside all of

weighed by '\emph{\textsuperscript{J} \textsuperscript{r}}

merits. one's prejudices against the weak and offensive person

ality of the author, and to forget the inconsistencies and contradictions of the work itself. The \emph{Emile} has always been accounted a work of great richness, power, and underlying wisdom, and each of its defects is more than balanced by a corresponding merit. Moreover, the most fundamental movements in modern educational progress - sociological, scientific, and psychological - may be said to have germinated through the \emph{Emile.} The Sociological Movements in Modern Education. - The most marked feature of the Rousselian education and the one most subject to criticism has been its ex

Reyoit from treme revolt against civilization and all social control. 'A state of nature is held to be the ideal condition, and all social relations are regarded as degenerate. The child is to be brought up in isolation by the laws of brute necessity and to have no social education until he is fifteen, when an impossible set of expedients for bringing him into touch with his fellows is devised. One should remember, however, that the times and the cause had\protect\hypertarget{ch19.xmlux5cux23para.276.1.0.box.151.236.1266.1392.q.60}{}{ he was able to emphasize the crying need of a natural development of man, and to tear down the effete traditions in educational organization, content, and methods. And many of the social movements in modern educational organization and content were made possible and even suggested by him, after having thus cleared the ground. He held that all members of society should be trained industrially so as to contribute to their own support and should be taught to be sympathetic and benevolent toward their fellows. Thus through him education has been more closely related to human welfare. The in- followed \textsc{Rou\^{} }dustrial work of Pestalozzi and Fellenberg, the moral aim socr\textsubscript{a}i\textsuperscript{stressed} of education held by Herbart, the `social participation' in \textsuperscript{activitiea}the practice of Froebel, and the present-day emphasis upon vocational education, moral instruction, and training of defectives and of other extreme variations, alike find some of their roots in the \emph{Emile.} In fact, the fallacy involved in Rousseau's isolated education is too palpable to mislead anyone, and those who have best caught his spirit and endeavored to develop his practice have in all cases most insistently stressed social activities in the training of children and striven to make education lead to a closer and more sympathetic cooperation in society. The Scientific Movement in Modern Education. - Moreover, since Rousseau repudiated all social traditions SSw'buteS and accepted nature as his only guide, he was absolutely P\textsuperscript{hasi}j?d ot\textgreater{}}

but extreme need of just so extreme a doctrine. Such radical in

doctnne . , - .

needed, dividualism alone could enable him to break the bondage

to the past. By means of paradoxes and exaggerations

\emph{\textsuperscript{c} \textsuperscript{J}} °' \emph{\textsuperscript{J}} servational

opposed to all book learning and exaggerated the value \textsuperscript{work}- of observation?"'He consequently neglected the past, and would have robbed the pupil of all the experience of his fellows and of those who had gone before. But he emphasized the use of natural objects in the curriculum and developed the details of nature study and observa-

tional work to an extent never previously undertaken. Partly as a result of this influence, schools and colleges have come to include in their course the study of physical forces, natural environment, plants, and animals. Therein Rousseau not only anticipates somewhat the nature study and geography of Pestalozzi, Basedow, Salzmann, and Ritter, but, in a way, foreshadows the arguments of Spencer and Huxley, and the modern scientific movement in education. The Psychological Movements in Modern Educative ufknowi-\textsuperscript{0}" \^{}\textsuperscript{on}' - A matter of even greater importance is Rousedge of chii- \textsubscript{se}au's belief that education should be in accordance with

dren, Rousseau

saw the need the natural interests of the child. Although his knowl

of studying

them. edge of children was defective, and his recommendations

were marred by unnatural breaks and filled with sentimentality, he saw the need of studying the child as the only basis for education. In the Preface to the \emph{Emile} he declares that ``the wisest among us are engrossed in what the adult needs to know and fail to consider what children are able to apprehend. We are always looking for the man in the child, without thinking of what he is before he becomes a man. This is the study to which I have devoted myself, to the end that, even though my whole method may be chimerical and false, the reader may still profit by my observation.'' As a result of such appeals, the child has become the center of discussion in modern training. Despite his limitations and prejudices, this unnatural and neglectful parent stated many details of child development with much force and clearness and gave an impetus to later reformers.

In this connection should especially be considered Rousseau's theory of stages of development. He makes\protect\hypertarget{ch19.xmlux5cux23para.278.1.0.box.138.233.1266.1232.q.60}{}{ a sharp division of the pupil's development into definite \^{}\^{}iMtur\^{} periods that seem but little connected with one another, \emph{\textsuperscript{ia}e-'} and prescribes a distinct education for each stage. This seems like a breach of the evolution of the individual, and the \emph{reductio ad absurdum} of such an atomic training is reached in his hope of rendering Emile warmhearted and pious, after keeping him in the meshes of self-interest and doubt until he is fifteen. But, as in the case of his attitude toward society, Rousseau takes an extreme view, and he has thereby shown that there are characteristic differences at different stages in the child's life, and that only as the proper activities are provided for each stage will it reach maturity or perfection. He may, therefore, be credited to a great degree with the increasing tendency to cease from forcing upon children a fixed method of thinking, feeling, and acting, and for the gradual disappearance of the old ideas that a task is of educational value according as it is distasteful, and that real education consists in overcoming meaningless difficulties. Curiosity and interest rather are to be used as motives for study, and Rousseau therein points the way for the Herbartians. It is likewise due to him primarily that we have recognized the need of physical activities Physical ac}

\emph{°} \textsuperscript{r}' tivities and

and sense training in the earlier development of the sense training. child as a foundation for its later growth and learning. To these recommendations may be traced much of the object teaching of Pestalozzianism and the motor expression of Froebelianism. Thus Rousseau made a large contribution to educational method by showing the value of motivation, of creating problems, and of utilizing the senses and activities of the child, and may be regarded as the father of the psychological movements in modern

education. He could not, however, have based his study of children and his advanced methods upon any real psychological foundation, for in his day the `faculty' psychology (see p.~182) absolutely prevailed. Instead of Sympathetic working out his methods from scientific principles, he

understanding \emph{\textsubscript{m}"} \textsubscript{\#} . .

of the child, obtained them, as did Pestalozzi afterwards, through his sympathetic understanding of the child and his ability to place himself in the child's situation and see the world through the eyes of the child.

The Spread of Rousseau's Doctrines. - Thus seeds of many modern developments in educational organization, method, and content, were sown by Rousseau, and he is

intellectual \textsubscript{seen} t\textsubscript{0} b\textsubscript{e} the intellectual progenitor of Pestalozzi,

progenitor of \textsc{Jt }o \emph{i}

modern re- Herbart, Froebel, Spencer, and many other modern

formers, but ' \ldots{}.. \ldots{} ,

influence upon reformers. But his principles did not take immediate mediate. hold on the schools themselves, although their influence

is manifest there as the nineteenth century advanced. In France they were apparent in the complaints and recommendations concerning schools in the lists of desired reforms \emph{(cahiers)} that were issued by the various towns, and afterward clearly formed a basis for much of the legislation concerning the universal, free, and secular organization of educational institutions. In England\^{}\_ since there was no national system of schools, little direct impression was made upon educational practice. But in America this revolutionary thought would seem to have had much to do with causing the unrest that gradually resulted in upsetting the aristocratic and formal training of the young and in secularizing and universalizing the public school system. The first definite attempt, however, to put into actual practice the naturalistic education of Rousseau occurred in Germany through

the writings of Basedow and the foundation of the fhrough\textsuperscript{te}Ba\textsuperscript{P}se'Philanthropinum,' and is of sufficient importance to \textsuperscript{dow}demand separate discussion.

Development of Basedow's Educational Reforms. - Johann Bernhard Basedow (1723-1790) was by nature the very person to be captivated by Rousseau's doc- \^{}-J\^{} \textsubscript{by} trines. He was talented, but erratic, unorthodox, tact- Rousseau's

doctrines.

less, and irregular in life. He had been prepared at the University of Leipzig for the Lutheran ministry, but proved too heretical, and, giving up this vocation, became a tutor in Holstein to a Herr von Quaalen's children. With these aristocratic pupils he first developed methods of teaching through conversation and play connected with surrounding objects. A few years after this, in 1763, Basedow fell under the spell of Rousseau's \emph{Emile,} which was most congenial to his methods of thinking and teaching, and turned all his energy toward educational reform. As in the case of Rousseau with education in France, he realized that German education of the day was sadly in need of just such an antidote as `nat- \^{}``jf''\^{}\^{} uralism' was calculated to furnish. The schoolrooms naturalism. were dismal and the work was unpleasant, physical training was neglected, and the discipline was severe. Children were regarded as adults in miniature (Fig. 25), and were so treated both in their dress and their education. The current schooling consisted largely of instruction in artificial deportment. The study of classics composed the entire intellectual curriculum, and the methods were purely grammatical. As a result, suggestions made by Basedow for educational improvement attained as great popularity as his advanced theological propositions had received abuse.

In 1768 by his \emph{Address on Schools and Studies, and their Influence on the Public Weal,} he called generally upon princes, governments, ecclesiastics, and others in power, to assist him financially in certain definite educational reforms. In addition to suggesting that the schools be made nonsectarian and that public instruction be placed under a National Council of Education, he proposed that, in contrast to the formal and unattractive training of the day, education should be rendered practical in content \emph{\textsc{Am\^{}s }}°and\textsuperscript{S} \textsuperscript{an}d playful\textsuperscript{m} method. To assist this reform, he planned hU\textsuperscript{d}t\textsuperscript{U}xt'book? \textsuperscript{to} b\textsuperscript{rul}S \textsuperscript{out} \textsuperscript{a} work on elementary education, which he described in outline. Great interest in his proposals was shown throughout Europe by sovereigns, nobles, prominent men, and others desiring a nonsectarian and more effective education, and a subsidy of some ten thousand dollars was speedily raised, to enable him to perfect his plans. Six years later, Basedow completed his promised text-book, \emph{Elementarwerk,} and the companion work for teachers and parents known as \emph{Methodenbuch.} The \emph{Elementarwerk} was accompanied by a volume containing ninety-six plates, which illustrated the subject-matter of the text, but were too large to be bound in with it. While in these manuals Basedow included many naturalistic ideas from Rousseau, he also embodied features from other reformers and even additions of his own. \emph{Elementarwerk} Text-books and Other Works. - The \emph{Elementarwerk} clearly combines many of the principles of Comenius as well as of Rousseau. It has, in fact, been often called `the \emph{Orbis Pictus} (see p.~170) of the eighteenth century,' and \emph{Metkoden-} \textsuperscript{anc}* gl\textsuperscript{ves} \textsuperscript{a} knowledge of things and words in the form \emph{\textsuperscript{buch}-} of a dialogue. The \emph{Methodenbuch,} while not following\protect\hypertarget{ch19.xmlux5cux23para.282.1.0.box.140.239.1250.1115.q.60}{}{ Rousseau completely, contains many ideas concerning natural training that are suggestive of him. In this study of the nature of children, the book makes some advance upon the Rousselian doctrine by finding that they are especially interested in motion and noise, although Basedow would have shocked Rousseau by being so much under the control of tradition as to suggest using these interests in the teaching of Latin. Later, Basedow, together with Campe, Salzmann, and others of his followers, also produced a series of popular story books especially adapted to the character, interests, and needs of children. These works are all largely filled with \^{}\textsubscript{k}\textsuperscript{u}\textsubscript{s}\textsuperscript{Ia}/\textsubscript{0}r\textsuperscript{tory} didactics, moralizing, religiosity, and scraps of scientific children. information. The best known of them is \emph{Robinson der Jungere} (Robinson Crusoe Junior), which was published by Campe in 1779. It seems to have been suggested by Rousseau's recommendation of \emph{Robinson Crusoe} as a text-book, and in turn a generation later it became the model for \emph{Der Schweizerische Robinson} (The Swiss Family Robinson) of Wyss, which has been so popular with children in America and elsewhere.}

Course and Methods of the Philanthropinum. - Eight years before this, however, Prince Leopold of Dessau had been induced to allow Basedow to found there a model school called the `Philanthropinum,' which should embody that reformer's ideas. Leopold granted him a generous salary, and three years later gave him an equip- Salary, equipment of buildings, grounds, and endowment. At first Basedow had but three assistants, but later the number was considerably increased. The staff then included \textsuperscript{teachers}several very able men, such as Campe, formerly .chaplain at Potsdam, and Salzmann, who had been a professor

at Erfurt. The underlying principle of the Philanthropinum was `everything according to nature.' The natural instincts and interests of the children were only to be directed and not altogether suppressed. They were to and pupils. jj\textsubscript{e} trained as children and not as adults, and the methods of learning were to be adapted to their stage of mentality. That all of the customary fashion and unnaturalness might be eliminated, the boys were plainly dressed and their hair cut short. cation\textsuperscript{rS}but\textsuperscript{edu} While universal education was believed in, and rich social distinc- \textsubscript{an}d p\textsubscript{00r} \textsubscript{a}lik\textsubscript{e} \textsubscript{were} to be trained, the traditional idea

tions. \emph{\textsuperscript{r}}

still obtained that the natural education of the one class was for social activity and leadership, and of the other for teaching. Consequently, the wealthy boys were to spend six hours in school and two in manual labor, while those from families of small means labored six hours and training\textsuperscript{*}' studied two. Every one, however, was taught handicrafts, - carpentry, turning, planing, and threshing, as suggested in the third book of the \emph{Entile,} and there were also physical exercises and games for all. On the intellectual side, while Latin was not neglected, considerable attention was paid to the vernacular and French. In keeping with the \emph{Elementarwerk,} Basedow planned a wide objective and practical course very similar to that suggested by Comenius. It was to give some account of man, including bits of anthropology, anatomy, and physiology; of brute creation, especially the uses of domestic animals and their relation to industry; of trees and plants with their growth, culture, and products; of minerals and chemicals; of mathematical and physical instruments; and of trades, history, and commerce. He afterward admitted that he had overestimated the\protect\hypertarget{ch19.xmlux5cux23para.284.1.0.box.136.234.1049.97.q.70}{}{ amount of content that was possible for a child, and greatly abridged the material.}

and wide objective course.

The most striking characteristic of the school, however, was its recognition of child interests and the consequently improved methods. Languages were taught by Jf\textsubscript{u}\textsuperscript{n}\^{}\textsuperscript{a}b\textsuperscript{es}\textsubscript{con} speaking and then by reading, and grammar was not versation and brought in until late in the course. Facility in Latin was acquired through conversation, games, pictures, drawing, acting plays, and reading on practical and interesting subjects (Fig. 26). His instruction in arithmetic, geometry, Progressive

, , . methods in

geography, physics, nature study, and history was fully other subjecu. as progressive as that in languages, and, while continuing Rousseau's suggestions, seems to anticipate much of the `object teaching' of Pestalozzi. Arithmetic was taught by mental methods, geometry by drawing figures accurately and neatly, and geography by beginning with one's home and extending out into the neighborhood, the town, the country, and the continent.

Influence of the Philanthropinum. - The attendance at the Philanthropinum was very small in the beginning, since the institution was regarded as an experiment, but eventually the number of pupils rose to more than fifty. Most visitors were greatly pleased with the school, p\^{}\textsuperscript{1} \textsuperscript{e}*P\textsuperscript{ec}' especially on account of the interested and alert appearance of the pupils. Kant declared that it meant ``not a slow reform, but a quick revolution,'' although afterward he admitted that he had been too optimistic. While it may not have served well for older pupils, it was certainly excellent in its stimulus to children under ten or \textsuperscript{Stimulus} \textsuperscript{for}

\emph{\textsuperscript{J}} younger

twelve, who can be reached by appeals to physical \textsc{P\textsuperscript{u}p\textsuperscript{!'s}- }activities and the senses better than by books. Basedow, however, proved temperamentally unfit to

Salzmann,

Similar Jnsti- direct the institution. Joachim Heinrich Campe (1746Campe, 1818), who first succeeded him, withdrew within a year to

found a similar school at Hamburg. Institutions of the same type sprang up elsewhere, and some of them had a large influence upon education. The most striking and enduring of these schools was that established in 1784 by Christian Gotthilf Salzmann (1744-1811) at Schnepfenthal under the patronage of the royal family of SaxeGotha. The natural surroundings - mountains, valleys, lakes - were most favorable for the purpose of the institution, and much attention was given to nature study, `lessons on things,' organized excursions, gardening, agricultural work, and care of domestic animals. Manual training, gymnastics, sports, informal moral and religious culture, and other features that anticipated later developments in education also formed part of the course. During the decade before the establishment of Salzmann's school, institutions embodying many of Basedow's ideas were also opened at Rechahn and his other Brandenburg estates by Baron Eberhard von Rochow (1734-1805). His schools were simply intended to improve the peasantry in their methods of farming and living, but, when this step toward universal education proved extraordinarily successful, Rochow advocated the adoption of a complete national system of schools on a nonsectarian basis.

In 1793 the Philanthropinum at Dessau was closed

permanently. Its teachers were scattered through

Europe, and gave a great impulse to the new education.

Becomes a fad, Ai\textsuperscript{1} unfortunate result of this popularity was that the

pHshed\textsuperscript{C}some Philanthropinum became a fad, and schools with this

good. name were opened everywhere in Germany by educa

and Rochow.

tional mountebanks. These teachers prostituted the system to their own ends, degraded the profession into a mere trade, and became the subject of much satire and ridicule. Nevertheless, the philanthropinic movement seems not to have been without good results, especially when we consider the educational conditions and the pedagogy of the times. It introduced many new ideas concerning methods and industrial training into all parts of France and Switzerland, as well as Germany, and these were carefully worked out by such reformers as Pestalozzi, Froebel, and Herbart. In this way there were embodied in education the first positive results of Rousseau's `naturalism.'

\textbf{Supplementary Reading}

Graves, \emph{In Modern Times} (Macmillan, 1913), chap. II; and \emph{Great Educators} (Macmillan, 1912), chaps. VII and VIII; Monroe, \emph{Text-book} (Macmillan, 190s), chap. X; Parker, S. C, \emph{History of Modern Elementary Education} (Ginn, 1912), chaps. VIII-X. The \emph{Emile} (Translated by Payne; Appleton, 1895) should be read, and the \emph{Elementarwerk} (Wiegandt, Leipzig, 1909) should be examined. A judicial description of the life and work of Rousseau is that by Morley, J. (Macmillan), while Davidson, T., furnishes an interesting interpretation of \emph{Rousseau and Education from Nature} (Scribner, 1902), but the standard treatise on \emph{The Educational Theory of Rousseau} (Longmans, Green, 1911) at present has been written by Boyd, W. A good brief account of \emph{Basedow: His Educational Work and Principles} (Kellogg, New York, 1891) is afforded by Lang, O. H. See also Barnard, H., \emph{American Journal of Education,} vol.~V, pp.~487-520.

\hypertarget{philanthropy-in-education}{%
\chapter{20. Philanthropy in Education}\label{philanthropy-in-education}}

\textbf{Outline}

\begin{itemize}
\tightlist
\item
  In England, during the eighteenth century, there were numerous attempts to provide education for the poor through charityschools. The most important factor in maintaining these institutions was the Society for the Promotion of Christian Knowledge.
\item
  Among other organizations, there sprang up a Society for the Propagation of the Gospel in Foreign Parts, which supported schools throughout the American colonies, except Virginia. Charity schools were also maintained in America by various other agencies.
\item
  An attempt was likewise made by Raikes of Gloucester, England, to establish Sunday schools, for training the poor to read, and these institutions spread throughout the British Isles and America.
\item
  A system of instruction through monitors, developed by Lancaster and Bell, while formal and mechanical, furnished a sort of substitute for national education in England, and, spreading throughout the United States, paved the way for state support, and greatly improved the methods of teaching.
\item
  `Infant schools' for poor children also grew up during the nineteenth century in France, England, and the United States, and found a permanent place in the national systems, but they soon became formalized and mechanical.
\item
  Philanthropic education proved a first step toward universal and national education.
\end{itemize}

Reconstructive Tendencies of the Eighteenth Century. - The eighteenth century cannot be regarded altogether as a period of revolution and destruction. While\protect\hypertarget{ch20.xmlux5cux23para.290.1.0.box.130.254.1263.798.q.60}{}{ such a characterization describes the prevailing tendencies, there were also social and educational forces that looked to evolution and reform rather than to a complete disintegration of society and a return to primitive living. Even in Rousseau, the arch-destroyer \^{}T and\^{}hT of traditions, we found many evidences of a recon- j\textsubscript{n}7\textsubscript{s}\^{}`\textsuperscript{£?}throi\textgreater{}\^{}' struction along higher lines, and such a positive movement was decidedly obvious in Basedow, Salzmann, and other philanthropinists. But in England reforms were especially apparent. In the land of the Briton, progress is proverbially gradual, and sweeping victories and Waterloo defeats in affairs of society and education are alike unwonted. The French tendency to cut short the social and educational process and to substitute revolution for evolution is out of accord ?"£especially}

in England.

with the spirit across the English Channel.

The Rise of Charity Schools in England. - And yet conditions in England at this time might well have incited people to revolution. Wages were low, employment was irregular, and the laboring classes, who num- wretched conbered fully one-sixth of the population, were clad in boring dass. rags, lived in hovels, and often went hungry. Opportunities for elementary education were rare. The few schools that remained after the Reformation had largely lost their endowments or had been perverted into secondary institutions, and had suffered from incompetent and negligent masters and from the religious upheaval of the times. It was as a partial remedy for this situation, that, toward the close of the seventeenth century, there sprang up a succession of `charity schools,' in which children of the poor were not only taught, but boarded and sometimes provided with clothes, and the as remedy.\textsuperscript{00}'

Foundation,

management,

books,

teachers,

and course.

boys were prepared for apprenticeship and the girls for domestic service. Probably about one thousand schools upon this general philanthropic basis had been established in England and Wales by the middle of the eighteenth century. Most of these had received substantial endowment, but numbers of them were maintained by private subscriptions.

The Schools of the S. P. C. K. - A factor that was even more important in opening charity schools was the `Society for the Promotion of Christian Knowledge' (often abbreviated to S. P. C. K.). This society was founded in 1698 by Reverend Thomas Bray, D. D., and four other clergymen and philanthropists. As a rule, its schools were established, supported, and managed by local people, but the Society guaranteed their maintenance, and assisted them from its own treasury whenever a stringency in funds arose. The S. P. C. K. also inspected schools, and advised and encouraged the local managers, and furnished bibles, prayer books, and catechisms at the cheapest rates possible. It made stringent regulations of eligibility for its schoolmasters, requiring, in addition to the usual religious, moral, pedagogical, and age tests, that they be members of the Church of England and approved by the minister of the parish. Each master was expected to teach the children their catechism, and purge them of bad morals and manners, besides training them in reading, writing, and elementary arithmetic. The pupils were, moreover, clothed, boarded, and at times even lodged.

The number of charity schools of the S. P. C. K. grew by leaps and bounds, and by the close of the first decade there were eighty-eight within a radius of ten\protect\hypertarget{ch20.xmlux5cux23para.292.1.0.box.129.228.1263.1556.q.60}{}{ miles of London. The gifts made had amounted to al- Development, most ten thousand pounds, and nearly one thousand boys and over four hundred girls had been sent out as apprentices. And before the middle of the eighteenth century the total number of these charity schools in England and Wales reached nearly two thousand, with about fifty thousand boys and girls in attendance. This increase in facilities for the education of the poor was not kindly received by many in the upper classes, who often felt that ``there is no need for any learning at all for the meanest ranks of mankind: their business is to labour, not to think.'' But the charity schools had also XSScT \textsuperscript{and} many warm supporters, and Addison even believed that as a result of them there would be ``few in the next generation who will not at least be able to write and read, and have not an early tincture of religion.'' The benefactions for these institutions continued to increase for nearly half a century, but by the middle of the eighteenth century popular interest had waned. The sub- decadence, scriptions began to fall off, the system of inspection and the teaching became less effective, and the schools ceased to expand. Nevertheless, the S. P. C. K. had succeeded in impressing the Church of England with a sense of re- \textsuperscript{and} influence, sponsibility for the establishment of a national school system upon a religious basis. Its schools were largely continued throughout the eighteenth century, and in most instances after 1811 were absorbed by the new educational organization of the English Church, the so-called `National Society' (see p.~239).}

Other British Charity Schools. - These institutions of the Church of England society may be regarded as typical of British charity schools in general. There were, how\protect\hypertarget{ch20.xmlux5cux23para.293.1.0.box.78.236.1269.43.q.60}{}{Nonconformist ever, also a dozen well-known foundations by non}

schools

conformists, including the `Gravel Lane School' of Southwark, London, which was started over a decade before the S. P. C. K. was organized. And an interesting type of philanthropic institution known as `cir'Circulating culating schools' was founded in Wales. These schools

schools.' \ldots{}. , .. „ . .

simply aimed to teach pupils to read the Bible in Welsh, and when this had been accomplished in one neighborhood, the school was transferred to another. But a much more important organization was the offshoot of the S. P. C. K., that arose chiefly to carry on charity schools in the American colonies. This association, the 'Society for the Propagation of the Gospel in

fhe\textsuperscript{U}s\textsuperscript{da}p°G °\textsuperscript{f} \textsuperscript{Forei}g\textsuperscript{n} Parts,' (commonly known as S. P. G.), was

founded by Dr.~Bray three years after the parent

society, but no schools were established for several

years.

The Charity Schools of the S. P. G - The first school

s. p. \textsc{G. }»chooi of the S. P. G. was opened in New York City in 1700

in New York

City, - under William Huddleston, who had been conducting

a school of his own there. It was intended that the new school should follow the plan of the charity schools in England, but, while free tuition and free books were granted from the beginning, it was not until many years later that the means of clothing the children gratuitously was provided. Under different masters and with varying fortunes, the school was supported by the society until 1783, when the United States had finally cut loose from the Mother Country and started on a career of its own. Meanwhile Trinity Church had come more and more to take the initiative in the endowment and support of the school, and since the withdrawal of the\protect\hypertarget{ch20.xmlux5cux23para.294.1.0.box.132.220.1232.98.q.60}{}{ society from America the institution has been known as now `Trinity 'Trinity Church School.' School.'}

Schools of the same type were active throughout the colonies in the eighteenth century. We possess more or less complete accounts of these institutions in New York and all the other colonies, except Virginia, where \textsuperscript{0ther} \textsuperscript{colonies}they were not believed to be needed. Except for size and local peculiarities, all of them closely resembled the school in New York City. The attendance ranged \textsuperscript{Attendanc}e. from eighteen or twenty pupils to nearly four times that number. Girls were generally admitted, and occasionally equalled or exceeded the boys in number. As a rule, children of other denominations were received on the same terms as those of Church of England members, and at times nearly one-half the attendance was composed of dissenters, but often those outside the Church were given secondary consideration, or the catechism was so stressed by the school that the dissenting children were withdrawn and rival schools set up. The character of the course of study in these charity schools is further £ooS!' \textsuperscript{and} indicated by the books furnished by the society. In packets of various sizes it sent over horn-books, primers, spellers, writing-paper and ink-horns, catechisms, psalters, prayer books, testaments, and bibles. There is also some evidence that secondary instruction was carried on intermittently in the various centers by the missionaries or by the schoolmasters in conjunction with their elementary work.

Throughout its work in the American colonies the S. P. G. met with various forms of opposition. The dissenters, \^{}\^{}'p\textsuperscript{0}\^{}*\textsuperscript{0} Quakers, and others were often openly hostile through fear of the foundation of an established national church

Its devotion and generosity,

and influence upon universal education.

Organization,

course, and

similar to that of England, and both sides displayed considerable sectarianism and bigotry. After 1750 the opposition to the society increased in bitterness and became more general, owing to the feeling that its agents were supporting the king against the colonists. Yet its patronage of schools was most philanthropic and important for American education in the eighteenth century. While it insisted upon the interpretation of Christianity adopted by the Church of England, it stood first and foremost for the extension of religion and education to the virgin soil of America. It carried on its labors with devoted interest and showed great generosity in the maintenance of schools, and the support of schools in the colonies by the S. P. G. must have exerted some influence toward universal education.

Charity Schools among the Pennsylvania Germans. - During the eighteenth century the efforts of the S. P. G. were supplemented by the formation of minor associations and the establishment of other charity schools in various colonies. Perhaps the most noteworthy instance was the organization in 1753 of `A Society for Propagating the Knowledge of God among the Germans,' and the maintenance of schools among the sects of Pennsylvania. These schools were managed by a general colonial board of six trustees, who visited the schools annually and awarded prizes for English orations and attainments in civic and religious duties. The course of study included instruction in ``both the English and German languages; likewise in writing, keeping of common accounts, singing of psalms, and the true principles of the holy Protestant religion.'' Twenty-five schools were planned, but probably there were never more than half\protect\hypertarget{ch20.xmlux5cux23para.296.1.0.box.132.228.1247.38.q.70}{}{ that number. The schools lasted only about a decade, as disappearance}

of S P K \emph{C}

the Germans soon came to feel that this English school- schools." ing threatened their language, nationality, and institu- tions.

The ' Sunday School' Movement in Great Britain. - A variety of charity school, quite different from those already mentioned, sprang up toward the close of the century under the name of `Sunday Schools.' To overcome the prevailing ignorance, vice, and squalor in the manufacturing center of Gloucester, England, Robert Foundation, Raikes in 1780 set up a school in Sooty Alley for the instruction of children and adults in religion and the rudiments. Six months later he started a new school in Southgate street, and soon had other schools established. He paid his teachers a shilling each Sunday to train the children to read in the Bible, spell, and write. This charity education, meager as it was, was attacked by °PP°\textsuperscript{sltl0n}' many of the upper classes, and was often viewed with suspicion by the recipients themselves. Yet the new movement had warm supporters among the nobility and \^{}\^{}\textsuperscript{cy}' \textsuperscript{and} such reformers as Wesley, and the schools soon spread to London, and then throughout England, Wales, Ireland, Scotland, and the Channel Islands. A Sunday School Society was founded in 1785, and within a decade distributed nearly one hundred thousand spellers, twenty-five thousand testaments, and over five thousand bibles, and trained approximately sixty-five thousand pupils in one thousand schools.

The `Sunday School' Movement in the United States. - The Raikes system of Sunday instruction was also soon introduced in America. The first school was or- ,. , ,

Individual

ganized m 1786 by Bishop Asbury at the house of Thomas centers

Crenshaw in Hanover County, Virginia, and within a quarter of a century a number of schools arose in various ar\^{}gmanent cities. Before long, permanent associations were also started to promote Sunday instruction. `The First Day or Sunday School Society' was organized at Philadelphia in 1791, and during the first two decades of the nineteenth century a number of similar societies for secular instruction on Sunday were founded in New York, Boston, Philadelphia, and elsewhere. In 1823 these associations were all absorbed into a new and broader organization, known ever since as the `American Sunday School Union.' At the start it published suitable reading-books, and furnished primers, spellers, testaments, and hymn-books to needy Sunday schools at a reasonable rate.

Value of the Instruction in `Sunday Schools.' - Both in Great Britain and the United States, however, the Sunday schools gradually tended to abandon their prepare\^{}'the1 secular instruction and become purely religious. At varsai\textsuperscript{0}educa- \emph{\^{}\textsuperscript{e}} \textsuperscript{sanle} \emph{\^{}\textsuperscript{rae}} *i\textsuperscript{e} teachers came to serve without pay \emph{\textsuperscript{tion}-} and to instruct less efficiently. And the value of the

secular teaching was not large at the best, as the work was necessarily limited to a few hours once a week. Raikes and all others interested in these institutions recognized their inadequacy as a means of securing universal education, and regarded them merely as auxiliary to a more complete system of instruction. But while a makeshift and by no means a final solution for national education, they performed a notable service for the times, and helped point the way to universal education. The Schools of the Two Monitorial Societies. - While philanthropic education started largely in the eighteenth\protect\hypertarget{ch20.xmlux5cux23para.298.1.0.box.143.247.1249.1065.q.60}{}{ century, some of the schools continued well into the nineteenth. This was especially the case with the `monitorial' system, started at Southwark in 1798. This district of London was thronged with barefoot and unkempt children; and Lancaster, the founder of the school, Lancaster undertook to educate as many as he could. His schoolroom was soon filled with a hundred or more pupils. In order to teach them all, he used the older pupils as assistants. He taught the lesson first to these `monitors,' and they in turn imparted it to the others, who were divided into equal groups. Each monitor cared for a single group. The work was very successful from the first, but Lancaster, attempting to introduce schools of this kind throughout England, fell so recklessly into debt that an association had to be founded in 1808 to continue the work on a practical basis. Within half a dozen years Lancaster withdrew from the organization, but the association, under the name of the `British and Foreign'\textsuperscript{3} and Foreign Society,' continued to flourish and found \textsuperscript{Society}; new schools.}

So successful was the Lancasterian work that the Church of England, fearing its nonsectarian influence upon education, in 1811 organized `The National Society for Promoting the Education of the Poor in the Principles of the Established Church.' This long-named association was to conduct monitorial schools under the management of Doctor Andrew Bell, who had experi- National' mented with the system in India before Lancaster opened \textsuperscript{Soaet}yhis school. Although they had formed no part of Bell's original methods, the Anglican catechism and prayer book were now taught dogmatically in the schools founded by the National Society. Bell proved an admirable

director, and a healthy rivalry sprang up between the societies. Value of the Monitorial System in England. - The

th\^{}two\^{}ys" pl\textsuperscript{ans} of the two organizations were similar, but differed

\textsuperscript{tems}- somewhat in details. Both used monitors and taught

writing by means of a desk covered with sand, but the system of Lancaster was animated by broader motives and had many more devices for teaching. It also instituted company organization, drill, and precision, and developed a system of badges, offices, rewards, and punishments. Monitorial instruction, however, was not

original``\textsuperscript{6}''° original with either Lancaster or Bell. It had long been used by the Hindus and others, although the work of the two societies brought it into prominence. It over

and mechani- emphasized repetition and recitation mechanics, and consisted of a formal drill rather than a method of instruction. Yet the monitorial schools were productive of some

Afforded sub- achievements. Most of them afforded a fair education in

stitute for national educa- th\textsubscript{e} elementary school subjects and added some industrial and vocational training. They also did much to awaken the conscience of the English nation to the need of general education for the poor. The British and Foreign and the National Societies afforded a substitute, though a poor one, for national education in the days before England was willing to pay for general education, and they became the avenues through which such appropriations as the government did make were distributed. In 1833 the grant of £20,000, constituting the first government aid to elementary education, was equally divided between the two societies (see p.~388), and this method of administration was continued as the annual grant was\protect\hypertarget{ch20.xmlux5cux23para.300.1.0.box.144.248.1265.799.q.60}{}{ gradually increased, until the system of public education was established. Likewise, in 1839, £10,000 for normal instruction was voted to the societies, and was used by the British and Foreign for its Borough Road Train- Training coling College, and by the National for St.~Mark's Training College. These were followed by several other training institutions, established by each society through government aid. In 1870, when the `board,' or public elementary, schools were at length founded, the schools of the British and Foreign Society, with their nonsec- foreign\textsuperscript{and} tarian instruction, fused naturally with them; but the \^{}bed\textsuperscript{S} ``but institutions of the National Society, though transferred \^{}m\^{}them*'' to school boards in a few cases, have generally come to \emph{\textsuperscript{xlves}}constitute by themselves a national system on a voluntary basis.}

Results of the Monitorial System in the United States. - In the United States the monitorial system was introduced into New York City in 1806. The' Society for the Establishment of a Free School,' after investigating the best methods in other cities and countries, decided to try the system of Lancaster (see p.~260). The method was likewise introduced into the charity schools of Phila- ., ,. .

, Adoption by

delphia (see p.~261). The monitorial system then spread New York and rapidly through New York, Pennsylvania, Massachusetts, Connecticut, and other States. It is almost impossible to trace the exact extent of this organization in the United States, but before long it seems to have affected nearly all cities of any size as far south as Augusta (Georgia), and west as far as Cincinnati. There are still traces of its influence throughout this region, - in Hartford, New Haven, Albany, Washington, and Baltimore, as well as in the places already mentioned (Figs. 27, 28,

Introduced into high schools and academies.

Increased school facilities

and 29). In 1818 Lancaster himself was invited to America, and assisted in the monitorial schools of New York, Brooklyn, and Philadelphia. A dozen years later the system began to be introduced generally into the high schools and academies. Through the efforts of Dr.~John Griscom, who had been greatly pleased with the monitorial high school of Dr.~Pillans in Edinburgh, a similar institution was established in New York City in 1825, and the plan was soon adopted by a number of high schools in New York and neighboring states. Likewise, the state systems of academies in Maryland and in Indiana, which became high schools after the Civil War, were organized on this basis. For two decades the monitorial remained the prevailing method in secondary education. Training schools for teachers on the Lancasterian basis also became common.

In fact, the monitorial system was destined to perform a great service for American education. At the time of its introduction, public and free schools were generally lacking, outside of New England, and the facilities that existed were meager and available during but a small portion of the year. In all parts of the country illiteracy was almost universal among children of the poor. This want of school opportunities was rendered more serious by the rapid growth of American cities. `Free school societies,' like that in New York City, formed to relieve the situation, came to regard the system of Lancaster, because of its comparative inexpensiveness, as a godsend for their purpose. And when the people generally awoke to the crying need of public education, legislators also found monitorial schools the cheapest way out of the difficulty, and the provision made for these schools

Fig. 27. - A monitorial school, with three hundred pupils and but one teacher.

gradually opened the road to the ever increasing expenditures and taxation that had to come before satisfactory schools could be established. Moreover, the Lancasterian schools were not only economical, but most effective, when the educational conditions of the times are taken into consideration. Even in the cities, the oneroom and one-teacher school was the prevailing type, and grading was practically unknown. The whole organization and administration were shiftless and uneconomical, and a great improvement was brought organiMt£\textsuperscript{\textgreater{}}n\textsuperscript{ed} about by the carefully planned and detailed methods of \textsuperscript{and} \textsuperscript{methods}Lancaster. The schools were made over through his definite mechanics of instruction, centralized management, well-trained teachers, improved apparatus, discipline, hygiene, and other features.

But while the monitorial methods met a great educational emergency in the United States, they were clearly mechanical, inelastic, and without psychological foundation. Naturally their sway could not last long, and as °hra\textsuperscript{P}M?ua?enlarged material resources enabled the people to make jj\textsubscript{l}°``t\textsuperscript{'}\textsubscript{in}\^{}\textsuperscript{nti}'' greater appropriations for education, the obvious de- proved. fects of the monitorial system became more fully appreciated and brought about its abandonment. Before the middle of the century its work in America was ended, and it gave way to the more psychological conceptions of Pestalozzi and to those afterward formulated by Froebel and Herbart.

The `Infant Schools' in France. - Another form of philanthropic education that came to be very influential during the nineteenth century and has eventually been merged in several national systems is that of the so-called `infant schools.' The first recorded instance of these

development in Paris:

part of national system.

institutions occurred late in the eighteenth century through the attempt of a young Lutheran pastor named OberHn\textsuperscript{ing} \textsuperscript{with} Oberlin to give an informal training to the small children in all the villages of his rural charge in northeastern France. This type of training was copied in Paris as early as 1801, but did not amount to much until its revival through the influence of a similar development in England a quarter of a century later. It then rapidly expanded, and in 1833 was adopted as part of the French national system of education. In 1847 \textsuperscript{a} normal school was founded to prepare directresses and inspectors for these institutions, and in 1881 they became known as `maternal schools,' and the present type of curriculum was adopted. Besides reading and writing, these schools have always included informal exercises in the mother tongue, drawing, knowledge of common things, the elements of geography and natural history, manual and physical exercises, and singing.

The • Infant Schools' in England. - Quite independently, though over a generation later than Oberlin, Robert Owen opened his `infant school' in 1816 at New Lanark, Scotland. He was a philanthropic cottonspinner, and wished to give the young children of his operatives a careful moral, physical, and intellectual training. From the age of three they were taught in this school for two or three years whatever was useful and within their understanding, and this instruction was combined with much singing, dancing, amusement, and out-of-door exercise. They were not ``annoyed with books,'' but were taught about nature and common objects through maps, models, paintings, and familiar conversation, and their "curiosity was excited so as to

Owen at New Lanark;

ask questions concerning them." To afford this informal training, Owen secured a ``poor simple-hearted weaver, named James Buchanan, who at first could scarcely read, write, or spell,'' but who, by following the instructions of Owen literally, made a great success of the system. But when Buchanan, with the consent of Buchanan's

"" ' school in

Owen, had been transferred to London, to start a sim- London, ilar school for a group of peers and other distinguished philanthropists, his lack of intelligence reduced the training to a mere mechanical imitation of the procedure he had learned at New Lanark. Unfortunately, this London school became the model for Samuel Wilderspin, who was destined to become the leading exponent of infant schools. The schools of Wilderspin, while \textsubscript{became} \textsubscript{model} retaining some of the principles and devices of Owen, for wilder

\textsuperscript{0} * \emph{\^{}} spin, - formal

were much more formal and mechanical. He thought and me- too highly of `books, lessons, and apparatus,' and con- founded instruction with education. He overloaded' the child with verbal information, depending upon the memory rather than the understanding. Before the child was six, it was expected that he had been taught reading, the fundamental operations in arithmetic, the tables of money, weights, and measures, a knowledge of the qualities of common objects, the habits of dif- ferent animals, the elements of astronomy, botany, and zoology, and the chief facts of the New Testament. Even the games were stereotyped, and the religious teach- ing most formal.

Wilderspin's first school was opened at Spitalfields, \textsubscript{Spreadof} London, and soon attracted a horde of visitors. He schools; then began lecturing upon the subject throughout the United Kingdom, often demonstrating his methods with

classes of children he had taken along, and organized Sodety-\textsuperscript{Sch001} \textsuperscript{in}f\textsuperscript{ant} schools everywhere. In 1824 an `Infant School Society' was founded and through it several hundred schools were established. A dozen years later an organization for training infant school teachers, known as Coi\textsuperscript{m}niT\textsuperscript{d} `The Home and Colonial School Society,' was founded Society; at London by Reverend Charles Mayo, D. D., and

others. This society undertook to graft Pestalozzianism upon the infant school stock. While the combination resulted in some improvement of the infant schools, and real object teaching and sense training were more emphasized than they had been, the spirit of Pestalozzi was largely lost, and there was too much imitation of the formal instruction of older children, and there was an evident attempt to cultivate infant prodigies. Through these agencies infant schools spread rapidly in Great Britain, and were adopted as a regular part of the Part of public public system, when it was established in 1870 (p.~388). And four years later a marked advance was made through merging in them some of the methods and games of the kindergarten.

`infant Schools' in the United States. - Schools open to all younger children also sprang up in the United States during the first quarter of the nineteenth century. For many years they were nowhere regarded as an essential part of the public school system, and were managed separately, but about the middle of the century they were generally united. In 1818 Boston made its ma\textsuperscript{St}°schoois' **\textsuperscript{rst} \textsuperscript{a}PP\textsuperscript{ro}P\textsuperscript{r}i\textsuperscript{a}tion for ``primary schools, to provide instruction for children between four and seven years of age.'' These schools were divided into four grades, beginning with the study of the alphabet and closing\protect\hypertarget{ch20.xmlux5cux23para.308.1.0.box.122.232.1048.422.q.70}{}{ with reading in the New Testament. Besides reading, writing, and spelling, sewing and knitting were taught the girls. A formal course and the monitorial method were employed until about 1840, when the primary schools became largely inoculated with the informal procedure of Pestalozzi. The primary schools were for a long time under a separate committee, but in 1854 the management was fused in a general city board.}

New York started an `Infant School Society' in 1827. This organization opened two `infant schools' for poor children between three and six years of age. One of 'Primary de

partments' in

these schools was located in the basement of a Pres- New York, byterian Church and the other in that of a monitorial institution belonging to the Public School Society (see p.~261). The Pestalozzian methods used in these infant schools greatly commended themselves, and in 1830 the Public School Society added them as `primary departments' in all their buildings, but under separate management. A committee was appointed in 1832 to examine the Society's schools and suggest improvements. Upon the recommendation of two of this committee, who had inspected education in Boston, primary schools were established in rented rooms in sufficient numbers to be within easy reach for the young children. The subject-matter and methods were likewise made less formal.

In 1827 three `infant schools' were also founded in Philadelphia and other centers of Pennsylvania through infant schools'

-r. , \textsc{Tt }\textsc{T, }, \emph{r . ,} , , in Philadelphia

Roberts Vaux. By 1830 the number of infant schools in the state had risen to ten, with two to three thousand pupils. As the numbers would indicate, the schools were largely organized upon the Lancasterian plan.\protect\hypertarget{ch20.xmlux5cux23para.309.1.0.box.73.228.1267.691.q.60}{}{ Two years later a model infant school was started in Philadelphia, and in 1834 six others were organized. By 1837 there were thirty primary schools in Philadelphia centera\textsuperscript{her} alone. Several other cities started infant schools early. Hartford began them in 1827, and Baltimore in 1829. These institutions were in most cases fostered by the leading men of the community, and the ultimate service performed for American education by this form of philanthropy was considerable. Among other improvements, the infant schools developed a better type of schoolroom, secured separate rooms for different classes, intSoughTnfant traduced better methods and equipment, encouraged a schools. movement toward playgrounds, and brought women}

into the city schools of the United States.

The Importance of Philanthropic Education. - Many other types of charity school arose during the eighteenth century both in Great Britain and America, but the chief movements have been described, and sufficient has been said to indicate the important part in education played by philanthropy. The moral, religious, Purpose, and economic condition of the lower classes had been

sadly neglected, and by means of endowment, subscription, or organized societies, a series of attempts was made to relieve and elevate the masses through education. As a result, charity schools of many varieties and more or less permanent in character arose in all parts of the British Isles, the United States, and even location, France. In many instances the pupils were furnished

course, with lodging, board, and clothes. The curriculum in

these institutions was, of course, mostly elementary. It generally included reading, spelling, writing, and arithmetic, while a moral and religious training was given\protect\hypertarget{ch20.xmlux5cux23para.310.1.0.box.141.245.1242.637.q.60}{}{ through the Bible, catechism, prayer book, and psalms, and sometimes through attendance at church under supervision of the master. Frequently industrial or vocational subjects were taught, or the pupils apprenticed to a trade or to domestic service. The course was usually most formal both in matter and method, but occasionally in the later types drawing, geography, nature study, physical exercises, and games were added, and the more informal methods of Pestalozzi or Froebel \textsuperscript{and} methods. were partially employed. Sometimes the training was especially intended for and adapted to children under the usual school age.}

These efforts to improve social conditions by means of philanthropic education encountered various sorts of Various sorts

\textsuperscript{r} \textsuperscript{r} of opposition.

opposition. Often the upper classes held that the masses should be kept in their place, and feared that any education at all would make them discontented and cause an uprising. The poor themselves, in turn, were often suspicious of any schooling that tended to elevate them, and were unwilling to stamp themselves as paupers. Moreover, the sectarian color that sometimes appeared in the religious training not infrequently repelled people of other creeds or kept the schools from receiving their children.

However, this philanthropic education may, in general, be considered a fortunate movement, although its greatest service consisted in paving the way for better things. In contrast to the negative phase of `naturalism,' it f\textsubscript{°}*\textsuperscript{v}\^{}\textsubscript{t}j\textsuperscript{t}on\textsubscript{a}T\textsuperscript{ay} represented a positive factor in the educational activities \textsuperscript{and} public ed

\emph{\textsuperscript{r} \textsuperscript{r} \textsubscript{m}} ucation.

of the century. Instead of attempting to destroy existing society utterly, it sought rather to reform it, and when the work of destruction gave opportunity for new\protect\hypertarget{ch20.xmlux5cux23para.311.1.0.box.283.244.1050.421.q.60}{}{ ideals, it suggested and even furnished a reconstruction along higher lines. Hence philanthropy in education exercised an important influence in the direction of universal, national, and public training for citizenship. It was in many of its forms merged in such a system in several countries, and in succeeding chapters references to the S. P. C K., S. P. G., Sunday, monitorial, and infant schools will naturally appear.}

\textbf{Supplementary Reading}

Graves, \emph{In Modern Times} (Macmillan, 1913), chap. III; and \emph{Great Educators} (Macmillan, 1912), chap. XII; Parker, \emph{Modern Elementary Education} (Ginn, 1912), pp.~101-107. Allen, W. O. B. and McClure, E., have presented \emph{The History of the S. P. C. K.} (Christian Knowledge Society, London, 1001), and Pascoe, C. F., \emph{Two Hundred Years of the S. P. G.} (Christian Knowledge Society, London, 1898), while Kemp, W. W., gives a detailed history of \emph{The Support of Schools in Colonial New York by the S. P. G.} (Columbia University, Teachers College Contributions, no. 56, 1913), and Weber, S. E., of \emph{The Charity School Movement in Pennsylvania} (Doctoral dissertation, University of Pennsylvania). Harris, J., furnishes a good description of \emph{Robert Raikes; the Man and His Work} (Dutton, New York, 1899); Salmon, D., of \emph{Joseph Lancaster} (Longmans, Green, 1004); Meiklejohn, J. M. D., of \emph{An Old Educational Reformer, Dr.~Andrew Bell} (Bardeen, Syracuse); and Salmon, D., and Hindshaw, W., of \emph{Infant Schools, Their History and Theory} (Longmans, Green, 1904).

\hypertarget{the-period-of-transition-in-american-education}{%
\chapter{21. The Period of Transition in American Education}\label{the-period-of-transition-in-american-education}}

\textbf{Outline}

\begin{itemize}
\tightlist
\item
  Between the `transplantation' period and that of the purely American conception of education was a distinctive stage in American education, - the `period of transition.'
\item
  During this period Virginia and the other Southern states began to develop sentiment for universal education, and started permanent school funds and `permissive' laws for common schools.
\item
  In the state of New York, appropriations were made for elementary education, but the public system was not really extended to the secondary field; while in New York City the way for universal education was prepared by quasi-public societies. In Pennsylvania, school districts were established at Philadelphia and elsewhere, but not until 1834 was the state system of common schools started. New Jersey and Delaware were even slower in getting their systems started.
\item
  The generous support of colonial education in Massachusetts was followed by a decline, and the control of schools was transferred from the towns to the districts. Academies were subsidized by the state and took the place of the grammar schools. A similar decline took place in the schools of the other New England states, except Rhode Island, which for the first time began to develop schools at public expense.
\item
  In the new states erected out of the Northwest Territory during this period there was a prolonged struggle to introduce common schools among those who had come from states not yet committed to this ideal, and state systems of education began to appear toward the close of the first quarter of the nineteenth century.
\item
  Thus before the educational awakening spread through the land, a radical modification had taken place in the European institutions with which America began its education.
\end{itemize}

\textbf{Evolution of Public Education in the United States.} - We may now return to our discussion of education in America. It has already been seen (chap. XVII) that the organization of schools in the various colonies was largely the result of educational ideals and conditions in the Mother Country. At first the schools of America closely resembled those of the European countries from which the colonists came, and the seventeenth century in American education is largely a period of `transplantation.' But toward the middle of the eighteenth century, as new social and political conditions were evolving and the days of the Revolution were approaching, there were evident a gradual modification of European ideals and the differentiation of American schools toward a type of their own. America has long stood, in theory at least, for equality of opportunity, and this conception of society is apparent in its views of education. The distinguishing characteristic of the American schools has throughout been the attempt of a free people to educate themselves, and, through their elected representatives, the people of the various states have come, in harmony with the genius of American civilization,~to initiate, regulate, and control their own systems of education. While the purely American conception of education cannot be fully discerned until almost the middle of the nineteenth century, there can for three quarters of a century before be clearly distinguished `a period of transition' from the inherited ideals to those of America to-day. This intervening stage of evolution covers roughly the last quarter century of colonial life and the first half century of statehood. To it we must now direct our attention.

\textbf{Rise of the Common School in Virginia.} - By the opening of the period, as we noted (p.~193), Virginia had voluntarily made a fair provision for secondary and higher education in various localities, but as yet no real interest in common elementary schools had been shown by the responsible classes. The nearest approach to such institutions was found in the plantation `field school.' Organized by a group of neighbors, these schools were supported by tuition fees and were not dependent upon any authority other than the good sense of the parents and pupils. But by the close of the Revolution a desire for genuine public education began to appear. The leader in the movement was the great statesman, Thomas Jefferson. As early as 1770, he first introduced into the~legislature a scheme of universal education. His bill proposed to lay off all the counties into small districts five or six miles square, to be called `hundreds.' Each hundred was to establish at its own expense an elementary school, to which every citizen should be entitled to send his children free for three years, and for as much longer as he would pay. The leading pupil in each school was to be selected annually by a school visitor and sent to one of the twenty `grammar' (i. e. secondary) schools, which were to be erected in various parts of the state. After a trial of two years had been made of these boys, the leader in each grammar school was to be selected and given a complete secondary course of six years, and the rest dismissed. At the end of this six-year course, the lower half of the geniuses thus determined were to be~retained as teachers in the grammar schools, while the upper half were to be supported from the public treasury for three years at the College of William and Mary, which was to be greatly expanded in control and scope.

This comprehensive plan for a system of common schools was, in the face of most discouraging opposition, constantly adhered to by Jefferson, although he did not live to see universal education an accomplished fact. He did, however, stimulate some movements toward this end. In 1796 the legislature passed an ineffective law whereby the justices of each county were permitted to initiate a school system by taxation, and in 1810 a~`literary fund' was established for public education.~When, in 1816, this fund had been increased to a million~dollars, those in charge of it recommended to the legislature the establishment of ``a system of public education, including a university, to be called the University of~Virginia, and such additional colleges, academies, and schools as should diffuse the benefits of education through the Commonwealth.'' This revision of Jefferson's suggestion did not immediately result in any legal steps toward universal education, except the appropriation in 1818 of \$45,000 from the income of the literary fund to have the poor children of each county sent to a proper school, but it did bring about in 1820 the foundation of the University of Virginia and a generous grant for the erection of a set of buildings. In the same year the effectiveness of the `permissive' law for common schools of 1796 and of the appropriation act of 1818 was somewhat strengthened by the division of the counties into districts, among which the appropriation for education of the poor was distributed and managed by special commissioners.

While this law marked one more step in advance, it was hampered by several of the features that in various states continually delayed the establishment of common schools at public expense. In the first place, it was based on the conception of public education as poor relief, rather than universal training for citizenship. It was~often viewed with hostility or indifference by the wealthy, cation, who felt that they were paying for that from which they received no benefit, and with pride and scorn by the poor, who refused to be considered objects of charity. Moreover, the sum distributed (\$45,000) was totally inadequate for over one hundred thousand children, and every variety of school, private as well as public, was subsidized without distinction. The system lacked a strong central organization, and the commissioners, often appointed by the county judges from the classes most opposed to the arrangement, were notoriously inefficient. The teachers also were generally incompetent, as it was practically impossible to persuade college or academy graduates to undertake the instruction of the poor. Nevertheless, under this apology for a people's common school, the state went on for a score of years, and there was a steady growth in the literary fund, the appropriations, the length of the school term, and the number of pupils who were willing to take advantage of such opportunities as it afforded. State officials of wide vision, moreover, sought in every way to improve the teaching corps and the defective administration. While the great majority of the school children still attended the denominational, private, and `field' schools (see p.~253), this~system of subsidies was educating public opinion for something better. By the close of the first half century of statehood, while Virginia was not yet ready to establish a complete system of public education, we shall later (seepp.327f.) find that the ground had been prepared for the development of common schools that was spreading throughout the country.

\textbf{Similar Developments in the Other Southern States.} - This advance toward the common school in Virginia is typical of the South. The development in Maryland was very similar to that of Virginia. The state began to move slowly toward universal education by subsidizing the education of the poor (1816), and by the passage of a `permissive' law for common schools in the counties (1825). In South Carolina an annual appropriation for `free schools' was started in 18.11. A law was passed establishing a number of schools in each election district equal to that of its members in the legislature and providing \$300 for each school. But these schools were largely regarded as pauper institutions, and, because legislative representation was based upon property, the distribution of the appropriation was very inequitable, for the inland parts of the state, which most needed assistance, received least. Yet the amount of appropriation gradually increased, and sentiment for universal education steadily developed. Within the first half dozen years of statehood, Georgia began the provision of land endowment for schools, and the organization of a state system under the title of the `University of Georgia.' While the value of the land was too small to establish a genuine system of public education so soon, before the close of the transition period, a permanent school fund had been started, and sentiment for public education had begun to grow. North Carolina made even earlier~progress toward common schools. The constitution of 1776 provided for the establishment of schools, and, by 1817, at the request of the legislature, Judge Archibald D. Murphy, a statesman with broad educational traditions, even formulated an elaborate plan for a complete system of public schools. This scheme failed, because it proposed to `maintain,' as well as educate, the children of the poor. But the suggestions of the Murphy committee shortly brought about the establishment of a `literary,' or common school fund (1825), the income of which was to be used for the support of public schools.

In the case of the other Southern commonwealths, which were admitted after the union had been formed, there was similarly a very gradual growth of sentiment and afterward~for universal~education. In every state there appeared an alliance between far-sighted statesmen and educators and the great middle class of citizens for the purpose of establishing common schools for all white children, and the old ecclesiastical and exclusive idea of education was beginning to fade. By the close of the first half century of national existence, a public system~had not actually materialized in any of the states, but most of them had begun to create `literary funds,' subsidize schooling for the poor, and enact `permissive' laws for establishing public schools. Except in Virginia and South Carolina, provisions had been made for general administration in state, county, and district; and in North Carolina the organization of a complete common school system awaited only a first hint of the great educational awakening (1835-1860). Moreover, most of the larger cities - ~Baltimore, Charleston, Louisville, Nashville, Memphis, Mobile, New Orleans - had already organized a regular system of public schools, and all of the older commonwealths had made some attempt at supporting a state institution of higher learning, which was virtually the head of a public school system. The various denominations had begun to found colleges in some numbers, but even these institutions were not so strictly ecclesiastical as William and Mary started out to be, and assumed a wider function than merely training for the ministry, while the aristocratic and classical `grammar' schools had largely given way to the `academies' (Fig. 32), which were nonsectarian, democratic, and more comprehensive in their curriculum.

\textbf{Evolution of Public Education in New York.} - After the English took possession of New York, we have seen (p.~195) how that territory lapsed into the \emph{laissez jaire} support of education. The upper classes of society largely sought their education abroad or through tutors and the clergy, although in 1754 King's College (now Columbia University) was founded, and during the century a number of secondary schools were organized and granted gratuities by the legislature. The few elementary schools that existed were either private or maintained by some church or philanthropic society. As already shown (pp.~234 ff.), this was the period distinguished for the schools founded by the Society for the Propagation of the Gospel. At the close of the Revolution, however, the various elements of the population had been welded together in the common struggle, and a sentiment for public education began to prevail over vested interests and sectarian jealousies. A series of broad-minded governors - the Clintons, Lewis, Tompkins, and Marcy - constantly reminded the legislature of its duty to establish common schools. In 1787 a system of public education was theoretically organized under the management of a Board of Regents, with the title of `The University of the State of New York,' but it did not include elementary schools. Two years later in each township were set apart for the endowment of~common schools, and in 1795 it was enacted that the sum of \$50,000 for five years should be distributed for the encouragement of elementary education in counties where the towns would raise by taxation half as much as the amount of their share. This arrangement was not carried on beyond the five years, but in 1805 the proceeds from 500,000 acres of land were appropriated for a com- mon school fund, which was not to be used until the interest reached \$50,000 per annum.

In 1812 further organization was enacted whereby a state superintendent of common schools was to be appointed, and the county unit replaced by a more democratic town and district basis. But it had been supposed that the state fund would provide for the entire support of the schools, and there still remained an obstinate opposition to local taxes. The towns, however, were gradually persuaded to raise the amount required to secure their share of the state donation. Much progress was brought about through the first superintendent, Gideon Hawley, and while, after eight years of service, he was removed by political manipulation and the office combined with the secretaryship of state, each of his successors undertook to distinguish the educational side of his administration by some marked advance or~improvement in the common schools. But for a generation the academies and colleges remained under supervision of the regents, and, except for state appropriations to academies, no one undertook to extend the public system into secondary and higher education. Moreover, the professional training of teachers in the academies~was encouraged by the state, and thereby the organization of normal schools was delayed. Hence, while New York started the first system of public education adjusted to the political and social conditions of the new nation, and probably had the most effective schools of the times, not until the great period of common school development (1835-1860) were its people fully willing to contribute for a general school system, make it entirely free, or develop it consistently in all directions.

\textbf{New York City.} - Meanwhile, an interesting development of educational facilities was taking place in New York City. In 1805 the opportunities offered in the private, church, and charity schools were seen by certain of the most prominent citizens to be totally inadequate for a city of seventy-five thousand inhabitants, and a `Free School Society' was founded to provide for the boys who were not eligible for these schools. The president was De Witt Clinton, afterward governor, and in 1806 the first school was opened, from motives of economy, upon the monitorial basis (see p.~241). The state fund did not reach a sufficient amount to be available until 1815, but special gifts were made to the school society from time to time by the legislature, the city, and private individuals, and there was a rapid increase in the number of the society's schools during the first quarter of a century. In 1826 the legislature authorized the organization to charge a small tuition fee and change~its name to the `Public School Society.' While the fee~system was soon found to injure the efficiency of the~work, and was abolished within six years, the new title~persisted, as it did not suggest pauperism in the way the~old name had. In 1828 the society was allowed the~benefit of a small local tax. For quite a time the work~of the association was unhindered, but in 1820-1825 a~vigorous effort was made to obtain a share of the state~appropriation for the sectarian schools of the Bethel~Baptist Church. This move was finally defeated, but~the Roman Catholics made a more successful protest~fifteen years later by indicating that the society, while~nominally nonsectarian, was really Protestant. To~settle this dispute, the legislature in 1842 established a~city board of education, and after eleven years the institutions of the Public School Society were merged in this~city system. Thus was the way prepared for a public~school system in New York City, and this development~was typical of the training of educational sentiment~through quasi-public societies that took place in Buffalo,~Utica, Oswego, and several other cities.

\textbf{Development of Systems of Education in Pennsylvania and the Other Middle States.} - The rise of public systems in the other Middle states was also gradual. In Pennsylvania, the state system slowly arose through a prolonged stage of `poor schools.' The new constitution~(1790) of the state declared: ``The legislature shall, as~soon as conveniently may be, provide by law for the~establishment of schools throughout the State, in such manner that the poor may be taught gratis.'' Men of broad vision, like Franklin, Benjamin Rush, and Timothy~Pickering, had striven hard to have popular education introduced, but the general sentiment of the times could not reach beyond providing free education for the poor. Moreover, although this moderate constitutional provision was a compromise, it was not until some years later (1802, 1804, and 1809) that the legislature passed acts to make it effective. Even then public institutions to fulfill the legislation were not established, but it was arranged that the tuition of poor children should be paid for at public expense in private, church, and neighborhood schools, and the proceeds of the sixty thousand acres of land appropriated for `aiding public schools' went to subsidize private institutions. But the idea of common schools continued to develop, and governors and other prominent men constantly called attention to the need of universal education. Philadelphia was the first municipality to be converted, and in 1818, under a special act of the legislature, it became `the first school district of Pennsylvania,' with the power to provide a system of education on the Lancasterian plan at public expense. After three or four years this special legislation was extended to five more `districts', and in 1824 a general law permitting the establishment of free schools in any community was enacted, though soon repealed.

Finally, in 1828, `the Pennsylvania Society for the~Promotion of Common Schools,' after demonstrating the~ineffectiveness of the `pauper school' law in a series school fund of memorials, succeeded in having a state school fund established, and in 1834, ``an act to establish a general system of education by common schools'' was passed. This law established a state system of schools under the general superintendency of the secretary of state. For this~system it appropriated \$75,000 per annum from the income of the state school fund, and permitted the wards, townships, and boroughs, which it constituted school districts, to share in this, provided they levied local taxes for schools. The Northern counties, settled mostly by New England colonists, and the Western portion of the state, with its large element of Scotch-Irish Presbyterians, ardently favored this encouragement of universal education, but the law was only `permissive' and was bitterly opposed by the Quaker and German inhabitants of `old' Pennsylvania, who feared that their own parochial schools would be replaced. The wealthy classes were also hostile to the new law, on the ground that they ought not to be taxed to educate other people's children. In a vigorous campaign to repeal the act, however, the opponents of the law, largely through the eloquent speech of Thaddeus Stevens, were defeated the following year (1835), and the desire to establish public schools was greatly increased in 1836 by the passage of a new law, which enlarged the annual appropriation to \$200,000, in which the school districts might participate only on condition of local taxation. Even then not more than one-half the districts took advantage of the opportunity, and it was several years before most of them claimed their share. Hence, while the battle was won by 1835, the consummation of public education in Pennsylvania did not take place until the great awakening of common schools had swept over the country. After the formation of the Union, New Jersey and~Delaware met with the same kinds of hindrance to the development of common schools as did Pennsylvania,~and they were even slower in getting a system established.

In both commonwealths a state school fund was started early in the nineteenth century, but it was not distributed for about a dozen years, and then it was used mostly for the education of paupers in subsidized private schools. Some `permissive' legislation for the organization of school districts and commissioners and the establishment of public schools was also passed, but it accomplished little before the middle of the century.

\textbf{Decline of Education in Massachusetts.} - In Massachusetts, on the other hand, efforts for the provision of universal training degenerated during the eighteenth century. The generous support of public education that had been started in 1647 was followed by a period of decline for about a century and a half. The causes of this decadence of local interest in education were rather complicated. In the first place, the complete domination~of Calvinism gradually disintegrated and was replaced by~a toleration of several creeds. The non-Puritans, who were constantly increasing in numbers, were obliged by the law of 1638 to preserve an outward conformity to the Calvinistic regime under penalty of banishment, but by 1662 a compromise was granted, whereby persons not conforming in every respect might be admitted to all church privileges, except communion, and the persecution of Quakers, Baptists, and other sects was largely abandoned. In 1670 came the successful secession of the Old South Church from the original church of Boston, as the result of a quarrel concerning this very compromise, and within a decade the Baptists were permitted to build a meeting-house in Boston. By 1692 recognition had been largely granted to all Protestant beliefs, and to be a `freeman,' or voter on all colonial questions, it was no longer necessary to be a member of a Puritan church. While every town was still required to support by tax an orthodox pastor, by 1728 the Episcopalians, Quakers, and Baptists were permitted to pay their assessments to their own ministers, and the alliance of the State with a despotic Church, which had made possible the system of public education, was largely broken.

Moreover, there was a decided lowering of intellectual standards upon the part of the colonists. The hard~struggle to wring a living from an unpropitious soil, and the disturbances due to wars, Indian skirmishes, and the difficulties of pioneer life greatly lessened their feeling of need for a literary training. Another reason for the educational decline was the dispersion of the population in the towns. At first, because of possible attacks by the Indians, a law forbade any dwelling to be built more than half a mile from the church and school, and not infrequently the school was equipped with a watchtower (Fig. 22). But, as the best land near the center was more and more taken up, the towns spread out in various directions, and the intervening hills, streams, swamps, and poor roads, together with the fear of Indians and wild animals, greatly hindered those on the outskirts in reaching the church and school of the town. As a result of all these conditions, the towns, most of which had been eager to establish schools even before being compelled to do so, began to seek various methods of evading the school law without incurring the fine. The minister was at times made the nominal schoolmaster,~or a teacher was even employed during the session of the `General Court' (i. e., legislature) and discharged upon adjournment. Laws were enacted against these subterfuges, greater vigilance was exercised, and the fine was increased first to £10 (1671) and then to £20 (1683), with a progressive increase where the number of families ran over one hundred (1712). Thus the fine came to be almost sufficient to support a schoolmaster, and it was made more and more unprofitable for a town to disobey the law.

Under these circumstances it became advantageous to many citizens, especially those at the center of a town, to have the entire support of the school come through general taxation rather than partially by means of tuition fees. But the people in the more distant portions of the town refused to vote a rate from which they themselves obtained no profit. They demanded that, in return for their taxes, the public school should be~brought nearer to them. Probably they were influenced in this stand by the fact that private `dame' schools,~and possibly elementary schools, had for some time been opened in various parts of the town conveniently near their homes. Another factor that may have aided in suggesting this solution was the legal recognition of various remote settlements within the town, known as `parishes' or `districts,' through the grant of self-government, separate church organizations, and other privileges similar to those of the town as a whole, though on a smaller scale. At any rate, we find that in the early part of the eighteenth century, wherever a rate was adopted as the sole means of school support, it was agreed that, instead of holding the town school for twelve months in the center alone, opportunities should be offered for a fraction of that period in various portions of the town. Usually the compromise at first took the form of having one town master teach in different districts through the year, and the result was known as a `moving' school. This necessitated holding the school in~a number of isolated communities, and the temple of learning often came at first to be located in a private house, usually in the kitchen. And although, in time, another room was added to the farm house for the accommodation of the school, the institution has since then been known as a `kitchen school' (Fig. 30). But, by a later development, when separate schools under different masters or mistresses came to be taught at the same time, the town school was said to be'divided.' Then in the winter,~when the big boys were out of the fields and came to school, the session was held in the center of the town, and usually required the brawn of a man. But in summer, when only the younger children could attend, schools were held in various parts of the town and were taught chiefly by women (Fig. 31). The divisions of the town that thus came to be recognized were allowed more and more control of their schools until they practically became autonomous. Before the time of the Revolution `divided schools' were recognized as a regular institution, and, together with other customs that had grown up during the eighteenth century, they were given legal sanction and denominated `district schools' in the law of 1780. By 1800 the districts were not only allowed to manage their own share of the town taxes, but were authorized to make the levy themselves; in 1817 they were made corporations and empowered to hold property for educational purposes; and in 1827 they were granted the right to choose a committeeman, who should appoint the teacher and have control of the school property.

Thus the year 1827 ``marks the culmination of a process which had been going on for more than a century, - the high-water mark of modern democracy, and the low-water mark of the Massachusetts school system.'' The district system did in its earlier stages bind the families of a neighborhood into a corporation whose intent was the most vital of human needs, - education, and the people came to feel the necessity of supporting it by their own generous contributions. But in the course of time the districts became involved in private~and petty political interests, and had but little consideration for the public good. The choice of the committeeman, the site, and the teacher caused much unseemly wrangling, and as each received only what it paid in, the poor district obtained only a weak school and that for but a short term. The increasing expense of the district system had also made it impossible for any except the larger towns to support the old-time `grammar' school, and this part of the old school requirements had fallen into disuse before the close of the eighteenth century. To meet the needs of secondary education, the policy of endowing `academies' (Fig. 32) with public lands, in Maine had gradually grown up, and this custom was legalized in 1797. Seven academies, - four in Massachusetts proper and three in the province of Maine, had originally been endowed with a township apiece, and some fourteen more had been chartered by towns at an early date, and empowered by the state to hold educational funds. By the time of the educational awakening there were some fifty of these private secondary institutions subsidized by the state, although managed by a close corporation. The first public high school~had been established in Boston (1821), but this type of secondary school had not begun to have any influence as yet. Into such a decadence had the liberally supported system of public education fallen, before the rapid development in common schools began and the influence of Horace Mann and other reformers was felt.

Fig. 32. - The first `academy,' founded by Benjamin Franklin at Philadelphia in 1750, and later developed into the University of Pennsylvania.

(Fig. 41)

\textbf{Developments in the Other New England States.} - The development of common schools in Massachusetts may be considered typical of New England in general, except Rhode Island. Connecticut similarly degenerated into a district system, which was recognized by law in 1794, and was destined later to constitute one of the greatest problems during the period of educational development (see pp.313 and 320). Vermont likewise made provision for town and district schools, and eventually established a state school fund and school commissioners, but this legislation was soon repealed, and the schools of the state were in a parlous condition when the awakening found them. New Hampshire and Maine also present very similar features. In Rhode Island the voluntary organization of education continued throughout the eighteenth century. In 1800 a law permitting each town to maintain `one or more free schools' was passed, but no municipality availed itself of this permission, except Providence, and the act was repealed in 1803. The basal state law for common schools was not passed until 1828, when at length \$10,000 was appropriated, and each town was required to supplement its share by such an amount as should annually be fixed in town meeting.

\textbf{The Extension of Educational Organization to the Northwest.} - It is thus evident that by the close of the first half century of the republic, there was everywhere slowly growing up a sentiment for public education. The development of common schools had, however, been greatly hindered in the Southern states by the separation of classes in an aristocratic organization of society. Yet the superior class had shown no lack of educational interest in their own behalf and had through the facilities offered reared a group of intellectual leaders, some of whom, like the far-sighted Jefferson, had caught the vision of universal education. The great diversity of nationality and creed in the Middle states, on the other hand, had fostered sectarian jealousies and the traditional practice of the maintenance of its own school by each congregation. This had proved almost as disastrous to the rise of a system of public schools, although Pennsylvania, and even more New York, had well begun the establishment of a public system. In both sections of the country public education was at first viewed as a species of poor relief, and the wealthy were unable to see any justice in being required to educate the children of others. As a result, the young `paupers' at times had their tuition paid in private schools, and these institutions were not infrequently allowed to share in public funds. The New England states, however, as a result of the homogeneity of their citizens, had early adhered to a system of public schools for all, organized, supported, and supervised by the people. While the efficiency of their common schools was eventually crippled by the grant of autonomy to local districts and the arising of petty private and political interests, they had initiated this unique American product, - a public system for all, dependent upon local support and responsive to local wishes.

This growth of a `common schools consciousness' was destined, as the result of a great educational awakening, to increase rapidly during the second quarter of the nineteenth century in the Middle and Southern, as well as the New England, states. But before describing this development further, it is important to see the effect of the ideals of these three sections of the country when introduced into a new part of the United States by emigrants from the older commonwealths. The new domain~referred to was those large tracts of unsettled territory, belonging, according to claims more or less overlapping, to six or seven of the original states, and finally (1781), in settlement of these disputes, ceded to the federal government, with the understanding that the territory should be `formed into distinct republican States.' After much discussion and various acts of Congress for half a dozen years, the famous `Ordinance of 1787' was passed for the government of this `Northwest Territory.' An earlier act (1785) had divided the entire territory into townships, six miles square, after the New England system, and of the thirty-six sections into which each township was subdivided, section sixteen was reserved for the support of public schools. A special contract also started the practice of providing two townships for the establishment of a university in each state. These provisions were later extended to the vast territory purchased from France in 1803 and known as ' Louisiana,' and to all the other territory afterward annexed to the United States. This federal land endowment gave an additional stimulus to the establishment of public education in the four commonwealths - Ohio, Indiana, Illinois, and Michigan - that were admitted from the Northwest Territory~before 1840. But the final system of public education in these new states took form slowly for various reasons. The settlers were poor; incessant Indian wars, the wilderness, wretched roads, and lack of transportation facilities tended to repel immigrants and leave the country sparsely settled; the large tracts of school land were slow in acquiring value, and, to attract settlers, were often leased at nominal rates or sacrificed at a small price; and social distinctions and sectarian jealousies persisted among the immigrants. As a whole, immigration from the earlier commonwealths had followed parallels of latitude, and the northern parts of Ohio, Indiana, and Illinois were occupied mostly by people from New England and New York, and the southern by former inhabitants of Virginia, Kentucky, Tennessee, Louisiana, and other states where the public school system was not yet as well developed. In Michigan, however, because of its northerly location, the great influx throughout the state had come from New York, New England, and Northern Ohio.

Consequently, the history of public education in the first three of the new states seems to be in each case largely a record of a prolonged struggle to introduce common schools among those of the people who had come from states not yet committed to this ideal, but Michigan, whose inhabitants had migrated from states where public education was in vogue, showed the germs of a public system even before statehood was conferred. The history of the common schools in Ohio, Indiana, and Illinois is very similar in general outline. Each one started off by claiming two townships of land for a university and the sixteenth section for schools, and the state constitution committed it to equal school opportunities for all. But not until the close of the first quarter of the nineteenth century was a system of common schools, with the organization of districts, appointment of school officers, and local taxation provided by the legislature. Even then the acts were largely `permissive,' the tax was not exacted from anyone who objected, and for some time various laws allowed public funds to be paid to existing private schools for the tuition of the poor. The complete system with a state superintendent was first organized in Ohio by 1836, but a similar stage of development was not reached by the other two states until after the great wave of common school development (1835-1860) had passed over the country. Michigan,~on the other hand, as early as 1817 established a `catholepistemiad,' which was to include a university and a system of schools of all grades, and a dozen years later in its revision of the school laws provided for a department of Education at the university and a territorial superintendency of schools. While under this law of 1829 tuition fees were to be required, except from the poor, by the first state constitution in 1837 the school lands were taken over from the wasteful management of the towns, and a public school was required to be open for three months in every district. The state superintendency was also established, and before 1840 Michigan was well started with a complete system of common schools.

\textbf{Condition of the Common Schools Prior to the Awakening.} - Thus, while some of the New England states, New York, and Ohio possessed the only definitely organized systems of public education, the movement for common schools had made some progress in all sections~of the country even before the educational awakening spread through the land. A radical modification had~taken place in the European institutions with which education in the United States began. To meet the demands of the new environment, education had become more democratic and less religious and sectarian. Wealth had become much greater and material interests had met with a marked growth. The old aristocratic institutions had begun to disappear. Town and district schools had been taking the place of the old church, private, and `field' schools, and in some of the cities the foundation for public education was being laid by quasipublic societies or even through local taxation. The academies (Fig. 32) had replaced the `grammar' schools, and the colleges had lost their distinctly ecclesiastical character. State universities were starting in the South and Northwest. All these evidences of the growth of democracy, nonsectarianism, and popular training in education were destined to be greatly multiplied and spread before long. Such an awakening will be found to be characteristic of the great development of common schools that took place in the decades around the middle of the nineteenth century. But, before pursuing the subject further, we must direct our attention to some new reforms in method and content that were being introduced by Pestalozzi into education in Europe and were destined to produce a great stimulus in the public systems of the United States.

\textbf{Supplementary Reading}

Graves, \emph{In Modern Times} (Macmillan, 1913), chap. IV; Parker. \emph{Modern Elementary Education} (Ginn, 1912), chap. XII. A general, but not always accurate account of the period has been contributed by Mayo, A. D., to the Report of the U. S. Commissioner of Education, 1893-94, XVI; 1894-95, XXVIII; 1893-96, VI and VII; 1897-98, XI; and 1898-99, VIII. For the special states, see Adams, H. B., \emph{Thomas Jefferson and the University of Virginia} (United States Bureau of Education, Circular of Information, 1888, no. 1); Boone, R. G., \emph{History of Education in Indiana} (Appleton, 1892), chaps. I - III, and V-VII; Johnston, R. M., \emph{Early Educational Life in Middle Georgia} (Report of the U. S. Commissioner of Education, 1894-95, XVI, and 1895-96, VII); Martin, G. H., \emph{Evolution of the Massachusetts Public School System} (Appleton, 1894), lect. III; Palmer, A. E., \emph{The New York Public School} (Macmillan, 1905); Randall, S. S., \emph{History of the Common School System of the Slate of New York} (Ivison, Blakeman, Taylor, New York, 1871) Second Period; Smith, C. L., \emph{History of Education in North Carolina} (U. S. Bureau of Education, Circular of Information, no. 2,1888); Smith, W. L., \emph{Historical Sketch of Education in Michigan} (Lansing, 1881), pp.~1-7, 39-49, and 57-78; Steiner, B. C, \emph{History of Education in Connecticut} (U. S. Bureau, Circular of Information, no. 2, 1893), and \emph{History of Education in Maryland} (U. S. Bureau, Circular of Information, no. 2, 1894), chaps. II-IV; Stock well, T. B., \emph{History of Public Education in Rhode Island} (Providence Press Co., Providence, 1876), chaps. II-V; Updegraff, H., \emph{The Origin of the Moving School in Massachusetts} (Columbia University, Teachers College Contributions, no. 17, 1907), chaps. V-X; Wickersham, J. P., \emph{History of Education in Pennsylvania} (Lancaster, Pennsylvania, 1886), chaps. XIII-XVII.

\hypertarget{observation-and-industrial-training-in-education}{%
\chapter{22. Observation and Industrial Training in Education}\label{observation-and-industrial-training-in-education}}

\textbf{Outline}

\begin{itemize}
\tightlist
\item
  Pestalozzi was the first prominent educator to develop the negative naturalism of Rousseau into positive reforms.
\item
  He desired to elevate the peasantry about him, and, failing in other expedients, undertook to accomplish this through a combination of industrial and intellectual training at Neuhof. This training he continued at Stanz, and began the development of his observational methods. In his work at Burgdorf, he was forced to suspend his industrial training, but he further developed his `A B C of observation,' and at Yverdon the method reached its culmination.
\item
  Like Rousseau, Pestalozzi conceived of education as a natural development of innate powers, and he extended its application to all children. In his method he held that clear ideas could be formed only by means of sense perceptions, and he undertook to analyze each subject into its simplest elements and develop it by graded exercises.
\item
  While not original, practical, or scientific, Pestalozzi made education the remedy for corruption in society, and started the modern methods in the elementary studies. Pestalozzian schools and methods spread rapidly through Europe and the United States.
\item
  The attempt to combine industrial training with intellectual, which Pestalozzi had to give up, was continued by his friend, Fellenberg, in his institutions at Hofwyl. Similar training wis developed throughout Europe. In the United States it stimulated the `manual labor' movement, and was later utilized as a solution for racial and other peculiar problems in education.
\end{itemize}

\textbf{Pestalozzi as the Successor of Rousseau.} - Having outlined the various phases of philanthropic education and surveyed the development of the common school in America, we may now turn again to the more immediate development of the movements that found their roots in Rousseau. It has been noted how Rousseau's `naturalistic' doctrines logically pointed to a complete demolition of the artificial society and education of the times. A pause at this point would have led to anarchy. If civilization is not to disappear, social destruction must be followed by reconstruction. Of course the negative attitude of the \emph{Emile} was itself accompanied by considerable positive advance in its suggestions for a natural training, but this advice was often unpractical and extreme and its main emphasis was upon the destruction of existing education. Hence the happiest educational results of Rousseau's work came through Pestalozzi, who especially supplemented that reformer's work upon the constructive side. Pestalozzi became the first prominent educator to develop the negative and somewhat inconsistent `naturalism' of Rousseau into a more positive attempt to reform corrupt society by proper education and a new method of teaching.

\textbf{Pestalozzi's Philanthropic and Industrial Ideals.} - Johann Heinrich Pestalozzi was born at Zurich in 1746. After the death of his father, he was brought up almost altogether by his mother. Through her unselfishness and piety, and the example of his grandfather, pastor in a neighboring village, Pestalozzi was inspired to relieve and elevate the degraded peasantry about him. He first turned to the ministry as being the best way to accomplish this philanthropic purpose, and later took up the study of law, with the idea of defending the rights of his people, but he was not able to succeed in either profession. Then, in 1769, he undertook to demonstrate to the peasants the value of improved methods of agriculture. He took up a strip of waste land at Birr, which he called \emph{Neuhof} (`new farm'), but within five years this experiment also proved a lamentable failure. Meantime a son had been born to him, whom he had undertaken to rear upon the basis of the \emph{Emile}, and the results, recorded in a \emph{Father's Journal}, suggested new ideas and educational principles for the regeneration of the masses. He began to hold that education did not consist merely in books and knowledge, and that the children of the poor could, by proper training, be taught to earn their living and at the same time develop their intelligence and moral nature.

\textbf{His Industrial School at Neuhof and the \emph{Leonard and Gertrude.}} - Hence the failure of his agricultural venture afforded Pestalozzi the opportunity he craved to experiment with philanthropic and industrial education. Toward the end of 1774 he took into his home some twenty of the most needy children he could find. These he fed, clothed, and treated as his own. He gave the boys practical instruction in farming and gardening on small tracts, and had the girls trained in domestic duties and needlework. In bad weather both sexes gave their time to spinning and weaving cotton. They were also trained in the rudiments, but were practiced in conversing and in memorizing the Bible before learning to read and write. The scholastic instruction was given very largely while they were working, and, although Pestalozzi had not as yet learned to make any direct connection between the occupational and the formal elements, this first attempt at an industrial education made it evident that the two could be combined. Within a few months there was a striking improvement in the physique, minds, and morals of the children, as well as in the use of their hands. But Pestalozzi was so enthusiastic over the success of his experiment that he greatly increased the number of children, and by 1780 was reduced to bankruptcy.

Nevertheless, his wider purpose of social reform by means of education was not allowed to languish altogether, for a friend shortly persuaded him to publish his views. His first production, \emph{The Evening Hour of a Hermit}, embodied most of the educational principles he afterward made famous, but he was advised to put his thought into more popular form, and soon wrote his highly successful story of \emph{Leonard and Gertrude} (1781). This work, with subsequent additions, gives an account of the degraded social conditions in the Swiss village of `Bonnal' and the changes wrought in them by one simple peasant woman. `Gertrude' reforms her drunkard husband, educates her children, and causes the whole community to feel her influence and adopt her methods. When finally a wise schoolmaster comes to the village, he learns from Gertrude the proper conduct of the school and begs for her continued cooperation. Then the government becomes interested, studies the improvements that have taken place, and concludes that the whole country can be reformed in no better way than by imitating Bonnal.

\textbf{His School at Stanz and Beginning of His Observational Methods.} - In 1798 he was given an opportunity to carry on his philanthropic and industrial ideals in education through the orphan home and school at Stanz, of which he was put in charge. Here he found it impossible to obtain any assistants, books, and materials, but he felt that none of these conventional aids could be of service in the work he desired to do. Hence he sought to instruct the children rather by experience and observation than by abstract statements and words (Fig. 33). This was the real beginning of his teaching through `observation,' and, while at Stanz he further developed his correlation of intellectual with manual training, his observational methods were thereafter destined to be more stressed. Religion and morals, for example, were never taught by precepts, but through instances that arose in the lives of the children he showed them the value of self-control, charity, sympathy, and gratitude. In a similarly concrete way the pupils were instructed in number and language work by means of objects, and in geography and history by conversation rather than by books. While they did not learn their natural history primarily from nature, they were taught to corroborate what they had learned by their own observation. About this method he said: ``According to my experience, success depends upon whether what is taught to children commends itself to them as true through being closely connected with their own observation. As a general rule, I attached little importance to the study of words, even when explanations of the ideas they represented were given.''

In connection with his observational method, Pestalozzi at this time began his attempt to reduce all perception to its lowest terms, `the A B C of observation, as he afterward called it. It was while at Stanz, for example, that he first adopted his well-known plan of teaching children to read by means of exercises known as 'syllabaries.' These joined the five vowels in succession to the different consonants, - `ab, eb, ib, ob, ub,' and so on through all the consonants. From the phonetic nature of German spelling, he was able to make the exercises very simple, and thus to furnish a necessary practice in basal syllables. In a similar way he hoped to simplify all education to such an extent that schools would eventually become unnecessary, and that each mother would be able to teach her children and continue her own education at the same time.

\textbf{Continuation of His Methods at Burgdorf, and \emph{How Gertrude Teaches Her Children.}} - From these experiments and concrete methods that Pestalozzi started at Stanz gradually developed all his educational contributions. But before the close of a year he was forced by circumstances to remove to Burgdorf. Here, on account of the social position of many of his pupils, he had to suspend his experiment of combining industrial with intellectual training, although, as will later be seen, his special efforts in this direction were greatly enlarged and perpetuated by Fellenberg. He now devoted himself to his `A B C of observation,' and further worked out and graduated his `syllabaries.' Language exercises were also given his pupils by means of examining the number, form, position, and color of the designs, holes, and rents in the wall paper of the school, and expressing their observations in longer and longer sentences, which they repeated after him. For arithmetic he devised charts upon which were placed dots or lines concretely representing each unit up to one hundred. By means of this `table of units' (Fig. 34), the pupil obtained a clear idea of the meaning of the digits and the fundamental processes in arithmetic. The children were also taught the elements of geometry by drawing angles, lines, and curves, and the development of teaching history, geography, and natural history by this method of observation was likewise continued.

Despite a want of system and errors in carrying out his method, Pestalozzi seems to have produced remarkable results from the start. Pupils poured in; a number of progressive teachers came to assist him; many persons of prominence visited the school and made most favorable reports upon its methods; and during the following three years and a half the Pestalozzian views on education were systematically developed and applied. While at Burgdorf also, he undertook a detailed statement of his method by the publication of his \emph{How Gertrude Teaches Her Children} (1801). This work does not mention Gertrude, but consists of fifteen letters to his friend, Gessner. Like all of Pestalozzi's works, it is quite lack ing in both plan and proportion, and is filled with repetitions and digressions, but the following portion of the summary of its principles, made by a biographer of Pestalozzi, may serve to give an idea of his educational creed:

``1. Observation is the foundation of instruction.''

``2. Language must be connected with observation.''

``3. The time for learning is not the time for judgment and criticism.''

``4. In each branch, instruction must begin with the simplest elements, and proceed gradually by following the child's development; that is, by a series of steps which are psychologically connected.''

``5. A pause must be made at each stage of the instruction sufficiently long for the child to get the new matter thoroughly into his grasp and under his control.''

``6. Teaching must follow the path of development, and not that of dogmatic exposition.''

\textbf{The `Institute' at Yverdon and the Culmination of the Pestalozzian Methods.} - As a result of political changes, Pestalozzi was obliged in 1805 to transfer his school to Yverdon. The `institute' here sprang into fame almost immediately, and increased in numbers and prosperity for several years. Children were sent from great distances, and teachers and visitors thronged there to learn and apply the new principles at home. The work of the institute formed a continuation and culmination of the observational methods started at Stanz and Burgdorf. The simplification introduced through the `syllabaries' and `table of units' was further elaborated. A `table of fractions' was also devised for teaching that subject concretely. It consisted of a series of squares, which could be divided indefinitely and in different ways. Some of the squares were whole, while others were divided horizontally into two, three, or even ten equal parts. There was further developed a `table of fractions of fractions,' or compound fractions, in which the squares were divided, not only horizontally, but vertically, so that the method of reducing two fractions to the same denominator might be self-evident.

Further, in order to draw and write, the pupil was first taught the simple elements of form. Objects, such as sticks or pencils, were placed in different directions, and lines representing them were drawn on the board or slate until all elementary forms, straight or curved, were mastered. The pupils combined these elements, instead of copying models, and were encouraged to design symmetrical and graceful figures. This also paved the way for writing. The children wrote on their slates, beginning with the easiest letters and gradually forming words from them, but soon learned to write on paper with a pen. Writing was, however, taught in connection with reading, although begun somewhat later than that study. Constructive geometry was also learned through drawing. The pupils were taught to distinguish, first vertical, horizontal, oblique, and parallel lines; then they learned right, acute, and obtuse angles, different kinds of triangles, quadrilaterals, and other figures; and finally discovered at how many points a certain number of straight lines may be made to cut one another, and how many angles, triangles, and quadrilaterals can be formed. To make the matter concrete, the figures were often cut out of cardboard or made into models.

In nature study, geography, and history the concrete observational work was likewise continued. Trees, flowers, and birds were viewed, drawn, and discussed. The pupils began in geography by acquiring the points of the compass and relative positions, and from this knowledge observed and described some familiar place. The valley of the Buron near at hand was observed in detail and modeled upon long tables in clay brought from its sides. Then the pupils were shown the map for the first time and easily grasped the meaning of its symbols. His ideas on geography, however, were more fully worked out by the scientist, Karl Ritter, who had already been trained in principles similar to Pestalozzi's in Salzmann's school at Schnepfenthal (see p.~228). Instead of the ``arbitrary and unmethodical collection of all facts ascertained to exist throughout the earth,'' which constituted the old `encyclopaedic' type of geography, Ritter presented a work based on principles indicated by the title, - \emph{The Science of the Earth in Relation to Nature and the History of Man.} Similarly, Pestalozzi's method was applied to music by his friend, Nageli, a noted Swiss composer, who began with the simplest tone elements and then combined and developed these progressively into more complex and connected wholes.

\textbf{Pestalozzi's Educational Aim and Organization.} - From the beginning of his work, Pestalozzi held that ``all the beneficent powers of man are due to neither art nor chance, but to nature,'' and that education should follow ``the course laid down by nature.'' So in all his works he constantly returns to the analogy of the child's development with that of the natural growth of the plant or tree. He even holds that ``the whole tree is an uninterrupted chain of organic parts, the plan of which existed in its seed and root,'' and that ``man is similar to the tree.'' Consequently, he defines education as ``the natural, progressive, and harmonious development of all the powers and capacities of the human being.'' This belief in the observance of development from within is in keeping with the naturalism of Rousseau, but that reformer viewed it chiefly from the negative side, and failed to make his educational doctrine concrete and explicit and to apply it to the school. Pestalozzi further modified and extended the Rousselian doctrine by recommending its application to all children, whatever their circumstances and abilities. Where Rousseau evidently had only the young aristocrat in mind in the education of Emile, Pestalozzi held that poverty could be relieved and society reformed only through ridding each and every one of his degradation by means of mental and moral development. Accordingly, he was the stanch advocate of universal education.

\textbf{His General Method.} - Pestalozzi's general method of giving free play to this natural development of the powers of all and so for reforming social conditions was to train his pupils through `observation.' He felt that clear ideas could be formed only by means of careful sense perceptions, and he was thoroughly opposed to the mechanical memorizing with little understanding that was current in the schools of the day. His method in general consisted in analyzing each subject into its simplest elements, or `A B C,' and developing it by graded exercises based as far as possible upon the study of objects rather than words. Yet Pestalozzi felt that ``experiences must be clearly expressed in words, or otherwise there arises the same danger that characterizes the dominant word teaching, - that of attributing entirely erroneous ideas to words.'' Accordingly, as shown in the summary of \emph{How Gertrude Teaches Her Children} (see p.~282), in all instruction he would connect language with `observation.' The special applications of this general method that were worked out by him and his followers in the most common subjects of the curriculum have been described in detail in the account of his work at Stanz, Burgdorf, and Yverdon, and do not require repetition here.

\textbf{The Permanent Influence of Pestalozzi.} - It is easy to exaggerate the achievements of this almost sainted reformer of Switzerland. Pestalozzi's methods were neither very original nor well carried out. His chief merit lay in developing and making positive the suggestions offered by Rousseau, and in utilizing them in the work of the schools. Even in this he failed somewhat in practicality and consistency. Moreover, Pestalozzi was groping and never possessed full vision. He did not grasp definite educational principles in a scientific way, but, like Rousseau, obtained his ideas of teaching from sympathetic insight into the minds of children. His writings for the most part record his empirical efforts at an effective training, and are revelations of methods of teaching in the concrete rather than the abstract. His works are also poorly arranged and inaccurate, and there was little organization or order in his schools.

But all these deficiencies are of small import when compared with Pestalozzi's influence upon society and education. In the eighteenth century caste ruled through wealth and education, while the masses, who supported the owners of the land in idleness and luxury, were sunk in ignorance, poverty, and vice. The schools for the common people were exceedingly few, the content of education was largely limited by ecclesiastical authority, and the methods were traditional and verbal. The teachers generally had received little training, and were selected at random. Ordinarily the pay was wretched, no lodgings were provided for the teacher, and he had often to add domestic service to his duties, in order to secure food and clothing. In the midst of such conditions appeared this most famous of modern educators, who never ceased to work for the reformation of society. As Voltaire, Rousseau, and others had held that the panacea for the corrupt times was rationalism, atheism, deism, socialism, anarchy, or individualism, Pestalozzi found his remedy in education. Like Rousseau, he keenly felt the injustice, unnaturalness, and degradation of the existing society, but he was not content to stop with mere destruction and negations. He saw what education might do to purify social conditions and to elevate the people by intellectual, moral, and industrial training, and he longed to apply it universally and to develop methods in keeping with nature.

Pestalozzi's achievements contained the germ of modern pedagogy, as well as of educational reform. It was he that stimulated educational theorists, instead of accepting formal principles and traditional processes, to work out carefully and patiently the development of the child mind and to embody the results in practice. From him have come the prevailing reforms in the present teaching of language lessons, arithmetic, drawing, writing, reading, geography, elementary science, and music. In harmony with his improved methods, Pestalozzi also started a different type of discipline. His work made clear the new spirit in the school by which it has approached the atmosphere of the home. He found the proper relation of pupil and teacher to exist in sympathy and friendship, or, as he states it, in `love.' This attitude, which appears so fully in his kindly treatment of the poor children at Neuhof and Stanz (Fig. 33), constituted the greatest contrast to that of the brutal schools of the times, and introduced a new conception into education.

\textbf{The Spread of Pestalozzian Schools and Methods through Europe.} - The `observational' methods of Pestalozzi and institutions similar to his were soon spread by his assistants and others throughout Europe. Strange to say, as a result of their familiarity with his weaknesses and the conservatism resulting from isolation, the Swiss were, as a whole, rather slow to incorporate the Pestalozzian improvements. In Zurich, however, Zeller of Wurtemberg, who had visited Burgdorf and had helped conduct a Pestalozzian training school, was early invited to give three courses of lectures in aid of the establishment of a teachers' seminary based upon the principles of Pestalozzi. Kriisi, after leaving the institute at Yverdon, also founded a number of schools and carried Pestalozzianism into various parts of Switzerland. And other disciples eventually started or reorganized schools in various parts of Switzerland.

But the Pestalozzian reforms in method secured their best hold upon Germany. The innovations were most remarkable in Prussia, and the elementary education there has come to be referred to as the `Prussian-Pestalozzian school system.' By the opening of the nineteenth century Pestalozzianism began to find its way into that state. In 1801 the appeal of Pestalozzi for a public subscription in behalf of his project at Burgdorf was warmly supported. In 1802 Herbart's account of \emph{Pestalozzi's Idea of an A B C of Observation} (see p.~337) attracted much attention. A representative was sent from Prussia to Burgdorf to report upon the new system in 1803. Meanwhile the Pestalozzian missionaries were fast converting the land. Plamann, who had visited Burgdorf, in 1805 established a Pestalozzian school in Berlin, and published several books applying the new methods to language, geography, and natural history. Zeller lectured to large audiences at Konigsberg, and organized a Pestalozzian orphanage there. A similar institution for educating orphans was opened at Potsdam by Turck. In 1808, two of Pestalozzi's pupils, Nicolovius and Silvern, were made directors of public instruction in Prussia, and sent seventeen brilliant young men to Yverdon to study for three years. Upon their return these vigorous youthful educators zealously advanced the cause. The greatest impulse, however, was given the movement by the philosopher, Fichte, who was ardently supported by King Frederick William III, and even more by the noble queen, Louise. They held that only through these advanced educational principles could a restoration of the territory and prestige lost to Napoleon at Jena be effected.

A similar spirit animated the other states of Germany, and Bavaria, Detmold, and other states early undertook to introduce the new principles. Everywhere in Germany the greatest enthusiasm prevailed among teachers, state officials, and princes. Thus in place of the reading, singing, and memorizing of texts, songs, and catechism, under the direction of incompetent choristers and sextons, with unsanitary buildings and brutal punishment, all Germany has come to have in each village an institution for training real men and women. Each school is under the guidance of a devoted, humane, and trained teacher, and the methods in religion, reading, arithmetic, history, geography, and elementary science are vitalized and interesting.

In France the spread of Pestalozzianism was at first prevented by the military spirit of the time and by the apathy in education, and later, when the reaction occurred, the schools came under ecclesiastical control and had little influence upon the people. Nevertheless, there were evidences of interest in the new doctrines. General Jullien came to Yverdon to study the methods, and issued two commendatory reports, which induced some thirty French pupils to go to Pestalozzi's institute. Chavannes also published a treatise upon the Pestalozzian methods in 1805. These efforts, however, had little effect upon education, and the Pestalozzian principles did not make much headway in France up to the revolution of 1830. After that time they rapidly became popular, especially through Victor Cousin. This famous professor, who was later minister of public instruction, issued in 1835 a \emph{Report on the State of Public Instruction in Prussia}, which showed the great merit of Pestalozzianism in the elementary schools of that country.

In England the influence of Pestalozzi was large, but the use made of his methods was not altogether happy. The private school opened by Mayo after his return from Yverdon employed object teaching in several subjects, and a popular text-book, entitled \emph{Lessons on Objects,} was written by his sister. This book of Elizabeth Mayo consisted of encyclopaedic lessons on the arts and sciences arranged in a definite series, and much beyond the comprehension of children from six to eight years old, for whom it was intended. Together with several texts of a similar sort, it had a wide influence in formalizing object teaching and spreading it rapidly. The Mayos were also interested in infant schools, and when they helped organize `The Home and Colonial School Society' in 1836, they combined the Pestalozzian methods with those of the infant school (see p.~246). Through the model and training schools of this society, formalized Pestalozzianism was extended through England and America.

\textbf{Pestalozzianism in the United States.} - Pestalozzianism began to appear in the United States as early as the first decade of the nineteenth century. It was introduced not only from the original centers in Switzerland, but indirectly in the form it had assumed in Germany, France, England, and other countries. The instances of its appearance were sporadic and seem to have been but little connected at any time. The earliest presentation was that made from the treatise of Chavannes in 1805 by William McClure. By this and other articles, McClure did much to make the new principles known in the United States, and in 1806 he induced Joseph Neef, a former assistant of Pestalozzi, to come to America and become his ``master's apostle in the New World.'' Neef maintained an institution at Philadelphia for three years and afterward founded and taught schools in other parts of the country. But his imperfect acquaintance with English and with American character and his frequent migrations prevented his personal influence from being greatly felt, and the two excellent works that he published upon applications of the Pestalozzian methods were given scant attention.

A large variety of literature, describing the new education, and translating the accounts of Chavannes, Jullien, Cousin, and a number of the German educationalists, was also published in educational journals, which were just beginning to appear in the United States (see p.~304). Returned travelers, like Professor John Griscom (see p.~305) published accounts of their visits and experiences at Yverdon and Burgdorf, such lecturers as the Reverend Charles Brooks began to suggest the new principles as a remedy for our educational deficiencies, and educational reformers, like the Alcotts, began to show the Pestalozzian spirit in their schools. Pestalozzi's objective methods and the oral instruction resulting from them were used in various subjects by a number of educators. For example, the methods advocated in arithmetic were introduced into America by Warren Colburn. He spread `mental arithmetic' throughout the country, and in his famous \emph{First Lessons in Arithmetic on the Plan of Pestalozzi,} published first in 1821, he even printed the `table of units' (Fig. 34). The Pestalozzi-Ritter method in geography was early presented in the United States through the institute lectures and text-books of Arnold Guyot, who had been a pupil of Ritter and came to America from Switzerland in 1848. The promotion of geographic method along the same lines was later more successfully performed by Francis Wayland Parker, who had studied with Guyot, in his training of teachers and his work on \emph{How to Teach Geography.} Colonel Parker has also had several successful pupils, who are to-day largely continuing the Pestalozzian tradition. The Pestalozzian method in music was brought into the Boston schools and elsewhere about 1836 by Lowell Mason. Mason, who was influenced by the works of Nägeli.

The most influential propaganda of the Pestalozzian doctrines in general, however, came through the account of the German school methods in the \emph{Seventh Annual Report} (1843) of Horace Mann (see p.~308), and through the inauguration of the `Oswego methods' by Dr.~Edward A. Sheldon. Mann spoke most enthusiastically of the success of the Prussian-Pestalozzian system of education and hinted at the need of a radical reform along the same lines in America. The report caused a great sensation, and was bitterly combated by conservative sentiment throughout the country, but the suggested reforms were largely effected. Dr.~Sheldon, on the other hand, caught his Pestalozzian inspiration from Toronto, Canada, where he became acquainted with the formalized methods of the Mayos through publications of the Home and Colonial School Society (see p.~291). He resolved to introduce the principles of Pestalozzi into the Oswego schools, of which he was at that time superintendent, and in 1861 secured from the society in London an instructor to train his teachers in these methods. There was some criticism of the Oswego methods on the ground of formalism, but as a whole they were pronounced a success, and in 1865 the Oswego training school was made a state institution. This was the first normal school in the United States where `object lessons' were the chief feature, but a large number of other normal schools upon the same basis sprang up rapidly in many states, and the Oswego methods crept into the training schools and the public systems of numerous cities. As a consequence, during the third quarter of the nineteenth century, Pestalozzianism, though somewhat formalized, had a prevailing influence upon the teachers and courses of the elementary schools in the United States.

\textbf{Pestalozzi's Industrial Training Continued by Fellenberg.} - Such was the wide influence of Pestalozzi upon education. But while throughout his work he continued to make new applications of his observational methods, his principle of combining industrial training with intellectual education, which he had begun so successfully at Neuhof and Stanz, could not be continued at Burgdorf. His pupils there came chiefly from aristocratic families and were not obliged to support themselves by manual labor. However, Pestalozzi still hoped to save enough of the income from the school payments of the rich to found a small agricultural school for the poor on this plan and connect it with the `institute,' and while this institution was never started, the opportunity for carrying out his aim came through his friend, Emanuel von Fellenberg (17 71-1844). Fellenberg belonged to a noble family of Berne, but, like Pestalozzi, he believed that an amelioration of the wretched moral and economic conditions in Switzerland should be accomplished by education. To secure the means for an experiment in this direction, he persuaded his father to purchase for him an estate of six hundred acres at Hofwyl, just nine miles from Burgdorf. Here Pestalozzi urged him to undertake his favorite idea of industrial education, and in 1806, with the aid of Zeller (see p.~289), who had been sent him by Pestalozzi, he opened a school to train teachers in the Pestalozzian method.

\textbf{The Agricultural School and Other Institutions at Hofwyl.} - Fellenberg especially desired, however, to combine Pestalozzi's observational work and his older principle of industrial training in an `agricultural institute' for poor boys. This plan was not fully realized until 1808, when he secured the enthusiastic Jacob Wehrli as an assistant. The work was so arranged that each old pupil, as fast as he was trained, took charge of a newer one as an apprentice, and the school from the first became a sort of family. The chief feature of the institute was agricultural occupations, including drainage and irrigation, but, from the requirements of farm life, it was natural to train also cartmakers, blacksmiths, carpenters, locksmiths, shoemakers, tailors, mechanics, and workers in wood, iron, and leather. Workshops for these industries were established upon the estate, and the pupils in the agricultural institute were enabled to select a training in a wide range of employments, without neglecting book instruction (Fig. 35). By this means, too, they could support themselves by their labor while being educated. Through the institute also, a considerable number of the pupils were trained to be directors of similar institutions, or to become rural school-teachers. Fellenberg thought it important that all who were to teach in the common schools should have a thorough acquaintance with the practical labor of a farm, the means of self-support, and the life and habits of the majority of their pupils.

But the work of Fellenberg did not stop there. From the beginning he had felt that the wealthy should understand and be more in sympathy with the laboring classes, and learn how to direct their work more intelligently. Hence he began very early an agricultural course for landowners, and many young men of the wealthy classes came to show a striking interest in his deep-soil ploughing, draining, irrigation, and other means of educating the poor. But these wealthier youths remained at the institute so short a time that he could not extend his ideals very widely. To retain them longer at Hofwyl, in 1809 he opened a `literary institute' which, besides the usual academic studies, used Pestalozzi's object lessons and strove to develop physical activities. Moreover, the pupils in the literary institute had to cultivate gardens, work on the farm, engage in carpentering, turning, and other mechanical occupations, and in many ways come into touch and mutual understanding with the poorer boys in the agricultural institute. The wealthy learned to dignify labor, and the poor, instead of envying those in the higher stations of life, became friendly and desirous of cooperating with them. Eventually there arose an independent community of youth, managing its own affairs outside of school, arranging its own occupations, games, and tours, choosing its own officers, and making its own laws. Within this little world was provided a training for society at large, with its various classes, associations, and corporations, which Fellenberg seems to have regarded as divinely ordained. Likewise, in 1823, a school for poor girls was opened by his wife, and four years later he started a `real,' or practical, school for the middle classes, which was intermediate between the two `institutes.'

\textbf{Industrial Training in the Schools of Europe.} - The educational institutions of Fellenberg (Fig. 36) were well managed and proved very successful, and the idea of education through industrial training spread rapidly. While, after the death of Fellenberg in 1844, the schools at Hofwyl gradually declined, various types of industrial education everywhere came to supplement academic courses, and extend the work of the school to a larger number of pupils. Thus the tendency of modern civilization to care for the education of the poor, the defective, and the delinquent through industrial training has sprung from the philanthropic spirit of Pestalozzi and his practical collaborator, Fellenberg, and has become apparent in all advanced countries. Industrial institutions rapidly increased in Switzerland, beginning in 1816 with the school in the neighboring district of Meykirch. In 1832 a cantonal teachers' association was formed at Berne, with Fellenberg as president and Wehrli as vice president, and every canton soon had its `farm school.' Industrial training was also introduced into most of the Swiss normal schools. In Germany the industrial work suggested by Pestalozzi and Fellenberg came into successful operation in many of the orphanages and most of the reform schools. Later, industrial education was taken up by the \emph{Fortbildungsschulen} (' continuation schools') of the regular system (see p.~420). At the reform and continuation schools of France industrial training has long formed the distinctive element in the course. Educators and statesmen of England likewise early commended the work of Fellenberg, and industrial training shortly found a foothold in various technical and reform schools of that country.

\textbf{Industrial Institutions in the United States.} - The industrial work of the Pestalozzi-Fellenberg system also began to appear in the United States about the close of the first quarter of the nineteenth century. After that, for twenty years or so, there sprang up a large number of institutions of secondary or higher grade with `manual labor' features in addition to the literary work. The primary object of the industrial work in these institutions was to enable students to earn their way through school or college and at the same time secure physical exercise. It was the first serious academic recognition of the need of a `sound mind in a sound body,' and did much to overcome the prevailing tendency of students toward tuberculosis and to furnish a sane substitute for the escapades and pranks in which college life abounded. The first of these manual labor institutions were established in the New England and Middle states between 1820 and 1830, but within a dozen years the manual labor system was adopted in theological schools, colleges, and academies from Maine to Tennessee. The success of this feature at Andover Theological Seminary, where it was begun in 1826 for `invigorating and preserving health, without any reference to pecuniary profit,' was especially influential in causing it to be extended. The `Society for Promoting Manual Labor in Literary Institutions,' founded in 1831, appointed a general agent to visit the chief colleges in the Middle West and South, call attention to the value of manual labor, and issue a report upon the subject. Little attention, however, was given to the pedagogical principles underlying this work. As material conditions improved and formal social life developed, the impracticability of the scheme was realized, and the industrial side of these institutions was given up. The physical exercise phase was then replaced by college athletics. By 1840-1850 most of the schools and colleges that began as `manual labor institutes' had become purely literary.

A further movement in industrial education has been found in the establishment of such schools as Carlisle, Hampton, and Tuskegee, which adopted this training as a solution for peculiar racial problems. But the original idea of Pestalozzi, to secure redemption through manual labor, has been embodied in American institutions since 1873, when Miss Mary Carpenter, the English prison reformer, visited the United States. Contract labor and factory work in the reformatories then began to be replaced by farming, gardening, and kindred domestic industries. At the present time, moreover, the schools for delinquents and defectives in the New England, Middle Atlantic, Middle West, and most of the Southern states, have the Fellenberg training, though without much grasp of the educational principles involved. Finally, there has also been a growing tendency in the twentieth century to employ industrial training or trade education for the sake of holding pupils longer in school and increasing the efficiency of the public system. In so far as it has tended to replace the more general values of manual training, once so popular, with skill in some particular industrial process, this modern movement represents a return from the occupational work started by Froebel to the philanthropic practice of Fellenberg and Pestalozzi.

\textbf{Supplementary Reading}

Graves, \emph{In Modern Times} (Macmillan, 1913), chap. V; and \emph{Great Educators} (Macmillan, 1912), chap. IX; Monroe, \emph{Textbook} (Macmillan, 1905), pp.~597-622; Parker, \emph{Modern Elementary Education} (Ginn, 1912), chaps. XIII-XVI. The \emph{Leonard and Gertrude} has been well arranged for English readers in the edition of Eva Channing (Heath, 1896) and \emph{How Gertrude Teaches Her Children} has been translated by Lucy E. Holland and Frances C. Turner (Bardeen, 1898). The standard English treatises on Pestalozzi are Guimps, R. de, \emph{Pestalozzi, His Aim and Work} (Appleton, 1890); Holman, H., \emph{Pestalozzi} (Longmans, 1908); Kriisi, H., \emph{Pestalozzi, His Life, Work, and Influence} (American Book Co., 1875); Pinloche, A., \emph{Pestalozzi and the Foundation of the Modern Elementary School} (Scribner, 1901), and, more recently, Green, J. A., \emph{Life and Work of Pestalozzi} (Clive, London, 1913) and \emph{Pestalozzi's Educational Writings} (Longmans, Green, 1912). Monroe, W. S., has furnished an interesting \emph{History of the Pestalozzian Movement in the United Slates} (Bardeen, 1907). \emph{The Institutions of De Fettenberg} were fully described by King, W. (London, 1842); and by Barnard, H., in his \emph{American Journal of Education,} vol.~Ill, pp.~591-596; XIII, 323-331; and XXVI, 359-368.

\hypertarget{development-of-public-education-in-the-united-states}{%
\chapter{23. Development of Public Education in the United States}\label{development-of-public-education-in-the-united-states}}

\textbf{Outline}

\begin{itemize}
\tightlist
\item
  During the second quarter of the nineteenth century a third period in the educational history of America, marked by further democratization and a great expansion of public education, appeared.
\item
  It began with an awakening generally known as `the revival of common schools,' which was most noticeable in New England. Here, owing to the attacks made upon him by reactionaries, Horace Mann was the most conspicuous reformer; while Henry Barnard, through his \emph{American Journal of Education,} enabled educators to look beyond the educational experience of America. But the influence of this awakening was also felt in every other section of the United States.
\item
  It was followed by a steady growth in universal education, state support and control, local supervision, and the organization of normal schools in New England and the Middle states.
\item
  In the Northwest, common school advocates overcame the opposition of settlers from states not committed to public education, and in the further expansion of the United States progress in common school sentiment has kept pace with the settlement of the country.
\item
  The South made considerable progress during the early years of the awakening, and while the Civil War crushed its educational facilities, the struggle for public education has since been won.
\end{itemize}

\textbf{The Third Period in American Education.} - Interest in the improved methods of Pestalozzi and other reformers that was manifesting itself everywhere in the United States during the second quarter of the nineteenth century seems to have been but one phase of a much larger movement. It was about this time that a third period in American education, which was marked by the development of democratic ideals and the extension of state systems of public schools, may be said to~have begun. During the period of `transition,' we found~(chap. XXI), half a dozen of the states had started an organization of common schools, and in a dozen others permanent school funds had been established, an influential minority of leading citizens were constantly advocating universal education, and public interest in the matter was evidently increasing- But the consummation of a regular system was still much hindered by sectarian jealousies, by the conception of public schools as institutions for paupers and the consequent custom of allowing private schools to share in public funds, by the unwillingness of the wealthy to be taxed locally for the benefit of other people's children, and, in New England, by the division of the system into autonomous districts and the interference of petty politics. Hence, while much progress had been made since the early days of `transplantation' of European ideals and institutions, there was still much need of the expansion and further democratization that now began to appear. Of the rapid development that took place during this final period of Americanization, much was accomplished before the middle of the nineteenth century, but educational progress continued through the final decade.

\textbf{Early Leaders in the Common School Revival.} - The educational awakening with which the beginning of this~third period seems to be marked, has been generally stop center known as `the common school revival.' It first became~evident during the latter part of the decade between~1830 and 1840, and had its storm center in Massachusetts~and Connecticut. While it greatly furthered the cause of public education everywhere, because of the decadence into which New England had fallen, the demand for an educational awakening was strongest there. In this revival the most conspicuous figure was probably Horace Mann, but there were several leaders in the field before him, many were contemporaneous, and the work was expanded and deepened by others of distinction long after he withdrew from the scene. For a score of years before Mann appeared, definite preparation for the movement had been in progress, and the labors of the individuals and associations engaged in these endeavors should be briefly noted. Many of the reformers seem to have\emph{~}recommended an improvement in methods through the creation of an institution for training teachers, thus anticipating one of the greatest achievements of Mann. Actual attempts at a private normal school were even made by the Reverend Samuel R. Hall at Concord, Vermont (1823), Andover, Massachusetts (1830), and Plymouth, New Hampshire (1837).

A number of educational journals, moreover, published articles on schoolbooks, the methods of Lancaster, Pestalozzi, Neef, and Fellenberg, the infant and Sunday schools, physical education, European school systems, and a variety of other timely topics and reforms. Among~these progressive publications were the \emph{American Journal of Education,} edited by William Russell from 1826-1830, and then continued from 1831 to 1839, as the \emph{American Annals of Education} under the editorship of William C. Woodbridge, and the \emph{Quarterly Register,} published 1828-1843 by the `American Educational Society.' The latest European ideas were also reported from first-hand observation by a number who had gone abroad to investigate. The most influential of these reports was \emph{A Year in Europe,} written in 1819 by Professor John Griscom (see p.~292), who was a lecturer before several New York associations, including the Public School Society. Almost as widely read were the reports of William C. Woodbridge in 1824, and of Professor Calvin E. Stowe of Lane Theological Seminary, Cincinnati, in 1836.

\textbf{Work of James G. Carter.} - All these movements indicate the educational ferment that was going on. But the predecessor of Mann, who accomplished most for the common schools, and influenced that reformer most directly, was James G. Carter (1795-1849). Carter (Fig. 37) was a practical teacher and wrote continually on the need of a training institution to improve instruction in the public schools. These appeals proved very successful, and earned him the title of `father of the normal schools.' After being elected to the legislature, he accomplished much by his zeal and skill in parliamentary tactics. In 1826 he secured an act by which each town~as a whole was required to choose a regular committee, instead of the ministers and selectmen, to supervise the schools, choose text-books, and examine, certify, and employ the teachers. But the effect of this enactment was largely lost the following year by allowing the districts, as a compromise, to choose a committeeman, who should appoint the teachers. In 1826 he placed~secondary education, then largely conducted by academies, more under public control through a law requiring each town of five hundred families to support a free English high school (Fig. 41), and every one of four thousand inhabitants to maintain a classical high school. Next, in 1834, Carter succeeded in getting a state school fund established from the proceeds of the sale of lands in the province of Maine and the state's claims against the federal government for military services. But his most fruitful victory was won in 1837, when he procured the passage of the bill for a State Board of Education, after it had been once defeated, by inducing the house to discuss it in `committee of the whole.'

\textbf{Horace Mann as Secretary of the Massachusetts Board.} - By reason of his merits as an educator, his persistent efforts in behalf of educational reform, and his advocacy of the bill, it was assumed by most people that Carter would be chosen secretary of the new board. To their surprise, a lawyer named Horace Mann (17961859), at that time president of the senate, was selected for the post, but the choice is now known to have been most fortunate. By both heredity and training Mann (Fig. 38) was suffused with an interest in humanity and all phases of philanthropy and education. He possessed a happy combination of lofty ideals, intelligence, courage, enthusiasm, and legislative experience, which equipped him admirably for leadership in educational reform. The law proposed for the new Board of Education numerous duties in the way of collecting and spreading information concerning the common schools and of making suggestions for the improvement and extension of public education, but it provided no real powers, and the permanence~and influence of the board depended almost wholly upon the intelligence and character of the new Secretary.

During his twelve years in the office, Mann subserved the interests of the state most faithfully. To awaken the people, he made an educational campaign through every portion of the state each year, but an even more effective means of disseminating his reforms was found in his series of \emph{Annual Reports.} These documents were by law to give information concerning existing conditions and the progress made in the efficiency of public education each year, and they deal with practically every educational topic of importance at the time. Sometimes they seem commonplace, but it must be remembered that they were not so then, and that the work of Mann did much to render them familiar. They vitally affected school conditions everywhere in New England, and were read with great interest in all parts of the United States, and even in Europe. He also published semimonthly the \emph{Massachusetts Common School Journal,} to spread information concerning school improvement, school law, and the proceedings of the State Board. But it consisted of only sixteen pages, and was not as valuable as some of the educational journals that had preceded it (see pp.~704 f.). Another medium in the improvement of educational facilities was Mann's general establishment of school libraries by state subsidy throughout Massachusetts. But probably the most permanent means of propagating his reforms came through securing the foundation of the first public normal schools in this country. Massachusetts was in 1838 induced to establish three schools, so located that all parts of the state might be equally served. The course in each school consisted in a review of the common branches from the teaching point of view, work in educational theory, and training in a practice school under supervision, and, while not largely attended, these institutions were a great success from the start.

The arduous and unremitting labors of Mann in instituting and promoting the various means of school reform made the greatest inroad upon his strength and financial resources. Moreover, he was for years violently assailed by reactionaries of all types. His controversy with the~Boston schoolmasters was especially sharp.~Mann's \emph{Seventh Annual Report} (1843) gave an account of~his visit to foreign schools, especially those of Germany, and praised with great warmth the `Pestalozzian' (see p.~289) instruction without text-books, the enthusiastic teachers, the absence of artificial rivalry, and the mild discipline in the Prussian system. The report did not stigmatize the conservatism of the Boston schools or bring them into comparison with those of Berlin, but the cap fitted. The pedagogues were disquieted, and proceeded to answer savagely. But when the smoke of battle had cleared away, it was seen that the leaders of the old order had been completely routed. A more insidious attack was that led by the ultra-orthodox. The~old schools of the Puritans, with their dogmatic religious teaching, had been steadily fading for more than a century before the new board had been inaugurated, but many narrow people were inclined to charge this disappearance to the reformer, whose liberal attitude in religion was well known. The assaults, however, were vigorously and successfully repelled by the Secretary. And while these controversies wore Mann out and probably led ultimately to his resignation, they had much to do with making his reputation as a great educator. They have even caused us at times to forget that he was but a striking figure in a general movement. Men like Carter were in the field long before him, and his co-worker, Barnard, served the cause of education for nearly half a century after Mann withdrew.

\textbf{The Educational Suggestions and Achievements of Mann.} - In surveying his educational positions, we find Mann's foremost proposition was that education should be universal and free. Girls should be trained as well as boys, and the poor should have the same opportunities as the rich. Public schools should furnish education of such a quality that the wealthy would not regard private institutions as superior. This universal education, however, should have as its chief aim moral character~and social efficiency, and not mere erudition, culture, and accomplishments. And morality, he felt, would not be accomplished by inculcating sectarian doctrines. Mann was, however, mainly a practical, rather than a theoretical reformer, and to the material side of education he gave serious attention. He declared that school buildings should be well constructed and sanitary. This matter seemed to him so important that~he wrote a special report upon the subject during his first year in office. He carefully discussed the proper plans for rooms, ventilation, lighting, seating, and other schoolhouse features, and insisted that the inadequate and squalid conditions which existed should be improved. As to methods, he maintained that instruction should methods, be based upon scientific principles, and not upon authority and tradition. He advocated the word method~of reading, in the place of the uneconomical, artificial, and ineffective method of the alphabet, and the Pestalozzian object methods and oral instruction were introduced by him. He held that the work should be guided by able teachers, who had been trained in a normal school, and should be imparted in a spirit of mildness and kindness through an understanding of child nature. In the matter of the studies to be pursued, Mann was inclined to be exceedingly practical. In discussing educational values, he failed to see any reason ``why algebra, a branch which not one man in a thousand ever has occasion to use in the business of life, should be studied by more than twenty-three hundred pupils, and bookkeeping, which every man, even the day laborer, should understand, should be attended to by only a little more than half that number.'' Similarly, he holds that of all subjects, save the rudiments, physiology should receive the most attention.

In order that these various reforms might be realized, Mann insisted frequently that the state should spare no labor or expense. But in a republic he felt that ``education can never be attained without the consent of the whole people.'' It was a general elevation of ideals, effort, and expenditure that he sought, and for which he began his crusade. And the general progress that resulted in this period covers a wide range. During his secretaryship the appropriations made for public education in Massachusetts were more than doubled, and the proportion of expenditure for private schools in the state was, in consequence, reduced from seventy-five to thirtysix per cent of the total cost of education. The salaries of masters in the public schools were raised sixty-two per cent, and, although the number of women teachers had grown fifty-four per cent, the average of their salaries also increased fifty-one per cent. The school attendance enormously expanded, and a full month was added to the average school year. When Mann's administration began, but fourteen out of forty-three towns had complied with the high school law of 1826, but, by the middle of the century, fifty new high schools had been established. The efficiency of supervision was largely increased by making the compensation of the town visiting committees, established through Carter, compulsory by law. The first state normal schools at last appeared, and~teachers' institutes, county associations, and public school libraries were given general popularity. Quite as marked was the improvement effected in the range and serviceability of the school studies, in text-books, methods of teaching, and discipline. Thus under the leadership of Horace Mann a practically unorganized set of schools, with diverse aims and methods, was welded into a well-ordered system with high ideals, and the people of Massachusetts renewed their faith in the common schools.

\textbf{Henry Barnard's Part in the Educational Awakening.} - But there was another important contribution to the awakening made by a New Englander, which was of a rather different nature from that connected with the influence of Horace Mann. Before that reconstruction of the common schools, which was responsible for the best elements in our national civilization, could be at all complete, it was necessary that America should have a better comprehension of what was being done in education elsewhere. The United States had for two centuries been undergoing a gradual transition from the institutional types transplanted from England and the Continent in colonial days, and was coming more and more to blossom out into democracy and the people's schools, but for a long time there was little knowledge of what was being done by the other countries that had by this time adopted similar ideals. Conceptions of universal and democratic education and of improved organization and methods had been slowly developing in Prussia and other German states, and had extended to France and elsewhere. A literature connected with the advanced theories of such reformers as Rousseau, the philanthropinists, Pestalozzi, and Fellenberg had likewise grown up in Europe. It was very important that America, now keenly alive to the need of educational reorganization, should become acquainted with all this, that the New World might secure the advantages of comparison, corroboration, and expansion of view from the work of older civilized peoples. Some reports on foreign education and translations of European treatises had already appeared (pp.~304 f.), but the time was now ripe for a more extensive and systematic exposition of European education and its application to popular education in America, and for a really capable scholar to bring these world views within the grasp of all classes of teachers and educational authorities. This literary representative of the awakening appeared at length in Henry Barnard (1811-1900), who is fully worthy of a place in the educational pantheon of America. Barnard (Fig. 39) made a brilliant record at Yale for general scholarship, and a position as assistant librarian during his last two years in college did much to afford him a wide grasp of~bibliography. After graduation, he obtained a valuable experience in teaching, and, by travelling extensively in America and Europe, formed a broad acquaintance with educational institutions, libraries, galleries, and social conditions in all the leading states and nations.

\textbf{Barnard as Secretary of the Connecticut State Board.} - Two years after Barnard's return to Connecticut, he began his part in the educational awakening as Secretary of the new State Board of Commissioners of Common Schools, and undertook to do a work similar to that of Mann in Massachusetts. Throughout the eighteenth century Connecticut schools had been among the most~efficient in the country, but since the income from the~Western Reserve lands had begun in 1798, and especially since this had been increased by the United States deposit fund in 1836, public education had steadily declined. A state tax was still maintained, but all local effort was paralyzed through lack of exercise. Another factor in producing this decline was connected with the transferal of the management of the common schools from the town to the `school society,' which was a species of district, almost identical with the parish of each Congregational church. The results of this ruinous policy had been revealed in an investigation made by the legislature, which showed that not one-half of the children of school age were attending the common schools, and that the~teachers were poorly trained and supervision was neglected. Barnard at once began to urge many reforms, and in his reports and the \emph{Connecticut Common School School journal Journal} made suggestions for a complete plan of public education. He also began the publication of his rich~collection of material bearing upon popular training at tionai material. home and abroad. But he was more a scholar and literary man than an educational statesman like Mann. He succeeded in getting the legislature to pass several reforms and a general revision and codification of the school laws, and in arousing several towns to amend their educational plans, although the crucial difficulty of the `school societies' could not be touched, and within four years the conservatives succeeded in legislating him out of office and in undoing all his reforms.

\textbf{Commissioner of Common Schools in Rhode Island.} - This gave Barnard an opportunity to pursue his favorite investigations, and for about a year and a half he was engaged in collecting material for a history of education in the United States. Then he was persuaded by the governor of Rhode Island to become the first Commissioner of Common Schools for that state. While he found in Rhode Island a better educational sentiment and less opposition than in Connecticut, the actual condition of the decentralized and individualistic schools was far worse (see p.~269). But, through his assemblies of Radical re- teachers and parents and his educational treatises, he soon began to convince the people of the unwisdom of district organization, untrained teachers, short terms, irregular attendance, poor buildings and ventilation, and meager equipment. He also continued to publish his collection of educational material through the foundation of the \emph{Rhode Island School Journal.} As a result of his efforts, when failing health compelled him to resign in 1849, the state no longer regarded wilfulness and personal opinion as praiseworthy independence, and he could honestly claim that Rhode Island had at the time one of the best school systems in the United States.

\textbf{State Superintendent of Schools in Connecticut.} -~But the \emph{clientele} that Barnard had built up in Connecticut continued his reforms and constructive work after his departure, and improved upon them. In 1851, they even succeeded in having him recalled virtually to his old duties. He was designated as State Superintendent of Common Schools, as well as Principal of the State Normal School, which had been established through the efforts of his adherents. The state had now learned its error in mingling politics with education, and Barnard was able to carry out his reforms unmolested. Through the normal school he sent out a great body of trained teachers. He revised the school code, checked the power of the `school societies,' consolidated and simplified the organization and administration of public education, made a more equitable distribution of the school fund, and encouraged local taxation. But his most distinctive work, as might be expected, was on the literary side. He prepared a valuable series of documents upon foreign education, normal schools, methods of teaching, school architecture, and other topics, and a long report upon \emph{The History of Legislation in Connecticut Respecting Common Schools up to 1838.}

\textbf{\emph{Barnard's American Journal of Education.}} - It was, too, during the last days of his Connecticut superintendency that Henry Barnard suggested the establishment of a national journal of education. He first broached the matter to the `American Association for the Advancement of Education' at its meeting in Washington, December, 1854. But the association soon found itself unable to pursue this enterprise for lack of financial~support, and in May of the next year Barnard began~the publication of the \emph{American Journal of Education} at his own expense. It was at first planned to run the journal for five years only, but, although the work was somewhat interrupted upon occasions by other duties, it continued for more than a generation, until at length thirty-one large octavo volumes, averaging about eight anr\textbar{}efiity-twoS hundred pages each, had been issued. In addition, fifty-two special treatises reprinted from articles in the journal brought the material together in a connected way. Besides giving nearly all his time to editing this \emph{magnum opus,} Barnard sank his entire fortune of \$50,000 in its publication. This great treasury of material includes every phase of the history of education from the earliest times down into the latter half of the nineteenth century. It furnishes accounts of all contemporaneous systems in Europe and America, descriptions of institutions for the professional training of teachers, and essays upon courses of study for colleges and technical schools, the education of defectives and delinquents, physical education, school architecture, great educators, and a large variety of other themes. While it is always most reliable in its treatises upon foreign education, of even greater value is its practical grasp of educational life in America from the beginning. It contains the greatest collection of interesting monographs upon the development of ideals and organization in the various states, and gives the most complete description in literature of the educational life of a nation.

\textbf{First United States Commissioner of Education.} - In 1867 Barnard was appointed the first United States Commissioner of Education. This office he had been constantly trying to have established ever since he had~found, as Secretary of the Connecticut Board, how absolutely lacking the federal government was in school statistics and documents. He hoped that, through the agency of the government, facilities might be secured to collect and publish trustworthy educational statistics, and to issue a library of independent treatises. The bureau was not created for many years, and then through the immediate initiative of another, but when Barnard was called to the commissionership, he organized the office practically upon the lines he had previously suggested. He suspended his \emph{Journal} and used the product of his investigations in the annual reports of the office.~He started that searching inquiry into the administration, management, and instruction of institutions of every grade, and into all educational societies, school funds, legislation, architecture, documents, and benefactions that has since been maintained by the Bureau of Education. However, within three years a change in politics brought a new incumbent into the commissionership, and Barnard gave his literary efforts once more to his beloved \emph{Journal.}

\textbf{Value of Barnard's Educational Collections.} - Hence, Barnard's real life work may be considered the collection of a great educational compendium. By temperament, native ability, and habit, he proved himself well fitted to be the leading representative of the literary side of~the awakening. Through his work American education awakening, was, in its period of greatest development, granted the opportunity of looking beyond the partial and local results of the first half century of national life. It was enabled to modify and adapt to its own uses the educational theories, practices, and organizations of the leading~civilized peoples, and to bring together for a comparative view sections and states that were widely separated. \emph{Barnard's American Journal of Education} was not intended to be a universal encyclopaedia of education, but often includes a condensation of important works or a presentation of highly scientific methods and profound philosophic systems in popular form. It was not possible, either, to classify and work out a connected and complete historical account, when there were no reliable records or collections of materials in existence. It was necessary that some one should first gather the information from newspapers, pamphlets, memorials, monographs, and plans, and publish it as it was found. In this way he accomplished a more valuable work than if he had published a systematic history of education in the United States.

\textbf{Educational Development in New England since the Revival.} - This great storehouse of information published by Barnard and the virile efforts of Mann and other practical leaders were but prominent evidences of The `revival' the progress that was at the time sweeping over the but ftsresuits entire country. The educational awakening of 183 5-1860 strikinghiNew was general and proved one of the most fruitful in history. England. jts influence was felt in every state, and it led to the third period of American education, which has been characterized by the expansion of public schools and state educational systems. During this period new ideals of democracy have come to be felt in American education, and a rapid advance has taken place in the evolution of that unique product, the American public school. In describing this development, we may turn first to New England.

In Massachusetts Horace Mann has been followed in the central administration by a succession of seven scholarly and experienced educators, who believed as firmly as he that all stages of education below the college should be open at public expense without let or hindrance to the richest and poorest child alike. Since the revival the state has seen a steady growth of sentiment for universal education and improved schooling, and never~again has such an upheaval of the educational strata been necessary. The income of the state school fund and additional appropriations have been steadily increased, their apportionment among the towns has been rendered more equitable from time to time, and an effort has constantly been made to distribute them in such a way as to encourage local effort and cooperation. The school term has been lengthened to ten months and the average attendance of pupils to seven years. The improvements in school buildings, sanitation, and equipment have steadily advanced. The district system died hard, and not until 1882 was it altogether forced out of existence.

Most of the academies, too, which proved such a hindrance to the development of public secondary education, gradually died or were merged in the public system as high schools. By means of state aid, it has been possible since 1903 for the smallest towns to afford a high school training for their children at public expense. Supervision has also become universal during the past quarter century. Springfield first introduced a superintendent of schools in 1841, Gloucester in 1850, Boston in 1851, and the other cities much later, but since 1888, through increasing state aid and the combination of~smaller towns into a district superintendency, expert supervision has become possible everywhere, and during the last decade it has been compulsory. The normal schools, which have now increased to ten, have brought about a striking improvement in teaching. It is practically impossible at present for an untrained teacher to secure a position in the elementary schools of Massachusetts, and, through a system of examinations and investigations, teachers of exceptional ability have, since 1896, been granted an extra weekly allowance by the state. Since the middle of the century, the state board has been permitted to appoint a number of agents, to assist in inspecting and improving the schools, especially in the smaller towns and rural districts.

The course of development since the awakening has been very similar in the other New England states. The successors of Barnard in the central administration both in Rhode Island and Connecticut have been skilled and earnest educators, and, while their reports lacked his literary touch, they were of rather more practical character. Until 1856, Connecticut made no attempt to return from the parish to the town organization. Even then, as well as later, legislation on the subject was `permissive,' and not until the twentieth century was the `school society,' or district system, given up in half of the towns. In Rhode Island, even after Barnard's reforms, almost one-third of the districts did not own their school buildings, owing to the survival of the method in use when the schools were private, but this condition has gradually been remedied. Likewise, the number of towns levying sufficient local taxes to secure a share in the state apportionment rapidly grew, and the state appropriation itself doubled and quadrupled within a generation. In Vermont, New Hampshire, and Maine, owing to insufficient wealth, infertility of soil, and sparseness of population, effective public education has been reached only by slow and cautious steps. But even these states have gradually centralized their educational administration through the abolition of the district system and the creation at various times of a state superintendent, a state commissioner, or a state board and secretary. This reorganization has been followed by increased state school funds and appropriations, more systematic statistics and reports from the schools, and great advances in universalizing and improving all stages of public education.

\textbf{Influence of the Awakening upon the Middle States.} - Although this awakened sentiment for education and progress in the common school has been most patent and spectacular in New England, it has not been peculiar to that part of the country. Nearly all of the other states~seem to have felt the influence of the awakening. In~close conjunction with the revival in New England, states, the movement appeared in New York, especially the western part, and was more or less evident in Pennsylvania, New Jersey, and Delaware. But because of its cosmopolitanism and the need of fusing so many different political, religious, and industrial traditions, the older parts of New York, where the school system had until the awakening been rather in advance of other states, did not progress as rapidly in the development of public education as Massachusetts and Connecticut. It had, however, by the time of the Civil War, succeeded in working over its heterogeneous people into a unified~civilization and in causing their children to be educated together for a common citizenship.

The most distinct advances during this period of final organization have been in the establishment of state normal schools, instead of subsidizing academies to train teachers, in the administration and supervision of the system, and in the methods of state support of education. The first state normal school was opened at Albany in 1844, and this pioneer institution has eventually been followed by ten others. In 1854 the state superintendency had once more been separated from the secretaryship of state, with which it had been combined for thirty-five years (p.~259). In 1856 local supervision was established through the appointment of school commissioners for the cities and villages. In the same year, a three-quarters of a mill tax was placed upon the property valuation of the state, and during the next dozen years many improvements were made in the disbursing and accounting of public funds. At length, in 1867, the long fight that had been made for entirely free education was successful. Until then nearly fifty thousand children had been deprived of all education, because their parents were too proud to secure payment of their tuition fees by confessing themselves paupers. It was during this era of progress, too, that New York City was, in 1842, allowed to place the direction of its schools in the hands of a board of education, elected by the people,~instead of giving over the city's share of the state funds to a quasi-public society, controlled by a close corporation. For eleven years, however, the Public School Society refused to give up its work, but by 1853 it decided to disband and merge its buildings and funds with those of the city school system (see p.~261).

Pennsylvania was slower than New York in showing the effects of the educational awakening, but the leaven was at work. While a number of progressive governors and other statesmen continually recommended the development of public education, and the `Pennsylvania Society for the Promotion of Common Schools' had been organized, the towering leader in this movement was Thomas H. Burrowes. As secretary of state and \emph{ex officio} superintendent of schools (1836-1838), as a public speaker and educational journalist (1838-1860), and as state superintendent (1860-1862), he constantly urged a complete system of public education, the establishment of normal schools, a separate state department of education, and the organization of state and county supervision. In 1849 the `permissive' feature of the law of 1834 was abolished, and the two hundred districts that had thus far refused to establish public schools were forced to do so under the new provisions. In 1854 a revised school law was passed, which, after twenty years, now made the state system of education complete. It established in the secretary of state's office a deputy superintendent of schools, who had virtually a separate department, and provided for county superintendents. Three years later the state educational department became absolutely independent under the care of a superintendent, and provision was made for a system of normal schools. These institutions were to be established at first~by private enterprise and without state subsidy. By 1877 there were ten in operation, largely maintained by the state. Three others have since been added, and the~state has begun to take over into its own hands the entire support and control of them all.

Educational progress in New Jersey also took some time to get under way, but when the reforms once started, they continued until an excellent system of common schools had been inaugurated. In 1838 the limitation of state funds to the education of the poor was removed, and the apportionment of the income from them was thereafter applied only to public schools. Since 1848, when a state superintendency was established, the development has been more rapid. County supervision has been introduced, state normal schools have been established at Trenton and Upper Montclair, and appropriations have been greatly increased. In 1911 a state commissioner of education with an efficient corps of deputies was provided. Delaware, on the other hand, failed to live up to the possibilities under her early `permissive' laws. Even the organization of `the friends of common school education' showed itself very conservative, and would not advocate the creation of a state superintendency or the establishment of state normal schools. In fact, Delaware did not organize a complete state system until after the war. Even then, while a state board and state superintendency were established in 1875, there were no county superintendents, and when county supervision was introduced in 1888, the state superintendency was abolished. It was not reestablished until 1912, but since then the state system has made evident progress.

\textbf{Public Education in the West.} - The budding of a common school system, which had just begun to appear in the new commonwealths of the Northwest before 1840, rapidly unfolded into full blossom during this educational springtime. Through this awakening the common~school advocates in Ohio, Indiana, and Illinois were~greatly aided in their struggle to overcome the opposition of~settlers from the states not committed to public education (see p.~272), and they were favored to some extent by accessions of emigrants from the home of the public school movement. During the decade just preceding the middle of the century, there was a decided elevation of public sentiment going on. Under the leadership of Samuel Lewis and Samuel Galloway in Ohio, Caleb Mills in Indiana, and Ninian W. Edwards in Illinois, the friends of public education had marshalled themselves for battle. Reports and memorials were constantly presented to the legislatures of these states, and public addresses in behalf of common schools were frequent in most large communities. A group of devoted schoolmen appeared, who were as successful in lobbying for good legislation as they were with institutes and public lectures. While reactions occasionally happened, like that in Ohio between 1840 and 1845, when the state superintendency was temporarily abolished, public education gradually came to be regarded as something more than merely free education for the poor, and public school funds were no longer granted as a subsidy to private institutions. After a quarter century of `permissive' laws, local taxation and free common schools were fully realized in all three states early in the fifties. The contest, of course, was not ended, as reactionary elements, with selfish, local, and sectarian interests, still remained, but their contentions have never again been more than partially successful. New features of the common schools, such as efficient teachers for the rural districts,~county supervision, state normal training, and free higher education in state universities, have gradually rendered the state systems more consistent and complete.

In Michigan, on the other hand, where there was not such a mixture of population, and a complete sympathy with the common school idea appeared, there was almost unhampered progress from the beginning of statehood. Under the first constitution (1837), there was provision made for a permanent school fund and for a local tax in every district, although the schools were partly maintained until 1869 by `rate bills' collected from the pupils. In accordance with the grant of two townships of land by Congress in 1826 for a university, the first legislature of the new state established the University of Michigan (Fig. 42), and its doors were open to students in 1841. It soon became the most prominent of the state universities. There was also provided a system of `branches' of the university, whereby a liberal grant was made for an academy in any county that would furnish suitable buildings and a sum equal to the appropriation from the state. As this proved a dissipation of the university funds, it was gradually stopped, and between 1852 and 1860 `union' and high schools were rapidly developed to supply the means of fitting for the university. In 1850 a state normal school was founded, and four ethers have since been added.

In all the other territory acquired or purchased by the United States in its westward expansion, the educational history has been very similar to that in the first states of the Northwest. Progress in common school sentiment has been made \emph{pari passu} with the settlement of the country. Each state, upon admission, has received its sixteenth section of school land and two townships for a university, and in the states admitted since 1848 the endowment of schools has been increased to two sections, while Texas, which had been an independent republic (1836-1845), stipulated before becoming a state that it should retain sole possession of its public lands, and has set aside for education nearly two and one-half millions of acres. Hence in the first constitution of each state, permanent school and university funds, together with a regular organization of the schools of the state, have generally been provided. In few cases have sectarian interests been able to delay or injure the growth of common schools in any of the later commonwealths, and the interpretation of public education as schools for the children of paupers has never seriously influenced the West.

\textbf{Organization of State Systems in the South.} - Thus through the awakening of common schools that occurred throughout the union from 1835 to i860 was the oldtime country and city district school of the North gradually lifted up to the present system of graded free elementary, secondary, and normal schools, together with city and state universities. But these results were not at first as fully realized in the South, because of the approach and precipitation of the dreadful internecine conflict that weighed down and finally prostrated the resources of that section. However, except for this impending calamity, the conditions in the South were not essentially different from those in any other section. During the earlier years of the awakening, and in some states up to the very verge of the Civil War, great progress in public education was noticeable. The attendance in the common schools, established in several states by `permissive' legislation, had been rapidly growing for a score of years, and there was an increasing body of prominent men desirous of enlarging popular education. During the early forties there were many efforts and suggestions for a system of public schools, and several conventions were held in the interest of such institutions. North Carolina actually established a state system in 1839. Tennessee (1838-1843) and Kentucky (1838) made less enduring efforts toward a similar organization, and as late as 1858 Georgia took a distinct step forward in this direction. Moreover, even in their secession conventions some states, like Georgia, adopted resolutions or constitutional amendments looking to the education of the people, and North Carolina in 1863, with the union army actually at its doors, undertook to grade the schools and provide for the training of teachers. But, in general, as the impending conflict drew near, attention to educational progress was forced to give way to the preservation of state and home, and after the war, which crushed and ravaged nearly every portion of the South, educational facilities had for the most part been totally wrecked.

Nevertheless, in the end the war served as a stimulus to common schools. It brought about a complete overturn of the old social and industrial order, and the South realized more fully than ever that it could arise from its desperate material and educational plight only through the institution of universal education. As early as 1865, school systems were organized in the border states, - Maryland, Kentucky, Missouri, and West Virginia, and even during the harsh and unhappy days of `reconstruction' (1867-1876), efforts were made in other states to build up systems of free public education. The organization of education became more thorough and mandatory than before the war. All children, white and colored, were to attend school between six and twentyone, and the term was to last from four to six months each year. Property and poll taxation were established for the support of the schools. A state superintendent and state board of education, county commissioners and a county board, and trustees in each district, were provided for. Text-book commissions were often established, and free books were granted to poor children. The foundation for a real system was thus laid.

This was a tremendous undertaking, and shows the greatest courage and executive ability upon the part of the South. Property had been diminished in valuation to the extent of nearly two billion dollars, and there were two million children to be educated. Moreover, under the reconstruction regime, the tax on property was often not collected, and the appropriations for education remained on paper. Indifference and inexperience were aggravated by the fear that `mixed' schools would be forced upon the white population by a reconstruction legislature or a Congress with millennial zeal in behalf of universal brotherhood. These obstacles, together with misdirected effort upon the part of Northern missionaries, and other serious interferences, for fully a decade constituted an enormous stumbling-block. Several factors, however, aided and encouraged the South in its efforts. Of these the most important was the foundation in 1867 of the Peabody Educational Fund of cationai Fund \$2,000,000, well characterized as ``a gift to the suffering~South for the good of the Union.'' This fund was placed in the management of the wisest and most sympathetic agents, who appealed to the higher sentiment of the communities and the states, and granted the assistance necessary to stimulate local effort in education. When the fund proved insufficient for the great task, the trustees pleaded with Congress for an additional subsidy, and made the whole country aware of the crying needs of education in the South. Through these appeals, more than ten million dollars from various sources have since been granted to the different grades of public education.

Despite the tremendous rally during the seventies, struggle won however, the struggle for public education in the South~was not won for twenty years, but complete systems of common schools have now at length been generally established. With the cessation of the reconstruction influence and the subsidence of the dread of mixed schools, attendance and appropriations have greatly increased, schools for the education of colored children have been furnished, and provision has been made for training and stimulating teachers of both races. Separate state institutions for higher education, cultural and vocational, have been established to furnish a broad education for both whites and negroes. Since 1890 there has been an ever increasing interest in improving the public school in all respects, and the expenditures and facilities for education have been constantly increasing.

\textbf{Development of the American System of Education.} - With its final development in the South during the last decade of the nineteenth century, the distinctly American public school system may be said to have been fully~elaborated. The educational ideals and institutions imported from Europe in the colonial period have gradually been modified and adapted to the needs of America. Schools have become public and free in the modern sense. The control of education has passed from private parties and even quasi-public societies to the state. The schools have likewise come to be supported by the state, and are open to all children alike without the imposition~of any financial obligation. In secondary education, and state unithe academies, which supplanted the `grammar' schools, ushed. first became `free academies' and made no charge for tuition from local patrons, though remaining close corporations, and then were in time replaced by the true American secondary institution, - the high school (Fig. 41). Colleges became largely non-sectarian, even when not nominally so, and state universities were organized in all except a few of the oldest commonwealths (Fig. 42). Thus has the idea of common schools and the right to use the public wealth to educate the entire body of children into sound American citizenship been made complete. Although the system is still capable of much improvement, it is expressive of American genius and development. It is simply the American idea of government and society applied to education. It is the educational will of the people expressed through the majority, and the resultant of the highest thinking and aspirations of a great nation made up of the most powerful and progressive elements from all civilized peoples.

\textbf{Supplementary Reading}

Graves, \emph{In Modern Times} (Macmillan, 1913), chaps. VI and VIII, and \emph{Great Educators} (Macmillan, 1912), chap. XIII; Parker, \emph{Modern Elementary Education} (Ginn, 1912), chap. XII. For the details of the life and work of Mann in brief form, read Hinsdale, B. A., \emph{Horace Mann and the Common School Revival} (Scribner, 1899), or the readable little work on \emph{Horace Mann the Educator} (New England Publishing Co., 1896) by Winship, A. E. Monroe, W. S., has briefly recounted \emph{The Educational Labors of Henry Barnard} (Bardeen, Syracuse, 1893), and a longer account of \emph{Henry Barnard} is that of Mayo, A. D., in \emph{Report of U. S. Commissioner of Education,} 1896-1897, vol.~I, chap. XVI. For the development of public education in the various parts of the country during this third period, see Martin, G. H., \emph{Evolution of the Massachusetts Public School System} (Appleton, 1894), lects. IV-VI; Steiner, B. C, \emph{History of Education in Connecticut (U. S. Bureau of Education, Circular of Information,} No.~2, 1893), chaps. III-V; Stockwell, T. B., \emph{History of Public Education in Rhode Island} (Providence Press Co., Providence, 1876), chaps. VI-X; Randall, S. S., \emph{History of the Common School System of the State of New York} (Ivison, Blakeman, Taylor, New York, 1871), third and fourth periods; Wickersham, J. P., \emph{History of Education in Pennsylvania} (Lancaster, Pennsylvania, 1886), chaps. XVII-XVIII; Mayo, A. D., \emph{The Development of the Common Schools in the Western States (Report of the U. S. Commissioner of Education,} 1898-99, vol.~I, pp.~357-450); Boone, R. G., \emph{History of Education in Indiana} (Appleton, 1892), chaps. IV and VIII-XXXIII; Smith, W. L., \emph{Historical Sketch of Education in Michigan} (Lansing, 1881), pp.~17-38, 49-57, and 78-109; Knight, E. W., \emph{The Influence of Reconstruction on Education in the South} (Columbia University, \emph{Teachers College Contributions,} No.~60,1913) and \emph{The Peabody Fund and Its Early Operation in North Carolina (South Atlantic Quarterly,} vol.~xiv. no. 2). Mayo, A. D., \emph{Education in the Several Status, Education oj the Colored Race,} and \emph{The Slater Fund (Report of the U. S. Commissioner of Education} 1894-95, XXX, XXXI, and XXXII).

Fig. 42. - The University of Michigan in 1855. (The oldest picture of the first prominent state university; established by the legislature in 1837, and opened in 1841.)

\hypertarget{development-of-education-practice}{%
\chapter{24. Development of Education Practice}\label{development-of-education-practice}}

\textbf{Outline}

\begin{itemize}
\tightlist
\item
  Of the two aspects to Pestalozzi's educational positions, Froebel stressed development from within, and Herbart development from without.
\item
  Through an early tutorial experience Herbart developed his pedagogy, but afterward invented an ingenious psychology upon which to base it. He undertook to show how the mind of the pupil is largely built up by the teacher, and he held to the moral aim of education. To accomplish this, he advocated `manysided interest,' and, while recognizing the value of both `historical' and `scientific' subjects, emphasized the former. But he also held that all subjects should be unified through `correlation,' and formulated the `formal steps of instruction.' The value of his work has been obscured by the formal interpretations of disciples, but he contributed greatly to the science of education. Herbartianism, developed by Ziller and others, spread throughout Germany; through the Herbart Society, it has greatly influenced educational content and methods in the United States.
\item
  Through his university environment, Froebel developed a mystic philosophy, but made it the basis of remarkable educational practices. He held to organic `unity' in the universe, and to the general method of `self-activity.' Besides this (1) `motor expression,' he also stressed (2) `social participation,' and attempted to realize both principles in (3) a school without books and set tasks, - the `kindergarten.' The training here has consisted chiefly in `play-songs,' `gifts,' and `occupations.' The chief weakness of Froebelianism is its mystic and symbolic theory, but it has comprehended the most essential laws of education at all stages. The kindergarten was spread through Europe largely by Baroness von Biilow, and through the United States by Elizabeth P. Peabody and others.
\item
  Few tendencies in educational practices to-day cannot be traced back for their rudimentary form to Herbart and Froebel, or their master, Pestalozzi.
\end{itemize}

\textbf{Froebel and Herbart as Disciples of Pestalozzi.} - In the discussion of observation and industrial training, we have noted the suggestions for improvement in educational practice that arose through Pestalozzi. While somewhat vague and based upon sympathetic insight rather than scientific principles, the positions of Pestalozzi not only left their direct influence upon the teaching of certain subjects in the elementary curriculum, but became the basis of the elaborate systems of Herbart and Froebel. These educators may be regarded as contemporary disciples of the Swiss reformer, who was born a generation before, but they continued his work along rather different lines. Each went to visit Pestalozzi, and it would seem from their comments upon what they saw that each found in the master the main principie which appealed to him and which he afterward developed more or less consistently throughout his work.

For there were two very definite aspects to Pestalozzi's positions, which may at first seem opposed to each other, but are not necessarily contradictory. On the one hand, Pestalozzi seems to have held that education should be a natural development from within; on the other, that it must consist in the derivation of ideas from experience with the outside world. The former point of view, which is apparent in his educational aim and definition of education (seep. 285), would logically argue that every characteristic is implicit in the child at birth in the exact form to which it is afterward to be developed, and that the teacher can at best only assist the child's nature in the efforts for its own unfolding. This attitude Pestalozzi apparently borrowed from the psychology implied in Rousseau's naturalism. The other conception, that of education as sense perception, which is evident in Pestalozzi's observational methods (see p.~286), depends upon the theory that immediate and direct impressions from the outside are the absolute basis of all knowledge, and holds that the contents of the mind must be entirely built up by the teacher. Some such naive interpretation has been common since speculation began, especially among teachers, and had been formulated before Pestalozzi's day by Locke, Hume, and others. In the main, Froebel took the first of these Pestalozzian viewpoints and rarely admitted the other, but the latter phase was developed by Herbart to the almost total disregard of the former. Hence we find that the one educator lays emphasis upon the child's development and activities, and the other concerns himself with method and the work of the teacher. The original contributions of both reformers to educational practice, however, were large, and are deserving of extended description.

\textbf{The Early Career and Writings of Herbart.} - Johann Friedrich Herbart (1776-1841) both by birth and by education possessed a remarkable mind, and was well calculated to become a profound educational philosopher. He came of intellectual and educated stock, and at the gymnasium and university displayed a keen interest in philosophy, Greek, and mathematics. Each of these subjects, too, was destined to play a part in his educational theories. Just before graduation (1797), however, Herbart left the university to become private tutor to the three sons of the governor of Interlaken, Switzerland, and during the next three years he obtained in this way a most valuable experience. The five extant reports that he made on the methods he used and on his pupils' progress reveal thus early the germs of his elaborate system. The youthful pedagogue seems to have recognized the individual variations in children, and to have shown a due regard for the respective ages of his pupils. He also sought, by means of his favorite work, the \emph{Odyssey}, to develop in them the elements of morality and a `many-sided interest.' This early experience, rather than his ingenious system of psychology and metaphysics, which he afterward developed in explanation, was the real foundation of his pedagogy, and furnished him with the concrete examples of the characteristics and individualities of children that appear in all his later works. He ever afterward maintained that a careful study of the development of a few children was the best preparation for a pedagogical career, and eventually made an experience of this kind the main element in his training of teachers.

While still in Switzerland, Herbart met Pestalozzi and was greatly attracted by the underlying principles of that reformer. He paid a visit to the institute at Burgdorf in 1799, and during the next two years, while at Bremen completing his interrupted university course, he undertook to advocate and render more scientific the thought of the Swiss educator. Here he wrote a sympathetic essay \emph{On Pestalozzi's Latest Writing, `How Gertrude Teaches Her Children,'} and made his interpretation of \emph{Pestalozzi's Idea of an A BC of Observation} (see p.~286). Next Herbart lectured on pedagogy at the University of Gottingen. The treatises he wrote there seem to have become more critical toward the Pestalozzian methods, and he no longer strives to conceal their vagueness and want of system. Sense perception, he holds with Pestalozzi, does supply the first elements of knowledge, but the material of the school course should be definitely arranged with reference to the general purpose of instruction, which is moral self-realization. This position on the moral aim of education he made especially explicit and complete in his work on \emph{The Science of Education} (1806).

\textbf{His Work at Konigsberg and Gottingen.} - In 1809 Herbart was called to the chair of philosophy at Konigsberg, and there established his now historic pedagogical seminary and the small practice school connected with it. The students, who taught in the practice school under the supervision and criticism of the professor, were intending to become school principals and inspectors, and, through the widespread work and influence of these young Herbartians the educational system of Prussia and of every other state in Germany was greatly advanced. In his numerous publications at Konigsberg, Herbart devoted himself chiefly to works on a system of psychology as a basis for his pedagogy. After serving nearly a quarter of a century here, he returned to Gottingen as professor of philosophy, and the last eight years of his life were spent in expanding his pedagogical positions. Here he issued the first edition of his \emph{Outlines of Educational Doctrine} (1835), which gives an exposition of his educational system when fully matured. It contains brief references to his mechanical metaphysics and psychology, but is a most practical and well-organized discussion of the educational process.

\textbf{Herbart's Psychology.} - Herbart's metaphysical psychology seems to have been an after-thought developed to afford a basis for the method of pedagogical procedure that he had worked out of his tutorial experience and his acquaintance with the Pestalozzian practice. But some explanation of this elaborate psychology may serve to make clearer his educational principles. For the most part he holds that the mind is built up by the outside world, and he is generally supposed to have left no place for instincts or innate characteristics and tendencies. With him the simplest elements of consciousness are `ideas,' which are atoms of mind stuff thrown off from the soul in endeavoring to maintain itself against external stimuli. Once produced by this contact of the soul with its environment, the ideas become existences with their own dynamic force, and constantly strive to preserve themselves. They struggle to attain as nearly as possible to the summit of consciousness, and each idea tends to draw into consciousness or heighten those allied to it, and to depress or force out those which are unlike. Each new idea or group of ideas is heightened, modified, or rejected, according to its degree of harmony or conflict with the previously existing ideas. In other words, all new ideas are interpreted through those already in consciousness. In accordance with this principle, which Herbart called `apperception,' the teacher can secure interest and the attention of the pupil to any new idea or set of ideas and have him retain it, only through making use of his previous body of related knowledge. Hence the educational problem becomes how to present new material in such a way that it can be `apperceived' or incorporated with the old, and the mind of the pupil is largely in the hands of the teacher, since he can make or modify his `apperception masses,' or systems of ideas.

\textbf{The Aim, Content, and Method of Education.} - Accordingly, Herbart holds that the purpose of education should be to establish moral and religious character. He believes that this final aim can be attained through instruction, and that, to determine how this shall furnish a `moral revelation of the world,' a careful study must be made of each pupil's thought masses, temperament, and mental capacity. There is not much likelihood of the pupil's receiving ideas of virtue that will develop into glowing ideals of conduct when his studies do not appeal to his thought systems and are consequently regarded with indifference and aversion. They must coalesce with the ideas he already has, and thus touch his life. But Herbart does not limit `interest' to a temporary stimulus for the performance of certain school tasks; he advocates the building up by education of certain broad interests that may become permanent sources of appeal in life. Instruction must be so selected and arranged as not only to relate itself to the previous experience of the pupil, but as also to reveal and establish all the relations of life and conduct in their fullness.

In analyzing this `many-sided interest,' Herbart holds that ideas and interests spring from two main sources, -- `experience,' which furnishes us with a knowledge of nature, and `social intercourse,' from which come the sentiments toward our fellowmen. Interests may, therefore, be classed as belonging to (1) `knowledge' or to (2) `participation.' These two sets of interests, in turn, Herbart divides into three groups each. He classed the `knowledge' interests as (a) `empirical,' appealing directly to the senses; (b) `speculative,' seeking to perceive the relations of cause and effect; and (c) `aesthetic,' resting upon the enjoyment of contemplation. The `participation' interests are divided into (a) `sympathetic,' dealing with relations to other individuals; (b) `social,' including the community as a whole; and (c) `religious,' treating one's relations to the Divine. Instruction must, therefore, develop all these interests, and, to correspond with the two main groups, Herbart divides all studies into two branches, - the (1) `historical,' including history, literature, and languages; and the (2) `scientific,' embracing mathematics, as well as the natural sciences. Although recognizing the value of both groups, Herbart especially stressed the `historical; on the ground that history and literature are of greater importance as the sources of moral ideas and sentiments. But, while all the subjects, 'historical' and `scientific,' are needed for a `many-sided interest,' and the various studies have for convenience been separated and classified by themselves, Herbart holds that they must be so arranged in the curriculum as to become unified and an organic whole, if the unity of the pupil's consciousness is to be maintained. This position forecasts the emphasis upon `correlation,' or the unification of studies, so common among his followers. The principle was further developed by later Herbartians under the name of `concentration,' or the unifying of all subjects around one or two common central studies, such as literature or history. But the selection and articulation of the subject matter in such a way as to arouse many-sidedness and harmony is not more than hinted at by Herbart himself. He specifically holds, however, that the \emph{Odyssey} should be the first work read, since this represents the interests and activities of the race while in its youth, and would appeal to the individual during the same stage. He would follow this with other Greek classics in the order of the growing complexity of racial interests depicted in them. This tentative endeavor of Herbart, in the selection of material for the course of study, to parallel the development of the individual with that of the race, was continued and enlarged by his disciples. It became especially definite and fixed in the `culture epochs' theory formulated by Ziller and others.

But to secure this broad range of material and to unify and systematize it, Herbart realized that it was necessary to formulate a definite method of instructing the child. This plan of instruction he wished to conform to the development and working of the human mind, and on the basis of what he conceived this activity to be, he mapped out a method with four logical steps: (1) `clearness,' the presentation of facts or elements to be learned; (2) `association,' the uniting of these with related facts previously acquired; (3) `system,' the coherent and logical arrangement of what has been associated; and (4) `method,' the practical application of the system by the pupil to new data. The formulation of this method was made only in principle by Herbart, but it has since been largely modified and developed by his followers. It was soon felt that, on the principle of `apperception,' the pupil must first be made conscious of the existing stock of ideas so far as they are similar to the material to be presented, and that this can be accomplished by a review of preceding lessons or by an outline of what is to be undertaken, or by both procedures. Hence Herbart's noted disciple, Ziller, divided the step of `clearness' into `preparation' and `presentation,' and the more recent Herbartian, Rein, added `aim' as a substep to `preparation.' The names of the other three processes have been changed for the sake of greater lucidity and significance by still later Herbartians, and the `five formal steps of instruction' are now given as (1) `preparation' (2) `presentation,' (3) `comparison and abstraction,' (4) `generalization,' and (5) `application.'

\textbf{The Value and Influence of Herbart's Principles.} - On all sides, then, as compared with Pestalozzi, Herbart was most logical and comprehensive. Where Pestalozzi obtained his methods solely from a sympathetic insight into the child mind, Herbart sought to found his also upon scientific principles. The former was primarily a philanthropist and reformer; the latter, a psychologist and educationalist. Pestalozzi succeeded in arousing Europe to the need of universal education and of vitalizing the prevailing formalism in the schools, but he was unable with his vague and unsystematic utterances to give guidance and efficiency to the reform forces he had initiated. While he felt the need of beginning with sense perception for the sake of clear ideas, he had neither the time nor the training to construct a psychology beyond the traditional one of the times, nor to analyze the way in which the material gained by observation is assimilated. Herbart, on the other hand, did create a system of psychology that, while fanciful and mechanical, worked well as a basis for educational theory and prao tice. In keeping with this psychology, he undertook to show how the ideas, which were the product of the Pestalozzian `observation,' were assimilated through `apperception,' and maintained the possibility of making all material tend toward moral development. This, he held, could be accomplished by use of proper courses and methods. In determining the subjects to be selected and articulated, he considered Pestalozzi's emphasis upon the study of the physical world to be merely a stepping-stone to his own `moral revelation of the world.'~While the former educator made arithmetic, geography,~natural science, reading, form study, drawing, writing, and music the object of his consideration, and is~indirectly responsible for the modern reforms in teaching these subjects, Herbart preferred to stress history, languages, and literature, and, through his followers, brought about improved methods in their presentation. He also first undertook a careful analysis of the successive steps in all instruction, and by his methodical principles did much to introduce order and system into the work of the classroom, although it is now known that his conception of the way in which the human mind works is hardly tenable.

A great drawback to the Herbartian doctrines is found in their formalization and exaggeration. For these~tendencies his enthusiastic and literal-minded followers, rather than Herbart himself, have probably been to blame. He was himself too keen an observer to allow his doctrines to go upon all fours. He is ordinarily credited by Herbartians with a psychology that takes no account of the innate characteristics of each mind, and holds that the mind is entirely built up by impressions~from the outside, but, while this is his main position, he~occasionally recognizes that there must be certain native~predispositions in the body which influence the soul in~one direction or another. This limitation of complete~plasticity by the pupil's individuality, and of the consequent influence of the teacher, causes him to perceive that ``in order to gain an adequate knowledge of each pupil's capacity for education, observation is necessary - observation both of his thought masses and of his physical nature.'' Again, while Herbart holds that every subject should, if possible, be presented in an attractive, interesting, and `almost playlike' way, he does not justify that `sugar-coated interest' which has so often put Herbartianism in bad odor. ``A view that regards the end as a necessary evil to be rendered endurable by means of sweetmeats,'' says he, ``implies an utter confusion of ideas; and if pupils are not given serious tasks to perform, they will not find out what they are able to do.'' Often, he realizes, ``even the best method cannot secure an adequate degree of apperceiving attention from every pupil, and recourse must accordingly be had to the voluntary attention, i. e., to the pupil's resolution.'' Moreover, `correlation' between different subjects, as well as between principles within the same subject, was advocated by Herbart, but he felt that the attempt to make such ramifications should not be unlimited. Further, while Herbart made some effort in shaping the course of study to parallel the development of the individual with that of the race, it was Ziller that erected this procedure into a hard and fast theory of `culture epochs.' But most common of all has been the tendency of his disciples to pervert his attempt to bring about due sequence and arrangement into an inflexible \emph{schema} in the recitation, and to make the formal steps an end rather than a means. Whereas, there is reason to believe that Herbart never intended that all these steps should be carried out in every recitation, but felt that they applied to the organization of any subject as a whole, and that years might even elapse between the various steps.

\textbf{The Extension of His Doctrines in Germany.} - At first the doctrines of Herbart were little known, but a quarter of a century after his death there sprang up two flourishing contemporary schools of Herbartianism. In its application of Herbart's theory, the school of Stoy for the most part held closely to the original form; but that headed by Ziller departed further and gave it a devdoiS\^{}and more extreme interpretation. Tuiskon Ziller (1817- popularized. 1882), both as teacher in a gymnasium and as professor at Leipzig, did much to popularize and develop Herbart's system. Through him was formed the Herbartian society known as the `Association for the Scientific Study of Education,' which has since spread throughout Germany. He it was that elaborated the doctrines of `correlation' and `concentration,' and first definitely formulated the'culture epochs' theory. ``Every pupil should,'' he writes, ``pass successively through each of the chief epochs of the general mental development of mankind suitable to his stage of development. The material of instruction, therefore, should be drawn from the thought material of that stage of historical development in culture, which runs parallel with the present mental stage of the pupil.'' All these principles Ziller worked out in a curriculum for the eight years of the elementary school, which he centered around fairy tales, \emph{Robinson Crusoe,~}and selections from the \emph{Old} and \emph{New Testaments.} He, moreover, developed Herbart's `formal stages of instruction' by dividing the first step and changing the name of the last.

Karl Volkmar Stoy (1815-1885), the founder of the other school, gave himself simply to a forceful restatement of the master's positions, but also established a most influential pedagogical seminary and practice school upon the original Herbartian basis at Jena. And eleven years later, Wilhelm Rein (1847- ), who had been a pupil of both Stoy and Ziller, succeeded the former in the direction of the practice school, and introduced there the elaborate development that had taken place since Herbart's time. He adopted Ziller's `concentration,' `culture epochs,' and other features, but made them a little more elastic by coordinating other material with the `historical' center in the curriculum. Through him Jena became known as the great seat of Herbartianism. Other Germans to develop the principles of Herbart have been Lange and Frick. The \emph{Apperception} of Karl Lange is an excellent combination of scientific insight and popular presentation. Otto Frick, director of the `Francke Institutions' at Halle (see p.~176), inclining more to the close interpretation of Stoy, devoted himself to applying Herbartianism to the secondary schools, and outlined a course for the gymnasium.

A throng of other German schoolmasters and professors have further adapted the doctrines of Herbart to school practice, and while their theories differ very largely from one another, from their common basis they are all properly designated `Herbartian.' As a result of this continuous propaganda, the content and methods of the school curricula in Germany have been largely modified. Herbart's emphasis upon the importance to the secondary schools of literary and historical studies as a moral training has been adapted to the elementary schools by the later Herbartians in the form of story and biographical material. History has consequently attained a more prominent place in the curriculum, and is no longer auxiliary to reading and geography. It is regarded as a means of moral development, and the cultural features in the~history of the German people are stressed more than the political. Ziller's plan for concentrating all studies about a core of history and literature, on the ground of thus producing `a moral revelation of the world' for the pupil, is in evidence everywhere. A twofold course, - Jewish history through Bible stories, and German history in the form of legends and tales, appears in every grade of the elementary school after the first two, and even in these lower classes there is some attempt to utilize literature as a moral training through the medium of fairy stories, fables, moral tales, \emph{Robinson Crusoe,} and the various stories of the philanthropinists (see p.~225).

\textbf{Herbartianism in the United States.} - Next to the land of its birth, the United States has been more influenced by Herbartianism than any other country. Before 1880 there were but few notices of Herbartianism in American educational literature, and not many appeared during the following decade. The movement was fostered largely by American teachers that were studying with Rein at Jena~during the last two decades of the century. Before 1890 nine Americans had taken their~degree there, and by the twentieth century more than fifty. These young men came back filled with the enthusiastic belief that Herbartian principles could supply a solution in systematic form for the many complicated problems with which American education was then grappling, and began at once to propagate their Northern faith. The movement centered chiefly in northern Illinois and was especially strong in the normal schools.

The staff of the State Normal University at this time included Charles DeGarmo, afterward professor of Education at Cornell, Frank M. McMurry, now of the Teachers College, Columbia University, and his brother, Charles A. McMurry, now of the faculty of the George Peabody College for Teachers; and the practice school at the Normal University was the first to be established upon Herbartian principles. The Schoolmasters' Club of Illinois gave much of its time to a discussion of Herbartian principles, and the first Herbartian literature in the United States was rapidly produced. During the last decade of the century there appeared large numbers of articles, textbooks, treatises, and translations, including \emph{The Method of the Recitation} and a variety of other works upon general and special methods by the McMurrys. In 1892 The Herbart Club was founded to promote a study of Herbartian principles and adapt them to American conditions, and during the first three years it spent its efforts in translating the words of Herbart The Herbart an(i in discussing Herbartian topics only. In 1895 the \emph{Year Book.} name of the club was changed to the Herbart Society for the Scientific Study of Education, many nonHerbartians were admitted, the scope of the discussions was enlarged, and the publication of a \emph{Year Book} was begun.

Then began the period of criticism and the formulation of American Herbartianism. The movement was vigorously opposed by many on the ground that it was Opposition, a foreign importation, was based upon absurd metaphysical presuppositions, or contained nothing new, but the disciples of Herbart stood valiantly by their guns. Although not always certain in their own minds, they endeavored to clear up all misunderstanding and confusion in the doctrines and to keep them practical through developing them in connection with actual experiments in teaching. They showed that the fanciful psychology of Herbart did not hold a determining place in his educational thought, and that it might be rejected, without affecting the merit of his pedagogy. One by one the doctrines were introduced in the order of their concreteness, - five formal steps, apperception, concentration,~interest - and little attempt was made to weave them into a single system. The critical season did not long endure, and the movement soon spread widely. By the close of the first year the Herbart Society had a membership of seven hundred, and the Herbartian principles were everywhere studied by local clubs and taught in schools and universities. In the report of the United States Commissioner of Education for 1894-1895, Dr.~Harris stated: ``There are at present more adherents of Herbart in the United States than in Germany.'' This, he believed, was due to the greater freedom of discussion that was allowed. The movement not only became an educational awakening, but it attained almost to the proportions of a cult. Moreover, many who hardly~considered themselves Herbartians undertook to modify and adapt the Herbartian principles, especially `correlation' and `concentration.' Francis W. Parker of Chicago, for example, among the phases of his educational practice (cf.~pp.~293 and 364), approached concentration so closely as to center the entire course of study around a hierarchy of natural and social sciences. And the Committees of Ten and Fifteen, appointed by the National Education Association to report upon secondary and elementary education respectively, showed a strong Herbartian influence in their recommendations of correlation.

Largely in consequence of the development of Herbartianism, an increased amount and larger utilization of historical material became general also in American elementary schools. A wide appreciation of the growth of morality, culture, and social life, rather than merely the development of patriotism, became the object in studying this subject. English and German history, as well as American, which alone was formerly taught, and sometimes Greek, Roman, and Norse, appear in the curricula of many elementary schools, and, instead of being confined to the two upper classes, historical material is often presented from the third grade up. Biographical and historical stories are largely employed in the lower classes, while in the upper some attempt is made to use European history as a setting for American. A similar development in the amount and use of literature also has appeared in the course of the elementary schools, partly as a result of the Herbartian influence. Instead of brief selections from the English and American writers, or the poorer material that formerly appeared in the school readers, complete works of literature have begun to be studied in the elementary curriculum, and a wide and rapid survey of the great English classics has been encouraged in the place of merely reading for the sake of oral expression. Even in the lowest grades some attempt to introduce the classics of childhood has been made.

While in these ways all elementary, and to some extent secondary, schools have been affected, Herbartianism pure and simple has largely been abandoned for less dogmatic methods. Even the Herbart Society has ceased to foster a propaganda, and has since 1901 dropped the first part of its name and been known as `The National Society for the Scientific Study of Education.' The later works of DeGarmo and Frank M. McMurry claim to be quite emancipated from Herbartianism. But, although professed Herbartians are now almost unknown in the United States, no other system of pedagogy, except that of Pestalozzi, has ever had so wide an influence upon American education and upon the thought and practice of teachers generally.

\textbf{Froebel's Early Life.} - Let us now turn to Froebel, the other great successor of Pestalozzi, and to his development and extension of the master's principle of `natural development.' Friedrich Wilhelm August Froebel (17821852) was born in a village of the Thuringian forest. He tells us that this environment started within him a search for the mystic unity that he believed to exist amid the various phenomena of nature, but it is more likely that this attitude was developed through a brief residence (1799-1800) at the University of Jena. The atmosphere about this institution was charged with the idealistic~philosophy, the romantic movement, and the evolutionary attitude in science. Froebel could not have escaped~the constant discussions upon the philosophy of Fichte~and Schelling. He must likewise have fallen under the spell of the Jena romanticists, - the Schlegels, Tieck, and Novalis. The advanced attitude in science at Jena may also have impressed the youth. While much of the science instruction failed to make clear that inner relation and mystic unity for which he sought, he must occasionally have caught glimpses of it in the lectures of professors belonging to the school of \emph{Natur-philosophie.}

\textbf{His Experiences at Frankfort, Yverdon, and Berlin.} - After leaving the university, Froebel was for four years groping for a niche in life. But he eventually (1805) met Anton Gruner, head of a Pestalozzian model school at~Frankfort, who persuaded him of his fitness for teaching and gave him a position in the institution. Here he undertook a systematic study of Pestalozzianism, and, through the use of modeling in paper, pasteboard, and wood with his pupils, he came to see the value of motor expression as a means of education. He then withdrew~to Yverdon and worked with Pestalozzi himself for two years (1808-1810). There he greatly increased his knowledge of the play and development of children, music, and nature study, which were to play so important a part in his methods. Next, he went to the University of Berlin to study mineralogy with Professor~Weiss, and through the work there he finally crystallized~his mystic law of `unity.' He became fully ``convinced~of the demonstrable connection in all cosmic development,'' and declared that ``thereafter my rocks and crystals served me as a mirror wherein I might discern mankind, and man's development and history.''

\textbf{The School at Keilhau.} - While at Berlin, he met his lifelong assistants, Langethal and Middendorf, and took them with him when he undertook the education of his five young nephews at Keilhau. Here he founded (1816) `The Universal German Institute of Education,' in which self-expression, free development, and social participation were ruling principles. Much of the training was obtained through play, and, except that the pupils were older, the germ of the kindergarten was already present. There was much practical work in the open air, in the garden about the schoolhouse, and in the building itself. The children built dams and mills, fortresses and castles, and searched the woods for animals, birds, insects, and flowers. To popularize the institute, Froebel published a complete account of the theory practiced at Keilhau in his famous \emph{Education of Man~}(1826). While this work is compressed, repetitious, and vague, and its doctrines had afterward to be corrected by experience, it contains the most systematic statement of his educational philosophy that Froebel ever made.

\textbf{Development of the Kindergarten.} - But the school at Keilhau was too radical for the times, and soon found itself in serious straits. Froebel then went to Switzerland~and, and for five years (1832-1837) continued his educational experiments in various locations there. While~conducting a model school at Burgdorf, it became obvious to him that ``all school education was yet without a proper initial foundation, and that, until the education of the nursery was reformed, nothing solid and worthy could be attained.'' The \emph{School of Infancy} of Comenius (see p.~171) had been called to his attention, and the educational importance of play had come to appeal to him more strongly than ever. He began to study and devise playthings, games, songs, and bodily movements~that would be of value in the development of small children, although at first he did not organize his materials into a system. Then, two years later, he returned to Germany, and established a school for children between the ages of three and seven. This institution was located at Blankenburg, two miles from Keilhau, one~of the most romantic spots in the Thuringian Forest, and was, before long, appropriately christened `Kindergarten' (i. e., garden in which children are the unfolding plants). Here he put into use the material he had invented in Switzerland, added new devices, and developed his system. The main features of this were the `play songs' for mother and child and the series of `gifts' and' occupations' (see pp.~3 58 f.). During his seven years in Blankenburg, he constantly expanded his material, and the accounts of these additions have been collected in~the works known generally as \emph{Pedagogics of the Kindergarten, Education by Development,} and \emph{Mother Play and Nursery Songs.}

While the kindergarten attracted considerable attention, Froebel's want of financial ability eventually compelled him to close the institution. After lecturing with much success for five years upon his system, he settled for the rest of his life near the famous mineral springs~at Liebenstein in Saxe-Meiningen. During this period~he obtained the friendship and support of the Baroness~von Marenholtz-Bulow, who brought a large number of people of distinction in the political and educational world to see his work in operation, and wrote most interesting \emph{Reminiscences} of Froebel's activities during the last thirteen years of his life. But owing to a confusion of his principles with the socialistic doctrines of his nephew, Karl, a decree was promulgated in Prussia by the minister of education, closing all kindergartens there. Froebel never recovered from this unjust humiliation, and died within a year.

\textbf{Froebel's Fundamental Concept of `Unity.'} - While Froebel's underlying principles go back to the developmental aspect of Pestalozzi's doctrines and even to Rousseau's naturalism, his conception of them, his imagery, and statement, seem to be a product of the idealistic philosophy, romantic movement, and scientific attitude of the day. These tendencies seem to have been assimilated by Froebel largely through his residence in Jena and Berlin. His conclusions as to educational theory and practice would have been possible as inferences from a very different point of view, but as he developed them logically and consistently with his metaphysical position, it may be of value to consider briefly the groundwork of the Froebelian philosophy. He regarded the `Absolute,' or God, as the self-conscious spirit from which originated both man and nature, and he consequently held to the unity of nature with the soul of man. His fundamental view of this organic unity appears in his general conception of the universe: ``In all things,'' says he, ``there lives and reigns an eternal law. This all-controlling law is necessarily based on an all-pervading, energetic, living, self-conscious, and hence eternal Unity. This Unity is God. All things have come from the Divine Unity, from God, and have their origin in the Divine Unity, in God alone. All things live and have their being in and through the Divine Unity, in and through God. The divine effluence that, lives in each thing is the essence of each thing.'' This fundamental mystic principle Froebel constantly~reiterates in various forms, and from it derives a number~subsidiary conceptions. These, however, play but~a small part in his actual practice, and scarcely require~consideration here.

\textbf{Motor Expression as His Method.} - But Froebel also holds that, ``while in every human being there lives humanity as a whole, in each one it is realized and expressed in a wholly particular, peculiar, personal, and unique manner.'' Thus he maintains that there is in every person at birth a coordinated, unified plan of his mature character, and that, if it is not marred or interfered with, it will develop naturally of itself. While he is not entirely consistent, and at times implies that this natural~development must be guided and even shaped, in the~main he reiterates Rousseau's doctrine that `nature is right,' and clearly stands for a full and free expression of the instincts and impulses. Hence he insists that ``education in instruction and training should necessarily be \emph{passive, following; not prescriptive, categorical, interfering."} But in his conclusion as to the proper method for accomplishing this `development,' Froebel naturally holds that it''should be brought about not in the way of dead imitation or mere copying, but in the way of living, spontaneous self-activity." By this principle of~`self-activity' as the method of education Froebel seeks not simply activity in response to suggestion or instruction from parents or teachers, but activity of the child in carrying out his own impulses and decisions. Individuality must be developed by such activity, and selfhood given its rightful place as the guide to the child's powers when exercised in learning. Hence with this~idea of development through `self-activity' is connected~his principle of `creativeness,' by which new forms and~combinations are made and expression is given to new~images and ideas. ``Plastic material representation in~life and through doing, united with thought and speech,''~he declares, ``is by far more developing and cultivating~than the merely verbal representation of ideas.''

\textbf{The Social Aspect of Education.} - His emphasis upon this psychological principle of motor expression under the head of `self-activity' and `creativeness' is the chief characteristic of Froebel's method. Rousseau had also recommended motor activity as a means of learning, but he had insisted upon an isolated and unsocial education for Emile, whereas Froebel stresses the social aspects of education quite as clearly as he does the principle of self-expression. In fact, he holds that increasing self-realization, or individualization through `self-activity,' must come through a process of socialization. The social instinct is primal, and the individual can be truly educated only in the company of other human beings. The life of the individual is necessarily bound up with participation in institutional life. Each one of the various institutions of society in which the mentality of the race has manifested itself - the home, the school, the church, the vocation, the state - becomes a medium for the activity of the individual, and at the same time a means of social control. As far as the child enters into the surrounding life, he is to receive the development needed for the present, and thereby also to be prepared for the future. Through imitation of cooperative activities in play, he obtains not only physical, but intellectual and moral training. Such a moral and~intellectual atmosphere Froebel sought to cultivate at Keilhau by cooperation in domestic labor, - `lifting, pulling, carrying, digging, splitting,' and through cooperative construction out of blocks of a chapel, castle, and other features of a village. Similarly, the kindergarten was intended to ``represent a \emph{miniature state} for children, in which the young citizen can learn to move freely, but with consideration for his little fellows.''

\textbf{The Kindergarten.} - Beside his basal principles of motor expression and social participation, Froebel made a third contribution to educational practice in advocating as a means of realizing these principles a school without books or set intellectual tasks, and permeated with play,~freedom, and joy. In the kindergarten, `self-activity' creativeness,' together with social cooperation, found complete application and concrete expression. The training there has always consisted of three coordinate forms of expression: (1) song, (2) movement and gesture, and (3) construction; and mingled with these and growing out of each is the use of language by the child. But these means, while separate, often cooperate with and interpret one another, and the process is connected as an organic whole. For example, when the story is told or read, it is expressed in song, dramatized in movement and gesture, and illustrated by a construction from blocks, paper, clay, or other material.

The \emph{Mother Play and Nursery Songs} were intended to exercise the infant's senses, limbs, and muscles, and, through the loving union between mother and child, draw both into intelligent and agreeable relations with the common objects of life about them. The fifty `play songs' are each connected with some simple nursery game, like `pat-a-cake,' `hide-and-seek,' or the imitation of some trade (Fig. 43), and are intended to correspond to a special physical, mental, or moral need of the child. The selection and order of the songs were determined with reference to the child's development, which ranges from almost reflex and instinctive movements up to an ability to represent his perceptions with drawings, accompanied by considerable growth of the moral sense. Each song contains three parts: (1) a motto for the guidance of the mother; (2) a verse with the accompanying music, to sing to the child; and (3) a picture illustrating the verse.

The `gifts' and `occupations' were both intended to~stimulate motor expression, but the `gifts' combine and rearrange certain definite material without changing the form, while the `occupations' reshape, modify, and transform their material. The emphasis in kindergarten practice has come to be transferred from the `gifts' to the `occupations,' which have been largely increased in range and number. Of the `gifts,' the first consists of a fast, box of six woolen balls of different colors. They are to be rolled about in play, and thus develop ideas of color, material, form, motion, direction, and muscular sensibility. A sphere, cube, and cylinder of hard wood compose the second `gift.' Here, therefore, are found a second, known factor in the sphere and an unknown one in the cube. A comparison is made of the stability of the cube with the movability of the sphere, and the two are harmonized in the cylinder, which possesses the characteristics of each. The third `gift' is a large wooden cube third, divided into eight equal cubes, thus teaching the relations of the parts to the whole and to one another, and~making possible original constructions, such as armchairs, benches, thrones, doorways, monuments, or steps. The three following `gifts' divide the cube in various ways so as to produce solid bodies of different types and sizes, and excite an interest in number, relation, and form. From them the children are encouraged to construct geometrical figures and `forms of beauty' or artistic designs. Beside the six regular `gifts,' he also added `tablets,' `sticks,' and `rings,' sometimes known as `gifts seven to nine.' This material introduces surfaces, lines, and points in contrast with the preceding solids, and brings out the relations of area, outline, and circumference to volume. The `occupations' comprise a long list of constructions with paper, sand, clay, wood, and other materials. Corresponding with the `gifts' that deal with solids, may be grouped `occupations' in clay modeling, cardboard cutting, paper folding, and wood carving; and with those of surfaces may be associated mat and paper weaving, stick shaping, sewing, bead threading, paper pricking, and drawing.

\textbf{The Value and Influence of Froebel's Principles.} - For one pursuing destructive criticism only, it would not be difficult to find flaws in both the theory and practice of Froebel. In the \emph{Mother Play} the pictures are rough and poorly drawn, the music is crude, and the verses are lacking in rhythm, poetic spirit, and diction (Fig. 43). But the illustrations and songs served well the interests and needs of those for whom they were produced, and Froebel himself was not insistent that they should be used after more satisfactory compositions were found. Other criticism of his material has been made on the ground that it was especially adapted to German ideals, German

Fig. 43. - \emph{Der Zimmermann} (The Carpenter).

(Reproduced by permission ol D. Appleton and Company frora the Eliot and Blow edition of FroebePs \emph{Mother Play.)}

children, and the relatively simple village life of Froebel's experience, and that it needs considerable modification~to suit other countries and the industrial organization of society to-day. Also the argument of `formal discipline' for care and accuracy in the use of the gifts, and the insistence upon the employment of every part of each gift upon all occasions in the exact order mentioned by Froebel, have been shown to violate the principles of modern psychology. His more liberal disciples, however, realize that it is the spirit of his underlying principles, and not the letter of his practice, that should be followed, and have constantly struggled to keep the kindergarten matter and methods in harmony with the times and the environment.

A more serious hindrance to the acceptance of Froebelianism has arisen from his peculiar mysticism and~symbolism. Since all things live and have their being in and through God and the divine principle in each is the essence of its life, everything is liable to be considered by Froebel as symbolic in its very nature, and he often resorts to fantastic and strained interpretations. Thus with Froebel the cube becomes the symbol of diversity in unity, the faces and edges of crystals all have mystic meanings, and the numbers three and five reveal an inner significance. At times this symbolism descends~into a literal and verbal pun, where it seems to a modern that Froebel can hardly be in earnest. Further, he holds that general conceptions are implicit in the child, and each of these can be awakened by `adumbration,' that is, by presenting something that will symbolically represent that particular `innate idea.' Thus, in treating the gifts and games, he maintains that from a ball the~pupils gather an abstract notion of `unity.' Moreover, because God is the self-conscious spirit that originated both man and nature, and everything is interconnected, he believes that each part of the universe may throw light on every other part, and constantly holds that a knowledge of external nature, - such as the formation of crystals, will enable one to comprehend the laws of the mind and of society.

Unfortunately, this mystic symbolism, vague and extreme as it is, is regarded by the strict constructionists among the kindergartners as the most essential feature in Froebelianism, and they expect the innocents in their charge to reveal the symbolic effect of the material upon their minds. There is no real evidence for supposing that such associations between common objects and abstract conceptions exist for children. But such an imaginary symbolic meaning may be forced upon an object by the teacher, and pupils in conservative kindergartens soon learn to adopt certain phrases and attitudes that imply such mystic meaning. This often tends to foster insincerity and sentimentalism rather than to inculcate abstract truth through symbols. Had Froebel possessed the enlarged knowledge of biology, physiology, and psychology that is available for one living in the twentieth century, it is unlikely that he would have insisted upon the symbolic foundations for his pedagogy. His excellent practice is heavily handicapped by these interpretations, and might as easily have been inferred from very different positions in modern psychology.

But Froebel has had a most happy effect upon education as a whole. In some respects he utilized features from other reformers. We can see that he adopted many of Pestalozzi's objective methods in geography, natural history, arithmetic, language, drawing, writing and reading, and constructive geometry; reiterated Rousseau's views upon the infallibility of nature; and advocated the physical training and excursions as a means of study that are stressed by both these reformers. In his use of stories, legends, fables, and fairy-tales, he paralleled his contemporary, Herbart, in his influence upon the curriculum. But in his emphasis upon motor expression and social~participation, together with his advocacy of a school~without books or set tasks, Froebel was unique, and made a most distinctive contribution to educational practice. And whenever the real significance of his principles has been comprehended, they have been recognized as the most essential laws in the educational process, and are valued as the means of all effective teaching.

Froebel himself never fully worked out his theories in connection with schooling beyond the kindergarten, but all stages of education have now come to realize the value of discovering and developing individuality by means of initiative, execution, and cooperation; and spontaneous activities, like play, construction, and occupational work, have become more and more the means to this end. For example, the `busy work,' `whittling,' `claymodeling,' `sloyd,' and other types of `manual training' have to a large degree sprung from the influence of Froebel. Uno Cygttaeus (1810-1888), who started the Cygnaeus manual training movement, owed his inspiration to Froebel and his own desire to extend the kindergarten occupations through the grades. As a result of his efforts, Finland in 1866 became the first country in the world to~adopt manual training as an integral part of the course in the elementary and teacher training schools. In 1874, through the visit of Otto Salomon (1849-1007) to Cygnaeus, Sweden transformed its sloyd from a system of teaching the elements of trades to the more educative method of manual training. This use of constructive and occupational work for educational purposes rather than for industrial efficiency soon spread throughout Europe, and was first suggested to the United States by the Centennial Exposition of 1876 at Philadelphia. Various types of modern educational theory and practice, especially those associated with experiments made in the United States, also reveal large elements of Froebelian influence. Among these might be included the work of Colonel Parker (Fig. 40) and of Professor John Dewey. The Froebelian emphasis upon motor expression, the social aspect of education, and informal schooling are evident throughout Parker's work in his elementary school, and are even extended so as to include speech and the language-arts. Similarly, Dewey's occupational work and industrial activities, which were used through the entire course of his `experimental school' in Chicago, although not copied directly from Froebel, closely approached the modified practice of the kindergarten (see pp.~430 f.).

\textbf{The Spread of Froebelianism through Europe.} - Directly after the death of Froebel, the kindergarten began to be spread through his devoted followers, especially Baroness von Bulow. By means of her social position and knowledge of modern languages, she was enabled to become his great apostle throughout Europe. Having failed to obtain a revocation of the edict against~the kindergarten (see p.~355) in Prussia, the baroness turned to foreign lands. She visited France, Belgium, Holland, Italy, Russia, and nearly every other section of Europe, and in 1867 was invited to speak before the ' Congress of Philosophers' at Frankfort. This distinguished gathering had been called to inquire into contemporary educational movements, and after her elucidation of Froebelianism, a standing committee of the Congress, known as the `Froebel Union,' was formed to study the system.~The propaganda was soon everywhere eagerly embraced. Kindergartens, training schools, and journals devoted to the movement rapidly sprang up. While the kindergarten was not generally adopted by the governments, it was widely established by voluntary means throughout Western Europe, and has since met with a noteworthy growth. Instruction in Froebelian principles is now generally required in most normal and teacher training institutions there. Sometimes, as in France and England,~it has been combined with the infant school movement, and has lost some of its most vital characteristics, but even in these cases the cross-fertilization has afforded abundant educational fruitage. Only in Germany, the native land of the kindergarten, has serious hostility to the idea remained. Kindergartens have, with few exceptions, never been recognized there as genuine schools or part of the regular state system. Even to-day the German kindergarten is regarded as little more than a day nursery or convenient place to deposit small children and have them amused.

\textbf{The Kindergarten in the United States.} - The development and influence of the kindergarten have been more marked in the United States than in any other~country. First attempts at a kindergarten in America were made shortly after the middle of the nineteenth century by educated Germans, who had emigrated to America because of the unsettled conditions at home. A more fruitful attempt was that of Elizabeth P. Peabody at Boston in the early sixties. Notwithstanding the immediate success of this institution and the evident enjoyment of the children, Miss Peabody felt that she had not succeeded in getting the real spirit of Froebel, and in 1867 she went to study with his widow, who had been settled in Hamburg for several years. Upon her return the following year Miss Peabody corrected the errors in her work and established a periodical to explain and spread Froebelianism. The remainder of her life was spent in interesting parents, philanthropists, and school boards in the movement, and a service was done for the kindergarten in America almost equal to that of Baroness von Biilow in Europe. In 1868 through Miss Peabody the first training school for kindergartners in the United States was established at Boston. A similar institution was opened in New York by 1872 in charge of Maria Bolte, who had also studied with Frau Froebel. The same year saw the beginning of Susan E. Blow's work in St.~Louis, where her free training school for kindergartners was opened. Another missionary effort began in 1876 through Emma Marwedel, who was employed to organize voluntary kindergartens and training classes throughout the chief centers of California. The kindergarten movement grew rapidly. Between 1870 and 1890 in all the leading cities of the country subscriptions for kindergartens were raised by various philanthropic agencies, and by the close of the century there were about five hundred such voluntary associations.

But private foundations are restrictive, and it was not until the kindergarten began to be adopted by school systems that the movement became truly national in the United States. Boston in the early seventies added a few kindergartens to her public schools, but after several years of trial gave them up on account of the expense.~The first permanent establishment under a city board~was made in 1873 at St.~Louis through the efforts of Miss Blow. Twelve kindergartens were organized at first, but others were opened as rapidly as competent directors could be prepared at Miss Blow's training school. Within a decade there were more than fifty public kindergartens and nearly eight thousand pupils in St.~Louis. San Francisco authorized the addition of kindergartens to the public schools in 1880; and between that date and the end of the century New York, Boston, Philadelphia, Buffalo, Pittsburgh, Rochester, Providence, Milwaukee, Minneapolis, and nearly two hundred other progressive cities made the work an integral part of their system. About twenty of the cities employed a special supervisor to inspect the work. Excellent training schools for kindergartners are now maintained by half a hundred public and quasi-public normal institutions.

\textbf{The Relative Influence of Pestalozzi, Herbart, and Froebel.} - It is now obvious how large a part in the development of modern educational practice has been played by Herbart and Froebel. There are few tendencies in the curricula and methods of the schools today that cannot in their beginnings be traced back to them, or to Pestalozzi, their master. But the reforms of all three find their roots in Rousseau (Fig. 44). His `naturalism' was continued by Pestalozzi (Fig. 45) in his `development' and `observation,' which were, in turn, further elaborated by Froebel and Herbart respectively (Figs. 47 and 46). Through his `observation' methods, Pestalozzi greatly improved the teaching of arithmetic, language work, geography, elementary science, drawing, writing, reading, and music, and, by means of Fillenberg's work, developed industrial and philanthropic training. As a result of Herbart's moral and religious aim, marked advances in the teaching of history and literature have taken place, and, largely through his carefully wrought educational doctrines, order and system have everywhere been introduced into instruction. From Froebel's mystic interpretation of `natural development' we have obtained the kindergarten training for a period of life hitherto largely neglected, the informal occupations, manual training, and other studies of motor expression, together with psychological and social principles that underlie every stage of education. Pestalozzi's reforms were felt in Europe throughout the first half of the nineteenth century, but did not have any wide effect upon the United States until after the `Oswego movement' in the sixties. The influence of Froebel appeared in Europe shortly after the middle of the century, and began to rise to its height in America about 1880. The Herbartian theory and practice became popular in Germany between 1865 and 1885, while the growth of Herbartianism in the United States began about five years after the latter date. Hence the development of modern educational practice, due to these three great reformers, falls distinctly within the period of the nineteenth century.

\textbf{Supplementary Reading}

Graves, \emph{In Modern Times} (Macmillan, 1913), chap. VII; \emph{Great Educators} (Macmillan, 1912), chaps. X and XI; Monroe, \emph{Textbook} (Macmillan, 1905), pp.~622-673; Parker, \emph{Modern Elementary Education} (Ginn, 1912), chaps. XVII and XVIII. Herbart's \emph{Science of Education} (translated by Felkin), and \emph{Outlines of Educational Doctrine} (translated by Lange and De Garmo, Macmillan, 1909), and Froebel's \emph{Education of Man} (translated by Hailmann; Appleton, 1894), \emph{Pedagogics of the Kindergarten} and \emph{Education by Development} (translated by Jarvis; Appleton, 1897 and 1899), and \emph{Mother Play} (translated by Eliot and Blow, Appleton, 1896), should be read at least cursorily. The best brief treatise on \emph{Herbart and Herbartianism} (Scribner, 1896) is that by De Garmo, C., a graphic description of \emph{The Herbarlian Psychology} (Heath, 1898) is given by Adams, J., in chap. III, and a history of \emph{The Doctrines of Herbart in the United States} as a doctoral dissertation (University of Pennsylvania) by Randels, G. B. A good account of \emph{Froebel and Education by Self-Activity} (Scribner, 1897) has been furnished by Bowen, H. C.; a conservative treatment of \emph{Kindergarten Education (Education in the United States,} edited by N. M. Butler, Monograph No.~1), by Blow, Susan E.; an interesting treatise on \emph{Kindergarten in American Education} (Macmillan, 1908), by Vandewalker, Nina C; and a critical account of \emph{The Psychology of the Kindergarten (Teachers College Record,} vol.~IV, pp; 377-408), by Thorndike, E. L.

\hypertarget{the-development-of-modern-systems-of-education}{%
\chapter{25. The Development of Modern Systems of Education}\label{the-development-of-modern-systems-of-education}}

\textbf{Outline}

\begin{itemize}
\tightlist
\item
  The leading states of Western Europe and of Canada have, during the past century and a half, organized systems of education, which may prove suggestive.
\item
  In Prussia, owing to a strong line of monarchs, state control has taken the place of ecclesiastical through a series of decrees and enactments. The people's schools are quite separate from the secondary schools. Three types of secondary institutions have developed, - the `gymnasium,' with the classics as staples; the `real-school,' with modern languages and sciences; and the `realgymnasium,' with its compromise between the other two. The universities have likewise been emancipated from ecclesiastical control.
\item
  In France, a highly centralized system has been developed. Napoleon united secondary and higher education in a single corporation; under Louis Philippe, an organization of elementary schools was made; and, during the third republic, elementary education has been made free, compulsory, and secular. The present secondary system - \emph{lyctes} and communal colleges - began with Napoleon, and has now been differentiated into several courses. One-half of the universities established by Napoleon were suppressed during the Restoration, but since 1896 there has been a university in each of the sixteen `academies,' save one.
\item
  In England national education has grown out of the conflict of a number of social elements. The sentiment for universal training appeared toward the close of the eighteenth century, but not until 1870 were `board schools' established. In 1899 a central Board of Education was created; and the Act of 1902, while permitting voluntary schools to share in the local rates, unified the system and established secondary education at public expense. During the nineteenth century also the classical and ecclesiastical monopoly in secondary and higher education was largely broken.
\item
  In Canada there have developed two types of educational control, - (1) the closely centralized system of public schools in Ontario, and (2) the public supervision of ecclesiastical schools In Quebec.
\end{itemize}

National Systems of Education in Europe and Canada. - In previous chapters (XVII, XXI, XXIII) we have witnessed the gradual evolution in America of state systems of universal education out of the unorganized and rather aristocratic arrangement of schools that had first been transplanted from Europe in the seventeenth century. But development of a centralized organization of public schools has not been confined to the United States. During the past century and a half, the leading powers of Western Europe and Canada have likewise organized state systems of education, similar in some respects to those of the American union. All of these states educatfonfree, have now established universal elementary education free oF1\^{}tuitous to all, although as yet in few instances are secondary g\^{}f\textsubscript{s}\textsuperscript{ar}\^{}\textsubscript{nd} schools also gratuitous, and only Canada has welded her France alone

, - , secularized.

elementary and secondary systems. France alone has completely secularized its system, but the public schools of the other nations, while still including religious instruction, have been emancipated from ecclesiastical control, and are responsible to the civil authorities. In all of them school attendance is compulsory. Yet the educational system in none of these countries is identical with that in the United States, but has been adapted in each case to the genius and social organization of the people concerned. Its characteristics must, therefore, be considerably modified, in order to be utilized or to prove suggestive to the United States or other nations, and can be understood only in the light of the educational history of the particular country to which it bewfen\^{}under- longs. For an intelligent appreciation of these modem toricaiiy\textsuperscript{8} school systems, we must, therefore, trace the gradual development to their present form in response to the changing ideals of successive periods.

The Beginning of State Control in Prussia. - We may look first at Germany. Up to the later years of the eighteenth century all stages of education in the various German states remained almost entirely under ecclesiastical control, but during this period the schools and universities were taken over by the state from the church, although the clergy still exercised a few prerogatives, and centralized national systems were gradually organized. Among these states of Germany the first and most influential in the organization of universal education was Prussia. While each of the others is characterized by an educational history and peculiarities of its own, this state may be taken as an illustration of the Rise of Prus- evolution of German school systems. The rise of Prussia,

sian education

due to eniight- educationally as well as politically, seems to have been \textsuperscript{cspo}'' due to the strong Hohenzollern monarchs, - despotic, but thoroughly awake to the interests of their people. Although for nearly two centuries state control of education was carried on more or less through the medium of the church, its development was well under way by the seventeenth century. While the `consistory,' or board of supervision, was still composed largely of the clergy, the schools were soon (1687) declared not to be simply church organizations, but to belong to the state, and some attempt was made to extend schools to the villages as well as cities. But the first noteworthy attempt to establish compulsory attendance occurred during the reign of Frederick William I. In 1717 that monarch decreed that, wherever schools existed, children £ornpu1S\textgreater{}TM \textsuperscript{f}°\textsuperscript{r} should be required to attend during the winter, and in \textbar{}*\^{}\^{}"Jfvn the summer whenever they could be spared by their Bam 1 in 1717; parents, which must be at least once a week. He also founded the first teachers' seminary at Stettin from his own private means (1735), and the next year had a definite law passed, making education compulsory for children from six to twelve years of age.

Educational Achievements of Frederick the Great. - His most important contribution, however, consisted in preparing the way for an educational movement that was to be greatly developed through his more able son, Frederick the Great. Frederick began by improving the administration of secondary education, and requiring that all vacancies on crown lands be filled by graduates from Hecker's normal school at Berlin. But the great step toward a national system was taken in 1763, when Frederick issued his \emph{General School Regulations for} W \emph{General}

\ldots{}\ldots{}. \emph{School Regula}

\emph{ble Country.} This decree required children to attend \emph{tions} decreed

school from five until thirteen or fourteen, and until in 1763, they ``know not only what is necessary of Christianity, fluent reading, and writing, but can give answer in every- thing which they learn from the school books prescribed and approved by our consistory.'' If any pupils should arrive at this state of proficiency before thirteen or four- teen, they could even then leave school only through the official certification of the teacher, minister, and

inspector. Provision was also made for the attendance of children who had to herd cattle or were too poor to pay the school fees. Sunday continuation schools were to be established for young people beyond the school age. Teachers must have attended Hecker's seminary and had to be examined and licensed by the inspector. This decree was two years later supplemented with by\textsuperscript{P}ReT/o/i\textsuperscript{?}TO5 \textsuperscript{s}TMil\textsuperscript{ar} \emph{Regulations for the Catholic Schools in Silesia, Schools\^{}*"} drawn up by Abbot Felbiger. The carrying out of the decree was, however, stubbornly opposed by many teachers, who could not meet the new requirements; by farmers, who objected to the loss of their children's time; and by the nobles, who feared the discontent and uprising of the peasants, in case they were educated. The execution of the regulation was still in the power of the clergy, and for some time it proved but little more than a pious wish. But Frederick strove hard to have it enforced, and it became the foundation for the more effective laws that have since become embodied in the Prussian school system.

Educational Influence of Zedlitz. - After 1771 the educational work of Frederick was substantially aided by the appointment of Baron von Zedlitz as head of the Department of the Lutheran Church and School Affairs. This great minister had been much impressed by Basedow's principles and experiments and by Rochow's application of the `naturalistic' training, and through him village schools were greatly strengthened and enriched, a regular normal school was opened at Halberstadt, and the humanistic ideal of secondary education revived.

(3) Establishment of Central Board of

Administra- under the reactionary monarch, Frederick William II,

ment of Cen- A year after Frederick's death Zedlitz succeeded, even

tral Board of"

in further developing the nationalization of education, tion under

Frederick Wil

In 1787 an \emph{Oberschulcollegium,} or central board of school Ham 11 in administration, was appointed instead of the former \textsuperscript{I7} \textsuperscript{7}' church consistories. However, while the organization was supposed to be made up of educational experts, and Zedlitz was actually made chairman, the membership was mostly filled from the clergy, and the king refused to extend its jurisdiction to the higher schools.

Despite the reactionary policy of the sovereign, the culmination of the attempts to establish a national nonsectarian system of education occurred during this reign. In 1794 there was published the \emph{General Code,} in which £f Ge\^{}\^{}aSe the chapter upon education declared unequivocally that\textsuperscript{in} '794; ``all schools and universities are under the supervision of the state, and are at all times subject to its examination and inspection.'' Teachers were, therefore, not to be chosen without the consent of the state, and where their appointment was not vested in particular persons, it was to belong to the state. Teachers of all secondary schools were to be regarded as state officials. No child was to be excluded from the schools because of his religion, nor compelled to stay for religious instruction when it differed from the belief in which he had been brought up.

Foundation of the Ministry of Education and Further Progress. - While this comprehensive code met with much opposition from the clergy and the ignorant masses, and the next king, Frederick William III, weakly yielded at first, the humiliation of Prussia by Napoleon (1803) brought the country to a realization of the need of a centralized organization of the school system. The \emph{Oberschulcollegium} was abolished, to get rid of the clerical\protect\hypertarget{ch25.xmlux5cux23para.451.1.0.box.85.234.1259.42.q.60}{}{ (5) Creation of domination that had crept in, and a Bureau of Educa}

a Bureau of , . \_ .

Education in tion was created as a section of the Department of the Inlater'became a terior in 1807. The Bureau was within a decade erected istryandthen into a separate Department or Ministry of Education, organized** Eight years later (1825) the state was divided into educational provinces; and a \emph{Schulcollegium,} or administrative board, with considerable independence, but subject to the minister, was established over each province. Since then there have been many further developments, and provinces themselves are now divided into `governments,' each of which has a' school commission' over it, and every government is divided into `districts,' whose chief officer is a `school inspector.' Under the district inspector are local inspectors, and each separate school also has a local board, to take charge of repairs, supplies, and other external matters.

Thus the supreme management of the schools has been gradually coming into the hands of the state for nearly two centuries. The decrees of 1717 and 1763, the establishment of the \emph{Oberschulcollegium} in 1787, the General Code promulgated in 1794, the foundation of a distinct civic administration of education in 1807, are the milestones that mark the way to state control. But, while the influence of the church has been constantly diminishing, many of the board members are ministers or priests and the inspectors come mostly from the clergy. Moreover, religious instruction forms part of the course in every school, although it is given at such an hour that any pupil may withdraw if. the teaching is contrary to the faith in which he has been reared. The secondary schools are largely interdenominational, but in elementary education there are separate schools\protect\hypertarget{ch25.xmlux5cux23para.452.1.0.box.111.220.1047.88.q.70}{}{ for Catholics and Protestants, alike supported by the state.}

The Elementary System. - Prussia, like most of the principal states of Europe, as a result of their educational history, has its elementary and secondary systems quite separate and distinct from each other (Fig. 48). The universities continue the work of the gymnasiums and real-schools, but these two latter institutions parallel the work of the \emph{Volksschiden} (people's schools), rather than \emph{Voiksschacn,} supplement it. The course of the secondary school ordinarily occupies the pupil from nine to eighteen years of age, while that of the elementary school carries him from six to fourteen, and after the first three years it is practically impossible to transfer from the elementary to the secondary system. A pupil cannot enter a gymnasium or real-school after completing the people's school, and the only further training he can obtain is that of the \emph{Fortbildungschulen,} or `continuation schools,' which supple- 'Continuation

schools *

ment the system (seep. 420). The people's schools are gratuitous and are attended mostly by the children of the lower classes, while the gymnasiums charge a tuition fee and are patronized by the professional classes and aristocracy. Hence the line between elementary and secondary education in Prussia is longitudinal and not latitudinal, as it is in the United States; the distinction is one of wealth and social status rather than of educational grade and advancement. There are also some \emph{Mittelschulen} (middle schools) for the middle \textsuperscript{and} \emph{urna}classes of people, who cannot send their children to the secondary schools, and yet can afford some exclusiveness. They have one more class than the people's schools, include a foreign language during

the last three years, and require teachers with a better training. The Secondary System. - The main types of secondary foSSSi,\^{} \textsuperscript{schools} \textsuperscript{in} Prussia are the \emph{Gymnasien} (see p.~114), with the classic languages as the main feature of their course, and the \emph{Realschulen,} or real-schools (see p.~176), characterized by larger amounts of the modern languages, mathematics, and the natural sciences. For more than a century after the first real-school was opened in Berlin by Hecker (1747), this type of institution had only six years in its course, and was considered inferior to the gymnasium. By the ministerial decree of 1859, however, two classes of real-schools were recognized, and those of the first class had a course of nine years, and included Latin, but not Greek. They were given full standing as secondary schools, and graduates were granted admission to the universities, except for the study of theology, medicine, or law. The course of the second class of these institutions contained no Latin, and was but six years in length. In 1882 the compromise character of the course of the first class of institutions led to their being designated as \emph{Realgymnasien,} while the second class in some instances had their work extended to nine years and became known as \emph{Oberrealschulen.} The graduates were then allowed the privilege of studying at the universities in mathematics and the natural sciences. Since 1901 the university courses have been thrown open to graduates of any of the three types of secondary schools, except that, to be eligible for theology, one must have completed the course of a gymnasium, and for medicine, the course of a real-gymnasium at least. Besides these schools that have been mentioned, in rural

\emph{Realgymna sien} and \emph{Oberrealschulen;}

districts where a complete course cannot be maintained, there are often secondary institutions that do not carry the student more than six years. These are known, according to the curriculum, as \emph{Progymnasien, Real-} \emph{progymnasien,} and \emph{Realschulen.} The first two classes are far less common than institutions with the longer course of the same character, but the \emph{Realschulen} are \textsuperscript{six}-y\textsuperscript{ear}

'courses;

nearly twice as numerous as the \emph{Oberrealschulen.}

Since these three types of secondary institutions are so distinct from each other (Fig. 48), it is evident that a parent is forced to decide the future career of his boy at nine years, long before his special ability can be known. If he once enters a real-school, he can never transfer to a gymnasium, because the Latin begins in the lowest class of the latter course, nor can he enter the gymnasium from the real-gymnasium, after twelve, since he has had no Greek. To overcome this objection, during the past quarter of a century efforts have been made to delay the irrevocable decision by grouping all three courses as one institution and making them identical as long as possible. In secondary schools of this new sort, French is usually the only foreign language taught for the first three years. Then the course divides, and one section takes up Latin and the other English. After two years more a further bifurcation takes place in the Latin section, and one group begins with Greek, while the other studies English. These institutions are known as \emph{Reformschulen} (Fig. 48), \emph{Mormchuien;} and the plan was first introduced at Frankfort in 1892. The `reform schools' are now growing rapidly, and there is evident an increasing tendency to postpone the choice of courses as long as possible. The three years of training preliminary to admission to a secondary school of

the \emph{Vorschule.}

Universities state institutions, but controlled by

charters and decrees.

\emph{Tecknische} \emph{Hochschulen.}

any type may be obtained through the people's or the middle schools. But there has also grown up, as an attachment of the secondary schools, a \emph{Vorschule} (preparatory school), to perform this function for pupils of the more exclusive classes.

Higher Education. - Like the other stages of education, the universities are now emancipated from ecclesiastical control, and may be regarded as part of the national system of education. The university is now coordinate and under the same authority with the church, for both are legally state institutions. Universities can, therefore, be established only by the state or with the approval of the state. In general, however, they are not controlled by legislation, but through charters and special decrees of the minister of education. As their income from endowments and fees is very small, they are for the most part supported by the state. They are managed internally by the rector and senate. The rector is annually chosen from their number by the full professors, with the approval of the minister, and the senate is a committee from the various faculties. The professors are regarded as civil servants with definite privileges, and they are appointed by the minister, although the suggestions of the faculty concerned are usually respected.

During the nineteenth century new institutions for the cultivation of science in application to practical and technological purposes have developed from technical schools of a more elementary character. While known as `technical high schools' \emph{(Tecknische Hochschulen),} they are institutions of higher learning, and exist side by side with the universities. They include schools of

engineering, mining, forestry, agriculture, veterinary medicine, and commerce.

Educational Development in France. - The development of a centralized system of education in France began almost a century later than in Germany. During the eighteenth and the early nineteenth century the different monarchic powers were not at all favorable to training the masses, and popular education was badly neglected. It required several revolutions in government and the establishment of a permanent republic, to break \emph{the} old traditions completely, and to make it evident that universal suffrage should be accompanied by universal education. Just after the middle of the eighteenth century the revolutionary spirit began to First agitation

\textsuperscript{0} * » \emph{\textsuperscript{r}} ° for elementary

manifest itself with the appearance of Rousseau's \emph{Emile} education dur

. . , , . . 111 ing the Revolu

(see p.~222), and, except for the traimng started by the tion. Christian Brothers (see p.~140), the first serious atten- tion was given to elementary education. Holland, to whom a general plan for reorganization had been com- mitted, recommended universal education and an ade- quate number of training schools for teachers. While his proposals were not adopted, they were the basis of much of the short-lived legislation that arose during the Revolution, and of the great principles of educational administration that have since been established.

Napoleon, from the beginning, endeavored to reorganize education upon a better basis, and when he had become emperor, ordered all the lycees, secondary colleges, and faculties of higher education to be united in a single corporation, dependent upon the state and known as the `University of France' (1808). This decree of Napoleon and

, . , . the University

centralization divided the country into twenty-seven of France.

administrative `academies,' each of which was to establish university faculties of letters and science near the principal lycees.

This organization, however, did not include elementary education, and little attempt was made to provide for schools of this grade before the reign of Louis Philippe. Upon the advice of his great minister of education, Skc\^{}rimar Guizot, that monarch organized primary education, schools began, requiring a school for each commune, or at least for a group of two or three communes, and starting higher primary schools in the department capitals and in communes of over six thousand inhabitants (1833). He also instituted inspectors of primary schools, and established department normal schools under the more effective control of the state authorities. The plan for higher primary schools was never fully realized, and the institutions of this sort that had been established disappeared during the second empire. The reactionary law of Falloux (1850) did not even mention these schools, but encouraged the development of denominational schools.

The Primary School System. - Guizot, however, had

given a permanent impulse to popular education, and

during the third republic foundations for a national sys

Under third tem of education have rapidly been laid. Schools have

republic pn- *'

mary system been brought into the smallest villages, and elementary 'education has been made free to all (1881) and compulsory between the ages of six and thirteen (1882). To provide trained teachers, every department has been

schools\textsuperscript{1} required to provide a normal school for teachers of each

sex; and two higher normal schools, one for men and one for women, to train teachers for the departmental normal schools, have been opened by the state (1882).\protect\hypertarget{ch25.xmlux5cux23para.460.1.0.box.102.243.1227.901.q.60}{}{ The higher primary schools have been reestablished and Higher priextended (1898), and `supplementary courses' offered tinuation for pupils remaining at the lower primary schools after \emph{\textsuperscript{x}} °°' graduation. The studies in the supplementary courses are technical, as well as general, and some of the higher primary schools have been established for vocational training rather than literary. In addition, there are continuation `schools of manual apprenticeship' in the various communes, subsidized by the state for industrial and agricultural education, and five large schools for training in special crafts have been organized in Paris. Institutions for children between two and six years of age became part of the primary system in the days of \^{}S*\textsuperscript{1} Guizot (1833), and half a century later the present name, \emph{ecoles maternelles} (see p.~244), was adopted (1881), although there have since been marked reforms made in the curriculum.}

Secularization of the school system has also gradually Secularization, taken place. First, the courses of study were secularized by the substitution of civic and moral instruction for religious (1881); next, the instructional force was secularized by providing that members of the clergy should no longer be employed in the public schools (1886), and by recognizing public school teachers as state officers (1889); and finally, the schools themselves were completely secularized by compelling the teaching orders to report to the state authorities (1902), and afterwards by closing the free schools directed by them (1904). Thus within a generation universal elementary education has been established in France and brought completely under state control.

The Secondary System. - As in Prussia, the secondary

Development of lycees and communal colleges.

Organization oflycees

and colleges.

school system of France does not connect with the primary, but is quite separate and distinct (Fig. 49). The training has, since the time of Napoleon, been furnished chiefly by the lycees and communal colleges. During the Restoration (1814-1830) and the reign of Louis Philippe (1830-1848) the lycees came to be called `royal colleges,' but, with the advent of the second republic (1848-1851), the Napoleonic name was restored and the curricula were completely reorganized. By this revision some elasticity was introduced into the last three years of the lycee by a bifurcation into a literary and a scientific course, and during the third republic further elections were introduced, until finally (1902) four distinct courses were established. In the leading lycees and colleges special preparation is also afforded for schools like the military institution of St.~Cyr or the Polytechnic of Paris; and in some there is a short course of three or four years in modern languages and sciences that in function closely approaches that of the German real-school.

The boys ordinarily begin the first `cycle' of the lycee or college at ten years of age, and while they may transfer from the primary system at this stage, in most lycees and colleges there are preparatory classes to train the pupil from six to ten. The second `cycle,' during which the differentiation in courses largely occurs, takes the pupil from fourteen to seventeen, and leads upon completion to the bachelor's degree. Education in a lycee or college is not gratuitous, but the income from tuition fees is so small as to cover but a small fraction of the cost, and the rest is contributed by the state. The communal colleges differ from the lycees in being local, and they are maintained by the communes, as well as by the state.\protect\hypertarget{ch25.xmlux5cux23para.462.1.0.box.103.252.1268.368.q.60}{}{ They have not the same standing, and the same attainments are not required of their professors. Until 1880 there were no lycees and communal colleges for girls, and convents and private schools furnished the only means Secondary bof female education. Even now the usual course in the for girls. public secondary institutions for girls is two years shorter than in those for boys.}

The Institutions of Higher Education. - More than one-half of the universities established in the various `academies' by Napoleon were suppressed as soon as andratoration the monarchy was restored. But about half a dozen °.\textsuperscript{f}'\textsuperscript{he} univer

•' sities.

were reopened in the reign of Louis Philippe, and were gradually improved by the addition of new chairs. Beginning in 1885, a number of decrees established a general council of faculties in each academy to coordinate the different courses and studies, and in 1896 a law was passed, which established a university in each of the sixteen `academies,' except one. These universities differ greatly in size, but all grant the \emph{license,} or master's degree, and the doctorate. The university degrees are Degrees, ordinarily conferred in the name of the state and carry certain definite rights with them, but of late years a new type of degree, `doctorate of the university,' is granted upon easier terms to foreigners more desirous of the degree than of its state privileges. In Paris, besides the university, there is the College of France, which still endeavors to foster freedom of thought (see p.~no), and a dozen other institutions of university grade, con- \emph{9°*TM *\&**}

. , , institutions.

nected with some special line, have been established.

Centralized Administration of the French Education. - The centralization of education is even more complete in France than in Germany. The supreme head of the\protect\hypertarget{ch25.xmlux5cux23para.463.1.0.box.95.217.1263.802.q.60}{}{ Duties of system is the minister of education. He is immediately assisted by three directors, one each for primary, secondrectors, ary, and higher education. A rector is in charge of each of the `academies,' except Paris, where the minister nominally holds the office and a vice rector performs the duties. The rector has authority over all three fields of education in his academy, but does not appoint the prefects, teachers. That office is performed by the prefect, or head of each civil department, upon the recommendation of the academy inspector. There is also a departmental council, presided over by the prefect, that appoints delegates in each canton, to take charge of the school premises and equipment. Further organization is effected through the maintenance of a complete corps of and inspectors, general, academy, and primary inspectors.}

Early Development of English Education. - In England the nationalization of education was delayed even longer than in France. This country was never controlled by enlightened despots, who could, as in Germany, force the growth of public educational sentiment, nor was it overwhelmed by the sweep of a great revolution, destroying, as in France, all opposition to popular progress. National education in England has gradually Slow evolution, grown out of the conflict of a number of elements represented in its society. It has been the product of a series of compromises among many different factors, - the church, state, economic conditions, private enterprise, and philanthropy. For several centuries education was regarded as a function of the church and family, and the nopory\textsuperscript{1} "\^{} sentiment for universal training was retarded by the attitude of the upper classes, who strove to keep the poor in ignorance and to maintain the educational control\protect\hypertarget{ch25.xmlux5cux23para.464.1.0.box.98.226.1243.738.q.60}{}{ of the church. This domination was first seriously challenged in the eighteenth century, and while the training then furnished through the Society for the Promotion of Christian Knowledge, the Sunday schools, and other philanthropic institutions (see pp.~232 ff.), was rather ?\^{}J\textsuperscript{n}S\textsuperscript{1Iopic} meager, these organizations, together with the `monitorial' instruction of the British and Foreign, and the National Societies (see pp.~240 f.), greatly advanced the cause of universal education. And toward the close of the century there began to appear a new point of view, especially with men like Bentham, Blackstone, Robert Owen, and Adam Smith, who advocated universal education, compulsory attendance, and a national system of schools.}

Educational Movements in the Nineteenth Century. - The theory of these great thinkers was somewhat in advance of the times, but, early in the nineteenth century, social changes began to favor better educational First signs of opportunities. The Factory Act (1802) provided for the obligatory training of apprentices; Mr.~Whithread introduced (1807) a bill to permit the civic officials of any township or parish to establish schools for the poor wherever none existed; and Brougham, while losing his bill for popular education (1820), previously secured two commissions of inquiry on school facilities. In 1832, the passage of a reform bill, which largely increased the suffrage, aroused Parliament to the need of educating the masses, and the next year the first grant, £20,000, J\^{}J\^{}IIjint was made for elementary education. This sum was to \textsc{\textsuperscript{!d} }1833. be used solely to aid in building schoolhouses for which subscriptions had been privately obtained, and so could be passed as a vote of `supply,' without referring it

to the House of Lords. For lack of a government organization of education, it was apportioned through the National and the British and Foreign Societies (see p.~240). Governmental activities constantly increased. In 1839 the annual grant was increased to £30,000 and allowed to be used for elementary education without restriction. In the same year, a separate comCommittee of mittee of the Privy Council was designated to adin i83g. minister the educational grants; and in 1856 a Vice

President was appointed to act as chairman of this educational committee. Then, in 1861, through another commission on popular education, it was arranged to base the grant to any school upon the results shown by resuits'iii 86 \^{} pupils in the governmental examinations. This'payment by results' was intended to increase efficiency, but, used as a sole means of testing, it soon proved narrowing and unfair, and had to be supplemented by the general opinion formed of each school by the inspectors. Yet it somewhat increased the efficiency of the work.

Agitation in behalf of universal education continued, and organizations like the `Lancashire Public School Association' of Manchester (1847) and `The League' of Birmingham (1869) spread rapidly through the manufacturing centers. And when the franchise was further extended in 1868, the necessity for preparing millions of the common people for new responsibilities in public in 1870 estab- affairs led in 1870 to the passage of the epoch-making `board schools', bill of William E. Forster. Under this act'board schools,' local'rates,`as or institutions in charge of a board chosen by the people of the community, were to be established wherever a deficiency in the existing accommodations required it. The 'voluntary,' or denominational schools, most of\protect\hypertarget{ch25.xmlux5cux23para.466.1.0.box.97.242.1054.693.q.60}{}{ which belonged to the Church of England, were to continue to share in the government grants upon equal terms with the new institutions, but the latter had also the benefit of local `rates.' Elementary instruction in all schools had to be open to government inspection, and the amount of the grant was partly determined by the report of the inspectors. The board schools were forbidden to allow ``any religious catechism or religious formulary, which is distinctive of any particular denomination;'' and religious instruction in either type of school had to be placed at the beginning or end of the school session, so that, under the `conscience clause' of the act, any scholar might conveniently withdraw at that time.}

well as grants.

This act of 1870 was, of course, the \emph{magna charta} of national education, and has become the basis of much school legislation. The compromise in the bill that allowed the voluntary schools, with their sectarian instruction, to continue receiving government support, however, prevented a logical and consistent system from being established. The dual system of elementary schools continued to be developed in a variety of enactments. Compulsory attendance laws were passed (1876, StSid\textsuperscript{Ulsory} 1880), the minimum age of exemption was set first at minimum age, eleven years of age, and then raised to twelve (1893, 1899), and an extra grant, to take the place of tuition \textsuperscript{freetmtion}. fees (1891), made it possible for most schools to become absolutely free. Finally (1899), there was created a central Board of Education, which assumed the functions Education! °\textsuperscript{f} of the Committee of Privy Council on Education and similar agencies for managing educational interests.

Subsequent Educational Movements. - Within a generation of existence the board schools met with a phenom\protect\hypertarget{ch25.xmlux5cux23para.467.2.0.box.296.236.1055.1721.q.60}{}{enal growth, and came to include about seventy per cent of the pupils. They were spending about half as much again upon each pupil as were the voluntary schools, and were able to engage a much better staff of teachers. This .extension of civil influence in education was bitterly opposed by the Established Church, and when the conservatives came into power through the assistance of the clergy, they passed the act of 1902, whereby the denominational schools were permitted to share also in the local rates. While under this act the administration of both board and voluntary schools was now centralized in the county and city councils, the immediate supervision of instruction in the individual schools was placed in the hands of a board of managers; and, despite their receipt of local taxes, the voluntary schools were required to have but two of their managers appointed by the council, and the other four were still selected by the denomination. Serious opposition to the enforcement of the new law arose among nonconformists and others, and coercive measures were taken by the government. The new act, however, while unfair to those outside the Church of England, tended to sweep away the dual system of public and church schools, since both were coming to rest upon a basis of public control and support. Since 1902 all elementary schools have been considered as part of one comprehensive system, and the board schools have been distinguished as `provided schools' and the voluntary as `nonprovided.' Moreover, under the legislation of 1902 steps were also taken to coordinate secondary with elementary education, and bring it somewhat within the public system. The board schools had early in their existence begun to develop upward into}\protect\hypertarget{ch25.xmlux5cux23para.468.1.0.box.114.234.1264.530.q.60}{}{ secondary education and before long had come to compete with the older grammar and public schools, but in 1900 the `Cockerton judgment' forbade the use of local rates for other instruction than elementary, and it remained for the new act to impose upon councils the duty to support instruction in subjects beyond the elementary Instruction \^{} work. The Board of Education was also empowered to pub\^{}Spen'se. inspect the work of the great public schools and other endowed secondary institutions, and to allow grants to all schools meeting the conditions of the Board.}

In 1902'voluntary' schools also allowed local rates,

but dual system swept away,

After the liberals returned to power, they continued the conservatives' policy of granting local rates to all elementary schools, and of bringing secondary education under public support and control. While the education bill of 1906, which was kept from passage by the House defeated?\textsuperscript{06} of Lords, did not recognize church schools as such, and insisted upon bringing them under the complete control of the public authorities, it made no attempt to return to the former dual system of schools and the isolation of secondary from elementary education. It still held also to religious, and, under safeguards, even to sectarian instruction in the elementary schools, and may yet be passed in a revised form. A voluntary committee for a `resettlement in English elementary education,' through the mediation of the President of the Board of Education and the Archbishop of Canterbury, has formulated a plan, which concedes the principle of public control and \textsuperscript{but} ?\textsuperscript{ew} r\textgreater{}\textsuperscript{lan}' support for all elementary schools and religious freedom schools under

1 .1 1 1 • \emph{e} 1 public control.

for teachers and pupils, but provides local option for the continuance of denominational schools. Thus, while England is not prepared to adopt a secular system, like that of France and the United States, and has not yet\protect\hypertarget{ch25.xmlux5cux23para.469.1.0.box.78.230.1265.476.q.50}{}{ fully articulated its secondary education with elementary, (Fig. 50), it is upon the high road to a complete centralization of school administration in the national government. During the nineteenth century also the classical and eccie\textsuperscript{S}sStica°\textsuperscript{d} ecclesiastical monopoly in secondary and higher educaSokenm \^{}on \textsuperscript{was} largely broken. All the older public and gramsecondary and \textsubscript{mar} schools (see pp.~412 f.) developed `modern sides,' and tion. during the Victorian era a number of new schools were}

founded, which gave considerable attention to the modern languages and the sciences from the start. A recognition of the scientific ideals began also to appear in the curriculum of Cambridge (1851) and Oxford (1853), and the theological requirements for a degree were dropped (1856). By the last quarter of the century actual laboratories had been introduced, and students were freed from all doctrinal tests at both universities. Moreover, new universities, better adjusted to modern demands and more closely related to the school systems and the civil government, began to arise in manufacturing centers. Since 1889 such municipal or `provincial' institutions as the Universities of Birmingham, Manchester, Leeds, Liverpool, and Bristol have sprung up, and the University of London, started as an examining body in 1836, has become a teaching institution.

Development of Education in the Dominion of Canada. - Canada developed schools in very early days. In the beginning education was cared for in the four provinces separately, and when the Dominion of Canada was finally formed (1867), the federal government left to each province the administration of public education within its borders. The same autonomy was extended

to the provinces that have since been admitted to the federation. Two types of educational control, - state \textsuperscript{Two} tyi*\textsuperscript{3}and ecclesiastical, have been developing from the first. The former method is best illustrated by the system of Ontario; and the latter by that of Quebec. Ontario d) Ontario and

(a) Quebec.

was settled mostly by emigrants from England, Scotland, and the United States, and practically all brought with them the concept of public control of education. The French Catholics of Quebec, on the other hand, naturally followed their traditions of parish schools.

The Public School System of Ontario. - The system of schools in Ontario began with the passage of its Common Schools Act in 1846. This was formulated after a careful study of the systems of Massachusetts, New York, and the European states, and included excellent elements from various systems and many original features of value. By 1871 this fundamental law had come to include free tuition, compulsory attendance, county inspection, and uniform examinations. In 1876 an even greater centralization of the provincial system was ef- Universal

\textsuperscript{0} \emph{\textsuperscript{r} \textsuperscript{J}} education, and

fected through substituting for the chief superintendent since 1870

,.. ,, \ldots{}. great centrah

a minister of education with much larger powers, and zation through bringing all stages of public education, - elementary, secondary, and higher schools, into much closer relationship. The minister now has many assistants, including an Advisory Council of Education; and he initiates and directs all school legislation, decides complaints and disputes, sets examinations for the high, elementary, model, and normal schools, prescribes the courses of study, chooses the text-books, and appoints the inspec- , , ,.

. . and subordi

tors. The system is further administered by subordinate nate authoriauthorities elected in the localities, whose duties are

Unification of the several stages of education.

Inspectors.

`Separate schools.'

clearly defined by law. The province is for educational purposes divided into counties, which are in turn divided into townships, and subdivided into sections and incorporated cities, towns, and villages. The central and local administrations are wisely balanced, and while the one determines scholastic standards through its professional requirements, the other establishes schools and appoints teachers.

The system of elementary schools, high schools, and universities, is fully unified, and the work of each stage fits into the others even more exactly than in the `ladder' system of the United States. The training of teachers is cared for through the departments of Education in the universities, the eight provincial normal schools, and a model school in each county. The teachers for secondary institutions are prepared at the universities, the normal schools grant a life certificate to teach in the elementary schools, while the model schools afford fourteen weeks of training for country teachers. The buildings, equipment, courses, and instruction of the high, elementary, and model schools are each reported upon by inspectors of assured scholarship and experience. Since 1863 permission has been granted to establish `separate schools' for any peculiar creed or race, wherever there are five families requesting it. This opportunity to have schools of their own faith has not been embraced by any save the Roman Catholics. Any one paying toward the support of a `separate school' is exempt from taxation for the regular public schools. Special provincial inspectors report upon these schools, but in the same way as for the public schools.

The System of Ecclesiastical Schools in Quebec. - \protect\hypertarget{ch25.xmlux5cux23para.474.1.0.box.93.248.1275.1452.q.60}{}{ The Ontario system may be considered typical of the educational administration in the various provinces of Canada, except Quebec. Every other province has j\^{}fJJJJ\^{} \textsubscript{t} sought uniformity of school provision and educational Ontario, standards through government control, although none of them grant their central official quite as much power as Ontario. Alberta and Saskatchewan likewise permit `separate schools,' and they existed in Manitoba until 1890. But the type of control in Quebec is very different from that of the other provinces. There in 1845 the in Quebec parish was by law made the unit of school administration. But seven years later government inspectors were estab- but since 1859 lished, and in 1859 a central organization was completed Public inwith a Council of Public Instruction. This authority is \textsuperscript{s} \textsuperscript{ruc} \textsuperscript{10n} composed of two divisions, a Roman Catholic and a Protestant, which sit separately and administer the schools of their respective creeds. The provincial super- \textsuperscript{and} superintendent of schools, appointed by the lieutenant gover- schools, nor, is \emph{ex officio} chairman of both divisions, but he can vote only with the division to which he belongs by religion. Each division makes regulations for the instruction and tests of its own schools, and appoints inspectors School support of its own faith. The proceeds from the general public school fund or from any educational legacies are divided in proportion to the Catholic and Protestant inhabitants, but the regular school rate may be assigned to whichever of the two school systems the taxpayer wishes.}

\textbf{Supplementary Reading}

Graves, \emph{In Modern Times} (Macmillan, 1913), chap. IX; Parker, \emph{Modern Elementary Education} (Ginn, 1912), chaps. X and XI. The following works throw light upon various phases of the respective\protect\hypertarget{ch25.xmlux5cux23para.475.1.0.box.288.231.1059.1275.q.50}{}{ countries: Nohle, E., \emph{History of the German School System (Report of the U. S. Commissioner of Education,} 1897-1898; vol.~I, pp.~2644); Paulsen, F., \emph{German Education} (Scribner, 1908); Russell, J. E., \emph{German Higher Schools} (Longmans, Green, 1896); Paulsen, F., \emph{The German Universities} (Macmillan, 1895; Scribner, 1906); Kandel, I. L., \emph{The Training of Elementary School Teachers in Germany} (Columbia University, \emph{Teachers College Contributions,} No.~31, 1910); Brown, J. F., \emph{The Training of Teachers for Secondary Schools in Germany} (Macmillan, 1911); Beard, Mary S., \emph{Hcoles maternelles of Paris (Great Britain, Board of Education, Special Reports on Educational Subjects,} vol.~VIII, no. 8); Farrington, F. E., \emph{French Secondary Schools} (Longmans, Green, 1910) and \emph{The Public Primary System of France} (Columbia University, \emph{Teachers College Contributions to Education,} no. 7, 1906); Smith, Anna T., \emph{Education in France (Reports of the United States Commissioner of Education,} 1890 to 1914, see tables of contents); Greenough, J. C, \emph{The Evolution of the Elementary Schools of Great Britain} (Appleton, 1903); Montmorency, J. E. G. de, \emph{State Intervention in English Education} (Macmillan, 1903); Sharpless, I., \emph{English Education in Elementary and Secondary Schools} (Appleton, 1892); Smith, Anna T., \emph{Education in England (Monroe Cyclopadia of Education,} vol.~II); Sandiford, P., \emph{The Training of Teachers in England and Wales} (Columbia University, \emph{Teachers College Contributions,} no. 32, 1910); Coleman, H. T. J., \emph{Public Education in Upper Canada} (Columbia University, \emph{Teachers College Contributions,} no. 15, 1909); Ross, G. W., \emph{The School System of Ontario} (Appleton, 1896); Smith, Anna T., \emph{Education in Canada (Monroe Cyclopadia of Education,} vol.~I).}

\hypertarget{the-scientific-movement-and-the-curriculum}{%
\chapter{26. The Scientific Movement and the Curriculum}\label{the-scientific-movement-and-the-curriculum}}

\textbf{Outline}

\begin{itemize}
\tightlist
\item
  During the past two centuries a great growth has taken place in the natural sciences. For a long time this development affected practical life very little, but during the nineteenth century the application of science to industrial problems has resulted in a host of inventions.
\item
  Because of the importance of the sciences to life, Spencer and others have urged the inclusion of them in the curricula of schools and colleges. While the content of the sciences has furnished the chief argument for this, many scientists have urged their value as formal discipline.
\item
  Instruction in the sciences has gradually been included in the higher, secondary, and elementary institutions of Germany, France, England, and the United States.
\item
  This marked scientific movement is allied with the psychological tendency in its improvement of method, and with the sociological in its emphasis upon human welfare.
\end{itemize}

\textbf{The Development of the Natural Sciences in Modern Times.} - We have already (chapter XV) witnessed the growth of the natural sciences and the beginning of their introduction into the curriculum toward the close of the seventeenth century. This tendency was also greatly stimulated by Rousseau, who, we have seen (pp.~218-222), may be held to advocate the scientific, as well as achievements the sociological and psychological movements. And during the past two centuries this development has become most rapid and extensive. The desire for scientific investigation steadily grew throughout the eighteenth and nineteenth centuries until its ideals, methods, and results became patent in every department of human knowledge. The strongholds of ignorance, superstition, and prejudice were rapidly stormed and taken through new discoveries or new marshallings of facts already discovered. But evident as this movement has been, it is scarcely possible here even to mention the more important scientific achievements, or to outline the broad sweep of progress in astronomy, geology, biology, physiology, chemistry, physics, and other sciences within a century. The Newtonian theory has been confirmed by the investigations of Lagrange and Laplace and by the discovery of Neptune by mathematical reasoning from the effects of its gravitation. Hutton's `Plutonic' theory of continents and Agassiz's hypothesis of a universal ice-age have been formulated; the doctrine of evolution of Darwin (Fig. 51) and Mendel's law of inheritance have been established; Liebig and others have thrown light upon the process of digestion and the functioning of the lungs and liver; atoms, molecules, and ions have been defined; Joule and Mayer have demonstrated the conservation of energy; and the periodic law of chemical elements has been discovered by Newlands.

\textbf{The Growth of Inventions and Discoveries in the Nineteenth Century.} - It should be noted, however, that the majority of these investigations were for a long time carried on outside the universities, and, owing to the almost proverbial conservatism of educational institutions, the natural sciences scarcely entered the course of study anywhere. In fact, these great discoveries at first seem not to have affected practical life in any direction. Huxley tells us that in the eighteenth century ``weaving and spinning were carried on with the old appliances; nobody could travel faster by sea or by land than at any previous time in the world's history, and King George could send a message from London to York no faster than King John might have done.'' But a little later, as he adds, ``that growth of knowledge beyond imaginable utilitarian ends, which is the condition precedent of its practical utility, began to produce some effect upon practical life.'' The nineteenth century will, on this account, always be known for its development of inventions and the arts, as well as of pure science. During this period science rapidly grew and took the form of applications to the problems of labor, production, transportation, communication, hygiene, and sanitation. The reaper, the sewing machine, the printing press, and the typewriter greatly reduced the cost of labor; the steamboat, locomotive, electric railway, telegraph, and telephone linked all parts of the world together; anthracite, friction matches, petroleum, and electric lighting and heating greatly enlarged the comforts of life; and stethoscopes, anaesthetics, antiseptics, and antitoxines added wonderfully to the span of human life.

\textbf{Herbert Spencer and \emph{What Knowledge is of Most Worth.}} - Because of these practical results, the vital importance of a knowledge of natural phenomena to human welfare and social progress was more and more felt throughout the century. It gradually became evident that the natural sciences were demanded by modern life and constituted elements of the greatest value in~modern culture and education. Many English and American writers began to maintain that an exclusive study of the classics did not provide a suitable preparation for life, and that the sciences should be included in the curriculum. This step was bitterly opposed by conservative institutions and educators. During a greater part of the century a contest was waged between the advocates of the classical monopoly and the progressives, who urged that the sciences should be introduced.

A representative argument for sciences in the course of study is that made by Herbert Spencer (Fig. 52) in his essay on \emph{What Knowledge Is of Most Worth.} He ventured to raise the whole question of the purpose of education. He held that ``to prepare us for complete living is the function which education has to discharge; and the only rational mode of judging of any educational course is, to judge in what degree it discharges such function. Our first step must obviously be to classify, in the order of their importance, the leading kinds of activity which constitute human life. They may be arranged into: 1. Those activities which directly minister to self-preservation; 2. Those activities which, by securing the necessaries of life, indirectly minister to selfpreservation; 3. Those activities which have for their end the rearing and discipline of offspring; 4. Those activities which are involved in the maintenance of proper social and political relations; 5. Those miscellaneous activities which make up the leisure part of life, devoted to the gratification of the tastes and feelings. The ideal of education is complete preparation in all these divisions. But failing this ideal, the aim should be to maintain a due proportion between the degrees of preparation in each, greatest where the value is greatest, less where the value is less, least where the value is least.''

Applying this test, Spencer finds that a knowledge of the sciences is always most useful in life, and therefore of most worth. He considers each one of the five groups of activities and demonstrates the need of the knowledge of some science or sciences to guide it rightly. An acquaintance with physiology is necessary to the maintenance of health, and so for self-preservation. Any form of industry or other means of indirect self-preservation will require some understanding of mathematics, physics, chemistry, biology, and sociology. To care for the physical, intellectual, and moral training of their children, parents should know the general principles of physiology, psychology, and ethics. A man is best fitted for citizenship through a knowledge of the science of history in its political, economic, and social aspects. And even the aesthetic or leisure side of life depends upon physiology, mechanics, and psychology as a basis for art, music, educational and poetry. Hence Spencer advocates a complete change~from the type of training that had dominated education since the Renaissance and calls for a release from the traditional bondage to the classics. Instead of Greek and Latin for `culture' and `discipline,' and an order of society where the few are educated for a life of elegant leisure, he recommends the sciences and a new scheme of life where every one shall enjoy all advantages in the order of their relative value. But Spencer uses the term `science' rather loosely, and seeks to denote the social, political, and moral sciences, as well as the physical and biological, as being `of most worth.' Hence~he does not deserve to be severely arraigned for his `utilitarianism,' as he has been so frequently. His `preparation for complete living' includes more than `how to live in the material sense only,' and with him education should contain such material as will elevate conduct and make life pleasanter, nobler, and more effective.

\textbf{Advocacy of the Sciences by Huxley and Others.} - Another great popularizer of the scientific elements in education, who also stressed the value of the sciences for `complete living' and social progress, was Thomas H. Huxley (Fig. 53). His use of English was vigorous and epigrammatic, and he showed great skill in bringing his conclusions into such simple language that the most unscientific persons could understand them. Especially in an address on \emph{A Liberal Education} before a `workingmen's college,' he has most forcefully depicted the value of the sciences and other modern subjects in training for concrete living, and ridiculed the ineffectiveness of the current classical education. He maintains that ``the life, the fortune, and the happiness of every one of us depend upon our knowing something of the phenomena of the universe and the laws of Nature. And yet this is what people tell to their sons:`At the cost of from one to two thousand pounds of our hard-earned money, we devote twelve of the most precious years of your life to school. There you shall not learn one single thing of all those you will most want to know directly you leave school and enter upon the practical business of life.'\,'' Instead of this, ``the middle class school substitutes what is usually comprised under the compendious title of the `classics' - that is to say, the languages, the literature, and the history of the ancient Greeks and Romans, and the geography of so much of the world as was known to these two great nations of antiquity.'' Thus ``the British father denies his children all the knowledge they might turn to account in life, not merely for the achievement of vulgar success, but for guidance in the great crises of human existence.''

Many other vigorous lecturers and writers entered into this reform of the curriculum. Opposition to the over-emphasis of languages, especially the classics, in the content of education was undertaken even earlier in the century by the distinguished phrenologist, George Combe. In his `secular' schools and in his work on \emph{Education,} he emphasized instruction in the sciences relating to moral, religious, social, and political life, as well as those bearing upon man's physical and mental constitution. After the middle of the century a number of men undertook to popularize the sciences in America by tongue and pen. One of the most effective of these was Edward L. Youmans, who collected and edited a set of lectures urging the claims of the various sciences under the title of \emph{Culture Demanded by Modern Life} (1867). He also founded the \emph{International Science Series} (1871) and the \emph{Popular Science Monthly} (1872). A service for the sciences, bearing more directly upon the educational world, was that performed by Charles W. Eliot (Fig. 54), President of Harvard. This he accomplished largely by an extension of the elective system and an emphasis upon science in the curriculum of school and college. In his description of `a liberal education,' he argues that ``the arts built upon chemistry, physics, botany, zoology, and geology are chief factors in the civilization of our time, and are~growing in material and moral influence at a marvelous rate. They are not simply mechanical or material forces; they are also moral forces of great intensity.''

\textbf{The Disciplinary Argument for the Sciences.} - Thus, in general, the writers and lecturers interested in the scientific movement held that a knowledge of nature was indispensable for human welfare and that the content of studies rather than the method was of importance in education. Many of them also expressed their dissent from the disciplinary conception of education urged by the classicists. Huxley, for example, parodies the usual Huxley paro- linguistic drill by stating: ``I could get up an osteological primer so arid, so pedantic in its terminology, so altogether distasteful to the youthful mind, as to beat the recent famous production of the head-master out of the field in all these excellences. Next, I could exercise my boys upon easy fossils, and bring out all their powers of memory and all their ingenuity in the application of my osteogrammatical rules to the interpretation, or construing, of those fragments.''

Yet the tradition of `formal discipline' and the belief~in faculties or general powers of the mind that might be~trained by certain favored studies and afterward applied~in any direction (see pp.~182f.) were too firmly rooted to~be entirely upset. Even the greatest of the scientists~seem to have been influenced by this notion and to have~attempted occasionally a defense of their subjects on~the basis of superiority in this direction. After Spencer~has made his effective argument for the sciences on the~ground that their `content' is so much more valuable~for the activities of life, he shifts his whole point of view,~and attempts to anticipate the classicists by occupying~their own ground. He admits that ``besides its use for guidance in conduct, the acquisition of each order of facts has also its use as mental exercise.'' As evidence of this, he undertakes to show that science, like language, trains the memory, and, in addition, exercises the understanding; that it is superior to language in cultivating judgment; that, by fostering independence, perseverance, and sincerity, it furnishes a moral discipline. A similar argument is made by Combe, when he maintains that ``it is not so much the mere knowledge of the details of Chemistry, of Natural Philosophy, or of any other science that I value, as the strengthening of the intellect, which follows from these studies.'' So Youmans declares that ``by far the most priceless of all things is mental power. Science made the basis of culture will accomplish this result.'' In fact, nearly every apologist for the natural sciences at some time or other has advocated these subjects from the standpoint of formal discipline, although the implied attitude toward the transfer of a generalized ideal is often in harmony with modern psychology (seep. 184).

\textbf{Fig. 51-54}

\textbf{Introduction of the Sciences into Educational Institutions; Germany.} - Contemporaneously with the growth of inventions and the cogent arguments and vigorous campaigns of advanced thinkers during the nineteenth century, training in the sciences was gradually creeping into educational practice. While the sciences began to work their way into institutions of all grades early in the eighteenth century, it was not until about the middle of the nineteenth that the movement was seriously felt in education. Even in Germany the first attempts at studying nature were made outside the universities in~the `academies of science.' We have seen (pp.~177 f.) that during the eighteenth century most of the Protestant universities had started professorships in the sciences. But it was not until the beginning of the second quarter of the nineteenth century that, in Liebig's laboratory at the University of Giessen, students first began to be taught through experiments, and it was after the middle of the century before this investigation work had generally replaced the formal science instruction in German universities. Since then the development of science in the higher education of Germany has been phenomenal. The \emph{Technische Hochschulen} (see p.~380) have also come to furnish instruction in all fields of applied science.

In German secondary instruction the realistic instruction of the pietists was brought by Hecker (see p.~176) to Berlin, where he started his famous \emph{Realschule} in 1747, and before the beginning of the nineteenth century similar institutions had spread throughout Prussia. Early in the nineteenth century the course of study in the gymnasiums of Prussia was considerably modified, and, as part of the compromise, some science was introduced. The movement later spread into the secondary education of states in South Germany, and, while the total amount of science was not large, it managed to hold its place in the gymnasial curriculum even during the reaction to absolutism between 1815 and 1848. But, as we have seen (p.~378), two types of real-schools were eventually recognized, - \emph{Realgymnasium} and \emph{Oberrealschule,} and they at present devote approximately twice as much time to the physical and biological sciences as do the gymnasia. Technical and trade schools, with scientific and mathematical subjects as a foundation for the vocational work, have also appeared as a species of secondary education in Germany (see p.~420). The first of these were opened in Nuremberg in 1823, but their rapid increase in numbers, variety, and importance has taken place since the middle of the century, and their development in organization and method has occurred within the past twenty-five years.

The scientific movement was also felt in the elementary schools of Germany during the early part of the nineteenth century. Science was considerably popularized by the schools of the philanthropinists (pp.~227 L), and was widely introduced into elementary education by the spread of Pestalozzianism in Prussia and the other German states (seep. 289f.). Before the close of the first quarter of the century the study of elementary science, - natural history, physiology, and physics, appeared in various grades; geography and drawing were taught throughout the course; and geometry was included in the upper classes of the \emph{Volksschulen.}

\textbf{France.} - Before the Revolution in France the higher and secondary institutions found little place for instruction in science. There was a chair of experimental physics at the College of Navarre of the University of Paris and at the Universities of Toulouse and Montpelier, and natural history was also taught at the more independent College of France, but, as a whole, education was dominated largely by humanism. However, with the establishment of the republic a new regime began in education, as in other matters, and science entered more largely into higher and secondary instruction. Most of the revolutionary proposals subordinated letters to science,~and in 1794 the republic founded a great central normal school, where the famous Laplace and Lagrange for a short time gave instruction in science. In 1802 Napoleon had included in the scientific course for the lycees natural history, physics, astronomy, chemistry, and mineralogy, and a definite advance in quantity and method of the scientific instruction in the secondary schools was made in 1814. On the ground that they were injuring classical studies, Cousin in 1840 had the sciences curtailed, but he was shortly forced to restore them upon an optional basis. A contest between the two types of studies was carried on in the lycees until 1852, when a bifurcation in the course put the two theoretically upon the same basis. The scientific course, however, has never been considered equal in prestige to the classical, although it has constantly increased in length and difficulty.

Some instruction in science has come to be given during the past forty years even in the elementary schools of France. In the lower primary schools the work is informal, and consists mostly of object lessons and first scientific notions. These are developed in connection with drawing, manual training, agriculture, and geography of the neighborhood and of France in general. Instruction becomes more formal in the `higher primary' schools, and includes regular courses in the natural and physical sciences and hygiene, as well as geography, drawing, and manual training. In the normal schools for primary teachers instruction in all the physical and biological sciences is even more thorough, and includes not only the facts and theories of general scientific importance, but it also emphasizes their applications to everyday life. For example, the flora and fauna of the neighborhood are studied in their special relation to agriculture.

\textbf{England.} - In England, several chairs in the natural sciences were established at Cambridge during the eighteenth century. But it was almost the middle of the nineteenth century before the biological sciences and the laboratory method of instruction were introduced, and not until toward the close of the century did science become prominent at Cambridge and Oxford. And the most marked promotion of the scientific movement in England has occurred within the past fifty years through the foundation of efficient municipal universities in munidpal such centers as Birmingham, Manchester, London, and Liverpool (see p.~392). For many years the laboratory instruction was given only in institutions outside the universities. Higher courses in science by the new methods were afforded through the foundation of the Royal School of Mines (1851), the Royal School of Naval Architecture and Marine Engineering (1864), and the Normal School of Science (1868), which were all combined in 1890 into a single institution known as the Royal College of Science, and in 1907, when the Technical College (founded 1881) of the City and Guilds of London Institute was also merged, the entire corporation became known as the Imperial College of Science and Technology. An agency that was instrumental in encouraging the advanced study of science, although it accomplished even more for elementary and secondary schools, was the national Science and Art Department. This science and organization was founded in 1858 to bring under a single management the science, trade, and navigation schools already existing, and to facilitate higher instruction in~science, and a few years later began to offer examinations and to grant certificates to teach science in the elementary schools. It was taken over by the national Board of Education, when that body was organized in 1899 (see p.~389).

In English secondary instruction the `academies,' in which science first appeared (pp.~157 f.), had before the close of the eighteenth century greatly declined, and the humanistic `public' schools and secondary institutions of a private character had as yet paid almost no attention to the sciences. In the first half of the nineteenth century an anti-classical campaign began, and, continuing with ever increasing force until the middle of the century, it brought about the foundation of numerous schools to embody the new ideals. Toward the close of 1848 the first `secular' school was opened by Combe (see p.~403) at Edinburgh, and included in its curriculum a study of geography, drawing, mathematics, natural history, chemistry, natural philosophy, physiology, phrenology, and materials used in the arts and manufactures. Similar institutions were organized at Glasgow, Leith, London, Manchester, Birmingham, Newcastle, Belfast, and many other cities of the United Kingdom. While short-lived, these schools did much to promote the introduction of sciences into secondary education that soon followed. Shortly after the middle of the century Rugby, and then Winchester, introduced science into the regular curriculum, and by 1868, as a result of the governmental investigation of the endowed schools, which showed an almost complete absence of science in the curricula, all the leading secondary schools schools,~began to establish a `modern side.' This course generally included physics and natural history, as well as modern languages and history, but it was most reluctantly organized by the institutions, and, while it has attained to great efficiency, it has never, except in a few schools, been accorded the same standing as the classical course. The Department of Science and Art also afforded much~encouragement to secondary instruction in the sciences by subsidizing schools and classes in physics, chemistry, zoology, botany, geology, mineralogy, and subjects involving the applications of science. Before its absorption into the Board of Education some ten thousand classes and seventy-five independent schools of secondary grade received assistance from this source.

The Department also gave aid to the study of science in elementary education. As early as the fifties, grants were made to establish work in elementary science, art, and design, but the educational value was for more than forty years subordinated to practical applications. And while, after the report by a Committee of the British Association in 1889,much aid was furnished for the equipment of laboratories, lecture rooms, and workshops, and an increase in the staff of instructors, for a decade no subjects except the rudiments were required in the elementary course, and such `supplementary' subjects as elementary science and geography, if taught, were given a special subsidy. But since 1900 this scientific work has been made compulsory in the elementary curriculum.

\textbf{The United States.} - In the colleges of the United States the courses show considerable evidence of science teaching by the eighteenth century. Harvard, Yale, Princeton, King's (afterward Columbia), Dartmouth, during the Union, and Pennsylvania had all come to offer work in~`natural philosophy' or `natural history,' which terms might then be used to cover physics, chemistry, geology, astronomy, botany, and zoology. However, before the Revolution physics seems to have been a subordinate branch of mathematical instruction, even less importance was attached to biology, and chemistry was only occasionally taught as an obscure and unimportant phase of physics. Laboratories and instruments of precision did not yet exist.

Since then whole fields of science have been discovered and defined, and others, like geology and astronomy, have been reclaimed from dogmatism, and science studies have slowly come into favor. Instruction in chemistry has grown up through a study of materia medica at the medical schools of Pennsylvania (1768), Harvard (1782), and Dartmouth (1798). A separate chair of chemistry was soon established at Princeton (1795), Columbia (1800), Yale (1802), Bowdoin (1805), South Carolina (1811), Dickinson (1811), and Williams (1812), and the movement continued until practically all the colleges had recognized it as an important branch of study. But while experiments were from the first performed as demonstrations by the instructors, it was generally not until almost the middle of the century that students were admitted at all to the laboratories. About the same time laboratories in physics began to be equipped with apparatus. Geology was included in the early professorship of chemistry at Yale, and was given a distinct chair upon the advent of James D. Dana about the middle of the century, while Amos Eaton taught it as a separate subject at Williams as early as 1825. Some attention was given to astronomy early in the century, although the instruments remained very ordinary and the methods~authoritative and prescriptive until the opening of the~observatories at Cincinnati (1844), Cambridge~(1846),~and Ann Arbor (1854). The biological sciences were even longer studied through mere observation rather~than investigation and experiment. Until Louis Agassiz~opened his laboratory at Harvard to students just after~the middle of the century, the courses were meager,~mostly theoretical and classificatory, and were given~entirely by lecture, without field or laboratory work.~Since then the development has been rapid.

But the greatest impulse was given to instruction in though evoiuscience through the publication of Darwin's \emph{Origin of} \emph{Species} (1859), and the dissemination of evolutionary doctrine through Asa Gray, professor of natural history at Harvard, and William B. Rogers, president of the Massachusetts Institute of Technology. The intellectual development ensuing also brought about the foundation of such new institutions as Cornell and Johns Hopkins, which emphasized the teaching of science as an unconscious protest against the exclusively classical training. Special scientific and technological schools~likewise began to arise. The Rensselaer Polytechnic Institute (1825) and the Lawrence Scientific School at Harvard (1847) had already been opened, but now similar schools of science, like Sheffield at Yale (i860), and the Massachusetts Institute of Technology (1862), sprang up in all parts of the country. In 1862 the Morrill Act of Congress appropriated lands in every state to promote education in agriculture, mechanic arts, and the natural sciences. These grants, which amounted at first to thirteen million acres, were subsequently~extended to new states as they were admitted, and the endowment was increased by the annual grants of money that were made under later acts. From these funds and private benefactions, further schools of science were started or old schools were strengthened in every state.

Through the academy movement (pp.~158 ff.) sciences were introduced into American secondary education. Sometimes these subjects were extended downward from the colleges, but often they had as yet been barely touched by the colleges. As the early high schools grew up, they continued the attention paid to the sciences by the academies. The first high school to appear, that at Boston in 1821 (pp.~268 f.), scheduled geography in the first year; navigation and surveying in the second; and natural philosophy and astronomy in the third. A similar emphasis upon science appeared during the first half of the century in all the secondary institutions, whether known as academies, high schools, union schools, or city colleges. In all cases, however, instruction was given mainly through text-books, and, while experiments were frequently used for demonstration by the teacher, there was no laboratory work for the students. Moreover, a tendency to overload the curriculum with sciences was much increased during the seventies by the demand of the legislatures in several states that candidates for teachers' certificates pass an examination in several sciences. The high schools and academies endeavored to furnish the necessary training to prepare for these examinations, and until toward the end of the century the courses in the sciences were numerous and of rather superficial character. Within the last twenty years, however, the schools have come to limit each student to a relatively few courses taught by thorough laboratory methods.

Except for geography, which appeared in the curriculum early in the century, the rudiments practically constituted the entire course of the elementary school until the time of Horace Mann. Largely through his efforts, physiology was widely introduced by the middle of the century. About a dozen years later the Pestalozzian object teaching began to come in through the Oswego methods, although it tended to become formalized. Thus materials in several of the sciences came to be used, and the pupils were required to describe them in scientific terms. Toward the close of the century the sciences came to be presented more informally by the method generally known as `nature study.' This movement quickly spread through the country, and has most recently appeared in the guise of agricultural instruction (see p.~424). Many states now require agriculture as a requisite for a teacher's certificate, and most normal schools have come to furnish a training in the subject.

\textbf{Interrelation of the Scientific with the Psychological and Sociological Movements.} - It is evident that there has been a marked scientific movement in the educational systems of all countries during the past two hundred years. The sciences began to appear in the curricula of educational institutions in the seventeenth and eighteenth centuries, but their rapid increase, and the use of laboratories and the scientific method in instruction, dated from the middle of the nineteenth. In some respects this scientific movement has been closely related~to the other modern tendencies in education, - the psychological and the sociological. The coincidence of~the scientific movement with the psychological on the question of formal discipline has been evident (pp.~183 f.). The influence of the development of the sciences upon educational method also constitutes part of the psychological movement. The sciences demanded entirely different methods of teaching from the traditional procedure. These innovations were worked out slowly by experimentation, and when they proved to be more in keeping with psychology, they reacted upon the teaching of the older subjects and came to be utilized in history, politics, philology, and other studies. A corresponding improvement in the presentation of the form, content, and arrangement of various subjects has taken place in text-books, and a radically different set of books and authors has been rendered necessary.

The scientific movement has even more points in common with the sociological. In its opposition to the disciplinarians and its stress upon content rather than form, the scientific tendency coincides with the sociological, although the former looks rather to the natural sciences as a means of individual welfare, and the latter~to the social and political sciences to equip the individual for life in social institutions and to secure the progress of society. But while the scientist usually states his argument in individual terms, because of his connection in time and sympathy with the individualism of the eighteenth and nineteenth centuries, the same writer usually, as in the case of Rousseau, Combe, Spencer, and Huxley, advocates the social, moral, and political sciences as a means of complete living. Similarly, the sociological movement has especial kinship with the economic and utilitarian aspects of the study of the sciences, for professional, technical, and commercial institutions have been evolved because of sociological as well as scientific demands. Again, the use of the sciences in education as a means of preparing for life and the needs of society overlaps the modern sociological principle of furthering democracy. Both tendencies lead to the best development of all classes and to the abandonment of artificial strata in society.

\textbf{Supplementary Reading}

Graves, \emph{In Modern Times} (Macmillan, 1913), chap. X; and \emph{Great Educators} (Macmillan, 1912), chap. XIV; Monroe, \emph{Textbook} (Macmillan, 1905), chap. XII; Parker, \emph{Modern Elementary Education} (Ginn, 1912), pp.~331-340. Popular accounts of the growth of science can be found in Buckley, Arabella B., \emph{A Short History of Natural Science} (Appleton), and Williams, H. S., \emph{Story of Nineteenth Century Science} (Harper). Spencer's \emph{Education} and Huxley's \emph{Science and Education} should be read. Further arguments for the study of science can be found in Coulter, J. M., \emph{The Mission of Science in Education (Science,} II, 12, pp.~281-293); Dryer, C. R., \emph{Science in Secondary Schools} (Prize Essay in \emph{The Academy,} May, 1888, pp.~197-221); Galloway, R., \emph{Education, Scientific and Technical} (Triibner, London, 1881); Norton, W. H., \emph{The Social Service of Science (Science,} II, 13, pp.~6445.); Pearson, K., \emph{Grammar of Science} (Macmillan, 1911), chap. I; Roberts, R. D., \emph{Science in the Nineteenth Century} (Cambridge University Press, 1901), chap. VII; Sedgwick, W. T., \emph{Educational Value of the Method of Science (Educational Review,} vol.~V, pp.~243ft.), and especially Youmans, E. L., \emph{Culture Demanded by Modern Life} (Appleton, 1867).

\hypertarget{present-day-tendencies-in-education}{%
\chapter{27. Present Day Tendencies in Education}\label{present-day-tendencies-in-education}}

\textbf{Outline}

\begin{itemize}
\tightlist
\item
  At the present time there is great progress in industrial, commercial, and agricultural training in the schools of Europe and America.
\item
  For a quarter of a century the educational systems of Europe have been giving attention to moral training, and of late there has been some discussion of the subject in the United States.
\item
  All the great nations now provide for the training of mental defectives, and for some time training has been afforded those defective in some sense organ.
\item
  The attempts at improved methods of teaching are witnessed by the study of industries in the experimental school of Dewey, by the formulation of a curriculum in terms of normal activities of other elementary schools, and by the `didactic apparatus' and the devices for learning the `three r's' of Montessori.
\item
  Methods of mental measurement are being devised for the elementary school subjects by Thorndike and others, and systems of measurement are being utilized in administration.
\item
  Darwin's theory of evolution has revolutionized our attitude, imagery, and vocabulary in education.
\item
  There is also a great variety of other educational movements in all grades of education.
\end{itemize}

\textbf{Recent Educational Progress.} - Because of the notable development of science and invention, which has been noted in the last chapter, the nineteenth century has often been referred to as the `wonderful' century. Such a term affords no better description of material achievement than of the remarkable progress that has taken place in education. Previous chapters have indicated the extent to which, through various movements, education has advanced and broadened in conception, but the near future of education will probably witness a much greater development. At the present time there are constant efforts at a modification and a reconstruction of education in the interest of a better adjustment of the individual to his social environment and of greatly improved conditions in society itself. It would, of course, be impossible to describe all of these movements even in the briefest manner, but some of the present day tendencies that appear most significant should now engage our attention.

The Growth of Industrial Training. - The movement that is perhaps most widely discussed to-day is the introduction of vocational training into the systems of education. There is now an especial need for this type Social reasons of training. Since the industrial revolution and the education. development of the factory system, the master no longer works by the side of his apprentice and instructs him, and the ambition of the youth can no longer be spurred by the hope that he may himself some day become a master. His experience is generally confined to some single process, and only a few of the operatives require anything more than low-grade skill. Nor, as a rule, will the employer undertake any systematic education of his workmen, when the mobility of labor permits of no guarantee that he will reap the benefit of such efforts, and the modern industrial plant is poorly adapted to supplying the necessary theoretical training for experts. Hence an outside agency - the school - has been called upon to assist in the solution of these new problems. To meet the demand for industrial education, all the principal states of Europe have maintained training of this sort for at least half a century, and the United States has in the twentieth century been making rapid strides in the same direction.

\textbf{Industrial Schools in Europe.} - In Germany, where this training is most effective, the work has for fifty years been rapidly developing through the \emph{Fortbildungsschulen} (see Fig. 55). The course in these schools at first consisted largely of review work, but the rapid spread of elementary schools soon enabled them to devote all the time to technical education. Training is now afforded not only for the rank and file of workmen in the different trades, but for higher grades of workers, such as foremen and superintendents. Girls are likewise trained in a wide variety of vocations. During the last twenty-five years there have also been developed continuation schools to furnish theoretical courses in physical sciences, mathematics, bookkeeping, drawing, history, and law. In North Germany there is a tendency to confine the courses to theoretical training, and leave the practical side to the care of the employers, but the South German states generally combine theoretical and practical work, and develop schools adapted to the industries of the various localities. Through the work of Kerschensteiner, Munich has even included an extra class in the elementary schools, to bridge the gap between school life and employment.

France goes still further, and, because of unsatisfactory conditions in apprenticeship, attempts to eliminate it altogether, and to furnish the entire industrial training through continuation schools articulating with the elementary system. The pupils are admitted at thirteen to the continuation schools (see p.~383) and obtain practice in the school workshops for three years. Woodwork is generally taught to the boys, but the other courses vary with local needs. Girls learn to make dresses, corsets, millinery, artificial flowers, and other industrial products. In England, grants were first made to evening industrial schools and classes in 1851, but twenty years later regular schools of science were organized, which had both day and evening sessions. In addition to these continuation schools, there have now been established higher elementary schools, which afford a four-year course in practical and theoretical science arranged according to local needs.

\textbf{Industrial Training in the United States.} - Industrial training first began to be offered in the United States during the latter half of the mneteenth century by means of a number of evening continuation schools. These were established through philanthropy in the larger cities, and included the Cooper Union and the Mechanics' Institute in New York; the Franklin Union and the Spring Garden Institute in Philadelphia; the Ohio Mechanics' Institute in Cincinnati; and the Virginia Mechanics' Institute in Richmond. The public schools at length followed this example, and of late years have organized evening classes in drawing, mathematics, science, and technical subjects. Day instruction was long delayed. It began in 1881 with the foundation of the New York Trade School, but at the end of twenty years there were only two others, - the Williamson Free School of Mechanical Trades near Philadelphia and the Baron de Hirsch Trade School in New York. Later the development was more rapid, and since 1906 several hundred day trade schools have been organized, mostly through public support, in the larger cities of the country. These schools are mostly for youths between sixteen and twenty-five, but `preparatory trade schools' for younger boys have also been started in New York, Massachusetts, and other states. Higher training to equip leaders for the industries has also come to be furnished through endowed secondary schools and technical high schools in a number of cities. A recent variety of vocational training is the `part-time' plan, by which students are given some theoretical and formal training in a regular high school or college, while they are obtaining their practical experience. This alternation of practical and theoretical training is sometimes carried on in a single institution, or even within a commercial establishment itself.

\textbf{Conditions Commercial Education in Europe and America.} - But the modern development of vocational training throughout the leading countries has not been confined to industrial lines. With the extension of the sphere of commerce and the development of its organization that have taken place in the nineteenth century, it has come to be recognized that preparation is essential for a business career. Only recently, however, has this training been felt to be a proper function of the schools, since for many years it was opposed by educators as sordid and commercializing, and by business men as unpractical and ineffective. Both classes have now been brought to realize the need of mutual support, and the rapid growth of commercial education indicates an appreciation of its usefulness.

Germany is generally admitted to lead in commercial education. The growth of this training has taken place since 1887, but there is now offered under state control a unified and thorough preparation for any line of business. Besides private continuation schools, in which a course of three years in modern languages and elementary commercial studies can be obtained, there have grown up both public secondary schools and university courses in which a thorough general education and theoretical work in commerce, as well as a practical and technical training, are provided (Fig. 55). England and France have been rather indifferent to commercial education, different. In both countries until very recently schools have been few, and the number of pupils in each has been small. But now continuation schools, free evening courses, and private classes have sprung up, and in a few large cities commercial schools of secondary and even higher grade have been established. In the United States commercial training began by the middle of the nineteenth century through private enterprise with classes in bookkeeping, and later with `business colleges.' Despite the name of the latter, the course is narrow and is generally shaped by pecuniary aims. During the last two decades of the nineteenth century high and normal schools began to offer commercial instruction, but until the twentieth century the courses were only tolerated as a necessary evil, and largely imitated those of the business colleges. Since then many cities have opened high schools of commerce, and university schools and colleges of commerce have arisen, and even a score of years before this development the Wharton School of Finance and Commerce was started at the University of Pennsylvania.

\textbf{Recent Emphasis upon Agricultural Training.} - A similar development has of late been taking place in agricultural education. France and Germany offer elementary instruction in agriculture, while the former has also introduced the subject into the normal schools, and the latter has established a secondary agricultural institution open to students at the close of their sixth year in the \emph{Realschule.} Through the feeling that the United States must become the great agricultural nation, and that the traditional methods of agriculture have been exceedingly wasteful, this country especially has been emphasizing that type of vocational education. The land grant colleges, first endowed by act of Congress in 1862, have greatly stimulated interest in the subject, and later Congress added other sources of revenue, and has recently furnished appropriations for instruction in the teaching of agriculture and for extension work in agriculture. Thus the way has been prepared for the introduction of the subject into the high school and grades. There are now at least one hundred agricultural high schools in the United States, and agriculture is taught as a branch of study in several thousand high and elementary school systems.

\textbf{Moral Training in the Schools To-day.} - But present day tendencies in education have to do with more than the material side of civilization. There is a growing sentiment in favor of moral instruction in the schools. There are many reasons why this need should be especially felt in the complex business life of to-day. When men work for impersonal corporations, sell products to people they never see, or intrust their welfare to officials whose names are scarcely known, one strong factor making for honesty and virtue, that of personal relations, is lost. Moreover, as a result of the weakening of old religious sanctions, the new conditions in large cities, and other causes, moral traditions are in need of being buttressed.

The educational systems of Europe have for a quarter of a century given more or less attention to moral training. In France this training has been purely secular and excluded all religious elements. But the education of England and Germany has always associated the teaching of morality with religion. In England, the `board' schools have furnished religious instruction of a nonsectarian character, but the religious training of the `voluntary' schools has occupied more time and has stressed the creed and denominational teaching of some church, usually the Church of England (see pp.~380 f.). The contest over religious teaching since the Act of 1902 (see p.~390) caused a self-constituted commission, with Michael E. Sadler as chairman, to investigate the subject of moral instruction, and in 1908-1909 it presented a large and illuminating report. In Germany the moral and religious instruction in all elementary schools is sectarian, and Catholic and Protestant schools are alike supported, wherever needed, at public expense. During the past decade there has been considerable discussion in the United States concerning moral education. In response to the demand for an investigation of the subject, a committee of the National Education Association in 1908-1909 made a report upon various phases of moral training, and recommended special instruction in ethics, not in the form of precepts, but through consideration of existing moral questions. In 1911 the Religious Education Association, whose convention in that year was devoted to moral training, gave in its \emph{Journal} a broad summary of the progress of moral education in the United States. The report reveals a wide difference of opinion and practice, but an evident tendency to trust other agencies than direct moral instruction. As a rule, state legislation seems as yet to have failed to provide a general system of training, but has confined itself to specific subjects, such as instruction in citizenship, the effects of alcohol and narcotics, and the humane treatment of animals.

\textbf{The Development of Training for Mental Defectives.} - One of the most patent evidences of the growth of the humane spirit in modern times is found in the universal attention now given to the education of mental defectives. This movement was given its greatest impulse through Edouard Seguin, who came to the United States in 1850 and developed his methods here. His general plan was to appeal to the mind through the senses by means of a training of the hand, taste and smell, and eye and ear. He used pictures, photographs, cards, patterns, figures, wax, clay, scissors, compasses, and pencils as his chief instruments of education. The stimulus he gave to the training of defectives has been epoch-making, and his `physiological' methods have remained the chief means of education. Although there has grown up a tendency to introduce intellectual elements into the training of the feeble-minded, the advantages of such a procedure are doubtful.

All the great nations now provide schools for the training of defectives. Germany has over one hundred institutions, with some twenty thousand pupils in them, although nine-tenths of them are not supported by the state, but are under church or private auspices. These schools generally stress manual education, but give some attention to intellectual lines, especially to speech training. There are but few schools for defectives in France, aside from the two near Paris and the juvenile department of the insane hospital at Bicetre, but these institutions largely follow the physical work formulated by Seguin. In London there is one excellent institution with two thousand pupils, where manual training consti- tutes almost the entire course. But there are five other schools so located as to serve the various parts of Eng- land, in which the training is rather bookish and emphasis is especially laid upon number work.

Thanks to the start given by Seguin, America has taken up the education of defectives more fully than any other country. Schools for the feeble-minded now exist in almost all the states, and there are some thirty-five or forty private institutions of considerable merit. Not far from twenty thousand defectives are being trained, although this is probably only about one-tenth of the total number of such cases in the country. The type of education differs greatly according to the institution, ranging from almost purely manual training to a large proportion of the intellectual rudiments, but in all the work is adapted to the various grades in such a way as to raise them a little in the scale of efficiency and to keep them as far as possible from being a burden to themselves and to society. Likewise, special clinics and investigations, like those of Lightner Witmer of the University of Pennsylvania and of H. H. Goddard of the Training School at Vineland (New Jersey), are greatly adding to our knowledge of the best methods for training defectives.

\textbf{Education of the Deaf and Blind.} - Persons defective in some sense organ, but otherwise up to the standard, have likewise for some time been receiving an education that will minimize the difficulty. There have been two chief methods for teaching the deaf. The manual or `silent' method of communication was invented by the Abbe de l'Epee in Paris during the latter part of the eighteenth century, and his school was adopted by the nation in 1791. The other method, the `oral,' by which the pupil learns to communicate through reading the movements of the lips, was started in Germany early in the eighteenth century, but was not employed to any great extent until the middle of the next century. Most countries now use the oral method exclusively, or in connection with the manual system. In the United States practically every commonwealth now has one or more schools for the deaf, and since 1864 even higher education has been furnished by Gallaudet College at Washington.

The first instruction of the blind through raised letters was given toward the end of the eighteenth century by Abbe Haiiy at Paris. While his schools, owing to his lack of judgment, were failures, the idea spread rapidly. Early in the nineteenth century there were one or more schools in each of the leading countries of Europe, and a generation later institutions of this sort were started in the United States. In schools for the blind or deaf, industrial training has in most instances been added to the intellectual (see p.~300), in order to fit every individual to be an independent workman in some line. Even pupils, both deaf and blind, like Laura Bridgeman and Helen Keller, have had their minds awakened through the sense of touch.

\textbf{Recent Development of Educational Method; Dewey's Experimental School.} - Nor has the past century witnessed any cessation of the attempts at improved methods of teaching. Various suggestions and systems have been put forward and many have had an important effect upon school procedure. It is impossible, however, to discuss any except a few of the more influential and prominent, and these can be considered but briefly. The occupational work of Professor Dewey and Colonel Parker's scheme of concentration have marked the Colonel growth of a body of educational theory and practice that contributions. places the methods of to-day far in advance of anything previously known. The combination and modification of Ritter, Herbart, and Froebel worked out by Parker have perhaps received sufficient attention (see pp.~293, 350, and 364), but we may at this point outline a little more fully the contributions made by John Dewey, who has probably been the leader in the reconstruction that has taken place in education almost since the twentieth century began.

The methods of Dewey were developed in an experimental elementary school connected with the University of Chicago and under his supervision from 1896 to 1903. The school did not start with ready-made principles, but sought to solve three fundamental educational problems. It undertook to find out (1) how to bring the school into Purpose closer relation with the home and neighborhood life; (2) how to introduce subject-matter in history, science, and art that has a positive value and real significance in the child's own life; and (3) how to carry on instruction in reading, writing, and figuring with everyday experience and occupation as their background ``in such a way that the child shall feel their necessity through their connection with subjects which appeal to him on their own account.'' The plan for meeting these needs was found largely in the study of industries. Since industries are most fundamental in the thought, ideals, and social organization of a people, these activities must have the most prominent place in the course of a school. ``The school cannot be a preparation for social life except as it reproduces the typical conditions of life.'' The means used in furnishing this industrial activity were evolved mainly along the lines of shopwork, cooking, sewing, and weaving, although many subsidiary industries were also used. These occupations were, of course, intended for a liberalizing, rather than a technical purpose, and considerable time was given to an historical study of them (Fig. 56). Dewey declares: ``The industrial history of man is not a materialistic or merely utilitarian affair. It is a matter of intelligence. Its record is the record of how man learned to think, to think to some effect, to transform the conditions of life so that life itself became a different thing. It is an ethical record as well; the account of the conditions which men have patiently wrought out to serve their ends.''

It can be seen how fully this plan is in accord with the real principles of social cooperation and expression of individual activities underlying the work of Froebel; and ``so far as these statements correctly represented Froebel's educational philosophy,'' Dewey generously grants that ``the school should be regarded as its exponent.'' But these industrial activities of the Chicago experimental school were not in the least suggested by Froebel's work, and were far more expressive of real life. They never became as stereotyped and external as the but not as gifts or even as the occupations of the kindergarten have generally been. Dewey is insistent that this training shall be carried on not for the purpose of furnishing facts or principles to be learned, but for enabling the child to engage in the industrial occupations in miniature. ``The school is not preparation for life: it is life.'' Hence this training is superior to the occupations of Froebel in that ``it maintains a balance between the intellectual and the practical phases of experience.'' Where Froebel has held to the construction of beautiful things in mechanical ways, Dewey emphasizes the ordinary activities and experiences of life, even though the expression of these be crude. The child should be ``given, wherever possible, intellectual responsibility for selecting the materials and instruments that are most fit, and given an opportunity to think out his own model and plan of work, led to perceive his own errors, and find how to correct them.'' Thus the work was never ``reduced to a mere routine or custom and its educational value lost.'' As a result, too, it was the consensus of opinion that ``while the children like, or love, to come to school, yet work, and not amusement, has been the spirit and teaching of the school; and that this freedom has been granted under such conditions of intelligent and sympathetic oversight as to be a means of upbuilding and strengthening character.''

\textbf{Other Experiments in Method.} - Hence, while the Chicago school is now at an end, the experiment in education developed there is still yielding abundant fruitage. It has stimulated similar undertakings elsewhere, and has been the largest factor in determining the theory and practice of the present day. Either as a result of Dewey's work or through independent thought, there has sprung up an important group of schools in which there is clearly an effort to bring boys and girls of elementary school age into more intimate relation to community life about Schools on them. Such are the Gary (Indiana) Public Schools, the Francis W. Parker School of Chicago, the Elementary School at the University of Missouri, the Pestalozzi-Froebel School of Berlin, the Abbotsholme School in Derbyshire (England), and a number of others.

A good illustration is afforded in the school developed University by Junius L. Meriam at Columbia, Missouri, although it has not been given much publicity. Its function is to help children do better in all those wholesome activities in which they normally engage. The school does not attend to the `three r's' as such, but specifically to parits purpose and ticular activities of children, including (1) play, (2) observation, (3) handwork, and (4) stories, music, and art. These four `studies,' representing real life, irrespective of the school, constitute the curriculum, and the `three r's' are studied only as they are needed. Their content, therefore, being used, as in life, in meeting real needs, is studied most effectively.

An experiment that has attracted widespread interest is that worked out in the Gary school system by William A. Wirt. While the achievement is mostly in the way of a remarkable organization and administration that have undertaken to make available ``all of the educational opportunities of the city all of the time for all of the people,'' the teaching has to some extent been carried on so as to reveal to the pupils ``that what they are doing is worth while.'' The school plant includes a playground, garden, workshop, social center, library, and traditional school, and it has been shown that these agencies, when properly organized, ``secure the same attitude of mind toward the reading, writing, and arithmetic that the child normally has for play.'' All the other schools that have been mentioned above make similar attempts to enable the children to get into closer touch with their environment. While each of them approaches the problems of elementary training from a different angle, they are all in harmony with the spirit of Dewey and present day theory.

\textbf{The Montessori Method.} - But probably the most spectacular development in educational procedure is that originating with Maria Montessori at Rome. Yet the Montessori method, except for some elements adapted from Seguin (see p.~426), is largely a combination of several of the concepts found in Rousseau, Pestalozzi, and Froebel, and fails to grasp the larger vision of education that appears in present-day theory, such as Dewey's. Like Rousseau and Froebel, Montessori holds fundamentally to the rightness of child nature and consequently to the liberty of the pupil, but she does not, like Dewey, realize that education is itself life and that the activities of real life should be utilized in training. Moreover, the sense training, which Montessori herself considers the most distinctive feature of her system, is neither original nor psychologically sound. Montessori began as a teacher of defectives, and her `didactic apparatus' and methods are largely borrowed from Seguin. Exercises of this sort are of great value in training defectives, but the assumption of their usefulness in the education of normal children is more doubtful. They are intended to train the senses to general powers and discriminations, and seem to be defended simply upon the basis of faculty psychology and the outworn theory of `formal discipline' (see p.~182 f.).

The feature of the Montessori method, however, that has attracted most attention is its apparent success with the formal elementary studies, especially the facility, enthusiasm, and speed with which it has enabled the pupils to learn to write. Montessori has carefully analyzed the process of writing and devised three exercises by which this art is unconsciously learned by three or four year old children in Italy. If this training can be applied to unphonetic languages, like the English, it may possibly be regarded as a contribution. It is evident, however, that Montessori lays too much stress upon the acquisition of the formal studies and starts them at too early an age. In this she fails to appreciate Froebel's great contribution of a school without books, and certainly does not realize, with Dewey, that the main purpose of education is to give a child some control of his social environment and that for this there are activities of more importance to child life than the school arts. Within a few years it will probably be difficult to understand \emph{the furore} that has been created by the Montessori methods.

\textbf{The Statistical Method and Mental Measurements in Education.} - One of the most significant of the present day movements is the application, especially in the United States, of scientific, statistical methods to problems of education. Statistics have long been used, though often without clearness or accuracy, in reports of school administration, but it remained for this century to apply to the various phases of education the same general technique and approximately the same precision as that long demanded by the physical and biological sciences. Quantitative, unambiguous statements are now sought and secured not only for the phenomena of attendance, retardation, expenditures, and the like, but also for the relative and absolute amounts of knowledge. As a consequence, emphasis has been placed upon the results of education rather than upon the declaration of intentions.

Probably the first scholar to apply the scientific principles of statistics to education was Edward L. Thorndike of Columbia University. In his \emph{Educational Psychology} he illustrates how a quantitative description of individual differences and of the factors that condition them is necessary to throw real light upon educational theory and practice, and in his \emph{Mental and Social Measurements} he presents the details of the method. Subsequently he maintained, in the face of much opposition, that scales, as objective and as impersonal as possible, should and could be devised for measuring variations in ability and changes that take place as a result of natural growth and instruction. Such scales, beginning at an ascertained zero and progressing by regular steps to a point near perfection, are, because of the complexity of their elements, difficult to construct, but they have been set forth more or less tentatively by various investigators for the measurement of achievement in handwriting (Fig. 57), arithmetic, English composition, spelling, drawing, freehand lettering, and reading respectively. Other scales to measure ability in the several high school subjects may be expected soon.

Studies are also being made in several universities to determine the relative importance of the numerous factors in methods of teaching. This is done by conducting experiments with hundreds or thousands of children to find out by the most accurate measurement yet devised the amount of progress in learning that is wholly due to the presence of some one factor of method in the technique of class-room exercises. Educational psychology has revealed the qualitative significance of many of the single elements in the very complex procedure that we have called a `method of teaching,' and this new type of research aims to determine the quantitative significance of each of these several elements of method as factors in the production of abilities. A. Duncan Yocum of the University of Pennsylvania has formulated a considerable number of tests, and, by preliminary experimentation, has determined the conditions under which they may with a high degree of accuracy be given to groups of students engaged in actual school work under ordinary class-room conditions. His students have made a number of tentative, but suggestive studies, which have not yet been published. Milo B. Hillegas of Columbia University and others are engaged on certain aspects of this general type of research. There is reason, therefore, to believe that we may sometime be able to measure with as much accuracy the efficiency of well-defined educational processes as we are now able to measure educational products. If this can be attained, the technique of class-room teaching and of educational supervision will begin to rest on a really scientific basis.

Moreover, by the use of the improved statistical method and of scales, studies of greatly increased value have been made of fatigue, retardation, elimination, and of other social and mental phenomena of individual children. And in 1911, with the reports of Paul H. Hanus of Harvard University and Ernest C. Moore of Yale University upon the school systems of Montclair and East Orange, New Jersey, there began to be instituted those measurements and consequent criticisms of whole school systems, known as `educational surveys.' These scientific reports have been extended to the educational work of a large number of cities and states throughout the Union. They are intended to enable school officers and patrons to comprehend with more definiteness the absolute, as well as the relative, achievements of their children.

\textbf{Education and the Theory of Evolution.} - A most characteristic influence in education to-day has come through the theory of evolution of Darwin (Fig. 51). This fruitful hypothesis came to be generally accepted during the last quarter of the nineteenth century as the guiding principle of education, and has constantly increased the illumination it has shed upon the educational process. It has given an entirely new meaning to education, and has greatly modified the course of study and revolutionized the method of approaching educational problems. It has wrought very much the same changes in the treatment of intelligence that it did in the biological sciences. Consciousness is no longer regarded as a fixed set of entities, but as a developmental process. Instead of classifying and cataloging mental processes in fixed groups, efforts are made to study their growth from the standpoint both of the race and of the individual. Studies of mental development in the race, begun by Darwin's \emph{Descent of Man,} which recognized `sexual' and `social selection,' as well as `natural selection,' have been continued by numerous investigators, and equally extensive researches have also been latterly made in genetic psychology, child study, mental development, and adolescence. Both observation and experimentation have been introduced into the study of mental processes. Even more revolutionary than this actual increase in cha e in knowledge, however, is the change that has taken place imagery and in the conception, imagery, and terminology of education. Writers upon education constantly employ the language of evolution. Educational discussions are now filled with such terms as `variation,' `selection,' `adjustment,' and `adaptation,' and such concepts dominate all educational thinking. If educational leaders of half a century ago could be present to-day at a gathering of educational thinkers, they would find themselves listening to what would seem to them almost a foreign language.

\textbf{Enlarging Conceptions of the Function of Education.} - Such are a few of the chief tendencies and advances that are being made in education to-day. There is also a great variety of other educational movements, almost too numerous to be mentioned. In the organization and administration of the public schools there is a decided tendency toward centralization in educational activities, corresponding to the centralization in industrial and political affairs. The United States Bureau of Education and the various State Departments of Public Instruction have had their functions much enlarged and their activities greatly increased. There are also such matters as the new procedure in school hygiene, arising from the modern attitude toward the prevention of disease; new health regulations, as a result of having so many children housed in the same buildings; medical inspection, open-air schools, and better nourishment; and new tendencies in school architecture. Likewise we find progressive legislation on compulsory school attendance; more extensive training of teachers; a rapid recognition of education as a profession; the organization of various professionaiitypes of teachers' associations; and the development of teaching. educational journalism. Secondary education is also being greatly extended and largely reorganized. `Junior Reorganizahigh schools,' combining the upper grades of the elemen- ary and higher tary school with the lower grades of the secondary school, and thus bridging the gap, are being widely introduced into American cities, and a variety of propositions for a six-year course are being seriously entertained. In connection with higher education there are such new tendencies as university extension, correspondence courses, summer sessions, university interest in the practical problems of the people, the correlation of the first two years of college with the secondary school, more flexible entrance requirements, an increasing number of fields of professional work, and, above all, the professional training of teachers through Departments of Education, Teachers Colleges, and Schools of Education. With this is connected the scientific study of Education, both in graduate courses and independent investigations.

Similar efforts to secure economy, guard health, improve method, and cause education to serve democratic ideals are everywhere apparent. Educational theory and practice are in a constant flux, and have entered upon a most distinctive epoch of experimentation, change, and improvement. While such a situation is not without its perils, and each proposal should be carefully scrutinized before acceptance, the present tendencies are in the main a sign of progress and life.

Supplementary Readings

Graves, \emph{In Modern Times} (Macmillan, 1913), chap. XI; Monroe, \emph{Textbook} (Macmillan, 1905), chaps. XIII-XIV. For the special tendencies mentioned, the following works may be consulted: Cooley, E. G., \emph{Vocational Education in Europe} (Chicago Commercial Club, 1912); Hanus, P. H., \emph{Beginnings in Industrial Education} (Houghton, Mifflin, 1908); Haskins, C. W., \emph{Business Education and Accounting} (Harper, 1904); Adler, F., \emph{Moral Instruction of Children} (Appleton, 1895); Palmer, G. H., \emph{Ethical and Moral Instruction in Schools} (Houghton, Mifflin, 1909); Goddard, H. H., \emph{Education of Defectives \{Monroe's Cyclopadia of Education);} Bell, A. G., \emph{Deaf Mute Instruction in Relation to the Work of the Public Schools;} Armitage, T., \emph{Education and Employment of the Blind} (Harrison \& Sons, London, 1886); Dewey, J., \emph{The School and Society} (University of Chicago Press, 1899), and \emph{Elementary School Record} (University of Chicago Press, 1000); Montessori, Maria, \emph{The Montessori Method} (Translated by Anne E. George, Stokes Co., New York, 1912); Kilpatrick, W. H., \emph{The Montessori Method Examined} (Houghton, Mifflin, 1914); Ayres, L. P., \emph{Measuring Educational Processes through Educational Results (School Review,} May, 1912); Strayer, G. D., \emph{Standards and Tests for Measuring the Efficiency of Schools} (Report of the Committee of the National Council of Education in the \emph{United States Bureau of Education Bulletin,} 1913, No.~13); Thorndike, E. L., \emph{The Measurement of Educational Products (School Review,} May, 1912).

\hypertarget{retrospect-and-prospect}{%
\chapter{28. Retrospect and Prospect}\label{retrospect-and-prospect}}

\textbf{Outline}

\begin{itemize}
\tightlist
\item
  Evolution in education may be interpreted from the standpoint of the development of individualism. Individualism was first fully recognized in the teachings of Christ, but was repressed during the Middle Ages. While it reappeared during the Renaissance, Reformation, and other movements, it soon lapsed, but a complete break from tradition occurred with Rousseau in the eighteenth century.
\item
  For a time individualism dominated, but education since then has endeavored to afford latitude to the individual without losing sight of the welfare of society.
\end{itemize}

The Development of Individualism. - The discussion of present day tendencies that has just been given, together with the account of educational evolution in the preceding chapters, serves to show how far modern times have progressed in the ideals and practice of education. This may perhaps be best appreciated from the standpoint of the development of individualism. To follow such an interpretation back to the beginning of the history of education, it may be stated that during the day of primitive man no real distinction was made between \^{}\^{}\^{}MHltic society and the individual, and practically all advance- tendencies

•" 1 during the days

ment was impossible, for no one looked much beyond of primitive the present. With the appearance of the transitional period in the Oriental countries, the individual had begun nations,

Roman civili zations,

Christian de velopment,

to emerge, but was kept in constant subjection to the social whole, for man was quite enslaved to the past.

Athenian and \^{}\textsuperscript{s} \textsuperscript{tne} J\textsuperscript{ew}ish; Athenian, and Roman civilizations developed, the beginnings of individualism were for the first time clearly revealed, and some regard was had for the future. Then, through the teachings of Christ, there came to be a larger recognition of the principle of individualism and the brotherhood of man. Owing to a necessity for spreading these enlarged ideals among a barbarous horde of peoples, individualism was repressed,

\emph{TM\textsuperscript{d}J\textsuperscript{he}} \textsuperscript{Middle} and throughout the Middle Ages the keynote was submission to authority and preparation for the life to come. The cultural products of Greece and Rome largely disappeared, and all civilization became restricted, fixed, and formal.

But the human spirit could not be forever held in bondage, and, after almost a millennium of repression and uniformity, various factors that had accumulated within the Middle Ages produced an intellectual awakening that we know as the `Renaissance.' Its vitality lasted during the fifteenth century in Italy and to the close of the sixteenth in the Northern countries, but by the dawn of the seventeenth century it had everywhere degenerated into a dry and mechanical study of the classics. This constituted a formalism almost as dense as that it had superseded, except that linguistic and literary studies had replaced dialectic and theology. A little later than the spread of the Renaissance, though overlapping it somewhat, came the allied movement of the `Reformation.' This grew in part out of the disposition of the Northern Renaissance to turn to social and moral account the revived intelligence and learning. Yet here also the revival failed in its mission, and the tendency to

the Renaissance,

the Reformation,

rely upon reason rather than dogma hardened into

formalism and a distrust of individualism. Again, in

the seventeenth century, apparently as an outgrowth of

the same forces, intellectual activity took the form of a

search for `real things.' The movement that culminated

in `sense realism' appeared, but this small and crude \textsuperscript{and} \textsuperscript{realism};

beginning of the modern scientific tendency was for some

decades yet held within limits. Associated with this

realistic tendency, on the religious and political sides

also appeared a quickening in such forms as `Puritanism'

and `Pietism,' which likewise degenerated eventually into tanism and

\emph{e} ,. . j 1 . Pietism;

a fanaticism and hypocrisy.

The Harmonization of the Individual and Society. - Thus the way was opened for the complete break with JUJJj \^{}°\textsuperscript{u}\textbar{}\textbar{}f\textsuperscript{au} tradition and authority that occurred in the eight- "trudge

\emph{\textsuperscript{J} \textsubscript{m}} ° tendency.

eenth century. This tendency, while in France at least most destructive and costly, was the inevitable result of the unwillingness to reshape society and education in accordance with changing ideals and conditions. Hence Rousseau undertook to shatter all educational traditions. But his recommendation of isolated education, so palpable in its fallacies, prepared the ground for the numerous social, scientific, and psychological tendencies (see pp.~218-222) that were destined to spring up in modern education and for the consequent improvement in the aim, organization, content, and method of education. Of course modern education has advanced infinitely beyond anything implied by Rousseau or even the later reformers of the past century, but it is out of his attempts at destruction that has grown this nobler structure. For a time individualism triumphed and

The present

ground authority under its heel, but when this extremity had been passed, the problem became how to harmonize the individual with society, and to develop personality progressively in keeping with its environment. Thus the nineteenth and twentieth centuries have put forth conscious efforts to justify the eighteenth and to bring out and develop the positions barely hinted at in its tendencies in negations. It is not alone the individual as such that

education .

seemtohar- has been of interest in the modern period, but more and dividual more the individual in relation to the social whole to

tho\^{}ofsodety. which he belongs, as only in this way can the value of his activities be estimated.

This is revealed in the works of those who followed

Rousseau, and especially in the attempts of recent

, , . . educational philosophers to frame a definition of educa

Recent defim- .

tions of educa- tion that shall recognize the importance of affording

tion show this. , . , , \ldots{}\ldots{} \ldots{}. . , , ,

latitude to the individual without losing sight of the welfare of the social environment in connection with which his efforts are to function. Thus Butler, though recognizing the individual factor, especially stresses the social by declaring education to be ``the gradual adjustment of the individual to the spiritual possessions of the race.'' Then he further declares: ``When we hear it sometimes said, `All education must start from the child,' we must add, `Yes, and lead into human civilization;' and when it is said on the other hand that `all education must start from a traditional past,' we must add, `Yes, and be adapted to the child.'\,'' And the balance between the two factors of the individual and society is even more explicitly preserved in Dewey's statement ``that the psychological and social sides are organically related, and that education cannot be regarded as a compromise between the two, or a superimposition of one upon the other.'' In the same way Bagley has made `social efficiency' the main aim in educating the individual to-day, and both elements are carefully considered by all modern writers in discussing educational values. Thus the central problem in education of the twentieth and succeeding centuries is to be a constant reorganization of the curriculum and methods Jona\^{}probiem of teaching, and this reconstruction must be such as to \textsuperscript{of} \textsuperscript{the} \textsuperscript{future}harmonize a due regard for the progressive variations of the individual with the welfare of the conservative institutions of society. It must include a continual effort to hand on the intellectual possessions of the race, but also to stimulate all individuals to add some modification or new element to the product. In this way there may develop unending possibilities for both the individual and society.

\textbf{Supplementary Reading}

Graves, F. P., \emph{History of Education before the Middle Ages} (Macmillan, 1909), chap. XII; \emph{History of Education during the Transition} (Macmillan, 1910), chap. XXIII; \emph{History of Education in Modern Times} (Macmillan, 1913), chap. XII; Monroe, P., \emph{Textbook in the History of Education} (Macmillan, 1905), chap. X.

\bibliography{book.bib,packages.bib}

\end{document}
